\documentclass[preprint,12pt]{elsarticle}

\usepackage{amsmath,amssymb,amsfonts,amsthm}
\usepackage{array}
\usepackage{booktabs}
\usepackage{graphicx}
\usepackage{url}
\usepackage{xcolor}
\providecommand{\nolinkurl}[1]{\url{#1}}

\journal{Computer Methods in Applied Mechanics and Engineering}
\emergencystretch=4em

\newtheorem{assumption}{Assumption}
\newtheorem{definition}{Definition}
\newtheorem{lemma}{Lemma}
\newtheorem{theorem}{Theorem}
\newtheorem{proposition}{Proposition}
\newtheorem{corollary}{Corollary}

\newcommand{\method}{Gauss6\slash FullVA}
\newcommand{\tfe}{TFE}
\newcommand{\figpath}{.}
\newcolumntype{P}[1]{>{\raggedright\arraybackslash}p{#1}}
\newcommand{\paperfig}[3]{%
  \begin{figure}[htbp]
    \centering
    \IfFileExists{#1}{%
      \includegraphics[width=0.90\linewidth]{#1}%
    }{%
      \fbox{\parbox{0.86\linewidth}{Missing figure: \nolinkurl{#1}}}%
    }
    \caption{#2}
    \label{#3}
  \end{figure}
}

\begin{document}

\begin{frontmatter}

\title{A Conditional Sixth-Order Lie-Group FullVA Integrator for Lower-Pair Mechanisms}

\author[inst1]{Jingquan Wang\corref{cor1}}
\ead{jingquanw@wisc.edu}
\cortext[cor1]{Corresponding author.}

\affiliation[inst1]{organization={University of Wisconsin--Madison},
  addressline={Department of Mechanical Engineering, 1513 University Avenue},
  city={Madison},
  state={WI},
  postcode={53706},
  country={United States}}

\begin{abstract}
Lower-pair mechanisms require Lie-group integrators that synchronize
position, velocity, and acceleration constraints. This paper proposes
\method{}, a three-stage Gauss collocation Lie-group integrator. The
contribution is a monolithic
FullVA stage object: stage kinematics, Newton--Euler balance, and lower-pair
position-, velocity-, and acceleration-level constraints are coupled in a
square nonlinear system with Lie-group endpoint reconstruction. The main
theorem proves a conditional sixth-order same-initial-state reported-grid
estimate for the selected smooth branch by transferring the
classical three-stage Gauss local defect through the implemented FullVA stage
root, endpoint-closure perturbation, and inexact-Newton perturbation. The
retained interfaces are compact-tube regularity, endpoint/stability,
implementation binding, and branch-selected solver scale
\(\eta_h^{\rm tube}\le c_\eta h^7\). The theorem's residual-value
input is discharged by a direct \(96+36\)
bridge on the lifted Gauss stage: the 96
non-dynamic FullVA rows have an \(O(h^7)\) row certificate and the 36
Newton--Euler rows vanish by direct substitution. This is neither an
empirical residual fit nor primitive/Taylor closure; the latter remains
non-active, with five open lift/reduction obligations and no accepted
\(162\)-subterm bound. On a
smooth cylindrical-chain benchmark, the implementation
records position/velocity orders 7.161/7.066 as branch-consistency
diagnostics. Four ASME-style
mechanisms exercise the same residual vocabulary: single- and double-pendulum
rows provide dynamic-order evidence, while four-link and slider-crank rows
provide closed-loop constraint/reaction consistency and coarse-window
mechanism diagnostics.
Fixed production tolerances remain finite-run diagnostics; residual-to-error
promotion, full external same-test reproduction, and complete source-paper
\tfe{} residual replacement are not claimed. The comparison to \tfe{} is kept
at the formal-order and diagnostic levels.
\end{abstract}

\begin{keyword}
Lie group integrator \sep index-3 DAE \sep multibody dynamics \sep
lower-pair mechanisms \sep Gauss collocation \sep temporal finite elements
\end{keyword}

\end{frontmatter}

\section{Introduction}

Implicit integration of constrained multibody systems must keep position,
velocity, and acceleration constraints synchronized while respecting the
geometry of finite rotations. This is especially important for absolute- or
lower-pair formulations, where a step update that looks accurate in position
can still carry endpoint velocity drift or inconsistent reaction dynamics.
Recent time finite element methods on Lie groups have shown that high-order
integration of frictional index-3 differential-algebraic equations is feasible
and practically valuable \cite{chaturvedi2026tfe}. The present work asks a
more specific numerical-methods question: can a Gauss-collocation Lie-group
FullVA residual be organized so that lower-pair position, velocity, and
acceleration constraints are enforced monolithically while retaining a
theorem-level conditional sixth-order claim on the retained smooth branch?

The resulting method is \method{}, a Gauss collocation FullVA integrator whose
sixth-order claim is tied to the branch smoothness, residual-defect,
endpoint/stability, implementation-binding, and nonlinear-solver contracts
stated below. Its novelty is not the existence
of Gauss collocation itself; it is the monolithic Lie-group FullVA residual
organization for lower-pair mechanisms, where Gauss stage kinematics,
$\mathrm{SO}(3)$ retraction variables, Newton--Euler dynamics, and
position/velocity/acceleration lower-pair rows are solved in one square
nonlinear system. Temporal finite elements enter this paper as a formal-order
comparator only after the method and theorem have been fixed.  Their role is
to mark the local \(m=3\) order-five formula target and the diagnostic
boundary for separate same-test reproduction, not to supply an implementation
superiority claim, a full-\tfe{} replacement, or a theorem premise.
Accordingly, the central mathematical object is the branch-selected
\method{} residual together with its retained-interface order theorem. The
comparison and reproducibility material states the formal-order, external
same-test, and source-package scope without modifying the theorem.
The proof organization follows the Lie-group constrained-BDF literature only
as an obligation ordering: fix the discrete object, insert the smooth branch,
prove the local defect, and then propagate it globally; all constants and
residual inputs used below are supplied by the FullVA lemmas and certificates
in this manuscript.
In particular, the Wieloch--Arnold BLieDF proof is invoked as a proof-order
template, not as an estimate source: its BDF order theorem, hidden-constraint
and multiplier estimates, and Lie-algebra recurrence constants are not
imported into the present Gauss/FullVA proof.

The contributions are as follows.
\begin{enumerate}
  \item a monolithic square 132-row Lie-group FullVA stage residual for
  lower-pair mechanisms;
  \item a complete one-step algorithm and conditional sixth-order
  same-initial-state reported-grid theorem obtained by
  local-defect-to-global-error transfer on the compact-tube Gauss-predictor branch
  under retained endpoint/stability, implementation-binding, and solver-scale interfaces,
  with finite-window smooth slopes 7.161/7.066 reported only as
  branch-consistency diagnostics rather than solver-scale proof discharge;
  \item method-side dynamic-order evidence for the single- and
  double-pendulum examples, and mechanism-coverage evidence for the four-link
  and slider-crank examples under the same FullVA vocabulary;
  \item a scope boundary that limits \tfe{} material to the local
  order-five \(m=3\) formula target and keeps implemented source-paper
  comparisons diagnostic until the same-test promotion criteria are met.
\end{enumerate}
A separate reproducibility appendix records external reproduction,
full-\tfe{}, primitive-route, and source-package boundaries. It
is included to delimit non-claims and reproduce the accepted rows, not as an
additional method contribution or theorem input.

Table~\ref{tab:intro-claim-hierarchy} fixes the hierarchy used throughout the
paper.  The proof contribution is the analytic residual-to-local-defect
transfer for the implemented \method{} map.  Numerical rows are downstream
evidence for the accepted path.  External same-test reproduction,
source-package, and full-\tfe{} replacement records are boundary evidence;
they keep stronger external claims out of the manuscript and are not premises of
the theorem.
The load-bearing proof is the direct local-defect-to-global-error argument for
that same-initial-state reported-grid theorem.

\begin{center}
\refstepcounter{table}\label{tab:intro-claim-hierarchy}
\small
\textbf{Table~\thetable: Claim hierarchy used in the manuscript.}\par\vspace{0.5em}
\begin{tabular}{P{0.23\linewidth}P{0.43\linewidth}P{0.22\linewidth}}
\toprule
Evidence layer & Role in the paper & Claim role \\
\midrule
Gauss6/FullVA residual and theorem &
Defines the 132-row Lie-group FullVA stage solve and the conditional
sixth-order same-initial-state reported-grid proof obtained from
compact-branch local-defect-to-global-error transfer. &
Main contribution. \\
Smooth cylindrical-chain and pendulum rows &
Check the accepted smooth path and provide method-side dynamic-order evidence
for the examples whose trajectory-order policy is accepted. &
Primary numerical evidence. \\
Four-link and slider-crank rows &
Exercise closed-loop lower-pair constraints, reaction consistency, and local
coarse dynamics under the same residual vocabulary. &
Mechanism coverage, not standalone asymptotic order proof. \\
\tfe{} and common-reference comparisons &
Place the method against the local \(m=3\) order-five formula target and
record bounded diagnostic comparisons. &
Formal-order and diagnostic comparison only. \\
External reproduction, source-package, and auxiliary primitive-route records &
Record why external same-test superiority, full-\tfe{} replacement,
primitive-route closure, and residual-to-error promotion are outside the
accepted theorem and numerical claim. &
Claim-scope and reproducibility boundary, not a method contribution or theorem
input. \\
\bottomrule
\end{tabular}
\end{center}

Throughout the paper, ``conditional order'' has this narrow meaning: a direct
132-row residual-defect and local-to-global consistency transfer under the
stated compact-tube, endpoint-closure, and branch-solver assumptions. It is
an analytic perturbation statement, not an empirical order fit: the smooth
trajectory slopes are reported only after the theorem-level estimates have
been stated. It is not a residual-to-error transfer theorem for mechanism-coverage
rows, not a proof that fixed production tolerances are asymptotically
sufficient, and not a full source-paper same-test \tfe{} reproduction.
The P7 residual-to-error boundary is therefore a downstream output nonclaim,
not one of the retained branch-domain hypotheses used to prove the
smooth-branch theorem.
The one-step map is defined only for transitions that remain on the selected
Gauss-predictor branch and satisfy the compact-tube, endpoint-closure, and
scaled Newton residual conditions in the theorem. Transitions outside that
domain are not assigned a theorem-level order by residual tables, fixed
tolerances, or external comparison rows.
The scaled Newton condition is a proof-norm residual condition imposed before
endpoint or trajectory errors are evaluated; it is not an assumed
local-defect or grid-error conclusion.

\noindent\textbf{Proof structure.}
The proof starts from the classical Gauss defect, passes through the
implemented 132-row residual bridge with the 96-row
\(O(h^7)\) non-dynamic certificate plus the direct Newton--Euler substitution
certificate, under which the 36 Newton--Euler rows vanish, adds the
\mbox{endpoint-closure and inexact-Newton perturbations}, and then applies
local-to-global transfer and the reporting map.\par\noindent
Section~\ref{sec:order} formalizes this handoff as a
four-term local-defect estimate on the selected branch: Gauss defect,
residual-bridge perturbation, endpoint-closure perturbation, and
inexact-Newton perturbation. The reported order statement then uses a fifth
step--the first-exit/discrete-Gronwall transfer followed by the reporting
map--rather than introducing another local-defect source. The Newton--Euler certificate discharges the 36
dynamic rows, while the 96-row non-dynamic certificate supplies the row-local
residual input under the fixed branch, norm, and implementation binding.
Residual-only mechanism rows, external comparison records, and primitive/Taylor
fragments are recorded separately because they are not terms in this
estimate. The theorem-facing tables below are reading maps for this
lemma chain: they identify hypotheses, row-local certificates, and forbidden
promotions, but the estimates themselves are supplied only by the stated
lemmas and equations. In particular, residual tables, solver logs, external
reproduction rows, and auxiliary primitive-route fragments are not theorem
inputs unless the corresponding theorem interface is explicitly closed.
The proof below separates domain hypotheses from proved algebraic residual
inputs.  The domain side fixes the compact branch, local inverse, endpoint
map, stability scale, implementation binding, and solver-scale envelope with
\(h\)-independent constants.  The residual side identifies the exact
implemented 132-row residual used by Newton, evaluates the lifted Gauss stages
in that residual, combines the 96-row non-dynamic certificate with the
36-row direct Newton--Euler substitution, and then passes through the
Kantorovich perturbation step.  Once those objects are fixed, the proof
derives the \(h^7\) local defect and the \(h^6\) same-initial-state
reported-grid estimate forward from the Gauss defect and the implemented
residual bridge.  Numerical evidence and appendix records do not substitute
for this local-defect argument.

\noindent\textbf{Reference proof usage.}
The non-\tfe{} proof-style reference is
Wieloch--Arnold~\cite{wieloch2021bdf}, not the \tfe{} comparison paper.  Its
role is limited to proof ordering for Lie-group constrained DAEs: define the
discrete Lie-group object, insert the smooth exact branch into that object,
prove the local truncation/defect estimate, and only then write the global
error recursion and reported-output consequence.  The present proof instantiates
that order with a different object, namely the row-scaled 132-row
Gauss/FullVA residual \(F_{A,h}\), the accepted Gauss-predictor branch, the
endpoint-output map, and the compact-tube solver envelope.  Thus the active
Taylor object is the full scaled residual map used by
Lemma~\ref{lem:stage-residual-defect}, not the separate 162 primitive
Newton--Euler subterm inventory.  No BLieDF constant, BDF \(p=k\) estimate,
hidden-constraint/multiplier estimate, or primitive FullVA subterm bound is
transferred into the theorem.
The reference paper's technical neighborhood and start-value assumptions are
not suppressed in this translation: their role is played here by the retained
compact-tube/branch domain, endpoint and solver-scale interfaces, and the
same-reduced-initial-state clause used in the local-to-global transfer.  These
domain clauses are fixed before diagnostics are read; finite residual tables,
smooth \(h\)-sweeps, or source-policy rows do not verify them retroactively.

\begin{table}[htbp]
\centering
\caption{How the non-\tfe{} reference proof is invoked.}
\label{tab:reference-proof-order-map}
\scriptsize
\setlength{\tabcolsep}{2pt}
\renewcommand{\arraystretch}{0.95}
\begin{tabular}{P{0.24\linewidth}P{0.39\linewidth}P{0.25\linewidth}}
\toprule
Reference-proof step & FullVA instantiation & Explicit non-transfer \\
\midrule
Fix the discrete Lie-group object before any estimate. &
Fix the row-scaled \(132\)-row residual \(F_{A,h}\), the accepted
Gauss-predictor branch, row order, proof norm, and endpoint-output map. &
No BLieDF equation or source-paper weak row is imported. \\
Insert the smooth exact branch into the fixed object. &
Insert the smooth Gauss lift \(Z_G\) into the implemented FullVA residual and
evaluate the certified \(96+36\) row split on that same branch. &
No hidden-constraint, multiplier-recursion, or primitive lift estimate is
borrowed. \\
Prove the local defect before global propagation. &
Combine the Gauss defect, the direct \(132\)-row residual bridge, endpoint
closure, and inexact-Newton perturbation into the four-term local-defect
bound. &
No finite \(h\)-sweep, residual table, or separate primitive/Taylor route closes a
separate theorem interface. \\
Only then derive the global reported error. &
Apply the same-branch local-to-global transfer and reporting-map estimate to
the accepted \method{} map. &
No residual-to-error, external-superiority, source-policy, or full-\tfe{}
claim follows without its separate theorem or source-policy criterion. \\
\bottomrule
\end{tabular}
\end{table}

\noindent
Table~\ref{tab:reference-proof-order-map} is the introductory
reader-facing map only.  The proof-level correspondence is
Table~\ref{tab:reference-proof-order-correspondence}, where the same ordering
discipline is tied to the FullVA lemmas and the explicit non-import rules are
restated at theorem-interface level.

\noindent\textbf{How to read the proof layers.}
The proof has one theorem route.  The direct Gauss/FullVA residual bridge
assembles the 96 non-dynamic rows and the 36 direct Newton--Euler rows into the
\(132\)-row stage-residual input used by
Theorem~\ref{thm:g6fullva-order}.  The primitive/Taylor material asks a
stronger term-by-term question for the same Newton--Euler rows using
primitive state, acceleration, and multiplier lifts; it is an alternative
sufficient route, not a premise of the theorem.  Residual and
external-comparison rows are treated separately because they would require
their own trajectory-error or source-policy promotion arguments.  Thus the
proof strength rests on the direct theorem route, while the auxiliary material
delimits what is deliberately not being claimed.

\noindent
The route is a one-way conditional implication, not a backfilled
interpretation of numerical evidence.  The compact-tube hypotheses choose the
branch and constants before any reported run is inspected.  On that fixed
object, the stage-row defect on \(Z_G\) supplies the residual input, endpoint
closure and inexact Newton add controlled perturbations, and the local-to-global
transfer then gives the same-initial-state reported-grid estimate.  Finite step
sweeps may identify the reported branch, but they do not discharge the
theorem-level solver hypothesis and do not promote residual-only mechanism or
external comparison rows.  The retained interfaces delimit where the theorem
may be invoked.

The remainder of the paper follows the evidence chain in that order.
Section~\ref{sec:method} defines the accepted one-step map and the complete
execution algorithm. Section~\ref{sec:validation} fixes the validation
policy, including the rule that external comparisons remain diagnostic unless
the same-test promotion criteria are satisfied.
Section~\ref{sec:cross-paper-boundary}
records the detailed same-test promotion boundary.
Section~\ref{sec:order} gives the conditional sixth-order theorem and separates
stage-residual inputs from theorem-domain interfaces. Sections~\ref{sec:numerical}
and~\ref{sec:four-examples} report the numerical evidence. The final
diagnostic and reproducibility sections document what has been checked without
turning those checks into broader claims.

\section{Related work}

The method sits at the intersection of geometric integration, constrained
multibody dynamics, collocation methods for differential-algebraic equations
(DAEs), variational integration, nonsmooth contact dynamics, and time finite
elements. The relevant literature separates into six clusters, each of which
leaves a different gap.

First, Lie-group integrators preserve the rotation manifold by updating in a
local tangent chart rather than advancing redundant orientation coordinates and
then repairing them. Runge--Kutta--Munthe-Kaas methods provide a general
high-order construction on manifolds \cite{munthe1998lie,munthe1999high}.
Multibody-specific work has adapted this geometric viewpoint to constrained
mechanics through Lie-group generalized-alpha schemes \cite{bruls2012lie},
BDF-type schemes on Lie groups \cite{wieloch2021bdf}, and quaternion-based
singularity-free rotational integration
\cite{terze2016singularity,terze2020aircraft}. These methods establish how to
evolve rotations geometrically, but they do not by themselves prescribe a
stage residual that simultaneously enforces lower-pair position, velocity, and
acceleration consistency in one Newton solve.

Second, constrained multibody systems in absolute or redundant coordinates are
index-3 DAEs. Their numerical order is affected not only by the nominal order
of the time integrator, but also by how the algebraic constraints and their
differentiated forms are enforced \cite{hairer1993solving}. The
absolute-coordinate Lie-group formulations of Kissel, Taves, Fang, Zhang, and
Negrut clarify the connection between small body-frame rotation increments and
the first-order sensitivities needed in Newton iterations
\cite{taves2020exponential,kissel2021dwelling,kissel2022absolute,fang2023halfimplicit,kissel2024velocitypartitioning}.
That line of work is important for constructing matrix rows for constraints
and reaction forces, and it also supplies a public-code benchmark suite
covering the single pendulum, double pendulum, slider-crank, four-link, and
scaling tests. The present paper uses the same body-frame perturbation
viewpoint to build a stage system in which the constraint row, velocity row,
and acceleration row are all part of the accepted residual; it cannot be read
as a fair external comparison until those published benchmark definitions are
matched exactly.

Third, high-order collocation and index-reduction methods are a standard route
to high-order DAE integration when the constrained flow is regular and the
algebraic variables are solved consistently
\cite{gear1985automatic,hairer1993solving}. Symplectic partitioned
Runge--Kutta methods also show how holonomic constraints can be preserved
within a structure-preserving Runge--Kutta framework \cite{jay1996symplectic}.
This places the present construction near constrained collocation, but with a
different emphasis on the lower-pair residual rows that are exposed to Newton's
method.
The three-stage Gauss path therefore provides a natural sixth-order reference
construction on the smooth reduced problem. What is less standard in a
lower-pair multibody setting is the concrete residual organization: the
accepted path here couples Gauss stage kinematics, Newton--Euler balance, Lie
retraction variables, and FullVA lower-pair rows into a square nonlinear
system. This is the level at which the manuscript's contribution is intended
to sit; it is not a claim that Gauss collocation itself is new.

Fourth, variational and symplectic integrators derive the discrete equations
from a discrete action principle and are valued for their momentum,
symplectic, and long-time energy behavior
\cite{marsden2001discrete,leyendecker2008variational,leok2012prolongation}.
Constrained variational integrators are therefore a close conceptual
neighbor: they treat holonomic constraints at the discrete level and can
reduce the constrained solve through null-space coordinates. Their objective,
however, is not the same as the FullVA residual used here. The present method
does not construct a discrete variational principle; it builds an
absolute-coordinate stage residual that enforces lower-pair position,
velocity, and acceleration consistency simultaneously and then evaluates the
resulting one-step map under the validation policy below.

Fifth, nonsmooth contact and friction solvers form a separate branch of
multibody dynamics. Time-stepping and complementarity formulations handle
unilateral contact, impact, and Coulomb friction without resolving every
event explicitly
\cite{stewart1996implicit,anitescu1997formulating,tasora2011matrixfree}.
Those methods are essential for large contact-rich systems, but their central
problem is a contact complementarity or cone-complementarity solve. The
validation problems studied here are bilateral lower-pair mechanisms with smooth
driving and reaction reconstruction; the paper therefore does not claim a
new nonsmooth contact integrator or a replacement for complementarity-based
frictional dynamics.

Sixth, time finite elements provide a different route to high-order direct
integration. Early space-time and time finite-element methods were developed
for structural dynamics and constrained rigid-body dynamics
\cite{borri1988rigid,hughes1988spacetime,hulbert1992time}. Later direct
integration families based on time finite elements and weighted residuals
showed high-order behavior and controllable numerical damping
\cite{kim2017higher,jog2016robust,bauchau2024time}. The frictional
Lie-group time finite-element work of Chaturvedi, Sandu, and Sandu extends
this direction to index-3 multibody DAEs with nonlinear friction and provides
the local paper-style temporal finite-element comparator used here
\cite{chaturvedi2026tfe}. Its Gauss--Lobatto family uses the order convention
$2m-1$, so its $m=3$ member is an order-five formula target, although DAE
order reduction can still appear in the reported pendulum experiments.

The present manuscript is intentionally different from a full source-paper
reimplementation. It reports method-side evidence for the accepted Gauss
collocation FullVA residual on lower-pair mechanisms and compares its
conditional sixth-order path against the local order-five temporal
finite-element formula target. Complete
substitution of all source-paper temporal finite-element weak rows remains
outside the accepted claim. The gap addressed here is narrower but concrete:
an implemented Lie-group lower-pair stage residual that carries position,
velocity, and acceleration consistency through the same nonlinear solve and is
checked on the cylindrical-chain benchmark plus four ASME-style mechanisms,
with dynamic-order, mechanism-coverage, and diagnostic evidence kept separate.

\section{Mathematical setting and definitions}

Consider a constrained multibody system with generalized position $q$,
velocity $v$, acceleration $a$, and Lagrange multipliers $\lambda$. The
absolute-coordinate index-3 form may be written abstractly as
\begin{align}
  M(q)\dot v + G(q,v) + \Phi_q(q)^T\lambda &= F(q,v,t), \label{eq:dae-dyn}\\
  \dot q &= \mathcal{K}(q)v, \label{eq:dae-kin}\\
  \Phi(q) &= 0. \label{eq:dae-constraint}
\end{align}
Equations~\eqref{eq:dae-dyn}, \eqref{eq:dae-kin}, and
\eqref{eq:dae-constraint} fix the
abstract constrained dynamics, kinematics, and holonomic constraint notation
used throughout the paper.
Here $\mathcal{K}$ includes the local Lie-group kinematic map for finite
rotations. For orientations, the step update is represented by a retraction
\begin{equation}
  Q_{n+1}=Q_n\,\mathrm{Exp}(\eta_{n+1}),
  \label{eq:retraction}
\end{equation}
with $\eta_{n+1}$ in the local rotation algebra. The implementation uses this
local representation to keep finite rotations on the group while solving a
Newton system in vector coordinates.

\begin{definition}[FullVA lower-pair residual]
In this manuscript, FullVA denotes full position-, velocity-, and
acceleration-level lower-pair enforcement. The FullVA residual is the square
nonlinear stage system that enforces lower-pair position constraints, their
velocity consistency equations, and their acceleration consistency equations in
the same Newton solve, together with the collocation dynamics and the
Lie-group update equations.
\end{definition}

\begin{definition}[Accepted method path]
The accepted method path is the \method{} one-step map generated by the
132-row FullVA residual in the implemented residual path. This is the only method path
used for the accepted dynamic-order rows and four-example
mechanism-coverage validation claims in this manuscript.
\end{definition}

\begin{definition}[Full-\tfe{} stage replacement]
In this study, full-\tfe{} stage replacement means replacing all 132
Gauss/FullVA stage rows inside Newton by paper-derived temporal
finite-element weak rows, without endpoint source data, terminal-row
replacement, output projection, or target-direction oracle information. This
is a stronger source-paper reproduction criterion and is not accepted by the
evidence reported here.
\end{definition}

\section{Nomenclature}

\begin{table}[htbp]
\centering
\caption{Main symbols used in the manuscript.}
\label{tab:nomenclature}
\begin{tabular}{P{0.18\linewidth}P{0.66\linewidth}}
\toprule
Symbol & Meaning \\
\midrule
$q$, $v$, $a$ & Generalized position, velocity, and acceleration variables. \\
$Q$ & Rotation variable on the Lie group. \\
$\eta$ & Local tangent-space rotation increment used in the retraction. \\
$\lambda$ & Lagrange multipliers for the lower-pair constraints. \\
$M(q)$ & Configuration-dependent mass and inertia operator. \\
$F(q,v,t)$ & Applied and generalized force vector. \\
$\Phi(q)$, $\Phi_q(q)$ & Constraint function and constraint Jacobian. \\
$\mathcal{K}(q)$ & Kinematic map from generalized velocity to configuration rate. \\
$h$ & Time-step size. \\
$Z$ & Stacked stage unknown vector in the nonlinear one-step solve. \\
$R_{\mathrm{G6FVA}}$ & Accepted \method{} residual solved inside Newton. \\
$m$ & Polynomial index for the paper-style Gauss--Lobatto \tfe{} comparator. \\
$p_{\mathrm{TFE}}$ & Expected order of the local paper-style temporal finite-element target. \\
\bottomrule
\end{tabular}
\vspace{-3pt}
\end{table}

\section{Accepted one-step integrator}
\label{sec:method}

Let $x_n=(q_n,v_n,a_n)$ denote the endpoint position, velocity, and
acceleration variables. The accepted method defines a one-step map
\begin{equation}
  x_{n+1}=\Psi_h^{\mathrm{G6FVA}}(x_n),
  \label{eq:onestep}
\end{equation}
Equation~\eqref{eq:onestep} is the method-level map induced by one accepted
FullVA stage solve, where the stage variables $Z$ solve the nonlinear residual
\begin{equation}
  R_{\mathrm{G6FVA}}(q_n,v_n,a_n,Z;h)=0 .
  \label{eq:residual}
\end{equation}
For the cylindrical-chain benchmark used for the main convergence evidence,
the residual has two bodies, two lower-pair joints, and three Gauss stages.
Each body stage stores translational position, local rotation increment,
translational velocity, angular velocity, translational acceleration, and
angular acceleration. This gives $18$ body unknowns per body per stage. Each
lower-pair joint contributes four multiplier unknowns per stage. Thus the
Newton stage vector has
\begin{equation}
  3\{2(18)+2(4)\}=132
  \label{eq:stage-count}
\end{equation}
scalar unknowns by Eq.~\eqref{eq:stage-count}, and the accepted residual is
square. The closed-loop
four-link and slider-crank examples use their own fully constrained
kinematic FullVA systems, but they preserve the same position, velocity, and
acceleration-level lower-pair residual vocabulary.

The residual is assembled from six row families. For stage $i=1,2,3$ with
Gauss abscissa $c_i$, weights $b_i$, and Runge--Kutta matrix $A_{ij}$, the
stage variables are denoted in Eq.~\eqref{eq:stage-variable-block}:
\begin{equation}
\label{eq:stage-variable-block}
  Z_i=(r_i,\eta_i,v_i,\omega_i,a_i,\alpha_i,\lambda_i),
\end{equation}
where $r_i$ is the translational position, $\eta_i$ is the local Lie-algebra
rotation increment, $v_i,\omega_i$ are translational and angular velocities,
$a_i,\alpha_i$ are translational and angular accelerations, and $\lambda_i$
collects lower-pair multipliers. The implementation evaluates rotations by
$Q_i=Q_n\mathrm{Exp}(\eta_i)$ and uses the right-trivialized differential of
the exponential map in the rotational kinematic rows.

The complete one-step equation can be read as the following square stage
system plus endpoint reconstruction.  With $t_i=t_n+c_i h$, define
$F^r_{b,i}=f_{b,i}+\Phi_{r,b}^T\lambda_i$,
$T^\theta_{b,i}=\tau_{b,i}+\Phi_{\theta,b}^T\lambda_i$, and
$A_{\Phi,i}=\Phi_q(q_i,t_i)a_i+\dot{\Phi}_q(q_i,v_i,t_i)v_i$.  Each stage
solves
\begin{equation}
\label{eq:g6fullva-expanded-stage}
\begin{aligned}
0 &= r_i-r_n-h\sum_{j=1}^3 A_{ij}v_j, &
0 &= \eta_i-h\sum_{j=1}^3 A_{ij}J_r^{-1}(\eta_j)\omega_j,\\
0 &= v_i-v_n-h\sum_{j=1}^3 A_{ij}a_j, &
0 &= \omega_i-\omega_n-h\sum_{j=1}^3 A_{ij}\alpha_j,\\
0 &= m_b a_{b,i}-F^r_{b,i},\\
0 &= J_b\alpha_{b,i}+\omega_{b,i}\times J_b\omega_{b,i}-T^\theta_{b,i},\\
0 &= \Phi(q_i,t_i)-g_i^{(0)},\qquad
0 = \Phi_q(q_i,t_i)v_i-g_i^{(1)},\\
0 &= A_{\Phi,i}-g_i^{(2)} .
\end{aligned}
\end{equation}
Here $g_i^{(0)},g_i^{(1)},g_i^{(2)}$ are the position-, velocity-, and
acceleration-level driven right-hand sides; they are zero for passive
lower-pair constraints.  Equation~\eqref{eq:g6fullva-expanded-stage} is the
method equation solved by Newton in Algorithm~1.  The next paragraphs unpack
the same equation by row family and define the endpoint map.

The first two row families are the collocation kinematic defects. In compact
notation,
\begin{align}
  r_i-r_n-h\sum_{j=1}^3 A_{ij}v_j &=0, \label{eq:trans-collocation}\\
  \eta_i-h\sum_{j=1}^3 A_{ij}J_r^{-1}(\eta_j)\omega_j &=0,
  \label{eq:rot-collocation}
\end{align}
where $J_r^{-1}$ is the inverse right Jacobian of $\mathrm{SO}(3)$. The next
two families enforce velocity collocation,
\begin{align}
  v_i-v_n-h\sum_{j=1}^3 A_{ij}a_j &=0, \label{eq:v-collocation}\\
  \omega_i-\omega_n-h\sum_{j=1}^3 A_{ij}\alpha_j &=0.
  \label{eq:w-collocation}
\end{align}
The fifth family is the Newton--Euler balance row,
\begin{align}
  m_b a_{b,i}-f_{b,i}-\Phi_{r,b}(q_i)^T\lambda_i &=0,
  \label{eq:trans-balance}\\
  J_b\alpha_{b,i}+\omega_{b,i}\times J_b\omega_{b,i}
  -\tau_{b,i}-\Phi_{\theta,b}(q_i)^T\lambda_i &=0,
  \label{eq:rot-balance}
\end{align}
for each body $b$. The sixth family is the FullVA lower-pair block. It
contains position, velocity, and acceleration-level constraint rows,
\begin{equation}
  \Phi(q)=0,\qquad
  \Phi_q(q)v=0,\qquad
  \Phi_q(q)a+\dot{\Phi}_q(q,v)v=0,
  \label{eq:fullva}
\end{equation}
with nonzero right-hand sides included for driven mechanisms. The lower-pair
library covers coordinate-difference (CD), dot-product-one (DP1),
dot-product-two (DP2), and distance (D) constraints. The method is therefore
a FullVA constrained stage solve, not a post-step position-only projection.

After Newton convergence, the endpoint is reconstructed by the Gauss weights,
\begin{align}
  r_{n+1} &= r_n+h\sum_{i=1}^3 b_i v_i, &
  v_{n+1} &= v_n+h\sum_{i=1}^3 b_i a_i, \label{eq:end-trans}\\
  Q_{n+1} &= Q_n\mathrm{Exp}\!\left(h\sum_{i=1}^3 b_i
      J_r^{-1}(\eta_i)\omega_i\right), &
  \omega_{n+1} &= \omega_n+h\sum_{i=1}^3 b_i \alpha_i .
  \label{eq:end-rot}
\end{align}
The Newton linearization solves
\begin{equation}
  \frac{\partial R_{\mathrm{G6FVA}}}{\partial Z}\Delta Z
  =-R_{\mathrm{G6FVA}},
  \label{eq:newton}
\end{equation}
with automatic differentiation used to assemble the dense reference Jacobian.
The sparse backend uses the same residual and Jacobian entries; it is tracked
as an implementation caveat rather than as a primary accuracy claim.

\subsection{Stage Jacobian structure}

This subsection records the block structure behind
Eq.~\eqref{eq:newton}. It is not a separate hand-coded Jacobian; the
implementation differentiates the same residual by automatic differentiation.
The formulas below specify the row-to-unknown dependencies used to verify the
AD Jacobian and to define the square Newton system.

For stage $i$, write the differential of the stage pose as
Eq.~\eqref{eq:pose-variation}:
\begin{equation}
\label{eq:pose-variation}
  \delta q_i=(\delta r_i,\delta\theta_i),\qquad
  \delta\theta_i=J_r(\eta_i)\delta\eta_i ,
\end{equation}
so that $Q_i=Q_n\mathrm{Exp}(\eta_i)$ is perturbed by the corresponding
right-trivialized rotation vector. Let $G_i=J_r^{-1}(\eta_i)$ and let
$D_\eta(G_i\omega_i)[\delta\eta_i]$ denote the directional derivative of
$J_r^{-1}(\eta_i)\omega_i$. The four collocation blocks have the exact
stage-coupling derivatives in Eq.~\eqref{eq:collocation-jacobian}:
\begin{align}
  \delta R^r_i
    &= \delta r_i-h\sum_{j=1}^3 A_{ij}\delta v_j,\\
  \delta R^\eta_i
    &= \delta\eta_i-h\sum_{j=1}^3 A_{ij}
       \left(D_\eta(G_j\omega_j)[\delta\eta_j]
       +G_j\delta\omega_j\right),\\
  \delta R^v_i
    &= \delta v_i-h\sum_{j=1}^3 A_{ij}\delta a_j,\\
  \delta R^\omega_i
    &= \delta\omega_i-h\sum_{j=1}^3 A_{ij}\delta\alpha_j.
  \label{eq:collocation-jacobian}
\end{align}
Thus the only off-diagonal stage couplings in the kinematic rows are the
Gauss matrix entries $A_{ij}$; all other nonlinear stage coupling enters
through dynamics and constraints evaluated at each stage.

For body $b$ at stage $i$, the translational and rotational balance
linearizations can be written
\begin{align}
  \delta R^{T}_{b,i}
    &= M_b\delta a_{b,i}
       -F_{q,b,i}\delta q_i-F_{\dot q,b,i}\delta\dot q_i
       -\delta(\Phi_{r,b,i}^{T})\lambda_i
       -\Phi_{r,b,i}^{T}\delta\lambda_i,\\
  \delta R^{R}_{b,i}
    &= J_b\delta\alpha_{b,i}
       +\mathcal G_{b,i}\delta\omega_{b,i}
       -T_{q,b,i}\delta q_i-T_{\dot q,b,i}\delta\dot q_i
       -\delta(\Phi_{\theta,b,i}^{T})\lambda_i
       -\Phi_{\theta,b,i}^{T}\delta\lambda_i .
  \label{eq:balance-jacobian}
\end{align}
Here Eq.~\eqref{eq:gyro-linearization-operator} defines the gyroscopic
linearization operator:
\begin{equation}
\label{eq:gyro-linearization-operator}
  \mathcal G_{b,i}\delta\omega
  =[\delta\omega]_\times J_b\omega_{b,i}
   +[\omega_{b,i}]_\times J_b\delta\omega ,
\end{equation}
and $F_{q},F_{\dot q},T_q,T_{\dot q}$ collect load derivatives, including the
velocity-dependent Brown--McPhee friction terms in the cylindrical-chain
benchmark~\cite{browmcphee2016friction}. The terms
$\delta(\Phi_{r,b,i}^{T})\lambda_i$ and
$\delta(\Phi_{\theta,b,i}^{T})\lambda_i$ are the geometric stiffness blocks
from the lower-pair Jacobian.

For a scalar lower-pair row $\Phi(q,t)=0$, the FullVA row linearization at a
stage is
\begin{align}
  \delta R^\Phi_i
    &= \Phi_q\,\delta q_i,\\
  \delta R^{\dot\Phi}_i
    &= D_q(\Phi_q\dot q-\nu)[\delta q_i]
       +\Phi_q\,\delta\dot q_i,\\
  \delta R^{\ddot\Phi}_i
    &= D_q(\Phi_q\ddot q+\dot\Phi_q\dot q-\gamma)[\delta q_i]
       +D_{\dot q}(\Phi_q\ddot q+\dot\Phi_q\dot q-\gamma)
        [\delta\dot q_i]
       +\Phi_q\,\delta\ddot q_i .
  \label{eq:fullva-jacobian}
\end{align}
The CD/DP1/DP2/D formulas in Section~\ref{sec:lower-pair-formulas} supply the
explicit $\Phi$, $\dot\Phi$, and $\ddot\Phi$ expressions used in these
directional derivatives.

With row groups ordered as Eq.~\eqref{eq:stage-row-group-order}
\begin{equation}
\label{eq:stage-row-group-order}
  (R^r,R^\eta,R^v,R^\omega,R^T,R^R,R^\Phi,R^{\dot\Phi},R^{\ddot\Phi})
\end{equation}
and unknown groups ordered as Eq.~\eqref{eq:stage-unknown-group-order}
\begin{equation}
\label{eq:stage-unknown-group-order}
  (\delta r,\delta\eta,\delta v,\delta\omega,\delta a,\delta\alpha,
  \delta\lambda),
\end{equation}
define $K_{ik}^{\rho,\chi}=\partial R^\rho_i/\partial \chi_k$. The nonzero
block pattern is summarized in Table~\ref{tab:stage-jacobian-block}. The
symbols $M$ and $J$ denote block-diagonal mass and inertia matrices, and the
$\Phi_\eta,\Phi_\omega,\Phi_\alpha$ columns are obtained from
Eq.~\eqref{eq:pose-variation}. This is the paper-level block pattern
corresponding to the 132-row dense AD Jacobian used by the accepted
cylindrical-chain solve.

\begin{table}[htbp]
\centering
\scriptsize
\setlength{\tabcolsep}{3pt}
\caption{Nonzero block pattern of the accepted stage Newton matrix.}
\label{tab:stage-jacobian-block}
\begin{tabular}{P{0.18\linewidth}P{0.27\linewidth}P{0.39\linewidth}}
\toprule
Row group & Direct stage blocks & Cross-stage or nonlinear blocks \\
\midrule
$R^r$ &
$K^{r,r}_{ii}=I$ &
$K^{r,v}_{ij}=-hA_{ij}I$. \\
$R^\eta$ &
$K^{\eta,\eta}_{ii}=I-hA_{ii}D_\eta(G_i\omega_i)$ &
$K^{\eta,\eta}_{ij}=-hA_{ij}D_\eta(G_j\omega_j)$ and
$K^{\eta,\omega}_{ij}=-hA_{ij}G_j$. \\
$R^v$ &
$K^{v,v}_{ii}=I$ &
$K^{v,a}_{ij}=-hA_{ij}I$. \\
$R^\omega$ &
$K^{\omega,\omega}_{ii}=I$ &
$K^{\omega,\alpha}_{ij}=-hA_{ij}I$. \\
$R^T$ &
$K^{T,a}_{ii}=M$, $K^{T,\lambda}_{ii}=-\Phi_r^T$ &
Load derivatives and geometric-stiffness blocks in
Eq.~\eqref{eq:balance-jacobian}. \\
$R^R$ &
$K^{R,\alpha}_{ii}=J$, $K^{R,\lambda}_{ii}=-\Phi_\theta^T$ &
Gyroscopic, load-derivative, and geometric-stiffness blocks in
Eq.~\eqref{eq:balance-jacobian}. \\
$R^\Phi$ &
$K^{\Phi,q}_{ii}=\Phi_q$ &
Rotation columns mapped by Eq.~\eqref{eq:pose-variation}. \\
$R^{\dot\Phi}$ &
$K^{\dot\Phi,\dot q}_{ii}=\Phi_q$ &
$D_q(\Phi_q\dot q-\nu)$ blocks. \\
$R^{\ddot\Phi}$ &
$K^{\ddot\Phi,\ddot q}_{ii}=\Phi_q$ &
$D_q$ and $D_{\dot q}$ blocks from Eq.~\eqref{eq:fullva-jacobian}. \\
\bottomrule
\end{tabular}
\end{table}

\begin{table}[htbp]
\centering
\caption{Operational definition of the accepted method path.}
\label{tab:method-path}
\begin{tabular}{P{0.24\linewidth}P{0.62\linewidth}}
\toprule
Item & Accepted definition \\
\midrule
Production method & \method{} one-step map. \\
Nominal order & Conditional sixth-order Gauss collocation path under the
retained interfaces, including solver-scale control, of
Theorem~\ref{thm:g6fullva-order}. \\
Stage residual & Square 132-row FullVA lower-pair nonlinear system. \\
Endpoint closure & Controlled by P2 local right-inverse and raw-defect
hypotheses; KKT/projection diagnostics are recorded separately. \\
Full-\tfe{} replacement & Not accepted; this stronger source-paper residual
replacement is outside the present claim. \\
\bottomrule
\end{tabular}
\end{table}

The endpoint closure certificate is part of the accepted \method{} evidence path. It
does not convert the method into a source-paper \tfe{} residual reproduction.
A full-\tfe{} replacement would instead replace all stage rows in
Eq.~\eqref{eq:residual} by paper-derived weak rows and would pass the same
rank, residual, endpoint, and convergence checks.

\begin{table}[htbp]
\centering
\caption{Algorithmic form of one \method{} step.}
\label{tab:algorithm}
\begin{tabular}{P{0.16\linewidth}P{0.70\linewidth}}
\toprule
Step & Operation \\
\midrule
1 & Build the lower-pair constraint graph, mass/inertia blocks, loads, and
FullVA stage unknown vector for $(q_n,v_n,a_n)$ and step size $h$. \\
2 & Assemble the 132-row residual
$R_{\mathrm{G6FVA}}(q_n,v_n,a_n,Z;h)$, including dynamics, Lie-group
kinematics, and position/velocity/acceleration constraint rows. \\
3 & Solve the nonlinear stage system by Newton iteration with automatic
differentiation for the residual Jacobian. \\
4 & Recover the endpoint by Gauss quadrature and the Lie-group retraction in
Eq.~\eqref{eq:retraction}. \\
5 & Check endpoint position, velocity, residual, rank, order, and
four-example quantities using the reproducibility checks summarized in
the reproducibility appendix. \\
\bottomrule
\end{tabular}
\end{table}

\subsection{Complete execution algorithm}
\label{sec:complete-algorithm}

The implemented \method{} path used for the cylindrical-chain and
four-example validation rows is the trajectory algorithm below.
The listing is deliberately written as an execution protocol: it starts from
mechanism data and ends with the error/order records used in the reproduced
results.

Operationally, for an example $e$ and a step size $h$, the accepted run is
the composition in Eq.~\eqref{eq:complete-execution-map}:
\begin{align}
  Z_{n,h,e}
    &=\operatorname{NewtonSolve}_{\tau_N,k_{\max}}
      \left(R_{\mathrm{G6FVA}}(x_{n,h,e},Z;h)=0;\,
      Z^{(0)}_{n,h,e}\right), \notag\\
  \widetilde{x}_{n+1,h,e}
    &=\mathcal E_h(x_{n,h,e},Z_{n,h,e}), \notag\\
  x_{n+1,h,e}
    &=\mathcal C_h(\widetilde{x}_{n+1,h,e}), \notag\\
  \mathcal R_e(h)
    &=\mathcal P_e\!\left(\{x_{n,h,e}\}_{n=0}^{N_h},
      \{D_{n,h,e}\}_{n=0}^{N_h-1},x_e^{\rm ref}\right).
  \label{eq:complete-execution-map}
\end{align}
Here $\mathcal E_h$ is the endpoint reconstruction in
Eqs.~\eqref{eq:end-trans}--\eqref{eq:end-rot}, $\mathcal C_h$ is the
accepted endpoint closure map, $D_{n,h,e}$ stores Newton and constraint
diagnostics, $x_e^{\rm ref}$ is the analytic or nested local reference
specified by the validation policy, and $\mathcal P_e$ is the postprocessing
map that writes the reproduced error, order, diagnostic, and comparison-scope
records.

\noindent\textbf{Algorithm 1 (complete accepted \method{} run).}
\begin{enumerate}
  \item Choose the validation example $e$, its mechanism data, its step-size
  set $H_e$, final time $T_e$, Newton tolerance $\tau_N$, Newton cap
  $k_{\max}$, Jacobian backend, and reference policy.
  \item Build the body mass and inertia blocks, lower-pair graph, joint
  frames, axes, CD/DP1/DP2/D constraint bases, driven right-hand sides, loads,
  and friction data.
  \item Construct the initial state $x_{0,h,e}$ on
  $\mathbb R^3\times\mathrm{SO}(3)$ for every body and project the initial
  velocity onto the velocity-level lower-pair constraints.
  \item For each $h\in H_e$, first check the accepted-grid condition
  $|T_e/h-\mathrm{round}(T_e/h)|\le\varepsilon_{\rm grid}$. Accepted order
  rows use only grids satisfying this condition; non-integer final-time grids
  are recorded as external-comparison diagnostics rather than promoted to reported
  order rows. This is a validation/reporting policy for the finite order
  rows, not an additional theorem condition; the proof theorem below states
  its grid-point estimate on the reported grid and does not require exact
  final-time divisibility. Set $N_h=\mathrm{round}(T_e/h)$, precompute the three-stage
  Gauss coefficients and the fixed FullVA stage-vector layout, and initialize
  empty trajectory and diagnostic records.
  \item For $n=0,\ldots,N_h-1$, form the predictor $Z^{(0)}_{n,h,e}$ from the
  endpoint state, constant-acceleration motion, zero initial Lie increments,
  and force-balance multiplier estimates.
  \item Assemble $R_{\mathrm{G6FVA}}(x_{n,h,e},Z;h)$ by evaluating stage
  poses, right-Jacobian terms, lower-pair positions, velocities,
  accelerations, Newton--Euler balance rows, and FullVA position-, velocity-,
  and acceleration-level constraint rows.
  \item Iterate Newton: evaluate the residual, form or apply
  $\partial R_{\mathrm{G6FVA}}/\partial Z$, solve
  Eq.~\eqref{eq:newton}, update $Z$, and stop only after the residual or
  update norm satisfies $\tau_N$; reject the step if $k_{\max}$ is reached.
  \item Reconstruct $\widetilde{x}_{n+1,h,e}$ with the Gauss endpoint formulas
  and the Lie-group retraction, then evaluate raw endpoint position and
  velocity constraint residuals.
  \item Apply the accepted endpoint closure map $\mathcal C_h$ when the
  velocity-level closure threshold is exceeded, and record raw residuals,
  corrected residuals, projection rank, condition number, and correction norm.
  \item Store $x_{n+1,h,e}$ and the diagnostics $D_{n,h,e}$, including Newton
  iterations, residual norms, linear-solve residuals, Jacobian or
  colored-derivative timings, endpoint constraint norms, stage acceleration
  constraint norms, and quaternion unit error.
  \item After all steps and all $h\in H_e$ are complete, compare each
  trajectory with $x_e^{\rm ref}$, compute the paper's error norms, error
  ratios, observed orders, mechanism-coverage rows, and reaction residual
  summaries.
  \item Write the CSV, JSON, plot, metadata, and comparison-scope records used in
  the manuscript. These postprocessing checks reproduce the evidence boundary;
  they do not alter the accepted numerical trajectory.
\end{enumerate}

All rows in Algorithm~1 are the accepted Gauss/FullVA rows in
Eq.~\eqref{eq:residual}. A full source-paper \tfe{} replacement would replace
the stage residual assembled in the residual-assembly step by paper-derived
weak rows; that stronger algorithm is explicitly outside the accepted claim in
this manuscript.
Algorithm--theorem tolerance boundary: Algorithm~1 records the executable
finite-run protocol, including the reported Newton tolerance \(\tau_N\).
Theorem~\ref{thm:g6fullva-order} uses the same implemented residual map only
after the accepted branch also satisfies the scaled proof policy
\(\eta_h^{\rm tube}\le c_\eta h^7\) in the fixed proof norm. A fixed
\(\tau_N\) run is therefore theorem-admissible only for those steps on which
its accepted branch-selected residuals meet that \(h^7\) bound; otherwise it
is finite-run diagnostic evidence and does not discharge P6.
The correspondence between Algorithm~1 and the supplied source package is
summarized in Appendix Table~\ref{tab:algorithm-source-map}; that table is a
reproducibility map, not an additional numerical claim.

\paperfig{Figure_11_method_stage_architecture.png}
{Accepted \method{} one-step architecture. The schematic separates the
132-row FullVA residual and conditional order boundary from non-claims: it is
not a full-\tfe{} replacement, does not require a default $10^{-4}$ step-size
sweep, and does not assert external superiority without same-test evidence.}
{fig:method-stage-architecture}

The reproducibility record is described in the reproducibility appendix. The
main text uses that record only to keep
the claim boundary precise; the scientific argument rests on the residual
definition, the conditional sixth-order theorem, and the numerical evidence below.

\noindent\textbf{Formal-order comparator guardrail.}

This short guardrail fixes only the scalar formal-order comparator used
later in the paper.  The comparison to Lie-group \tfe{} methods is
formal-order-level unless same-test implementation evidence is provided.  The
source-paper comparison used in this manuscript is therefore local and
explicit.  It is placed before validation only to define the comparator
convention; it is not benchmark evidence, not a theorem input, and
not a claim that the \tfe{} implementation in the source paper has been
reproduced.  The main evidence chain remains method definition, conditional
FullVA proof, accepted dynamic-order rows, and mechanism-coverage rows.
The paper-style temporal finite-element family indexed by polynomial degree $m$
uses the expected order convention
\begin{equation}
  p_{\mathrm{TFE}}(m)=2m-1 .
  \label{eq:tfe-order}
\end{equation}
The $m=3$ Gauss--Lobatto temporal finite-element target therefore has expected
order five. The source paper should not be read as reporting uniform fifth-order
DAE convergence in every experiment: in its pendulum study it observes that the
$m=3$ TFE implementation can lose order on the DAE problem, even though the
underlying Kim--Reddy TFE construction has the expected $2m-1$ formula order
\cite{kim2017higher,chaturvedi2026tfe}. This manuscript therefore uses the
paper-style formula target $p_{\mathrm{TFE}}(3)=5$ as the comparator boundary,
not a weakened observed-order target. The accepted \method{} path has a
conditional method-order claim of six. The manuscript-level order comparison
is the formal-order comparison of that conditional \method{} path against the
local order-five target. It is not a claim that the present implementation is
faster or more accurate than the source-paper \tfe{} implementation on the
original source-paper tests.

\begin{table}[htbp]
\centering
\caption{Comparison boundary used in the manuscript.}
\label{tab:claim-boundary}
\begin{tabular}{P{0.29\linewidth}P{0.55\linewidth}}
\toprule
Question & Boundary \\
\midrule
What is claimed? &
Under the stated compact-branch, endpoint/stability, implementation-binding,
residual-defect, and solver-scale contracts, \method{} has a conditional
sixth-order Gauss-collocation FullVA path. Its
formal-order contrast is order six versus the local order-five $m=3$
Gauss--Lobatto \tfe{} formula target; no implementation superiority follows
from this contrast. \\
What is the comparator? &
A local paper-style temporal finite-element formula target with expected order
$2m-1=5$; source-paper implementation superiority requires same-test evidence. \\
What is not claimed? &
Complete source-paper residual reproduction and accepted independent
full-\tfe{} stage replacement; faster or more accurate implemented
source-paper \tfe{} performance on the original tests is also not claimed. \\
Where is the reproducibility record? &
The reproducibility appendix lists the source package, comparison-scope files,
and read-only checks used to reproduce the reported evidence. \\
\bottomrule
\end{tabular}
\end{table}

\noindent\textbf{External-comparison guardrail.}

The comparison policy used in this paper is deliberately conservative. A
statement about implemented accuracy or cost against a source-paper method is
accepted only after a same-test row satisfies all entries in
Table~\ref{tab:same-test-rule}. Until then, the \tfe{} comparison remains the
formal-order comparison in Eq.~\eqref{eq:tfe-order}, and work/precision plots
are candidate or common-reference diagnostics.

\begin{table}[htbp]
\centering
\caption{Guardrail for promoting a baseline row to an implemented same-test comparison.}
\label{tab:same-test-rule}
\small
\begin{tabular}{P{0.24\linewidth}P{0.48\linewidth}P{0.18\linewidth}}
\toprule
Promotion item & Required evidence & Evidence boundary \\
\midrule
Mechanism and parameters &
The source-paper frictionless pendulum or revolute-friction benchmark is run
with the same masses, geometry, loads, friction law, final time, and output
variables. & Open for external same-test promotion. \\
Methods &
The same script or verified harness reports Newmark--$\beta$, trapezoidal,
\tfe{} $m=1,2,3$, and \method{} rows. & Candidate diagnostics only. \\
Grid and solver policy &
The rows use the same step sizes, same reference solution, same Newton
tolerance policy, and same accepted-grid handling. & Not fully same-test closed. \\
Error norms &
Position, velocity, orientation, angular-velocity, constraint, and reaction or
multiplier errors are evaluated with the same norm definitions. & Partly
available in common-reference rows. \\
Work metric &
Wall time, Newton iterations, Jacobian assembly time, and linear-solve time are
reported so that an accuracy claim is not separated from cost. & Diagnostic
compendium only. \\
\bottomrule
\end{tabular}
\end{table}

Consequently the manuscript does not say that \method{} is faster or more
accurate than an implemented source-paper \tfe{} method. The stronger sentence
would require the Table~\ref{tab:same-test-rule} promotion criterion to close on
the same benchmark. The accepted claim is narrower: \method{} is a
conditional sixth-order Lie-group Gauss-collocation FullVA path under the
retained branch-smoothness, residual-defect, endpoint/stability,
implementation-binding, and solver-scale
interfaces later stated in Theorem~\ref{thm:g6fullva-order}.

\section{Validation protocol}
\label{sec:validation}

The validation protocol organizes the numerical evidence into three layers:
accepted dynamic-order evidence for the \method{} branch, mechanism-coverage
evidence for lower-pair examples, and diagnostic comparison/reproducibility
boundaries. It also fixes the error norms and reference policies used in the
paper, so that the numerical tables are not merely unsupported assertions.
For a trajectory sampled at accepted output times
$t_k$, translational errors use the maximum-in-time Euclidean or infinity
norms recorded in the CSV files. Orientation errors use the Lie-group
distance
\begin{equation}
  d_R(Q_h,Q_{\mathrm{ref}})=
  \left\|\log(Q_{\mathrm{ref}}^TQ_h)\right\|_2 .
  \label{eq:rotation-error}
\end{equation}
Equation~\eqref{eq:rotation-error} defines the rotation component of the
reported trajectory error.
Angular-velocity errors use the corresponding vector norm. The observed order
reported in the tables is the least-squares slope of
$\log e(h)$ against $\log h$ over the stated step sizes. The fit protocol is
fixed before the reported row is inspected: step window, reference policy,
reported components, error norm, and inclusion or exclusion of floor-limited
rows are the policy stated for that evidence row. Recomputing a slope after
dropping a coarse or fine row, swapping position, velocity, or orientation
components, changing a norm, or moving a row between dynamic-order,
mechanism-coverage, and diagnostic status defines a different evidence row.
Such a post-hoc row may be recorded as a new diagnostic, but it cannot upgrade
the theorem claim or replace the predeclared row in the acceptance policy. A
result can support the main manuscript claim only if it satisfies the following
acceptance criteria.

\begin{enumerate}
  \item The method identifier is \method{}, and the order claim is the
  conditional sixth-order path of Theorem~\ref{thm:g6fullva-order}, not an
  unconditional property of fixed-tolerance or source-policy rows.
  \item The smooth cylindrical-chain h-sweep uses at least three step sizes and
  reports position and velocity orders above the local order-five comparator.
  \item The four ASME-style examples are accepted under the same method-side
  policy, but their roles differ: only the single and double pendula are used
  as accepted dynamic-order examples, while the closed-loop examples provide
  mechanism coverage, constraint closure, reaction residuals, and local coarse
  dynamic evidence.
  \item The source-paper comparator remains a local $m=3$ Gauss--Lobatto
  temporal finite-element target with expected order five.
  \item Any claim against a source-paper \tfe{} implementation requires
  same-test evidence with matched parameters, step sizes, tolerances,
  reference solution, error norm, and work metric; this manuscript does not
  make that source-paper superiority claim.
  \item Complete source-paper residual reproduction and independent
  full-\tfe{} stage replacement remain outside the accepted method claim.
  \item Routine reproducibility checks are documented separately in
  the reproducibility appendix.
\end{enumerate}

\begin{table}[htbp]
\centering
\caption{Reproducibility policy for the numerical evidence.}
\label{tab:experiment-policy}
\small
\begin{tabular}{P{0.22\linewidth}P{0.18\linewidth}P{0.25\linewidth}P{0.21\linewidth}}
\toprule
Evidence block & Step sizes & Reference & Primary reported quantities \\
\midrule
Smooth cylindrical chain &
$0.04,0.02,0.01$ &
$h=0.005$ trajectory &
Pos./vel. errors; constraints; Newton iterations. \\
Sharp-friction cylindrical chain &
$0.04,0.02,0.01$ and refined diagnostics &
$h=0.005$ or finer reference by diagnostic check &
Order reduction, endpoint constraints, refinement cost. \\
Driven single pendulum &
$0.2,0.1,0.05$ &
Analytic $\theta(t)$ trajectory &
Position, velocity, orientation, and residual errors. \\
Double pendulum &
$0.04,0.02,0.01$ &
Nested local FullVA $h=0.005$ reference &
Orientation and angular-velocity errors plus endpoint constraints. \\
Four-link and slider-crank &
0.02, 0.01, 0.005 &
Local kinematic FullVA $h=0.001$ reference &
Constraints; rank; reaction residuals. \\
\bottomrule
\end{tabular}
\end{table}

The paper uses three evidence types throughout the results. A dynamic-order row
is used to support an order statement for the accepted one-step map. A
mechanism-coverage row checks that the same FullVA position, velocity, and
acceleration vocabulary closes a lower-pair mechanism and produces consistent
reaction dynamics, but it is not by itself an asymptotic order row. A diagnostic
row clarifies a boundary condition, failure mode, or external-comparison gap and
is never promoted to a theorem or superiority claim. This separation is
important for the four-link and slider-crank examples, where the closed-loop
constraint/reaction checks and coarse-window trajectory traces are reported
only as coverage/consistency diagnostics; they do not become accepted
asymptotic order proof, residual-to-error evidence, or source-policy
comparison closure.

The accepted cylindrical-chain h-sweep is reported with endpoint velocity
closure, and the closure itself is checked explicitly rather than hidden inside the error
table. Two endpoint records are retained. The first is the actual projected
endpoint path used for the main finite-window order row. The second is a
square endpoint KKT check in which endpoint velocity variables, least-squares
stationarity, and endpoint velocity constraints are solved as an augmented
closure residual. Table~\ref{tab:endpoint-closure-scaling} reports the smooth
branch scaling. The endpoint functional \(E\) used in
Lemma~\ref{lem:endpoint-closure} is the velocity-level/KKT closure subsystem
actually corrected by \(\mathcal C_h\), not every raw endpoint diagnostic in
the table. The raw endpoint position defect is retained as a separate
constraint monitor: it remains small and approximately sixth order over the
reported window, but it is not the \(\|E(z_u)\|\le C_{E,\mathrm{raw}}h^7\)
hypothesis. The raw endpoint velocity defect and the KKT correction norm both
have finite-window fitted slopes near seven over the reported window, and the
corrected endpoint velocity residual is at roundoff scale. This evidence is a
finite-run consistency check for the endpoint-closure perturbation scale in
Lemma~\ref{lem:endpoint-closure}; it is not a proof of that lemma's retained
P2 endpoint-ball, raw-defect, or uniform right-inverse hypotheses, and it
neither closes the full-\tfe{} replacement nor the source-policy boundaries.
The sharp-branch coarse h-sweep remains a limitation, not theorem evidence.

\begin{table}[htbp]
\centering
\scriptsize
\setlength{\tabcolsep}{3pt}
\caption{Smooth cylindrical-chain endpoint closure scaling diagnostic.}
\label{tab:endpoint-closure-scaling}
\begin{tabular}{P{0.08\linewidth}P{0.16\linewidth}P{0.16\linewidth}P{0.16\linewidth}P{0.16\linewidth}P{0.10\linewidth}}
\toprule
$h$ & Raw pos. defect & Raw vel. defect & KKT correction & Closed vel. defect & Newton iters. \\
\midrule
0.04 & $5.883\times10^{-12}$ & $9.517\times10^{-11}$ & $7.926\times10^{-11}$ &
$2.431\times10^{-16}$ & 12 \\
0.02 & $9.196\times10^{-14}$ & $7.427\times10^{-13}$ & $6.395\times10^{-13}$ &
$1.573\times10^{-16}$ & 24 \\
0.01 & $1.635\times10^{-15}$ & $6.505\times10^{-15}$ & $5.484\times10^{-15}$ &
$3.634\times10^{-16}$ & 48 \\
Obs. order & 5.907 & 6.918 & 6.909 & roundoff limited & -- \\
\bottomrule
\end{tabular}
\end{table}

\section{Conditional sixth-order theorem}
\label{sec:order}

\noindent
Under the retained compact-branch, endpoint, stability,
implementation-binding, and solver-scale interfaces stated below, the theorem
proves a sixth-order same-initial-state reported-grid error bound by a
conditional local-defect-to-global-error argument on the accepted compact-tube
Gauss-predictor branch. It starts from the standard
three-stage Gauss order on a smooth reduced constrained flow and then accounts
for the perturbations introduced by the implemented FullVA residual, endpoint
closure, and Newton tolerance. It is not a derivation of the complete
source-paper temporal finite-element weak form.
The branch-selected map on the accepted Gauss-predictor domain is the only
theorem object in this section.
The theorem below is a conditional numerical-analysis theorem: under the fixed
compact-branch, endpoint, implementation, stability, and solver-scale
hypotheses on that branch, the \(h^6\) same-initial-state reported-grid
bound follows from a proved \(h^7\) local defect, not from fitted slopes or
downstream diagnostic records.

\paragraph{Proof spine}
The load-bearing proof route is a direct residual-to-local-defect argument.
After the smooth reduced Gauss stage \(Z_G\) is lifted into the implemented
stage variables, the proof evaluates the fixed row-scaled \(132\)-row residual
at that same \(Z_G\).  The 96 position-, velocity-, acceleration-, endpoint-,
and lower-pair rows are supplied by the non-dynamic FullVA certificate, and
the 36 Newton--Euler rows are supplied only as a same-branch base-point zero
block by the D5 direct-substitution route:
D1/D2 balance identities, the D3 virtual-work multiplier-wrench identity,
D4 smooth-force lift, and D6 row-ordering/scaling/AD binding together give
zero dynamic residual on \(Z_G\).
Together these two row families, under the retained branch, norm, and
implementation-binding tuple, give the \(O(h^7)\) implemented residual input
consumed by Lemma~\ref{lem:stage-residual-defect}: the single T2 input is
\(\|F_{A,h}(Z_G;y)\|\le C_Rh^7\) before Taylor expansion in
\(\Delta=Z-Z_G\).  The D5 identity supplies only the base-point dynamic block
of that same residual; it does not supply derivative bounds, primitive
subterm bounds, or an additional local-defect term.  This is the active route
used by the theorem.  The primitive/Taylor decomposition of the dynamic rows is
organized as a separate candidate sufficient route for verifying the same
dynamic block; until its primitive subterm bounds are proved on the same
compact branch, it stays outside the theorem and contributes no additional
summand to the local-defect estimate below.

\paragraph{Proof contract}
The theorem below should be read as the branch-local implication in
Eq.~\eqref{eq:proof-contract-implication}, with all constants fixed before
numerical evidence is used:
\begin{equation}
\label{eq:proof-contract-implication}
\begin{gathered}
  \mathcal I_h(K)
  \quad\wedge\quad
  \mathsf R_{96}(Z_G)=O(h^7)
  \quad\wedge\quad
  \mathsf D_{36}(Z_G)=0
  \quad\wedge\quad
  \eta_h^{\rm tube}\le c_\eta h^7 \\
  \Longrightarrow
  \max_{0\le n\le N_h}
  \|{\mathcal R}(y_n)-{\mathcal R}(y(t_n))\|
  \le C_{qv}h^6 .
\end{gathered}
\end{equation}
The residual part of Eq.~\eqref{eq:proof-contract-implication} is the
same direct bridge later used in the theorem's residual-certificate slot:
\begin{equation}
\label{eq:direct-bridge-contract}
\begin{gathered}
  \mathsf B_h^{96+36}(Z_G;y)
  :=
  \bigl\{\mathsf R_{96}(Z_G)=O(h^7),\ 
  \mathsf D_{36}(Z_G)=0\bigr\},\\
  \mathsf B_h^{96+36}(Z_G;y)
  \Longrightarrow
  \|F_{A,h}(Z_G;y)\|\le C_Rh^7 .
\end{gathered}
\end{equation}
Equation~\eqref{eq:direct-bridge-contract} is the early notation for the
\(\mathsf R_{\rm dir}\) certificate fixed in
Eq.~\eqref{eq:residual-certificate-slot}; it is not a second residual route
and it is not combined with the primitive/Taylor replacement template.
Here \(\mathcal I_h(K)\) denotes the retained compact-tube, inverse,
endpoint, stability, implementation-binding, and proof-norm interfaces on the
accepted Gauss-predictor branch.  The predicate
\(\mathsf R_{96}(Z_G)=O(h^7)\) is the proved non-dynamic FullVA row
certificate, and \(\mathsf D_{36}(Z_G)=0\) is the proved direct
Newton--Euler row identity on the lifted Gauss stage.  Both are formula-level
statements on the same branch-selected, row-scaled residual object, not
finite-run logged residuals, and they are fixed before the stage-root
perturbation is invoked.  Thus the theorem does
not assume the implemented residual defect as a black box: the
\(96+36\) row bridge supplies it before the Taylor/Kantorovich step.  At the
same time, \(\mathcal I_h(K)\) and the solver-scale predicate remain domain
and branch hypotheses; they are not inferred from observed slopes, solver
logs, external-comparison diagnostics, or the separate primitive/Taylor route.
The theorem closure below refers only to the direct residual-bridge/Kantorovich
argument under the stated branch, endpoint, implementation, and solver-scale
hypotheses; it does not prove a separate theorem for fixed production
tolerances, a theorem converting mechanism residual histories into trajectory
errors, or the separate primitive/Taylor route.
The Taylor step used in this implication is the full-map
Taylor/Kantorovich step for the fixed \(132\)-row residual \(F_{A,h}\), not a
term-by-term expansion of the separate primitive Newton--Euler inventory.  The
direct D5 identity supplies the dynamic block of \(F_{A,h}(Z_G;y)\) at the
base point before that full-map Taylor step is applied.  A separate primitive
\(162\)-term proof, if supplied on the same branch and proof norm, would
therefore be an alternative sufficient certificate for that dynamic-block
residual value only; it would not replace the
full-map inverse, endpoint closure, solver-scale condition, or local-to-global
slots of the theorem.

\paragraph{Residual-bridge handoff}
The theorem uses the direct dynamic-row identity,
together with the already certified 96 non-dynamic rows, to form the
\(132\)-row residual bridge at \(Z_G\).  The compact smooth branch,
local inverse, endpoint ball/right inverse, implementation binding, and
branch-solver scale are theorem-domain hypotheses rather than empirical
closures; they are fixed before the local defect is transferred.  The
P7 residual-to-error route is deliberately not
used: residual-only mechanism rows, reaction rows, common-reference rows, and
external-comparison rows cannot enter the sixth-order theorem unless a
separate transfer theorem is proved.  This is the sense in which the theorem
is strong but conditional: the local residual-defect input is proved for the
implemented FullVA residual, while the domain and solver-scale hypotheses are
stated instead of being hidden in numerical evidence.

\paragraph{Fixed theorem object}
Before any numerical row is read, fix the compact proof tube \(K\), the
reduced chart \(y\), the proof residual norm, the row ordering and scaling of
the implemented \(132\)-row residual, and the Gauss-predictor branch.  For
\(y\in K\), let \(\varphi_h(y)\) be the exact reduced constrained flow,
\(Z_G(y,h)\) the isolated three-stage Gauss stage in the reduced chart,
\(F_{A,h}(Z;y)\) the implemented \method{} residual written in the same
unknowns and row scaling, \(Z_A(y,h)\) the local zero selected from the
Gauss-predictor branch, \(\widetilde Z_A(y,h)\) the inexact Newton iterate,
\(\mathcal E_h\) the Lie endpoint reconstruction, \(\mathcal C_h\) the
accepted endpoint-closure map, and
\(\mathcal P_h=\mathcal C_h\circ\mathcal E_h\).  The one-step map whose order
is proved is Eq.~\eqref{eq:g6fva-onestep-map}:
\begin{equation}
\label{eq:g6fva-onestep-map}
  \Psi_h^{\mathrm{G6FVA}}(y)=\mathcal P_h(\widetilde Z_A(y,h)).
\end{equation}
The proof evaluates the residual bridge at the lifted Gauss stage \(Z_G\),
uses the local branch inverse to obtain \(Z_A\), then applies endpoint closure
and the inexact-Newton perturbation.  Thus the proof object is fixed
analytically before \(h\)-sweeps, Newton logs, residual tables, mechanism
coverage rows, external comparison rows, or source-package records are interpreted.
Those records can identify or corroborate an accepted finite branch, but they
cannot define the map, choose constants, or supply a separate theorem input.
This fixes the theorem object before numerical evidence is introduced,
including the retained P2 interface: the retained domain hypotheses and
discharged residual inputs are stated explicitly before the lemmas are used.

\paragraph{Data for theorem invocation}
The theorem is invoked only after the following four inputs are fixed on the
same accepted branch.  First, the compact tube, reduced chart, endpoint map,
row ordering, row scaling, proof norm, and Gauss-predictor branch define the
one-step object and its local inverse constants.  Second, the direct
residual bridge supplies the stage-residual value at the lifted Gauss stage:
the 96 non-dynamic rows contribute \(O(h^7)\), and the 36 Newton--Euler rows
vanish by direct substitution.  Third, the endpoint-closure and stability
interfaces provide the right-inverse and local-to-global constants used after
the stage perturbation.  Fourth, the P6 solver-scale envelope
\(\eta_h^{\rm tube}\le c_\eta h^7\) is retained for the selected Newton
branch.  These inputs are not interchangeable. The theorem input is not a
separate primitive/Taylor subterm closure, a residual-to-error promotion, a
fixed-tolerance log, or a source-policy comparison row; those records are
outside this invocation unless their own theorem interfaces are separately
proved.

\noindent\textbf{Positive and retained theorem inputs.}
The invocation contains two logically different kinds of theorem input.  The
positive residual input is the proved \(96+36\) row bridge at \(Z_G\): the
P4 row-local estimate and the P5 direct dynamic-row identity together provide
the residual value consumed by Lemma~\ref{lem:stage-residual-defect}.  The
retained interfaces are the compact chart/tube, local inverse, endpoint,
stability, implementation-binding, and P6 solver-envelope hypotheses; they
define the admissible branch and constants, but they are not outputs of P5 or
of numerical evidence.  The proof uses the positive residual input only after
the retained interfaces have fixed the branch and norm, and it uses the
retained interfaces only as hypotheses.  No retained interface is promoted to
a proved residual row, and no proved residual row is reused as a proof of an
interface.

\noindent\textbf{Load-bearing proof skeleton.}
The theorem has five analytic moves, all on the same accepted branch: (i) the
classical three-stage Gauss defect gives the reduced \(C_Gh^7\) local
truncation term; (ii) the \(96+36\) direct residual bridge plus the
Kantorovich inverse gives the accepted stage-root perturbation \(C_Ah^7\);
(iii) the endpoint-closure right-inverse gives \(C_Eh^7\); (iv) the retained
P6 branch-solver scale gives the inexact-Newton term \(C_Nc_\eta h^7\); and
(v) the first-exit/discrete-Gronwall transfer plus the reporting map converts
the resulting one-step \(C_{\rm loc}h^7\) bound into the \(h^6\)
same-initial-state reported-grid estimate.  The surrounding tables are only
signposts for this sequence:
they identify the common theorem object, constant dependencies, and excluded
claims, but they are not additional local-defect estimates and not substitutes
for any of the five analytic moves.

\paragraph{One-step perturbation decomposition}
For a fixed retained-tube state \(y\) and step \(h\), the proof uses the
four maps in Eq.~\eqref{eq:onestep-perturbation-maps} on the same branch:
\begin{equation}
\label{eq:onestep-perturbation-maps}
\begin{aligned}
  \Psi_h^G(y)&=\mathcal E_h(Z_G(y,h)),\\
  \Psi_h^A(y)&=\mathcal E_h(Z_A(y,h)),\\
  \Psi_h^E(y)&=\mathcal P_h(Z_A(y,h)),\\
  \Psi_h^{\mathrm{G6FVA}}(y)&=\mathcal P_h(\widetilde Z_A(y,h)).
\end{aligned}
\end{equation}
After the smooth FullVA lift identifies the reduced Gauss endpoint with
\(\mathcal E_h(Z_G)\), the four comparisons in the theorem are, in order:
Gauss truncation \(\Psi_h^G-\varphi_h\), accepted stage-root perturbation
\(\Psi_h^A-\Psi_h^G\), endpoint closure \(\Psi_h^E-\Psi_h^A\), and inexact
Newton perturbation \(\Psi_h^{\mathrm{G6FVA}}-\Psi_h^E\).  Each term is
evaluated for the same \(y\), same \(h\), same compact tube, same row-scaled
residual, and same endpoint-output map.  Thus the local \(O(h^7)\) defect is
one triangle inequality on a single branch-selected one-step object, not a
patchwork of Gauss, endpoint, Newton, or diagnostic quantities from different
runs.

\paragraph{Reference-paper proof-method alignment}
The proof organization follows the Lie-group constrained-BDF convergence
analysis of Wieloch and Arnold \cite{wieloch2021bdf}, not the temporal
finite-element comparator used later for formal-order context.  The
constrained-BDF proof paper and the comparator paper therefore have disjoint
roles: the former organizes the proof-obligation order, whereas the latter supplies
only the order-five formula target and external comparison boundary used
outside the theorem.  The constrained-BDF reference
first fixes the Lie-group update and increment variables, then proves a
Taylor-expansion local truncation estimate, defines global configuration error
in the Lie algebra, uses Baker--Campbell--Hausdorff perturbation estimates for
the exponential composition, and finally closes the constrained problem through
holonomic-constraint differences, a multiplier estimate, and a coupled error
recursion.  We invoke the reference as a proof-obligation ordering template, not
as a literal proof transfer: the BLieDF hidden-constraint multiplier estimate
and coupled multistep recursion are not repeated below.  More specifically,
the reference is invoked only for the Section~6.2/Theorems~2--4 style of proof
sequencing: Theorem~2 local truncation estimates, Theorem~3 configuration and
velocity error recursion, and Theorem~4 hidden-constraint difference
approximation, multiplier-error estimate, and coupled-recursion closure.
It supplies no FullVA residual Taylor
constant, no D5 primitive lift bound, and no one-step Gauss theorem for the
present residual.  Adapted to one-step
Gauss collocation, the stage vector \(Z\), FullVA derivative rows, lower-pair
constraints, Newton--Euler rows, row scaling, proof norm, and endpoint
reconstruction are fixed first as one algebraic object.  The formal Gauss
order is used only for the reduced smooth Gauss branch that supplies
\(\Psi_h^G-\varphi_h\).  FullVA enters the next proof slot through the fixed
\(132\)-row residual bridge evaluated at \(Z_G\): the direct-route
\(O(h^7)\) stage-residual input is discharged as the residual value of that
bridge, while the compact-tube inverse, branch, endpoint, and solver-scale
interfaces remain theorem hypotheses before any \(h\)-sweep, residual table,
work/precision row, or source-policy row is interpreted.
This is the manuscript's reference-style Taylor slot T2: fix the discrete
object, insert the smooth branch into the fixed \(F_{A,h}\), then use the
full-map Taylor/Kantorovich perturbation.  The T3 primitive Taylor inventory is
a separate alternative certificate specification for the dynamic-block residual
value; because its primitive lift hypotheses are not theorem inputs here, it
is not the Taylor layer required by the accepted theorem route.
If a future primitive/Taylor proof is completed, its mathematical role is only
to replace the dynamic-row residual-value certificate in the same \(132\)-row
residual bridge before the T2 full-map Taylor/Kantorovich step; it would not
create a second Taylor layer, add a second residual summand, or strengthen the
accepted theorem by combination with the direct certificate.
Thus the reference paper informs the order of proof
obligations---define the Lie-group reconstruction, prove the local defect,
transfer the defect through the constrained algebraic solve, and only then
report numerical diagnostics after the theorem boundary is fixed---while the
load-bearing sixth-order estimate here is supplied by the
direct residual-bridge and perturbation lemmas below, not by importing the
BDF \(p=k\) theorem.
The reference-paper correspondence is therefore structural and stepwise, not
an estimate transfer.  The reference local-truncation step corresponds only
to the proof slot occupied below by the four-term local defect sum; its
hidden-constraint and multiplier steps correspond only to the proof slots
occupied here by the retained compact FullVA chart, endpoint, lower-pair
right-inverse, and branch-stability interfaces; the Newton--Euler row
identities are separate FullVA residual-bridge inputs, not imported
constraint/multiplier estimates. Its final coupled error recursion
corresponds only to the accepted-branch stability recurrence used
after Lemma~\ref{lem:local-global-transfer}.  The estimates themselves are
not reused: every constant in Theorem~\ref{thm:g6fullva-order} is fixed in
Assumption~\ref{ass:regularity} or proved by the FullVA lemmas below.  The
present proof uses the same discipline--insert the smooth
exact branch into the fixed discrete equations before proving global
error--but replaces the multistep BDF recursion by a one-step
Gauss-collocation residual bridge on the \(132\)-row FullVA system.
In particular, the reference Appendix's coupled recursion for stacked
history variables such as \(E_y,E_z\), and its zero-stability norms for
the BDF companion matrices \(T_y,T_z\), are not theorem inputs here.  They
mark the proof slot occupied in this manuscript by the retained one-step
stability hypothesis for \(\Psi_h^{\mathrm{G6FVA}}\) and the scalar/vector
Gronwall recurrence for the reduced-chart error \(e_n\).  This replacement
is legitimate only because the FullVA stage multipliers, lower-pair
constraints, endpoint map, and Newton--Euler rows have already been consumed
inside the same \(132\)-row stage-root object before the recurrence is
written; it does not prove multiplier-history, reaction-output, or
residual-output order.
The reference theorem also carries multistep start-value assumptions and a
separate multiplier-error recursion. Those are not silently omitted here:
there is no \(2k-1\) start-up layer for the one-step accepted map, and the
lower-pair multipliers are stage unknowns inside \(F_{A,h}\), consumed through
the same residual bridge and endpoint perturbation rather than reported as a
separate global-error component. A separate trajectory-order theorem for
multiplier histories or residual-only mechanism rows would require additional
outputs and assumptions; it is not part of
Theorem~\ref{thm:g6fullva-order}.
This is an output-scope choice backed by the retained theorem interfaces, not
a deletion of the constrained part of the proof.  The recurrence below is
written for the reduced-chart state \(y\) after the lower-pair algebraic
variables have been absorbed into the accepted stage-root map on the compact
branch.  That absorption is admissible only under the same compact inverse,
endpoint right-inverse, branch-ball, and stability hypotheses that appear in
Assumption~\ref{ass:regularity}.  It therefore supports the stated
position--velocity grid estimate, but it does not by itself give a multiplier
history, reconstructed reaction-force, or residual-output order estimate.  To
claim those outputs one would have to augment the reporting map and prove a
separate multiplier/reaction stability or residual-to-error theorem on the
same branch.

\noindent\textbf{Reference proof-order correspondence.}
Table~\ref{tab:reference-proof-order-correspondence} records how the
Lie-group constrained-BDF proof architecture of Wieloch and
Arnold~\cite{wieloch2021bdf} is invoked in this manuscript: as an ordering
discipline, not as a theorem transfer.
This table is the canonical reader-facing map.  The following
no-estimate-transfer and translation-map blocks are guardrails on that same
map, not additional proof routes or theorem inputs.
\begin{table}[htbp]
\centering
\caption{Reference proof-order correspondence.  The reference proof is used
only to organize the FullVA proof obligations; none of its BLieDF estimates,
BDF order theorem, multiplier-history estimates, or source-paper comparison
claims are imported.}
\label{tab:reference-proof-order-correspondence}
\scriptsize
\setlength{\tabcolsep}{2pt}
\renewcommand{\arraystretch}{0.88}
\begin{tabular}{P{0.25\linewidth}P{0.34\linewidth}P{0.31\linewidth}}
\toprule
Reference proof step & Present FullVA analogue & What is not imported \\
\midrule
Fix Lie-group increments and reconstruction before the error proof &
Fix the stage vector \(Z\), SO(3) retractions, FullVA rows, row scaling,
proof norm, and endpoint reconstruction before any estimate. &
No BLieDF multistep update, no hidden-coordinate convention, and no
post-hoc residual-row selection. \\
Taylor local truncation estimate for the fixed discrete equations &
Insert the smooth lifted Gauss stage \(Z_G\) into the implemented
\(132\)-row residual, use Lemma~\ref{lem:full-132-row-residual-bridge} to
obtain the \(O(h^7)\) residual defect from the 96-row non-dynamic certificate
and the D5 direct-substitution certificate, and then use
Lemma~\ref{lem:stage-residual-defect} to transfer that residual defect to the
accepted local stage-root and endpoint perturbation. &
No BLieDF Taylor algebra, hidden-constraint or multiplier estimate, BDF
\(p=k\) theorem, source-paper/\tfe{} estimate, FullVA residual constant, or
primitive D5 Taylor subterm is imported. No primitive D5 Taylor subterm is
certified by the reference correspondence. The primitive/Taylor route is a
separate sufficient specification for verifying the dynamic-block residual
value; no primitive D5 subterm bound is consumed by the accepted theorem
unless that separate route is independently proved on the same branch and
reassembled with the same 96-row non-dynamic block into a same-tuple
132-row residual certificate. \\
Constraint and multiplier estimates before global propagation &
Retain the compact FullVA chart, lower-pair right inverse, endpoint closure,
and branch stability as theorem-domain interfaces; the Newton--Euler row
identities enter separately as FullVA-specific residual-bridge inputs. &
No source-policy row, residual table, or fixed-tolerance run supplies these
interfaces. \\
Multistep start values and multiplier-error recursion &
For the one-step accepted map, use the single reduced initial state and
absorb stage multipliers into the \(132\)-row stage solve and endpoint
perturbation chain. &
No BLieDF start-up theorem, multiplier-history order theorem, or residual-row
trajectory-error theorem is imported. \\
Coupled error recursion and Lie-algebra error transfer &
Use the accepted-branch stability recurrence and the reporting-map
perturbation after the local defect has been proved for the same branch. &
No BDF \(p=k\) theorem, no source-paper superiority claim, and no
residual-to-error transfer theorem for mechanism-coverage rows. \\
\bottomrule
\end{tabular}
\end{table}

\noindent\textbf{Reference no-estimate-transfer table.}
The correspondence above is intentionally one-way: it fixes the order in
which proof obligations are discharged, but it supplies no numerical,
asymptotic, or hidden-constraint estimate for the present method.  The
following table states which reference estimates are not imported and gives
the replacement argument used in this proof for each corresponding role.
\begin{table}[htbp]
\centering
\caption{No-estimate-transfer table for the Wieloch--Arnold constrained-BDF proof.}
\label{tab:reference-no-estimate-transfer-ledger}
\scriptsize
\setlength{\tabcolsep}{2pt}
\renewcommand{\arraystretch}{0.9}
\begin{tabular}{P{0.24\linewidth}P{0.31\linewidth}P{0.35\linewidth}}
\toprule
Reference ingredient & Forbidden estimate transfer & Replacement in this proof \\
\midrule
BLieDF local-truncation algebra &
No BLieDF Taylor constant, BDF \(p=k\) conclusion, or multistep defect bound
is a premise for Theorem~\ref{thm:g6fullva-order}. &
The local \(h^7\) input is the three-stage Gauss defect plus the implemented
132-row residual bridge evaluated at the lifted Gauss stage. \\
BCH perturbation and Lie-algebra recursion &
The BCH constants and global Lie-algebra recursion constants from the
reference paper are not imported. &
SO(3) chart and reporting constants are retained compact-tube constants in
Assumption~\ref{ass:regularity} and the reporting-map lemmas. \\
Hidden-constraint and multiplier estimates &
The reference multiplier-history estimate and hidden-constraint estimate are
not used to control lower-pair multipliers or endpoint closure here. &
Lower-pair and endpoint effects enter through the FullVA residual rows,
endpoint/KKT closure lemma, and retained right-inverse interfaces. \\
Reference correspondence table &
The table is not a theorem input, not a certificate for primitive Taylor
subterms, and not a source-policy or full-\tfe{} replacement argument. &
It is only an ordering rule enforcing the proof order:
fixed object, smooth-branch insertion, local defect, and local-to-global
propagation. \\
\bottomrule
\end{tabular}
\end{table}

\noindent\textbf{Reference-to-FullVA translation map.}
The reference proof is used only as a partial order on obligations, and
each translated obligation below is discharged or retained by a present
manuscript object.  R1 fixes the discrete Lie-group object before expansion;
here this is the accepted Gauss-predictor branch, \(F_{A,h}\), \(Z_G\),
\(\mathcal E_h\), \(\mathcal C_h\), the row scaling, and the fixed proof norm.
R2 inserts the smooth exact branch into that same object; here this produces
the Gauss defect \(C_Gh^7\) and the full residual value
\(\|F_{A,h}(Z_G;y)\|\le C_Rh^7\) from the direct \(96+36\) bridge.
R3 converts the fixed-object defect to a nearby accepted algebraic solution;
here the conversion is Lemma~\ref{lem:stage-residual-defect} together with
the endpoint-closure and inexact-Newton perturbation lemmas, using constants
\(C_A\), \(C_E\), and \(C_N\) fixed on the compact tube.  R4 applies a
stability recursion only after the local defect is fixed; here this is the
first-exit/local-to-global transfer and the reporting-map consequence.  This
translation is one-way: every right-hand object is proved or retained
independently in the present FullVA proof.  If any right-hand lemma or retained
interface failed, the constrained-BDF reference theorem would not repair it.
The reference therefore supplies a proof-order dependency graph, not a theorem
premise, source-policy substitute, primitive Taylor bound, solver-policy
discharge, or residual-to-error promotion.

\noindent\textbf{Reference constrained-DAE slot map.}
The constrained part of the reference proof has two distinct DAE slots: the
local truncation estimate for the fixed Lie-group equations and the later
hidden-constraint/multiplier estimate obtained from constraint differences
before the coupled recursion is closed.  In the present proof these slots are
not merged.  The local-defect slot is occupied by the Gauss defect plus the
full \(132\)-row residual bridge and
Lemma~\ref{lem:stage-residual-defect}.  The constrained-DAE slot is occupied
by the retained compact FullVA chart, lower-pair row certificate, endpoint
KKT/right-inverse hypotheses, and branch-stability interfaces, together with
the fact that multipliers are stage unknowns in \(F_{A,h}\), not separate
reported histories.  Consequently the reference multiplier argument is not
imported as a multiplier-order theorem, and the FullVA proof does not bypass
the constraint layer: it records exactly where those DAE hypotheses are
consumed before Lemma~\ref{lem:local-global-transfer} and the reporting-map
estimate are applied.  The slot analogy is not symmetric: the BDF
hidden-constraint estimate cannot replace the FullVA lower-pair
position/velocity/acceleration row certificate or endpoint KKT interface, and
the FullVA direct residual bridge cannot be read as proving the reference
paper's multiplier-history theorem.

\noindent\textbf{Reference-style constrained-slot contract.}
At the theorem-interface level the reference analogy is the following
four-slot contract, with every slot discharged or retained by a present
FullVA object.  C1, the local-error slot, is filled by the smooth Gauss
defect, the implemented \(132\)-row residual bridge, and
Lemma~\ref{lem:stage-residual-defect}.  C2, the constrained algebraic slot,
is filled by the lower-pair position-, velocity-, and acceleration-row
certificate, the retained compact chart/right-inverse interfaces, the
endpoint KKT closure, and the fact that multipliers are solved as stage
unknowns inside \(F_{A,h}\).  C3, the one-step perturbation slot, is filled
by the Kantorovich stage-root, endpoint, and inexact-Newton perturbation
lemmas on the same compact branch and proof norm.  C4, the global-recursion
slot, is filled only after C1--C3 by the accepted-branch stability recurrence
and reporting-map estimate.  No slot imports a BLieDF constant,
multiplier-history estimate, BDF zero-stability theorem, finite residual log,
or source-policy row.  Conversely, if any of C1--C4 failed for the present
FullVA object, the Wieloch--Arnold theorem would not close that slot for us;
the contract is a proof-order guard, not an extra theorem premise or an
estimate-transfer device.

\noindent\textbf{Typed-slot causality rule.}
The C1--C4 labels above are not merely a narrative outline.  They fix the
causal order in which estimates may be consumed.  A later slot may use the
output of an earlier slot only after the earlier slot has been proved or
retained on the same FullVA object; it may not be read backward to prove that
earlier slot.  Thus the global recursion in C4 cannot certify the local
residual bridge in C1, a finite reported order row cannot certify the P6
solver-scale envelope in C3, and a mechanism residual or reaction table cannot
certify the P7 residual-to-error theorem outside the stated output scope.
Likewise, the separate primitive \(162\)-term Taylor route cannot borrow the
T2 full-map Kantorovich conclusion to prove its own primitive lift
hypotheses.  This slot-local causality is the present proof's operational
use of the constrained-BDF reference discipline: it preserves a forward
chain from fixed discrete object, to smooth-branch insertion, to local defect,
to global/reporting consequence, without importing any estimate from the
reference theorem or from downstream diagnostics.

\begin{lemma}[Reference-order contract and no-estimate transfer]
\label{lem:reference-order-contract}
The use of the Wieloch--Arnold constrained-BDF proof in this section is a
one-way proof-order contract.  It allows the following typed ordering rule
and no other theorem input: first fix the discrete Lie-group object, then
insert the smooth branch into that fixed object, then prove the local defect
for that object, then propagate the accepted-branch error.  This is an
expository ordering rule, not a mathematical premise of
Theorem~\ref{thm:g6fullva-order}; deleting the reference comparison would
leave the theorem dependent only on the FullVA assumptions and lemmas stated
here.  The reference theorem's technical-neighborhood and start-value clauses
are not imported either; in the present FullVA theorem they are replaced by the
retained compact-tube/branch domain, endpoint/right-inverse, solver-scale, and
same-reduced-initial-state clauses, all fixed before any diagnostic table is
read.  In the present
FullVA theorem these four slots are instantiated respectively by the fixed
row-scaled residual \(F_{A,h}\) and endpoint map, the lifted Gauss stage
\(Z_G\), Lemmas~\ref{lem:full-132-row-residual-bridge} and
\ref{lem:stage-residual-defect} together with the endpoint and inexact-Newton
perturbation lemmas, and finally Lemma~\ref{lem:local-global-transfer} plus
the reporting-map consequence.  No BLieDF local-truncation constant, BDF
\(p=k\) theorem, BCH recurrence constant, hidden-constraint estimate,
multiplier-history estimate, start-value theorem, primitive Newton--Euler
Taylor subterm bound, source-policy row, solver log, or residual-to-error
transfer theorem is imported from the reference paper.
\end{lemma}

\begin{proof}
Each estimate used in Theorem~\ref{thm:g6fullva-order} is typed by the
object on which it is proved.  The constrained-BDF reference estimates are
estimates for a BLieDF multistep update, its Lie-algebra increments, its
hidden constraints, and its multiplier recursion.  The estimates needed here
are estimates for the one-step Gauss/FullVA stage vector \(Z\), the
row-scaled \(132\)-row map \(F_{A,h}\), the endpoint-output map
\(\mathcal P_h\), and the accepted compact-tube branch.  These objects have
different unknowns, residual rows, endpoint maps, start-up structure, and
global recursions, so a numerical constant or asymptotic conclusion attached
to the BLieDF object cannot be substituted for a FullVA theorem interface.

The only invariant common to the two proofs is the order in which typed
objects must be fixed and used.  The present proof fixes \(F_{A,h}\), its row
ordering, row scaling, proof norm, endpoint map, and branch before any
Taylor or stability estimate is invoked.  It then inserts \(Z_G\) into that
same object and obtains the local residual defect from the \(96+36\) direct
bridge.  The residual defect is converted to a local stage-root and endpoint
defect by the compact-tube Taylor/Kantorovich and endpoint lemmas, and only
after those local estimates are available is the stability recursion applied.
Thus the reference proof is used only to order obligations; the actual
assumptions, constants, and conclusions are exactly those stated or proved for
the FullVA object in this manuscript.

Consequently, if any FullVA slot above failed, the Wieloch--Arnold theorem
would not close it: the BDF proof is not a replacement for the P1/P2 compact
branch interfaces, the P6 solver-scale envelope, the P7 residual-to-error
transfer theorem, or the separate primitive/Taylor lift obligations.  This
proves both directions needed for the contract: the reference-order analogy
prevents estimates from being used out of order, and the no-estimate-transfer
rule prevents reference-paper constants or conclusions from being used as
missing FullVA hypotheses.
\end{proof}

\noindent\textbf{Reference-order Taylor discipline.}
The Taylor-related lesson from the Wieloch--Arnold proof is structural
ordering discipline only; no Taylor estimate, primitive lift bound, or
multiplier estimate is imported.  Their Theorem~2 first inserts the analytical
solution into fixed BLieDF equations and proves the local defect estimates;
their Theorem~4 then uses constraint differences and a multiplier estimate
before closing the coupled recursion.  In the present one-step FullVA proof,
the analogue of the first slot is fully occupied by T1 and T2 below: the
smooth Gauss defect and the Taylor/Kantorovich expansion of the fixed
\(132\)-row map \(F_{A,h}\).  The reference proof separates this local-defect
slot from the constrained multiplier/coupled-recursion slot.  In this
manuscript, T1 and T2 fill only the local-defect slot for the fixed FullVA
object; the later Gronwall step is a reduced-chart propagation step and
imports no BLieDF multiplier, hidden-constraint, or coupled-recursion
estimate.  A literal primitive Newton--Euler Taylor proof would still have to
prove, inside the present stage variables and row scaling, the non-circular
root/direct primitive obligations \(P_{\mathrm{state}}\),
\(P_{\mathrm{acc}}\), and \(P_{\lambda}\).  The actual
\(P_{\lambda}\) rate is not borrowed from the reference paper: in the
primitive dependency graph its PL4 propagation is instantiated only after
independent \(P_{\mathrm{state}}\) and \(P_{\mathrm{acc}}\) lift inputs are
available, and only then may the downstream \(P_{\mathrm{geom}}\) and
\(P_{\mathrm{gyro}}\) reductions be used.  Thus the active
T2 full-map Taylor/Kantorovich estimate is established only for the fixed
\(132\)-row residual map \(F_{A,h}\), where ``established'' means the direct T2
residual-defect input is discharged for that theorem object only, while the
separate primitive subterm
Taylor proof remains conditional at the analogous present FullVA
lift/multiplier interfaces; the reference paper only motivates the ordering
discipline, not those bounds.

\paragraph{Taylor-layer separation}
There are three distinct Taylor-layer slots in this manuscript.  T1 is the
standard smooth Gauss collocation local-truncation layer: in the reduced chart
the exact smooth branch gives the \(C_Gh^7\) defect for
\(\Psi_h^G-\varphi_h\).  T2 is the theorem-level FullVA residual-map layer:
after the implemented residual, row ordering, row scaling, proof norm, and
branch have been fixed, Lemma~\ref{lem:stage-residual-defect} expands the
single \(132\)-row map \(F_{A,h}\) in the full stage increment
\(\Delta=Z-Z_G\), uses the bridged value
\(\|F_{A,h}(Z_G;y)\|\le C_Rh^7\), and applies the compact-tube
Kantorovich argument to transfer that formula defect to the accepted local
stage root.  T3 is the separate primitive Newton--Euler Taylor route: it
decomposes the same 36 dynamic rows into state, acceleration, multiplier,
geometry, and gyroscopic lift subterms.  Only T1 and T2 are load-bearing in
Theorem~\ref{thm:g6fullva-order}.  T3 is a separate primitive
residual-certificate specification for the same 36 dynamic rows:
it is a stricter alternative factorization of the dynamic-row residual value
and is useful as a finer verification route, but it is not the local
truncation proof used by the theorem and it is not imported into the active
residual-bridge estimate.  The primitive subterm-bound inventory is therefore
a conditional specification for that separate sufficient route, not evidence
against the direct theorem route and not an additional summand in \(C_{\rm loc}\).
Equivalently, the conditional Taylor templates and the actual primitive-route
Taylor bounds are two separate mathematical statuses: the first gives a
row-by-row reduction plan under named lift hypotheses, while the current
absence of certified actual primitive-route bounds prevents those conditional
templates from being spliced into the direct residual-bridge proof.
There is also no partial promotion by selecting apparently benign subterms:
the theorem consumes one norm bound for the fixed dynamic residual block, and
the primitive route contributes no termwise dynamic-block bound until the
common lift and multiplier inputs are proved on the same compact tube.
The active theorem already consumes the direct dynamic-row identity through
the full 132-row residual bridge.
Thus the Taylor analysis is not being omitted or replaced by a numerical
record.  The active Taylor argument is Lemma~\ref{lem:stage-residual-defect}:
Taylor's formula with integral remainder expands the implemented \(132\)-row
residual about \(Z_G\), the \(O(h^7)\) residual value is supplied by the
direct bridge, and the compact-tube inverse absorbs the quadratic remainder.
Concretely, the Taylor variable is the full stage increment
\(\Delta=Z-Z_G\) of the scaled residual map \(F_{A,h}\); the input is the
single norm bound \(\|F_{A,h}(Z_G;y)\|\le C_Rh^7\) after the row-block bridge
has already supplied the 96-row non-dynamic \(O(h^7)\) block and the 36
Newton--Euler zero rows.  It is not a term-by-term Taylor expansion of the
162 primitive Newton--Euler subterms.
The primitive Newton--Euler route is a stricter factorization of one block of
that residual value, not a prerequisite for the theorem-level Kantorovich
stage-root perturbation.  This matches the reference-paper cadence at the
level needed here: insert the smooth branch into the fixed discrete equations,
prove the local defect in those equations, and only then propagate the
accepted-branch error.
The typed Taylor objects are therefore Eq.~\eqref{eq:taylor-layer-typed-objects}:
\begin{equation}
\label{eq:taylor-layer-typed-objects}
\begin{aligned}
\mathsf T_1&:\quad
  \|\Psi_h^G(y)-\varphi_h(y)\|\le C_Gh^7,\\
\mathsf T_2&:\quad
  \mathcal I^{\mathrm{T2}}_{\rm inv}(K)\wedge
  \|F_{A,h}(Z_G;y)\|\le C_Rh^7
  \Longrightarrow
  \|\Psi_h^A(y)-\Psi_h^G(y)\|\le C_Ah^7,\\
\mathsf T_3&:\quad
  (P_{\mathrm{state}},P_{\mathrm{acc}},P_{\lambda},
  P_{\mathrm{geom}},P_{\mathrm{gyro}})\\
&\quad\Longrightarrow
  \mathcal T_{\mathrm{D5}}^{162}=O(h^7)
  \quad\hbox{for all primitive D5 subterms}.
\end{aligned}
\end{equation}
In this \(\mathsf T_3\) shorthand, \(P_{\mathrm{tube}}\) is the already
closed compact-tube constant slot; the five displayed antecedents are the
remaining open primitive lift/reduction inputs.
Here \(\mathcal I^{\mathrm{T2}}_{\rm inv}(K)\) denotes the retained full-map
inverse, branch ball, and second-derivative bounds used by
Lemma~\ref{lem:stage-residual-defect}; the T2 implication in
Eq.~\eqref{eq:taylor-layer-typed-objects} is not a
residual-only implication without those compact-tube hypotheses.
Theorem~\ref{thm:g6fullva-order} consumes \(\mathsf T_1\) and
\(\mathsf T_2\).  A completed \(\mathsf T_3\) would be a separate sufficient
route only to the dynamic-row part of the same \(C_Rh^7\) bridge.  It would
not replace the full-map inverse, endpoint, inexact-Newton, or local-to-global
steps, and failure to complete \(\mathsf T_3\) is not the negation of the
direct bridge.  It only rules out treating primitive subterm bounds whose
hypotheses have not been established as theorem inputs.
Conversely, a separate completed T3 certificate would not be an appendix shortcut that
retroactively changes the present proof: it would be a replacement residual
certificate for the dynamic block, after which the same full-map inverse,
endpoint, inexact-Newton, and local-to-global steps would still have to be
invoked on the same branch.
Equivalently, a separate T3 proof could only change the provenance of the
36-row dynamic-block contribution to the same \(C_Rh^7\) residual input.  It
would replace the direct P5 zero-row certificate for that block, not add a
second local-defect term, not bypass the T2 full-map inverse/Kantorovich
step, and not discharge the endpoint, P6 solver-scale, P7 residual-to-error,
or global-transfer interfaces.

\begin{table}[htbp]
\centering
\caption{Taylor-layer dependency map for the accepted proof route.}
\label{tab:taylor-layer-consumption-ledger}
\scriptsize
\setlength{\tabcolsep}{2pt}
\renewcommand{\arraystretch}{0.92}
\begin{tabular}{P{0.10\linewidth}P{0.28\linewidth}P{0.26\linewidth}P{0.28\linewidth}}
\toprule
Slot & Fixed object and estimate & Theorem use & Explicit non-use \\
\midrule
T1 &
Reduced smooth Gauss collocation defect in the retained chart,
\(\|\Psi_h^G(y)-\varphi_h(y)\|\le C_Gh^7\). &
Positive local-defect summand before endpoint and solver perturbations. &
Does not choose the FullVA branch, residual row scaling, endpoint map, or
solver residual policy. \\
T2 &
Full implemented \(132\)-row residual map \(F_{A,h}\) in the stage increment
\(Z-Z_G\), with \(\|F_{A,h}(Z_G;y)\|\le C_Rh^7\). &
The load-bearing Taylor/Kantorovich step converting the bridged residual value
to the accepted stage-root perturbation. &
Does not decompose the dynamic block into primitive subterms and does not
infer any primitive lift estimate. \\
T3 &
Separate \(162\)-subterm primitive Newton--Euler Taylor inventory under
\(P_{\mathrm{state}},P_{\mathrm{acc}},P_{\lambda},
P_{\mathrm{geom}},P_{\mathrm{gyro}}\). &
No theorem use in the accepted route; it is only a separate alternative
dynamic-block residual-certificate route for \(C_Rh^7\).
It can replace the direct dynamic-block certificate only after the five
antecedent primitives are proved on the same branch and the resulting
dynamic block is assembled with the same 96-row non-dynamic certificate into
the same \(132\)-row residual bridge. &
Cannot be added to T2, used to lower constants, remove the solver-scale
condition, prove a residual-to-error theorem for reported mechanism rows, or
claim primitive-route closure before those antecedents are proved. \\
\bottomrule
\end{tabular}
\end{table}

Taylor-layer consumption rule: the reference-style ordering is fixed object
first, smooth-branch insertion second, local defect third, and propagation
last.  In this manuscript that means T2 consumes exactly one same-branch
dynamic-block residual certificate before the full-map Taylor/Kantorovich step
is invoked.  In the accepted theorem instance that certificate is the direct
P5 substitution certificate together with the 96-row non-dynamic certificate,
giving \(\|F_{A,h}(Z_G;y)\|\le C_Rh^7\).

The D5 primitive route remains a replacement-certificate route, not a second
Taylor layer.  For the results claimed here it is only a template-level
alternative: it supplies conditional templates for \(162/162\) subterms but
no certified primitive Taylor bounds and five open lift/bilinear primitives.
A future completed T3 proof on the same branch
and proof norm could replace the dynamic-block certificate only inside the
same \(132\)-row residual bridge that feeds the T2 chain; it would not be
added to T2, lower constants, discharge P6/P7, or shorten the endpoint and
global-transfer steps.  Definition~\ref{def:t3-primitive-route-certificate-template}
and Lemma~\ref{lem:t3-replacement-certificate-use-rule} spell out this
replacement certificate as a same-object closure target and use rule.
The canonical status card for this separate route is Eq.~\eqref{eq:t3-nonclosure-status-card}:
\begin{equation}
\label{eq:t3-nonclosure-status-card}
\mathsf S_{\rm T3}
=
\left(
\begin{array}{l}
\text{conditional primitive templates }162/162,\\
\text{certified primitive Taylor bounds }0/162,\\
\text{open same-object primitive inputs }5,\\
\text{theorem use: replacement only, not additive}
\end{array}
\right).
\end{equation}
Here \(\mathsf S_{\rm T3}\) records the current status of the separate T3
route, not a new theorem assumption.  Every later reference to primitive/Taylor
nonclosure, replacement-only use, or non-additivity is tied to
Eq.~\eqref{eq:t3-nonclosure-status-card} unless a future manuscript version
replaces it by a discharged same-object T3 certificate.

\begin{definition}[T3 primitive-route certificate template]
\label{def:t3-primitive-route-certificate-template}
For a fixed theorem instance, a completed T3 primitive-route certificate would
be the same-object tuple in Eq.~\eqref{eq:t3-certificate-tuple},
\begin{equation}
\label{eq:t3-certificate-tuple}
  \mathcal T_{\rm T3}
  =(K,\chi,Z_G,F^{\rm dyn}_{A,h},\Pi_{\rm D5},\tau_h,
    \|\cdot\|_{\rm proof},C_{\rm state},C_{\rm acc},C_{\lambda},
    C_{\rm geom},C_{\rm gyro},\Gamma_{\rm D5},h_0),
\end{equation}
chosen before any finite residual log, observed slope, direct residual-bridge
output, or source-policy row is read.  Here \(K\) and \(\chi\) are the same
compact proof tube and chart used by Theorem~\ref{thm:g6fullva-order},
\(Z_G\) is the lifted smooth Gauss stage, \(F^{\rm dyn}_{A,h}\) is the
36-row Newton--Euler block of the implemented residual map, \(\Pi_{\rm D5}\)
is the fixed primitive D5 inventory map, and
\(\tau_h=(\tau_{1,h},\ldots,\tau_{162,h})\) is the ordered primitive
Taylor-subterm vector.  The tuple must certify, uniformly for
\(0<h\le h_0\) on the same branch, the two root lifts in
Eq.~\eqref{eq:t3-root-lift-bounds},
\begin{equation}
\label{eq:t3-root-lift-bounds}
  P_{\rm state}:\|\delta S\|\le C_{\rm state}h^7,\qquad
  P_{\rm acc}:\|\delta A\|\le C_{\rm acc}h^7,
\end{equation}
with \(P_{\rm acc}\) in the unweighted acceleration norm used by the D5
Newton--Euler rows.  It must then propagate \(P_{\lambda}\),
\(P_{\rm geom}\), and \(P_{\rm gyro}\) only forward from those root lifts,
produce all \(162\) bounds in Eq.~\eqref{eq:t3-subterm-bounds},
\begin{equation}
\label{eq:t3-subterm-bounds}
  \|\tau_{j,h}(Z_G;y)\|\le C_jh^7,\qquad 1\le j\le 162,
\end{equation}
and use the fixed finite assembly bound in Eq.~\eqref{eq:t3-assembly-bound},
\begin{equation}
\label{eq:t3-assembly-bound}
  \|F^{\rm dyn}_{A,h}(Z_G;y)\|_{\rm proof}
  \le \Gamma_{\rm D5}\sum_{j=1}^{162}\|\tau_{j,h}(Z_G;y)\|
\end{equation}
to obtain a dynamic-block residual certificate.  Only after this dynamic-block
certificate is assembled with the same 96-row non-dynamic certificate, row
scaling, proof norm, endpoint map, and branch rule could it replace the direct
P5 dynamic-block certificate in the T2 full-map Taylor/Kantorovich proof.
The finite D5 inventory is the D5 term-budget record for this replacement
certificate: it records 162 Taylor subterms with zero term-level bounds
currently certified on the separate primitive/Taylor route, and the present
theorem therefore asserts no term-level \(O(h^7)\) bounds for those subterms.
The tuple is therefore ordered like the Wieloch--Arnold constrained-BDF proof:
fix the discrete object, insert the smooth branch, prove local primitive
defects in that object, assemble them, and only then enter the residual-to-root
and global-error steps.  It is not allowed to use
Lemma~\ref{lem:stage-residual-defect}, the direct D5 zero-row certificate, the
already assembled \(132\)-row residual estimate, a reported grid error, a
finite Newton log, or a source-policy row as an input to the primitive lifts.
This certificate template is not discharged here: the available status is
\(162/162\) conditional inventory coverage and \(0/162\) actual primitive
Taylor bounds for the separate T3 route.
\end{definition}

\begin{lemma}[T3 replacement-certificate use rule]
\label{lem:t3-replacement-certificate-use-rule}
Fix one theorem instance and the same compact branch, proof norm, row scaling,
endpoint map, lifted Gauss stage \(Z_G\), and implemented residual
\(F_{A,h}\) as in Theorem~\ref{thm:g6fullva-order}.  If a completed T3
primitive-route certificate satisfying
Definition~\ref{def:t3-primitive-route-certificate-template} were available,
its only admissible theorem-level use would be to replace the direct P5
dynamic-block certificate before the full \(132\)-row residual bridge is
formed.  It could then be combined with the same \(96\)-row non-dynamic
certificate, branch rule, row scaling, and proof norm to produce a new
same-tuple residual-value bound for \(F_{A,h}(Z_G;y)\).  It could not be
appended to the direct certificate, averaged with it, used to lower constants,
remove the P6 solver-scale envelope, prove a P7 residual-to-error boundary,
or bypass the endpoint, inexact-Newton, or local-to-global steps.  Until such
a T3 tuple is proved, the primitive inventory supplies no theorem input.
Equivalently, the active residual bridge and any future T3 certificate have a
one-slot replacement type: they may produce the same full-residual input, but
neither route is an inverse for the other route's internal hypotheses.
\end{lemma}

\begin{proof}
The definition of a T3 certificate has the same fixed object as the accepted
FullVA theorem instance and its output is only a dynamic-block residual
certificate for the Newton--Euler rows evaluated at \(Z_G\).  The theorem-level
residual bridge has a different type: it requires one dynamic-block
certificate and the existing \(96\)-row non-dynamic certificate to obtain a
single full residual-value estimate for the row-scaled map \(F_{A,h}\).  Thus
a completed T3 route would have the same input type as the direct P5 dynamic
certificate and could occupy the same slot, but it would not create a second
local-defect summand.

After that slot is filled, Lemma~\ref{lem:stage-residual-defect} still
requires the full-map inverse and Taylor/Kantorovich hypotheses on the same
branch, and the proof still requires the endpoint-closure, inexact-Newton,
and local-to-global transfer lemmas.  None of those hypotheses is a dynamic
primitive subterm bound.  The single residual-certificate slot in
Eq.~\eqref{eq:residual-certificate-slot} therefore makes the two routes
alternative provenance for the same full residual input, not cumulative
evidence.  Since the present manuscript records no completed T3 tuple and no
actual primitive Taylor bounds, the primitive inventory remains a conditional
replacement route only.
The non-inference used here can be written as
\begin{equation}
\label{eq:t3-no-reverse-inference}
  \begin{aligned}
  \mathsf R_{\rm dir}(Z_G;y)
  &:=
  \bigl(\|F_{A,h}(Z_G;y)\|\le C_Rh^7\bigr)
  \ \hbox{or}\ \bigl(\mathsf D_{36}(Z_G)=0\bigr),\\
  \mathcal P_{\rm T3}
  &:=
  \bigl(P_{\mathrm{state}},P_{\mathrm{acc}},P_{\lambda},
  P_{\mathrm{geom}},P_{\mathrm{gyro}}\bigr)
  \wedge
  \{\|\tau_{j,h}(Z_G;y)\|\le C_jh^7\}_{j=1}^{162},\\
  \mathsf R_{\rm dir}(Z_G;y)&\not\Longrightarrow \mathcal P_{\rm T3}.
  \end{aligned}
\end{equation}
Equation~\eqref{eq:t3-no-reverse-inference} is not an additional assumption;
it is the type separation just proved.  A full residual-value estimate is the
input required by the T2 full-map perturbation lemma, whereas the displayed
primitive lift and subterm bounds are sufficient antecedents for one possible
future T3 proof of the dynamic-block part of that input.
\end{proof}

\noindent\textbf{Primitive-route obligation ladder.}
If the separate T3 route is pursued, its proof obligations have the same
forward-only ordering discipline as the Wieloch--Arnold constrained-BDF proof: define
the fixed discrete object, insert the smooth branch, prove local defects in
that object, and only then propagate them through the error recursion.  For the
present FullVA dynamic block the corresponding implication chain is
Eq.~\eqref{eq:primitive-route-obligation-ladder}:
\begin{equation}
\label{eq:primitive-route-obligation-ladder}
\begin{aligned}
  (P_{\mathrm{state}},P_{\mathrm{acc}})
  &\Longrightarrow P_{\lambda}
  \Longrightarrow (P_{\mathrm{geom}},P_{\mathrm{gyro}})\\
  &\Longrightarrow
  \{\tau_{j,h}(Z_G;y)=O(h^7)\}_{j=1}^{162}\\
  &\Longrightarrow
  \|F^{\mathrm{dyn}}_{A,h}(Z_G;y)\|=O(h^7).
\end{aligned}
\end{equation}
All arrows in this ladder are one-way sufficient implications on the same
compact branch, with the same stage variables, row scaling, and proof norm as
the theorem-level residual map.  The last arrow is the finite-sum implication
proved below; the earlier arrows are exactly the primitive lift and reduction
obligations listed as open in the T3 status map.  Therefore a primitive-route
proof for that route cannot start from the assembled 132-row residual, from the direct D5
zero-row certificate, from a finite grid order, or from the Wieloch--Arnold
theorem and then infer the primitive lifts backward.  Conversely, once all
arrows in this ladder were proved, the output would still be only a sufficient
dynamic-block residual certificate to be inserted into the existing T2
full-map perturbation chain.

\noindent\textbf{Taylor accounting rule.}
The useful methodological constraint inherited from the reference proof is
not merely that a Taylor expansion appears, but that the scheme, variables,
comparison solution, and constants are fixed before the expansion is used.
In the present proof, after \(K\), the reduced chart, the row-scaled residual
\(F_{A,h}\), the proof norm, the Gauss-predictor branch, and the endpoint map
are fixed, all Taylor constants in the load-bearing lemmas are compact-tube
constants of that fixed object.  They are not calibrated from \(h\)-sweeps,
Newton logs, source-policy rows, or mechanism residual diagnostics.  The
Taylor variable for T2 is \(\Delta=Z-Z_G\), and the D5 direct certificate is a
base-point residual identity at \(Z_G\): the zero Newton--Euler rows enter
before the full-map Taylor expansion, not as consequences of the stage-root
perturbation.  By contrast, a primitive D5 Taylor proof would have to prove
the \(P_{\mathrm{state}}\), \(P_{\mathrm{acc}}\), and \(P_{\lambda}\) lift
bounds and their downstream geometry/gyroscopic reductions independently
before invoking the theorem-level local-defect result; it cannot use the T2
conclusion to prove those primitives.  This is the anti-circularity rule that
keeps the reference-style Taylor discipline strong while keeping the separate
primitive route separate from the accepted direct route.

\noindent\textbf{Full-map Taylor sufficiency.}
The load-bearing Taylor estimate is a statement about the differentiable
family \(F_{A,h}:Z\mapsto\mathbb R^{132}\) after row scaling and proof-norm
selection, not about a coordinate projection onto primitive symbolic
subterms.  The finite-dimensional residual inverse in
Lemma~\ref{lem:stage-residual-defect} maps the full bridged residual value
\(R_h=F_{A,h}(Z_G;y)\) to the stage increment on the same branch.  Therefore
the proof needs a uniform full-map inverse, a uniform full-map second
derivative bound, and the bridged value \(\|R_h\|\le C_Rh^7\); it does not
need a right inverse from the 132 implemented rows to the 162 primitive
Newton--Euler subterms.  Such a primitive projection would be a
strictly stronger verification route.  It is deliberately not part of the theorem
instance, and no theorem conclusion is used to infer it.
The direct D5 zero-row identity is consumed only in the base-point value
\(R_h\): it supplies the dynamic block of
\(F_{A,h}(Z_G;y)\) before the Taylor expansion is taken.  It is not used as a
derivative estimate for \(D_ZF_{A,h}\), not used as a bound on
\(D_Z^2F_{A,h}\), and not used to choose the Kantorovich ball.  Those
quantities remain full-map compact-tube inputs of
Lemma~\ref{lem:stage-residual-defect}, fixed before the direct residual
bridge is inserted.

This gives the active Taylor sufficiency statement used by the theorem:
once the fixed scaled map \(F_{A,h}\) has uniform \(J_h^{-1}\) and
\(D_Z^2F_{A,h}\) bounds on the compact branch, and once the direct bridge
gives \(\|F_{A,h}(Z_G;y)\|\le C_Rh^7\), Taylor's formula for the full map is
enough to produce \(\|Z_A-Z_G\|=O(h^7)\).  A term-by-term primitive proof of
the 162 Newton--Euler subterms would be a different way to certify part of
the same base-point residual value; it is not an additional hypothesis of this
full-map Taylor argument.

\begin{lemma}[Conditional finite-sum Taylor implication for the D5 route]
\label{lem:d5-conditional-taylor-finite-sum}
Let \(\tau_{j,h}(Z_G;y)\), \(1\le j\le 162\), denote the primitive
Newton--Euler Taylor subterms recorded in the D5 primitive-route
specification, grouped by
the fixed row ordering into the 36 dynamic residual rows.  Suppose that, on
the same compact branch, with the same row scaling and proof norm used for
\(F_{A,h}\), the five primitive lift obligations in
Eq.~\eqref{eq:d5-primitive-open-inputs}
\begin{equation}
\label{eq:d5-primitive-open-inputs}
P_{\mathrm{state}},\quad P_{\mathrm{acc}},\quad P_{\lambda},\quad
P_{\mathrm{geom}},\quad P_{\mathrm{gyro}}
\end{equation}
are proved with constants independent of \(h\), and that each recorded
primitive subterm satisfies Eq.~\eqref{eq:d5-primitive-subterm-bound},
\begin{equation}
\label{eq:d5-primitive-subterm-bound}
  \|\tau_{j,h}(Z_G;y)\|\le C_jh^7,\qquad 1\le j\le 162 .
\end{equation}
Then the D5 dynamic block assembled from those primitive subterms satisfies
Eq.~\eqref{eq:d5-primitive-dynamic-block-bound},
\begin{equation}
\label{eq:d5-primitive-dynamic-block-bound}
  \|F^{\mathrm{dyn}}_{A,h}(Z_G;y)\|
  \le C_{\mathrm{D5},\mathrm{prim}}h^7
\end{equation}
with \(C_{\mathrm{D5},\mathrm{prim}}\) depending only on the finite row
assembly operators, the fixed norm-equivalence constants, and
\(\sum_{j=1}^{162} C_j\).  This is a sufficient primitive-route implication
for the dynamic-row part of the theorem's stage-residual input.  It is not an inverse
projection from the already assembled \(132\)-row residual estimate back to
primitive subterms, and it is not additive proof evidence once the direct D5
zero-row certificate has already supplied the same dynamic block.
\end{lemma}

\begin{proof}
Each \(\tau_{j,h}\) is already an expanded primitive Taylor summand: product
and bilinear Taylor reductions, multiplier-geometry reductions, and
gyroscopic mean-value terms have been assigned to the finite D5 inventory
before this lemma is invoked.  After the branch, row ordering, row scaling,
and proof norm have been fixed, assembling those already-expanded summands
into the 36 dynamic rows is only a fixed finite row aggregation.  On a fixed
finite-dimensional row space, all block norms are equivalent; hence there is a constant
\(\Gamma_{\mathrm{D5}}\), independent of \(h\) and of the accepted step
index, such that the dynamic block assembled from the primitive vector
\(\tau_h=(\tau_{1,h},\ldots,\tau_{162,h})\) obeys
Eq.~\eqref{eq:d5-finite-aggregation-bound},
\begin{equation}
\label{eq:d5-finite-aggregation-bound}
  \|F^{\mathrm{dyn}}_{A,h}(Z_G;y)\|
  \le \Gamma_{\mathrm{D5}}\sum_{j=1}^{162}\|\tau_{j,h}(Z_G;y)\|.
\end{equation}
The assumed primitive bounds give
Eq.~\eqref{eq:d5-finite-sum-output},
\begin{equation}
\label{eq:d5-finite-sum-output}
  \|F^{\mathrm{dyn}}_{A,h}(Z_G;y)\|
  \le
  \Gamma_{\mathrm{D5}}\Big(\sum_{j=1}^{162}C_j\Big)h^7,
\end{equation}
so \(C_{\mathrm{D5},\mathrm{prim}}
=\Gamma_{\mathrm{D5}}\sum_{j=1}^{162}C_j\) proves the displayed dynamic-block
bound.  The argument is one-way: a small assembled residual does not identify
or bound each primitive subterm, because the finite assembly may contain
cancellations or many-to-one projections.  Therefore this lemma records the
conditional primitive-route implication only.  In the accepted theorem the
dynamic block is already supplied by the direct D5 substitution certificate,
so this conditional primitive implication is not combined with that direct
certificate to strengthen \(C_R\) or to remove any retained interface.
\end{proof}

\begin{lemma}[Taylor-route separation]
\label{lem:taylor-route-separation}
For Theorem~\ref{thm:g6fullva-order}, the local Taylor object is the
implemented scaled residual map \(F_{A,h}\) on the full stage variable
\(Z\), expanded about the lifted smooth Gauss stage \(Z_G\).  In the theorem
invocation this Taylor object is used only after
Lemma~\ref{lem:full-132-row-residual-bridge} supplies the full-residual
estimate in Eq.~\eqref{eq:taylor-route-full-residual-input},
\begin{equation}
\label{eq:taylor-route-full-residual-input}
  \|F_{A,h}(Z_G;y)\|\le C_Rh^7
\end{equation}
on the accepted compact tube.  The separate primitive Newton--Euler Taylor
route is a different proof factorization: it would first prove uniform
\(O(h^7)\) lift bounds for \(P_{\mathrm{state}}\), \(P_{\mathrm{acc}}\), and
\(P_{\lambda}\), then use those bounds to close \(P_{\mathrm{geom}}\) and
\(P_{\mathrm{gyro}}\), and only after that certify the 162 primitive dynamic
subterms.  Consequently, once
Lemma~\ref{lem:full-132-row-residual-bridge} independently supplies the full
residual estimate, an open primitive lift obligation cannot invalidate the
theorem-level Taylor/Kantorovich route; conversely the theorem-level route
does not certify any primitive subterm bound.  The only projection used by
the theorem is the fixed finite-dimensional block projection separating the
96 non-dynamic rows from the 36 Newton--Euler rows before the full residual
norm is formed; it is not a projection from the full residual estimate back
onto the 162 primitive symbolic terms.
Any primitive-route proof for that route would therefore be a replacement proof
for the dynamic-block portion of the same full-residual
hypothesis, not an additional theorem layer or a shortcut around the full-map
perturbation lemma.
\end{lemma}

\begin{proof}
The reference BLieDF proof first fixes the reconstruction and discrete
equations, inserts the smooth solution into that fixed scheme, estimates the
local truncation defect, and only then derives the Lie-algebra/coupled error
recursion.  The analogue used here fixes the SO(3) retractions, row scaling,
FullVA stage variables, endpoint reconstruction, and accepted branch before
any Taylor estimate is invoked.  In the theorem invocation, after
Lemma~\ref{lem:full-132-row-residual-bridge} has supplied the \(C_Rh^7\)
residual input from the 96 non-dynamic rows and the 36 direct Newton--Euler
rows on the same branch, Lemma~\ref{lem:stage-residual-defect} applies
Taylor's formula with integral remainder to the single smooth map
\(F_{A,h}\) in the full increment \(\Delta=Z-Z_G\).
The inverse and quadratic-remainder constants are compact-tube constants of
this full map, so the stage-root perturbation uses only that full-residual
defect and the retained branch/tube hypotheses.
This inverse use is local-stage, not residual-to-error: it maps the
fixed \(132\)-row stage residual perturbation near \(Z_G\) to a nearby
stage root on the accepted Newton branch before endpoint reporting.  It does
not map reported mechanism residual histories, reaction residuals, or
constraint-monitor rows to trajectory errors.  Such an output-level transfer
would require the separate P7 fixed-operator theorem and an
\(h\)-independent constant \(C_{\rm P7}\) for the reported residual/error
objects.  Therefore closing the T2 full-map Taylor/Kantorovich step supplies
only the local \(C_Ah^7\) stage-root term in the four-term defect sum; it
does not prove or replace the P7 boundary theorem.
The contraction proof is therefore a full-map Taylor argument: it expands
\(F_{A,h}(Z_G+\Delta;y)\), inverts \(D_ZF_{A,h}(Z_G;y)\) on the accepted
branch, and controls the integral quadratic remainder in the same 132-row
proof norm.  It never projects the resulting stage increment back into
primitive subterm errors, and it never treats aggregate cancellation in the
132-row residual as evidence that a primitive term-by-term Taylor proof
is closed.

The primitive route rearranges only the proof of the 36 dynamic-row residual
value.  Its dependency graph starts from the stage state, acceleration, and
multiplier lifts; the geometry and gyroscopic reductions are downstream of
those lifts.  The primitive-route reduction records rules for the
162 subterms but not the required uniform lift bounds, so no primitive subterm
is used as a theorem input.  This is a route-separation statement rather than
a gap in the active theorem: the theorem consumes the full-residual bridge,
while the primitive route remains a separate finer implication record for one
input to that bridge.
\end{proof}

\begin{assumption}[Accepted-path regularity and perturbation scale]
\label{ass:regularity}
On the smooth cylindrical-chain benchmark over a fixed interval
$0\le t\le T$, the active lower-pair constraint Jacobian has constant full row
rank, the constrained mass matrix restricted to the tangent space is
nonsingular, and the index-3 FullVA equations admit a smooth reduced
Lie-group chart $y$ with a $C^7$ vector field $\dot y=f(y,t)$.  Write
\(K\) for the retained compact reduced-state proof tube on this branch; the
stage neighborhoods used below are contained in the corresponding compact
stage tube over \(K\).  This is notation for the retained theorem domain, not
a posteriori selection from the reported grids. The
implemented three-stage accepted residual is considered on the
Gauss-predictor branch neighborhood in the accepted compact proof tube, with
the stage Jacobian at the lifted Gauss predictor uniformly invertible and
with the local derivative bounds required by
Lemma~\ref{lem:stage-residual-defect}. The existence of the local accepted
root \(Z_A\) is then the conclusion of that Kantorovich perturbation lemma,
not an additional assumption here. This is a compact-tube inverse condition
on the Gauss-predictor branch; finite rank probes or pointwise solved-stage
Jacobian checks are not used as a compact-tube inverse proof.
Compact-inverse quantifier discipline: the compact tube, Gauss-predictor branch,
proof residual norm, and inverse constants are fixed before any \(h\)-sweep,
rank probe, solver log, or observed convergence row is inspected. Shrinking
\(h_0\) after these objects are fixed is allowed; enlarging the inverse
constant or changing the tube after diagnostics changes the retained P2
interface rather than proving it.
The inexact-Newton
residual-to-stage-error estimate is also part of this local branch contract:
it uses the strong local residual inverse stated in
Lemma~\ref{lem:inexact-newton}, equivalently uniform invertibility of the
averaged Jacobian along the segment from the accepted root to the computed
iterate under the small-neighborhood perturbation bound. This condition is
not inferred from pointwise Jacobian invertibility alone. The final accepted
one-step maps on this branch also satisfy the stability scale in
Eq.~\eqref{eq:assumption-stability-scale}:
\begin{equation}
\label{eq:assumption-stability-scale}
  \|\Psi_h^{\mathrm{G6FVA}}(y)-\Psi_h^{\mathrm{G6FVA}}(\bar y)\|
  \le (1+C_s h)\|y-\bar y\|
\end{equation}
for points in the compact proof tube and sufficiently small \(h\), with a
nonnegative stability constant \(C_s\ge0\) after enlarging the bound if
needed. This is the local-to-global stability input used below; ordinary
compact-tube Lipschitz continuity alone is not enough for the stated
fixed-time global order transfer. The exact accepted trajectory remains a
positive distance from the boundary of the same compact tube, and the
local-to-global transfer uses the tube-retention bootstrap in
Lemma~\ref{lem:local-global-transfer}.
This is a conditional compact-tube hypothesis, not an invariant-region proof
for remote nonlinear roots or residual-only mechanism rows. The
endpoint-closure map has a same-branch closed endpoint-anchor family on the
accepted branch, written locally as \(z_\ast=z_\ast(y,h)\) with
\(E(z_\ast)=0\), and satisfies the endpoint-closure local right inverse
condition, together with the endpoint-ball margin, uniformly in \((y,h)\) as in
Lemma~\ref{lem:endpoint-closure}; this is a compact-tube splitting assumption
and is not inferred from a single endpoint Jacobian right inverse without the
small-neighborhood perturbation bound. The smooth FullVA lift used in
Lemma~\ref{lem:fullva-stage-lift} exists on the same tube. The
endpoint-closure defect satisfies the hypothesis of
Lemma~\ref{lem:endpoint-closure}. Inexact Newton solves are treated by
Lemma~\ref{lem:inexact-newton}: for the asymptotic order statement the
accepted residual tolerance must scale no larger than $O(h^7)$. The fixed
tolerances used in the numerical section are reported separately as finite-run
solver data, not as an asymptotic solver-error proof. More precisely, the
theorem uses a per-step branch-selected residual bound
\(\eta_h^{\rm tube}\le c_\eta h^7\) for every branch-selected solve to
which the compact-tube one-step map and stability estimates are applied. On a
reported trajectory this policy is recorded as the maximum
\(\eta_h^{\rm reported}=\max_{0\le n<N_h}\eta_{h,n}\le c_\eta h^7\) over the reported
branch-selected Newton steps \(0\le n<N_h\), where the corresponding output
grid satisfies \(t_n=n h\), \(0\le n\le N_h\), and \(N_hh\le T\). It does
not assert that arbitrary Newton initializations, fixed tolerances, or remote
nonlinear roots select that branch; the theorem branch is initialized by the
Gauss predictor on the same accepted branch. Dense and sparse backends
evaluate the same residual and Jacobian entries. The proof row-scaling map
and residual norm are fixed before \(h\), \(N_h\), backend choice, and
diagnostic logs; they introduce no hidden \(h\)-dependent row weights beyond
the explicit \(h\)-factors written in the residual formulas.
\noindent\textbf{Backend-norm compatibility rule.}
A backend termination record is theorem-grade P6 evidence only after a
predeclared \(h\)-uniform conversion to the proof residual norm of the same
row-scaled map \(F_{A,h}\) has been fixed on the accepted branch. In
particular, if a backend reports a norm or proxy
\(\eta_{h,n}^{\rm be}\), the proof may consume it only through an
\(h\)-independent constant \(C_{\rm be}\) satisfying
Eq.~\eqref{eq:backend-proof-norm-compatibility}:
\begin{equation}
\label{eq:backend-proof-norm-compatibility}
  \|F_{A,h}(\widetilde Z_{A,n};y_n)\|_{\rm proof}
  \le C_{\rm be}\eta_{h,n}^{\rm be}
\end{equation}
for the same branch-selected iterate, residual rows, and row scaling used by
Lemma~\ref{lem:inexact-newton}. Without this pre-run compatibility
certificate, backend norms, relative residuals, correction norms, and mixed
scaled/unscaled residual summaries remain diagnostics; they cannot supply
\(\eta_h^{\rm tube}\), choose \(c_\eta\), or absorb a solver term into
\(C_{\rm loc}\).
The reported maximum over accepted steps is theorem-grade only after this
backend-to-proof-norm conversion and the compact-branch solver envelope have
both been fixed before the run; taking a maximum over finite solver logs
afterward records a consistency diagnostic, but it does not create the
uniform P6 envelope required by Lemma~\ref{lem:inexact-newton}.
The theorem retains this as a uniform solver-scale hypothesis. Fixed-tolerance
runs and backend stopping records provide finite diagnostics only unless they
are prebound to the proof norm and the \(h^7\) envelope on the same compact
branch.
\end{assumption}

\paragraph{Discharged residual input, not an assumption}
The implementation-specific residual step used by the theorem is supplied
after Assumption~\ref{ass:regularity} has fixed the branch, tube, chart,
proof norm, row ordering, and row scaling. Lemma~\ref{lem:full-132-row-residual-bridge}
then gives the discharged pre-perturbation residual-defect input
\(\|F_{A,h}(Z_G;y)\|\le C_Rh^7\) for the implemented accepted residual.
The implementation part of the theorem is not a
compact-tube regularity assumption; it is the 96-row non-dynamic certificate
together with the D5 direct-substitution certificate for the 36
Newton--Euler rows. After the mathematical row identities are expanded into
the implemented row ordering, the lifted Gauss stage has the defect bound in
Eq.~\eqref{eq:discharged-132-row-residual-input}:
\begin{equation}
\label{eq:discharged-132-row-residual-input}
  \begin{aligned}
  \|F_{A,h}(Z_G;y)\|
  &\le C_{\rm blk}\|F_{\rm nd,h}(Z_G;y)\|
     + C_{\rm blk}\|F_{\rm dyn,h}(Z_G;y)\| \\
  &\le C_R h^7,
  \qquad F_{\rm dyn,h}(Z_G;y)=0 .
  \end{aligned}
\end{equation}
This displayed handoff says that the lifted Gauss stage has a defect bounded by
$C_R h^7$; it is the point where the proved 96-row non-dynamic certificate and
the P5 direct Newton--Euler certificate enter the proof. It is a one-step
accepted residual defect statement only: a
consumed row-local certificate on \(Z_G\), not a compact-tube regularity
hypothesis, not an empirical residual fit, and not a promotion of mechanism
residual-to-error rows, source-policy rows, or fixed-tolerance production
residuals into theorem interfaces.

\paragraph{Proof skeleton}
The theorem is proved by a perturbation chain. The lemmas and final proof carry
the estimates; the tables below only make the interfaces and nonuse rules
explicit. First, the classical three-stage Gauss
construction on the smooth
reduced chart gives a local reduced-flow truncation term \(C_Gh^7\). Second,
Lemma~\ref{lem:fullva-stage-lift} lifts the reduced Gauss stage to the
FullVA variables, and Lemma~\ref{lem:stage-residual-defect} consumes the
proved 132-row residual-defect input: the 96 non-dynamic rows and
the 36 Newton--Euler rows give \(\|F_{A,h}(Z_G;y)\|\le C_Rh^7\), hence the
accepted local stage root changes the endpoint by at most \(C_Ah^7\). Third,
Lemma~\ref{lem:endpoint-closure} bounds the endpoint closure correction by
\(C_Eh^7\), with the closed endpoint-anchor family, endpoint-ball constants,
and local right inverse kept explicit. Fourth, Lemma~\ref{lem:inexact-newton} converts the
branch-selected solver residual into a perturbation of the full endpoint-output
map \(\mathcal P_h=\mathcal C_h\circ\mathcal E_h\), bounded by
\(C_N\eta_h^{\rm tube}\), and the theorem assumes
\(\eta_h^{\rm tube}\le c_\eta h^7\), so this term is \(C_Nc_\eta h^7\).
Finally, Lemma~\ref{lem:local-global-transfer} applies the \(1+C_sh\)
stability scale and tube-retention bootstrap, giving the fixed-time
reduced-chart bound with the constant in
Eq.~\eqref{eq:proof-skeleton-global-constant}:
\begin{equation}
\label{eq:proof-skeleton-global-constant}
  C_{\rm red}=(C_G+C_A+C_E+C_Nc_\eta)\Gamma_s(T).
\end{equation}
Lemma~\ref{lem:reporting-map-norm-consequence} then converts this
reduced-chart estimate to the position--velocity quantities reported in the
numerical section by a compact-tube chart-to-output derivative bound. The
proof-interface tables that follow do not add another proof source.  They
only locate the assumptions, the consumed row-local certificates, and the
diagnostic rows that are deliberately excluded from the theorem input set.
The load-bearing residual-defect step is therefore upstream of all empirical
trajectory evidence: the implemented 132-row residual is evaluated at the
lifted exact Gauss stage \(Z_G\), before the accepted stage root, endpoint
closure, inexact Newton iterate, local-to-global transfer, and reporting-map
conclusion are invoked. Observed \(h\)-sweeps, residual/reaction tables, and
source-policy rows cannot be used to justify this stage-local certificate.

\paragraph{Constant hierarchy for the local-defect chain}
The constants used in Theorem~\ref{thm:g6fullva-order} are fixed in the same
order as the proof objects.  The row-local non-dynamic certificate gives
\(C_{\rm nd}\).  The fixed finite-dimensional block norm and row scaling give
\(C_{\rm blk}\), hence \(C_R=C_{\rm blk}C_{\rm nd}\) after the D5 direct
Newton--Euler rows give \(F_{\rm dyn,h}(Z_G;y)=0\).  The compact-tube inverse
constants in Lemma~\ref{lem:stage-residual-defect} then give
\(C_Z=2MC_R\) and \(C_A=M_EC_Z\).  Lemma~\ref{lem:endpoint-closure} supplies
the endpoint-closure constant \(C_E\), and Lemma~\ref{lem:inexact-newton}
supplies \(C_N=2M_AM_N\).  Therefore the one-step constant consumed by the
global theorem is Eq.~\eqref{eq:local-defect-chain-constant}:
\begin{equation}
\label{eq:local-defect-chain-constant}
  C_{\rm loc}=C_G+C_A+C_E+C_Nc_\eta ,
\end{equation}
where \(C_G\) is the classical three-stage Gauss local-defect constant on the
same smooth reduced flow.  Finally Eq.~\eqref{eq:reported-output-constant}
records the constants after the global stability and reporting-map steps:
\begin{equation}
\label{eq:reported-output-constant}
  C_{\rm red}=C_{\rm loc}\Gamma_s(T),\qquad
  C_{qv}=C_{\mathcal R}C_{\rm red},
\end{equation}
where \(C_{\mathcal R}\) is the compact-tube Lipschitz constant of the
reported \((q,v)\) map.  This hierarchy is not an additional assumption: it
records the proof's quantifier order.  Each constant may depend on the fixed
compact tube, chart, row scaling, inverse margins, endpoint right-inverse,
stability factor, and \(c_\eta\), but none is selected from fitted slopes,
finite residual logs, separate primitive-route status,
source-policy rows, backend timing, or reproducibility package state.
In particular, \(C_R\) is fixed before the inexact-Newton term is introduced:
P6 enters only through the separate \(C_Nc_\eta\) contribution to
\(C_{\rm loc}\), and P7 residual-only mechanism diagnostics do not enter
either constant.

Same-object proof invariant: every estimate in the local-defect chain is
taken for the same transition index, same accepted compact-tube branch, same
reduced chart, same implemented 132-row residual, and same fixed proof norm.
The Gauss truncation term, residual-to-stage perturbation, endpoint closure
term, and inexact-Newton term are \mbox{not assembled from different runs}, different
row scalings, or different Newton roots. This invariant is what permits the
four \(O(h^7)\) terms to be added in one triangle inequality before the
local-to-global estimate is invoked.

\paragraph{Order-accounting checkpoint}
The theorem uses the standard one-step/global order accounting: the accepted
branch-selected map has a local defect bounded by \(C_{\rm loc}h^7\), and
the same-initial-state reduced-chart reported-grid estimate is
\(C_{\rm red}h^6\) because \(N_h=O(h^{-1})\) transitions are accumulated
through the stability recurrence.  The reporting-map estimate for \((q,v)\)
has the same reported-grid \(h^6\) order.  Therefore the word sixth-order in
the theorem refers to the same-initial-state reported-grid \(h^6\) bound,
while the \(h^7\) quantities are local-defect and residual-envelope inputs.
The observed smooth-window slopes
7.161/7.066 are downstream consistency evidence for the selected branch; they
are not a seventh-order theorem, not a global \(O(h^7)\) error bound, and not
an additional solver-policy discharge.  A stronger global-order statement
would require a new theorem with its own local defect, stability, solver, and
reporting-map estimates, not just larger finite-window fitted slopes.

\paragraph{Analytic inputs versus numerical evidence}
The theorem is therefore not an empirical order fit. Before any numerical
\(h\)-sweep is cited, each load-bearing estimate in the proof is quantified
uniformly on the accepted compact proof tube: the Gauss local defect, the
132-row FullVA residual defect on \(Z_G\) together with the local stage-root
inverse, the endpoint-closure perturbation, the inexact Newton tolerance
scale, and the \(1+C_sh\) stability/tube-retention transfer. Numerical
dynamic-order rows are used only after this analytic chain, as consistency
evidence for the smooth accepted path. Mechanism-coverage rows verify that
the same FullVA lower-pair vocabulary, reaction reconstruction, and closure
checks are exercised on closed-loop examples, but they do not replace a
residual-to-error transfer theorem. Diagnostic rows, source-policy rows, finite
solver probes, and the local \tfe{} comparator likewise cannot be inserted
as analytic inputs. If one analytic input is not discharged, the corresponding
statement remains a retained condition; it is not promoted by a small
residual table or by an observed order fit.

\paragraph{Theorem status}
The word conditional is used in the standard numerical-analysis sense: the
result is a full local-defect and fixed-time global-error theorem once the
listed compact-branch, implementation-identity, and solver-scale hypotheses
are in force. It is not a weak empirical claim. The proof contribution in the
present manuscript is to make the 132-row \method{} residual enter that
theorem through a square FullVA stage-root perturbation, including the
Newton--Euler rows by the P5 direct route, and then to carry endpoint closure
and inexact Newton as explicit \(O(h^7)\) perturbations before the
local-to-global transfer. What remains retained is the automatic satisfaction
of the compact-tube and branch-solver hypotheses for every possible
initialization or fixed-tolerance run; these are domain-of-validity
conditions, not substitutes for the theorem's error estimate.
Thus the conditional qualifier is attached to the domain on which the
branch-selected map is defined and stable.  It is not attached to the
implemented 132-row residual-defect input itself: that input is supplied by
the certified non-dynamic rows and the D5 direct Newton--Euler route before
the endpoint, Newton, and global-transfer estimates are applied.
The resulting theorem is therefore conditional but not evidentiary: retained
interfaces are quantified domain and solver-scale hypotheses with fixed
constants, while discharged interfaces are algebraic inputs proved in the
manuscript.  The conclusion follows for every step size in that fixed domain;
it is not a finite-sweep acceptance rule and it is not weakened by keeping
residual-to-error and source-policy promotion outside the theorem.

\paragraph{Proof-scope invariant}
Throughout the proof section, ``proved'' refers only to a row-local or
algebraic certificate after the branch, norm, row order, and implemented
residual map have already been fixed.  ``Retained'' refers to a quantified
theorem interface that is assumed on that same compact branch; it is not
empirical evidence and it is not treated as discharged.  Thus the P4 binding
interface is retained, while P4 also supplies a proved 96-row non-dynamic
certificate.  The theorem can consume that certificate only after P1--P3 have
fixed the branch, norm, row ordering, and implementation binding, so the P4
theorem interface is not a free-standing
compact-tube discharge.  P5 is the only direct-route dynamic-row interface
discharged at the theorem-input level.  The P7 residual-to-error boundary,
residual-only mechanism rows, and
source-policy rows are neither theorem inputs nor theorem conclusions here;
they are explicit scope boundaries.  Consequently an open primitive/Taylor or
P7 entry cannot weaken the direct residual-bridge estimate, and a closed row-local
certificate cannot be reused to prove P6, prove P7, establish source-policy
readiness, or establish a full-\tfe{} replacement unless a separate theorem
proves that implication.

\noindent
To make the proof interfaces explicit, Assumption~\ref{ass:regularity} is used through separate
interfaces rather than as a single opaque hypothesis.  The theorem-interface
list is P1--P6.  The final entry is displayed separately as the P7
residual-to-error boundary because it is an output nonclaim boundary, not a
theorem premise; the P7 shorthand is used only so the proof tables,
residual tables, and reproducibility records point to the same transfer-scope
rule.
For reader traceability, \mbox{P7 output nonclaim/residual-to-error boundary}
denotes this displayed output boundary and not an additional theorem premise.
\begin{enumerate}
  \item[P1.] smooth reduced-chart and FullVA-lift regularity on the accepted
  trajectory tube;
  \item[P2.] uniform local invertibility of the accepted stage residual
  Jacobian at the lifted Gauss predictor on the compact-tube branch,
  together with the local derivative bounds used to obtain the accepted root,
  the
  endpoint-closure local right inverse with a closed anchor and endpoint-ball
  margin, the strong local residual inverse used by
  Lemma~\ref{lem:inexact-newton}, and
  the one-step stability scale \(1+C_s h\) with tube-retention bootstrap;
  \item[P3.] row identity between the mathematical 132-row residual and the
  residual differentiated and solved by the implemented residual path;
  \item[P4.] two distinct subslots: the branch/norm/implementation-binding
  side is retained as a theorem interface, while the 96-row non-dynamic
  implementation-defect certificate on the lifted Gauss branch is proved and
  consumed only through the P1--P3 branch, norm, and implementation binding
  fixed for the theorem;
  \item[P5.] the discharged direct Newton--Euler implementation-defect
  certificate that supplies the 36 dynamic rows on the lifted Gauss branch;
  \item[P6.] the theorem-level nonlinear-solver scale
  \(\eta_h^{\rm tube}\le c_\eta h^7\) uniformly on the compact-tube
  branch-selected solver hypothesis, with the reported-grid maximum
  \(\eta_h^{\rm reported}=\max_{0\le n<N_h}\eta_{h,n}\) as its trajectory specialization,
  distinct from fixed-tolerance production runs.  The theorem retains this as
  a uniform solver-scale hypothesis; fixed-tolerance runs provide finite
  diagnostics only.  Changing the solver scale changes the theorem hypothesis
  and the local-defect constant, so this finite diagnostic status does not
  close P6;
  \item[Output boundary (P7 boundary).] a separate residual-to-error
  promotion condition for residual-only mechanism rows. P7 is not used to
  prove the accepted smooth-branch theorem below; it is kept as an explicit
  non-claim boundary for four-link,
  slider-crank, and source-policy rows.  No residual-to-error transfer theorem
  is assumed here; such a theorem would require a fixed reported residual/error
  operator on the same reported grid and an \(h\)-independent transfer
  constant, so the seven transfer obligations remain outside this theorem.
\end{enumerate}
The notation is a compact proof-domain partition, not a list of empirical
acceptance tests.  Interpreted as a theorem instance, P1/P2/P3, the P4 binding side, and
P6 are the retained admissibility interfaces: they fix the compact branch,
proof norm, local inverse, endpoint closure, stability scale, implementation
path, and solver scale on which the argument is allowed to run.  P5 is different: it is the
proved stage-local Newton--Euler identity on \(Z_G\), supplying the 36
dynamic rows by direct substitution.  The P4 notation has two theorem roles,
distinct from the four-layer theorem-output reading guide below:
the branch/norm/implementation-binding side is part of the retained admissibility
interface, while the 96-row \(O(h^7)\) non-dynamic residual estimate itself is
proved and discharged.  P4 supplies the proved 96-row non-dynamic part, but
only after the same branch, norm, and implemented
residual map have been fixed; it is not a standalone 132-row residual proof
without the P5 dynamic handoff.  The PC1--PC2 local residual route therefore
refers to this particular direct residual-defect route under the retained P6
solver-scale condition, while PC4 records the
residual-to-error exclusion.  The output boundary (P7 boundary) is a
declared transfer-scope clause: residual-only mechanism rows require a separate
residual-to-error theorem before they can become trajectory-order evidence.
The primitive dynamic symbolic oracle, the primitive/Taylor symbolic-defect
route, external same-test claims, and reproducibility package claims likewise remain
outside the theorem-level PC2 route until their own implications are proved.

With that admissibility map fixed, the proof is a one-way conditional
implication rather than a backfilled interpretation of numerical evidence.
The compact-tube hypotheses choose the branch and constants before \(h\); the
direct residual bridge supplies the stage-row defect on \(Z_G\); endpoint
closure and inexact Newton add controlled \(O(h^7)\) perturbations; and the
local-to-global transfer then gives the \(O(h^6)\) position--velocity grid
estimate.  Finite step sweeps may identify the reported branch and check
consistency, but they do not prove automatic branch selection for arbitrary
initializations, do not discharge the theorem-level solver hypothesis from fixed
production tolerances, and do not promote residual-only mechanism or
external comparison rows through a residual-to-error theorem.  Those retained
interfaces delimit where the theorem may be invoked; they are not hidden
residual terms and they are not inferred from the theorem conclusion.
Using the theorem-interface labels below:
P1/P2/P3, the retained P4 binding side, and P6 describe the branch,
implementation, and solver-scale interfaces under which the theorem is
invoked; P4 also contains the proved 96-row non-dynamic certificate; P5 is the proved stage-local Newton--Euler identity;
the direct residual bridge is used under the retained solver-scale condition,
and the residual-to-error clause records only a transfer-scope exclusion; and the
primitive dynamic symbolic oracle and primitive/Taylor symbolic-defect route
remain outside the active route. The implication direction is fixed
throughout, so no stage-row, solver-scale, or residual-to-error premise is
inferred backward from observed error.

\noindent\textbf{Theorem invocation criterion.}
Theorem~\ref{thm:g6fullva-order}
may be invoked only after retained P1/P2/P3 interfaces, the P4 binding
interface, and the P6 solver-scale interface are fixed and the discharged P5
direct-route certificate is bound on the same accepted
compact branch. P1--P2 supply the smooth FullVA lift, compact-tube inverse and
derivative bounds, endpoint-closure, and stability/tube-retention hypotheses;
P3 binds the implemented residual path, row ordering, row scaling, and proof
norm; P4 contributes the proved 96-row non-dynamic certificate only under that
binding; P5 contributes the 36 Newton--Euler rows through the direct
substitution route; and P6 supplies the
branch-selected solver scale \(\eta_h^{\rm tube}\le c_\eta h^7\). The theorem
output is only the reported-grid order-six position--velocity error on that
same branch. It does not prove a solver-policy theorem, a P7
residual-to-error transfer theorem, source-policy reproduction, closed-loop
residual promotion, or a full-\tfe{} stage replacement. Therefore numerical
support rows may corroborate the accepted branch after the theorem has been
invoked, but they cannot backfill a separate invocation condition.

\noindent\textbf{Retained/discharged input rule.}
The retained domain
hypotheses are the compact-branch regularity, endpoint, stability, inverse,
implementation-binding, and P6 solver-scale interfaces; the certified
stage-residual input is assembled only after those retained interfaces are
fixed, from the 96-row P4 non-dynamic certificate and the discharged P5
direct-substitution certificate for the 36 Newton--Euler rows. Thus the
theorem does not merely assume the accepted implementation residual defect as
an unproved premise; it consumes the implemented direct-route certificate for
that input while retaining the compact-tube and solver-scale interfaces.
Thus the direct PC2 residual-defect input is discharged under the retained
P1/P2/P3 interfaces, the P4 binding interface, and the P6 solver-scale
interface: the exact Gauss defect, the 96-row non-dynamic lift, the 36-row
Newton--Euler direct substitution, endpoint closure, inexact Newton
perturbation, and local-to-global transfer are chained in this order. The
theorem remains conditional on those domain and solver interfaces; it does
not complete the primitive/Taylor route, prove a P7 residual-to-error boundary,
or establish source-policy or full-\tfe{} readiness.

\noindent\textbf{Assumption admissibility note.}
Assumption~\ref{ass:regularity}
does not assume the theorem conclusion.  In particular, P2 supplies a
compact-tube stability scale and tube-retention bootstrap for the
branch-selected map; it is independent of, and logically prior to, the
discrete Gronwall step.  It is \mbox{not an assumed local defect estimate}, not an
assumption that \(\Psi_h^{\mathrm{G6FVA}}\) is already \(O(h^7)\)-close to
the exact flow, and not inferred from observed sixth-order slopes,
residual/reaction tables, or successful Newton logs.  The local \(O(h^7)\)
defect is proved later from Gauss truncation, the certified 132-row
stage-root perturbation, endpoint closure, and the retained P6 inexact-Newton
scale.  The stability/tube-retention input only permits that one-step defect
to be iterated on the same compact tube; it does not choose \(C_{\rm loc}\),
does not prove a theorem-level solver-policy result, and does not prove a
P7 residual-to-error boundary.

\noindent\textbf{Non-vacuity/admissibility note.}
Assumption admissibility note: Assumption~\ref{ass:regularity} and the
retained solver-scale hypotheses define the theorem domain rather than a
shortcut to the theorem conclusion.
The retained
P1--P3 conditions, P4 binding condition, and P6 solver-scale condition define
a local same-branch domain of applicability, not the error estimate itself.
They require a smooth compact tube,
Gauss-predictor branch selection, local inverse/endpoint/stability structure,
implementation-path binding, and a branch-solver residual scale fixed before
the asymptotic limit. They do not assume the \(O(h^7)\) local defect, the
\(O(h^6)\) grid estimate, or the observed position--velocity orders. Once
that domain is fixed, the proof still derives the local \(h^7\) defect from
Gauss truncation, the proved 132-row direct residual bridge, endpoint
closure, and the retained P6 solver-scale term, and then derives the global
grid bound by the stability/tube-retention transfer. Finite numerical rows may
identify the reported branch and check consistency, but they are not part of
this admissibility claim and do not turn the conditional theorem into an
unconditional source-policy result.

\begin{table}[htbp]
\centering
\caption{Assumptions and theorem scope.}
\label{tab:p-interface-satisfaction-ledger}
\scriptsize
\setlength{\tabcolsep}{2pt}
\renewcommand{\arraystretch}{0.92}
\begin{tabular}{P{0.10\linewidth}P{0.17\linewidth}P{0.47\linewidth}P{0.18\linewidth}}
\toprule
Interface & Theorem role & Meaning & Explicit exclusion \\
\midrule
P1--P2 & Domain interface &
Smooth FullVA lift, compact-tube inverse, endpoint-ball/right-inverse, and
tube-retention stability remain theorem interfaces. &
Finite sweeps, rank probes, and endpoint checks diagnose the selected branch
only. \\
P3--P4 & Implementation interface; P4 row certificate proved &
Runtime formula/AD binding identifies the implemented direct residual path,
and P4's row-local certificate supplies the proved 96-row non-dynamic part of
PC2. &
Primitive symbolic oracle completion remains a separate verification route, not an input
to the active direct-route 132-row bridge; P3/P4 also require the separate P5
dynamic-row handoff. \\
P5 & Stage-local identity proved &
The D5 direct-substitution certificate proves the 36 Newton--Euler residual
rows vanish on \(Z_G\). &
This is stage-local and does not promote finite trajectories, source-policy
rows, or residual-to-error claims. \\
P6 & Solver-scale condition &
The theorem uses the compact-tube branch-solver hypothesis
\(\eta_h^{\rm tube}\le c_\eta h^7\). &
Finite solver probes do not prove a theorem-level solver-policy theorem or
turn fixed tolerances into an asymptotic proof. \\
P7 boundary & Residual-output nonclaim &
Residual-to-error promotion for mechanism-coverage and source-policy rows
requires a separate theorem. &
Residual tables remain diagnostics until all seven residual-to-error
obligations close. \\
Stage-residual bridge / solver-scale hypothesis / residual exclusion &
The direct bridge supplies the local residual route under the retained
solver-scale hypothesis; residual-to-error promotion is not used. &
This separation does not promote theorem-domain interfaces into proved theorem
inputs. \\
\bottomrule
\end{tabular}
\end{table}

\noindent\textbf{Proof-domain partition.}
The Proof-domain partition is therefore explicit: P5 is the only stage-local dynamic identity proved by
direct substitution.  P1/P2/P3, the retained P4 binding side, and P6 describe
the branch, implementation, and solver-scale interfaces on which the theorem
is invoked, with P4's row-local certificate proving the 96 non-dynamic rows
used in the PC2 stage-residual condition.  P7 remains a separate
residual-to-error boundary rather than a local-defect theorem
premise.  PC1--PC2 supply the local residual route under the retained P6
solver-scale hypothesis, while PC4 records the residual-to-error exclusion. These
interfaces do not turn domain conditions or the P7 boundary into
a proved solver-policy theorem,
residual-to-error promotion, source-policy rows, or reproducibility package readiness
claims.  Thus PC2 is discharged only as a route-level stage-residual
statement: the
accepted stage-residual input is the 96-row non-dynamic certificate together
with the 36-row P5 direct Newton--Euler certificate. It is not a statement
that P4 by itself has become a full 132-row theorem input, and it is not a primitive
symbolic-oracle or source-paper residual-replacement closure.
No residual-to-error theorem is assumed in this route. The seven
residual-to-error obligations remain outside the local-defect proof, so small
mechanism residuals and reaction rows cannot enter PC2 or P6 as theorem
inputs.

\begin{table}[htbp]
\centering
\caption{Route and interface separation table for
Theorem~\ref{thm:g6fullva-order}.}
\label{tab:pcp-close-requirement-separation-ledger}
\scriptsize
\setlength{\tabcolsep}{2pt}
\renewcommand{\arraystretch}{0.9}
\begin{tabular}{P{0.12\linewidth}P{0.24\linewidth}P{0.27\linewidth}P{0.29\linewidth}}
\toprule
Route requirement & Satisfaction mode & What it supplies & What remains outside it \\
\midrule
PC1 balance closure &
D1--D2 weak-balance identities and row binding. &
The Newton--Euler balance identities used by the direct residual route. &
It does not prove compact-tube regularity, source-policy rows, or package
readiness. \\
PC2 stage-residual condition &
96-row \(O(h^7)\) non-dynamic certificate plus the 36-row P5
direct-substitution zero-residual certificate. &
The accepted 132-row stage-residual defect on the smooth Gauss lift. &
It does not complete the primitive symbolic route or prove residual-to-error transfer. \\
PC3 solver-scale retention &
P6 is retained as the theorem condition
\(\eta_h^{\rm tube}\le c_\eta h^7\). &
The inexact-Newton perturbation is absorbed into the local \(h^7\) defect. &
It does not prove a solver-policy theorem or convert finite tolerances into
an asymptotic proof. \\
PC4 residual demotion &
Residual-to-error promotion is not used for the accepted theorem. &
Residual-only four-link, slider-crank, and source-policy rows stay diagnostic. &
It does not prove a separate residual-to-error transfer theorem or promote
residual rows to trajectory-error evidence. \\
\bottomrule
\end{tabular}
\end{table}

\noindent\textbf{Route and interface separation.}
The separation table has a narrow mathematical role: it separates the PC1--PC2
local residual route together with the retained P6 solver-scale condition from
the PC4 residual-to-error exclusion.  PC3 is therefore a retained-condition
reference, not an empirical proof and not an extra certificate.  PC4 is
recorded only as a retained transfer-scope exclusion: residual-to-error
promotion is not made.  The term \mbox{PC4 residual-to-error exclusion}
denotes precisely this non-use of residual-to-error promotion.  These modes do
not prove P1/P2/P3 theorem interfaces,
the retained P4 binding interface, or the P6 solver-scale interface, do not
prove a solver-policy theorem, do not prove a separate residual-to-error
transfer theorem, and do not supply source-policy reproduction evidence or
reproducibility package readiness.

\begin{lemma}[Route-to-local-defect sufficiency on the same accepted branch]
\label{lem:route-local-defect-sufficiency}
Suppose P1/P2 fix the compact reduced-chart tube, P3/P4 fix the implemented
residual path and the 96-row non-dynamic certificate, P5 supplies the 36-row
Newton--Euler direct substitution, and P6 supplies the retained branch-solver
scale on the same accepted branch. Then the PC1--PC2 local residual route,
together with the retained P6 solver-scale hypothesis, supplies exactly the
local-defect input used by
Theorem~\ref{thm:g6fullva-order}: there is a compact-branch constant
\(C_{\rm loc}=C_G+C_A+C_E+C_Nc_\eta\), independent of \(h\), the reported
grid length, backend, finite solver logs, and source-policy rows, such that
Eq.~\eqref{eq:route-local-defect-bound} holds:
\begin{equation}
\label{eq:route-local-defect-bound}
  \|\Psi_h^{\mathrm{G6FVA}}(y)-\varphi_h(y)\|
  \le C_{\rm loc}h^7
\end{equation}
for every accepted branch-selected transition in the retained proof tube.
PC4 is not a local-defect ingredient; it is the retained transfer-scope exclusion
for residual-only mechanism and source-policy rows.
\end{lemma}

\begin{proof}
The proof uses the same proof-obligation ordering discipline as the reference
convergence analysis, not an estimate transfer: first insert the smooth branch
into the already fixed discrete equations, then estimate the one-step local
defect, then move that defect through the constrained algebraic solve and
endpoint map, and only afterwards apply the stability recursion.  The
variables and estimates differ from the BLieDF proof--this is a one-step
Gauss FullVA residual rather than a multistep hidden-constraint recursion--so
only the causal order of the proof is shared.
Under P1/P2, the smooth reduced constrained flow and the fixed compact tube
give the classical three-stage Gauss local defect in
Eq.~\eqref{eq:route-suff-gauss-defect}:
\begin{equation}
\label{eq:route-suff-gauss-defect}
  \|\Psi_h^G(y)-\varphi_h(y)\|\le C_Gh^7
\end{equation}
with \(C_G\) chosen from reduced-flow regularity before any implementation
diagnostic is read.  On the same transition define
\(\Psi_h^G=\mathcal E_h(Z_G)\),
\(\Psi_h^A=\mathcal E_h(Z_A)\),
\(\Psi_h^E=\mathcal C_h(\Psi_h^A)\), and
\(\Psi_h^{\mathrm{G6FVA}}=\mathcal P_h(\widetilde Z_A)\), with
\(\mathcal P_h=\mathcal C_h\circ\mathcal E_h\).  These maps use the same
state \(y\), branch, chart, endpoint-output map, and proof norm; no table row
or computed trajectory is used to identify them.
Endpoint-map identity rule: the same local closure map \(\mathcal C_h\),
endpoint norm, endpoint-anchor family, and endpoint ball are used in
\(\Psi_h^E=\mathcal C_h(\mathcal E_h(Z_A))\) and
\(\Psi_h^{\mathrm{G6FVA}}=\mathcal C_h(\mathcal E_h(\widetilde Z_A))\).
The inexact-Newton term therefore compares two outputs of one fixed map
\(\mathcal P_h\), not an unclosed exact endpoint against a closed computed
endpoint and not two separately fitted projection operators.  If the
endpoint correction rule, endpoint norm, or endpoint anchor were changed
between the exact-stage and inexact-stage terms, the four-term triangle would
refer to different one-step maps and would require a separate equivalence
lemma before the constants could be added.
The PC2 condition supplies the implemented residual estimate in
Eq.~\eqref{eq:route-suff-pc2-residual}:
\begin{equation}
\label{eq:route-suff-pc2-residual}
  \|F_{A,h}(Z_G;y)\|\le C_R h^7
\end{equation}
for the implemented 132-row residual on the lifted Gauss stage.  The
certificate combines the 96 non-dynamic rows identified by P3 and proved by
P4 with the 36 Newton--Euler rows discharged by P5, all in the same chart,
row order, row scaling, and proof norm fixed by P1--P3 plus the retained P4
binding side. Equivalently, after the
fixed block norm and row-scaling map are chosen, write
\(F_{\rm nd}=F_{\mathrm{nd},h}(Z_G;y)\) and
\(F_{\rm dyn}=F_{\mathrm{dyn},h}(Z_G;y)\). Then
Eq.~\eqref{eq:route-suff-block-residual} gives the fixed block-norm bridge:
\begin{equation}
\label{eq:route-suff-block-residual}
 \begin{aligned}
  \|F_{A,h}(Z_G;y)\|
  &\le C_{\rm blk}\bigl(\|F_{\rm nd}\|+\|F_{\rm dyn}\|\bigr) \\
  &\le C_{\rm blk}(C_{\rm nd}h^7+0)\le C_Rh^7 .
 \end{aligned}
\end{equation}
Here \(F_{\rm nd}\) is the 96-row non-dynamic block and
\(F_{\rm dyn}=0\) is the 36-row Newton--Euler block supplied
by the D5 direct-substitution certificate. This displayed bridge is
stage-local at \(Z_G\), before Newton termination, endpoint closure, and
local-to-global transfer; it is not a finite-probe, primitive/Taylor, or
residual-to-error promotion statement. Lemma~\ref{lem:stage-residual-defect}
therefore gives the exact accepted-stage perturbation \(C_Ah^7\).  Lemma~\ref{lem:endpoint-closure}
adds the endpoint contribution \(C_Eh^7\), and
Lemma~\ref{lem:inexact-newton} adds
the solver term in Eq.~\eqref{eq:route-suff-inexact-solver-term}:
\begin{equation}
\label{eq:route-suff-inexact-solver-term}
  C_N\eta_h^{\rm tube}\le C_Nc_\eta h^7
\end{equation}
under the retained P6 solver-scale condition.  Thus the local sum consumed by
the theorem is Eq.~\eqref{eq:route-suff-local-sum}:
\begin{equation}
\label{eq:route-suff-local-sum}
  C_Gh^7+C_Ah^7+C_Eh^7+C_Nc_\eta h^7,
\end{equation}
with constants chosen on the compact branch before the \(h\)-limit.  PC4 is a
negative route rule: no residual-only mechanism, source-policy, or reaction row
is inserted into that decomposition.  Explicitly, the same-branch triangle
inequality is
\begin{equation}
\label{eq:route-suff-triangle}
\begin{aligned}
  \|\Psi_h^{\mathrm{G6FVA}}(y)-\varphi_h(y)\|
  &\le \|\Psi_h^G(y)-\varphi_h(y)\|
      +\|\Psi_h^A(y)-\Psi_h^G(y)\|\\
  &\quad+\|\Psi_h^E(y)-\Psi_h^A(y)\|
      +\|\Psi_h^{\mathrm{G6FVA}}(y)-\Psi_h^E(y)\|\\
  &\le (C_G+C_A+C_E+C_Nc_\eta)h^7 .
\end{aligned}
\end{equation}
Equation~\eqref{eq:route-suff-triangle} is a single same-branch triangle
inequality; it does not splice
quantities from different runs, endpoint policies, row scalings, or Newton
roots. Read as a sufficiency implication, this
statement is not an equivalence or a necessity theorem for all possible
sixth-order routes. It does not choose constants from finite slopes, Newton
logs, source-policy rows, or reproducibility package records. Hence every term used by the
local-defect proof is either a theorem-domain interface or a proved
row-local certificate on the same accepted branch, and the route is sufficient
for the stated conditional theorem without promoting residual-only evidence.
\end{proof}

\begin{table}[htbp]
\centering
\caption{Theorem-interface and output-boundary anchors for
Theorem~\ref{thm:g6fullva-order}.}
\label{tab:theorem-assumption-anchor-ledger}
\scriptsize
\setlength{\tabcolsep}{2pt}
\renewcommand{\arraystretch}{0.9}
\begin{tabular}{P{0.10\linewidth}P{0.28\linewidth}P{0.18\linewidth}P{0.34\linewidth}}
\toprule
Interface & Manuscript anchor & Theorem role & Consumption rule \\
\midrule
P1 & Assumption~\ref{ass:regularity} and
Table~\ref{tab:p1p2-compact-tube-obligation-ledger} &
Domain interface &
Defines the smooth reduced-chart and FullVA-lift tube consumed by every local
estimate. It is not discharged by finite trajectories or row checks. \\
P2 & Assumption~\ref{ass:regularity}, Lemma~\ref{lem:endpoint-closure},
Lemma~\ref{lem:inexact-newton}, Theorem~\ref{thm:g6fullva-order}, and
Table~\ref{tab:p1p2-compact-tube-obligation-ledger} &
Domain interface &
Provides compact-tube inverse, endpoint-ball/right-inverse, stability, and
branch-solver hypotheses; rank probes and endpoint diagnostics do not close it. \\
P3 & Tables~\ref{tab:dynamic-proof-closure-matrix} and
\ref{tab:p3p4-implementation-defect-obligation-ledger} &
Implementation interface &
Binds the implemented residual/Jacobian path; primitive symbolic
oracle completion remains a separate non-active certificate track and is not
used by PC2; this is also not source-paper residual replacement. \\
P4 & Lemma~\ref{lem:fullva-stage-lift} and
Lemma~\ref{lem:stage-residual-defect}, plus
Table~\ref{tab:p3p4-implementation-defect-obligation-ledger} &
Implementation interface; row-local certificate proved &
Supplies the proved 96 non-dynamic rows only after the accepted branch,
proof norm, and implementation binding have been fixed; the full 132-row
defect also requires P5. \\
P5 & Tables~\ref{tab:newton-euler-obligations} and
\ref{tab:newton-euler-row-target-map}, plus
Lemma~\ref{lem:stage-residual-defect} and
Table~\ref{tab:p5-direct-route-satisfaction-ledger} &
Stage-local identity proved &
Supplies the 36 Newton--Euler rows through the D5 direct route on \(Z_G\);
it does not promote P6, P7, or source-policy readiness. \\
P6 & Lemma~\ref{lem:inexact-newton} and
Theorem~\ref{thm:g6fullva-order}, plus
Table~\ref{tab:p6-solver-policy-obligation-ledger} &
Solver-scale condition &
Consumes only the compact-tube \(\eta_h^{\rm tube}\le c_\eta h^7\)
solver condition; finite solver probes are consistency diagnostics and do
not instantiate theorem-level solver-policy proof. \\
P7 & Table~\ref{tab:proof-traceability} and
Table~\ref{tab:p7-residual-to-error-obligation-ledger} &
Output nonclaim boundary &
Records the residual-to-error boundary; residual, reaction, and
common-reference/source-policy rows are scope records only. No residual-to-error
transfer theorem is assumed here; seven transfer obligations remain outside
the theorem. \\
\bottomrule
\end{tabular}
\end{table}

\noindent\textbf{Assumption-anchor map.}
The proof records where the compact-tube interfaces,
the implemented residual path, the direct Newton--Euler identity, the
solver-scale condition, and the residual-to-error boundary enter the
manuscript before any dependency is consumed. Anchor presence is not a proof
discharge: it does not prove compact-tube regularity, the solver hypothesis,
primitive symbolic-oracle work, residual-to-error promotion, source-policy
rows, full-\tfe{} replacement, or reproducibility package completeness.  Its
purpose is to make the theorem domain explicit without replacing the
local-defect proof, which still proceeds through the residual bridge,
endpoint perturbation, inexact Newton perturbation, and fixed-time transfer.

\begin{table}[htbp]
\centering
\caption{Theorem input-output flow for
Theorem~\ref{thm:g6fullva-order}.}
\label{tab:theorem-dependency-consumption-ledger}
\scriptsize
\setlength{\tabcolsep}{2pt}
\renewcommand{\arraystretch}{0.9}
\begin{tabular}{P{0.15\linewidth}P{0.29\linewidth}P{0.28\linewidth}P{0.20\linewidth}}
\toprule
Dependency & First consumed by & Later theorem output & Explicit non-use \\
\midrule
P1/P2 compact-tube data &
Endpoint closure, local stage-root perturbation, inexact Newton, and
local-to-global transfer lemmas. &
Accepted-branch local defect, tube retention, and reduced-chart grid estimate. &
Finite sweeps, endpoint checks, rank probes, or remote nonlinear roots. \\
P3/P4 implementation path &
Formula-row/AD binding and the 96-row non-dynamic certificate before the
stage-residual perturbation step. &
Implemented direct residual path for the non-dynamic part of the 132-row
defect. &
Primitive symbolic oracle closure, source-paper residual replacement, or
dynamic-row closure by P4 alone. \\
P5 direct dynamic rows &
D5 direct-substitution identity evaluated on the lifted Gauss stage \(Z_G\). &
The 36 Newton--Euler rows needed for the direct stage-residual route. &
Residual control at inexact Newton iterates, finite trajectories, or
residual-to-error promotion. \\
P6 branch-solver scale &
Inexact Newton endpoint contribution \(C_Nc_\eta h^7\). &
The accepted local defect remains \(O(h^7)\) under the retained
compact-tube solver hypothesis. &
Fixed production tolerances or finite solver probes as theorem-level
solver-policy proof. \\
P7 residual-to-error boundary &
Only the theorem-use and exclusion tables, not the order proof. &
Residual-only mechanism rows stay diagnostic until a separate transfer theorem
closes. &
Four-link, slider-crank, source-policy, or reaction residual tables as
trajectory-error evidence. \\
PC1--PC2 route checks plus retained P6 and PC4 boundary &
The conditional direct residual-defect route and residual-to-error exclusion. &
PC1--PC2 close the local residual route under the stated theorem interfaces
and retained P6 solver-scale hypothesis; PC4 records residual-to-error exclusion only. &
Promotion of domain interfaces into proved theorem inputs. \\
\bottomrule
\end{tabular}
\end{table}

\noindent\textbf{Theorem input-output flow.}
The displayed inputs record the proof-input flow
together with explicit exclusions. P1/P2
are consumed by the compact-tube and local-perturbation lemmas; P3/P4
identify the implemented residual path; P5 supplies the direct
Newton--Euler rows; P6 enters only through the retained branch-solver
scale; and P7 is recorded only in the exclusion column as a
residual-to-error boundary, not as a theorem premise. The
input-output flow does not prove a theorem-level solver-policy result,
a residual-to-error transfer, source-policy readiness, full-\tfe{} readiness,
or external reproducibility package claims.
In short, the input-output table lists theorem inputs and explicit exclusions only; it is not
an additional proof argument.

\noindent\textbf{Route separation.}
The PC1--PC2 labels are local route checks for the accepted residual-defect
proof route under the retained P6 solver-scale hypothesis, while PC4 labels
the residual-to-error exclusion; they are not replacements for
retained P1/P2/P3 theorem interfaces, the retained P4 binding interface, the
P6 solver-scale interface, the discharged P5 direct-route certificate, or the
P7 output boundary and do not assert that every
auxiliary proof route has closed.  The
active route uses the implemented residual/Jacobian path, the non-dynamic
implementation certificate, and the D5 direct-substitution certificate at the
lifted Gauss stage \(Z_G\).  The primitive symbolic/Taylor decomposition is
retained only as a separate residual-certificate route.  Likewise, finite
AD-expansion probes and scaled-tolerance probes check the recorded
implementation path and solver scale on finite windows, while
Theorem~\ref{thm:g6fullva-order} uses only the explicit branch-selected
\(\eta_h^{\rm tube}\le c_\eta h^7\) condition for its asymptotic claim.
Thus PC3 is interpreted only through the retained P6 theorem condition:
the finite scaled-tolerance probes are consistency checks, not a discharge of
the theorem-level nonlinear-solver hypothesis.
P6 is the retained theorem-level solver-policy condition.
P6 solver boundary: finite one-step, short-trajectory, and
\(h\)-scaled tolerance probes are finite diagnostics of branch-scale consistency only. P6
remains the explicit theorem condition
\(\eta_h^{\rm tube}\le c_\eta h^7\) on the compact-tube branch-selected
solver hypothesis; these probes do not prove a theorem-level solver-policy theorem, do not turn fixed tolerances into an asymptotic proof, and do not
assert global Newton convergence or remote-root selection.

\begin{table}[htbp]
\centering
\caption{P6 solver-policy interface table for
Theorem~\ref{thm:g6fullva-order}.}
\label{tab:p6-solver-policy-obligation-ledger}
\scriptsize
\setlength{\tabcolsep}{2pt}
\renewcommand{\arraystretch}{0.9}
\begin{tabular}{P{0.18\linewidth}P{0.26\linewidth}P{0.25\linewidth}P{0.23\linewidth}}
\toprule
Retained object & Consumed by & Current diagnostic & Non-discharge \\
\midrule
Branch residual scale \(\eta_h^{\rm tube}\le c_\eta h^7\) &
Theorem~\ref{thm:g6fullva-order} and Lemma~\ref{lem:inexact-newton}. &
Finite one-step and short-trajectory scaled-tolerance probes. &
No finite window proves a theorem-level solver-policy theorem for every
accepted step. \\
Proof residual norm and implementation binding &
The residual norm used in Lemma~\ref{lem:inexact-newton} after the accepted
chart, row ordering, and residual scaling are fixed. &
Implementation residual logs bind to this norm only through the accepted
implementation path and compact-tube norm equivalence. &
Logged residuals do not create a second residual budget, change the proof
norm, or prove P6 as a uniform solver theorem. \\
Stopping-proxy rejection &
Only the fixed proof-norm inequality
\(\|F_{A,h}(\widetilde Z_A;y)\|\le c_\eta h^7\) is consumed. &
Correction-norm, relative-residual, iteration-cap, or backend success flags
are usable only after a predeclared conversion to that same absolute residual
bound on the accepted branch. &
No proxy stopping rule, machine-precision floor, or adaptive backend flag
proves P6 by itself. \\
Strong local residual inverse &
The inexact-Newton perturbation contribution to the local defect. &
AD Jacobian probes and branch-selected Newton diagnostics. &
These diagnostics do not assert global Newton convergence or remote-root
selection. \\
Same accepted branch and initial state &
The PC3 perturbation route and accepted reduced-chart branch. &
Recorded reduced-chart initialization and Gauss-predictor branch selection. &
Arbitrary Newton initializations and remote nonlinear roots remain outside
the theorem. \\
Fixed production tolerances &
Finite reported diagnostics and reproducibility checks only. &
Tolerance-regime sweeps and finite solver residual summaries. &
Fixed tolerances are not promoted into an asymptotic proof. \\
Uniform all-transition solver-policy theorem &
Additional theorem evidence needed only to prove P6 as a separate uniform
all-transition solver-policy result; the present theorem remains conditional on P6. &
The finite one-step and short-trajectory probes are recorded diagnostics,
not theorem evidence. &
The uniform all-transition scaled trajectory solver theorem is not proved
here and is not needed to state the conditional theorem. It would be needed
only to replace P6 as a retained hypothesis by a proved solver theorem. \\
\bottomrule
\end{tabular}
\end{table}

\begin{definition}[P6 solver-policy certificate template]
\label{def:p6-solver-policy-certificate-template}
For a fixed theorem instance, a theorem-level P6 solver-policy certificate
would be the same-branch tuple in Eq.~\eqref{eq:p6-certificate-tuple},
\begin{equation}
\label{eq:p6-certificate-tuple}
  \mathcal T_{\rm P6}
  =(\mathcal N_h,\mathcal S_h,\|\cdot\|_{\rm proof},c_\eta,h_0,
    r_N,M_N,\mathcal B_N).
\end{equation}
and this tuple is fixed before any residual log or error table is inspected.  Here
\(\mathcal N_h\) is the nonlinear-solver map initialized from the accepted
Gauss predictor, \(\mathcal S_h\) is the stopping rule, and
\(\|\cdot\|_{\rm proof}\) is the same proof residual norm used in
Lemma~\ref{lem:inexact-newton}.  The tuple must certify the following
uniform facts on the same compact branch for every \(0<h\le h_0\) and every
accepted transition in the theorem domain:
\begin{enumerate}
  \item the Newton iterates selected by \(\mathcal N_h\) remain in the
  branch ball \(\mathcal B_N\) of radius \(r_N\) attached to the strong local
  residual inverse bound \(M_N\);
  \item the stopping rule \(\mathcal S_h\) is stated in, or is converted in
  advance to, the same absolute proof residual norm;
  \item the terminated iterate \(\widetilde Z_{A,n}\) satisfies
  Eq.~\eqref{eq:p6-certificate-template-bound},
  \begin{equation}
  \label{eq:p6-certificate-template-bound}
    \|F_{A,h}(\widetilde Z_{A,n};y_n)\|_{\rm proof}
    \le c_\eta h^7,\qquad 0\le n<N_h ;
  \end{equation}
  \item the constants \(c_\eta,h_0,r_N,M_N\) are independent of \(h\), of the
  reported grid length, and of finite residual/error tables.
\end{enumerate}
If the implementation stops in a backend residual norm rather than directly in
\(\|\cdot\|_{\rm proof}\), the conversion must be part of
\(\mathcal T_{\rm P6}\) before the run.  In particular, the certificate must
include an \(h\)-independent constant \(\kappa_N\) on the same branch ball
such that Eq.~\eqref{eq:p6-backend-to-proof-norm-conversion} holds,
\begin{equation}
\label{eq:p6-backend-to-proof-norm-conversion}
  \|F_{A,h}(Z;y)\|_{\rm proof}
  \le
  \kappa_N\,
  \|r_h^{\rm backend}(Z;y)\|_{\rm backend},
  \qquad Z\in\mathcal B_N .
\end{equation}
Only after Eq.~\eqref{eq:p6-backend-to-proof-norm-conversion} is proved in
advance can a backend stopping target be interpreted as the proof-norm bound
in Eq.~\eqref{eq:p6-certificate-template-bound}.  Without this predeclared
norm conversion, a logged backend residual remains a finite diagnostic and
does not instantiate P6.
Such a certificate would promote P6 from a retained branch-solver hypothesis
to a discharged solver-policy theorem for that fixed instance.  No such
certificate is proved in this manuscript.  The one-step scaled-tolerance
probe, short-trajectory probe, tolerance-regime sweep, fixed production
tolerances, and Newton logs are finite diagnostics only; they do not
instantiate \(\mathcal T_{\rm P6}\) unless the branch, norm, stopping
conversion, and constants above were fixed before the run.
\end{definition}

\noindent\textbf{P6 solver-scale interpretation.}
The displayed separation is not a
solver-policy assumption discharge. It separates
finite one-step probes, finite short-trajectory probes, tolerance-regime
sweeps, AD Jacobian probes, and branch-selected Newton diagnostics from the
retained theorem-level solver-policy condition. The proof residual norm is
fixed before reading implementation logs; those logs enter only through the
accepted implementation-path binding and an \(h\)-independent compact-tube norm
equivalence. It does not prove P6, does not promote fixed tolerances into
asymptotic proof, and does not assert global Newton convergence, remote-root
selection, a P7 residual-to-error boundary, source-policy/full-\tfe{}
readiness, or reproducibility package readiness evidence.

\noindent\textbf{P6 admissibility boundary.}
The retained
\(\eta_h^{\rm tube}\le c_\eta h^7\) condition is a
\mbox{residual-tolerance contract, not a restatement of the theorem conclusion}.  It is stated in terms
of the computable 132-row residual norm after chart, scaling, row-ordering,
and branch conventions are fixed, so an \(h\)-dependent stopping rule can be
applied before endpoint error or observed order slopes are measured. For a
trajectory-level invocation, the same \(c_\eta\) bounds the maximum selected
branch residual over all accepted transitions before the local-to-global
bootstrap is applied; it is not a stepwise constant fitted after individual
Newton logs are inspected. The
condition is not inferred from observed slopes, fixed-tolerance logs, or
post-hoc trajectory errors, and it does not assume access to the exact
trajectory. What remains retained is the compact-tube branch-membership and
strong-local-inverse hypothesis that turns this computable residual bound
into the \(C_Nc_\eta h^7\) perturbation term in
Lemma~\ref{lem:inexact-newton}.
The \mbox{compact-tube branch-membership and strong-local-inverse hypothesis}
remains retained.

\noindent\textbf{Solver-envelope quantifier discipline.}
The scalar solver-envelope constant, proof
residual norm, row scaling, and branch rule are fixed before the step-size
window and before any residual logs are read. Increasing that constant after
inspecting logs changes the theorem hypothesis and the local-defect constant;
it is not a proof of P6 and does not convert finite diagnostics into
theorem-level solver-policy proof.

\noindent\textbf{P6 evidence-grade rule.}
P6 has only one theorem-grade input: a predeclared same-branch residual
envelope in the fixed proof norm,
\(\eta_h^{\rm tube}\le c_\eta h^7\), with the branch rule, row scaling,
proof norm, and \(c_\eta\) fixed before the accepted transition is reported.
Finite one-step probes, short-trajectory probes, tolerance-regime sweeps,
fixed-tolerance production logs, and observed order slopes are
diagnostic-grade evidence only.  They may check whether a reported run is
consistent with a predeclared P6 instance, but they cannot define
\(c_\eta\), change the proof norm or row scaling, certify branch-ball
membership, or upgrade a fixed tolerance into theorem-level solver-policy
proof.  Thus a solver log can specialize a P6 hypothesis that has already
been fixed; it cannot create that hypothesis after the fact.

\noindent\textbf{P6 conditional-instantiation checkpoint.}
The retained P6 condition has an operational sufficient instantiation, but
that instantiation is not the same as the current fixed-tolerance evidence.  A
trajectory instantiates P6 for Theorem~\ref{thm:g6fullva-order} only after the
proof norm, row scaling, branch rule, and \(c_\eta\) have been fixed before
the run, consistently with
Definition~\ref{def:p6-solver-policy-certificate-template}; every accepted transition certifies the same Gauss-predictor branch
and the branch-ball/strong-local-inverse hypotheses of
Lemma~\ref{lem:inexact-newton}; and the nonlinear solve enforces the
proof-norm stopping target in Eq.~\eqref{eq:p6-proof-norm-stopping-target}:
\begin{equation}
\label{eq:p6-proof-norm-stopping-target}
  \|F_{A,h}(\widetilde Z_{A,n};y_n)\| \le c_\eta h^7,\qquad 0\le n<N_h .
\end{equation}
Under those already fixed choices, the reported-grid specialization
\(\eta_h^{\rm reported}=\max_{0\le n<N_h}\eta_{h,n}\le c_\eta h^7\) is a consequence of the
solver hypothesis and can be consumed by Lemma~\ref{lem:inexact-newton}.  This is
a sufficient instantiation rule, not a hidden empirical discharge: it does
not prove a theorem-level solver-policy theorem for the current fixed-tolerance
runs, does not choose \(c_\eta\) from residual ratios, does not prove global
Newton convergence or remote-root exclusion, and does not prove P7,
source-policy, full-\tfe{}, or reproducibility package readiness.
A transition obtained only after changing the proof norm, row scaling, branch
rule, or \(c_\eta\) in response to the computed residual log is therefore not
an instantiation of P6 for this theorem.
Likewise, P1 remains a smooth reduced-chart and FullVA-lift existence
condition on the compact proof tube: finite \(h\)-sweeps and local
mechanism-coverage rows diagnose branch plausibility but do not prove or
remove the theorem's uniform smooth-lift hypothesis.
Similarly, P2 remains a uniform compact-tube inverse, endpoint-ball, and
tube-retention interface: finite rank probes, endpoint-Jacobian checks, and
ordinary Lipschitz estimates may diagnose the branch but do not eliminate
those hypotheses.

\noindent\textbf{P1/P2 compact-tube boundary.}
P1 and P2 remain retained
compact-tube theorem interfaces. Finite \(h\)-sweeps, local mechanism rows,
finite rank probes, endpoint-Jacobian checks, and ordinary Lipschitz estimates
diagnose the selected branch only; they do not prove or remove smooth-lift
regularity, uniform compact-tube invertibility, endpoint-ball/right-inverse
hypotheses, or tube-retention stability.

\noindent\textbf{Compact-inverse quantifier discipline.}
The compact tube, Gauss-predictor branch,
proof residual norm, and inverse constants must be selected before numerical
diagnostics are read. Shrinking the admissible step-size window after those
objects are fixed is part of the theorem hypothesis; refitting the tube or
inverse constant from diagnostics is not proof closure.
A rank probe, solver log, per-step projection log, and endpoint norm may
diagnose the selected run, but none of them fixes the retained tube, inverse
constant, or right-inverse margin after the asymptotic theorem instance has
been chosen.

\begin{table}[htbp]
\centering
\caption{P1/P2 compact-tube interfaces for
Theorem~\ref{thm:g6fullva-order}.}
\label{tab:p1p2-compact-tube-obligation-ledger}
\scriptsize
\setlength{\tabcolsep}{2pt}
\renewcommand{\arraystretch}{0.9}
\begin{tabular}{P{0.18\linewidth}P{0.26\linewidth}P{0.25\linewidth}P{0.23\linewidth}}
\toprule
Retained object & Consumed by & Current diagnostic & Non-discharge \\
\midrule
P1 smooth lift and chart tube &
Gauss-stage lift, local defect, and reporting-map estimates. &
Accepted smooth-path construction and finite branch diagnostics. &
No finite trajectory row proves a uniform smooth FullVA lift on the whole
compact proof tube. \\
P2 stage inverse on the Gauss-predictor branch &
Stage-root perturbation and accepted-branch consistency. &
Formula-row/AD binding and finite rank probes diagnose the selected branch. &
Rank probes do not prove a uniform compact-tube inverse or remote-root
exclusion. \\
P2 endpoint-ball and local right inverse &
Lemma~\ref{lem:endpoint-closure} and the endpoint contribution to the
one-step defect. &
Endpoint anchor and local correction diagnostics diagnose the branch. &
Endpoint checks do not discharge the closed-anchor, ball-margin, or
right-inverse hypotheses. \\
P2 stability and tube-retention bootstrap &
Lemma~\ref{lem:local-global-transfer}. &
Observed smooth-window order and bounded local rows diagnose plausibility. &
Ordinary Lipschitz estimates and finite \(h\)-sweeps do not prove tube
retention for every retained state. \\
P2 strong local residual inverse &
Lemma~\ref{lem:inexact-newton}. &
AD Jacobian probes and branch-selected Newton diagnostics identify the
implemented local path. &
They do not replace the compact-tube residual inverse or the separate P6
\(\eta_h^{\rm tube}\le c_\eta h^7\) policy. \\
\bottomrule
\end{tabular}
\end{table}

\begin{definition}[P1/P2 compact-branch certificate template]
\label{def:p1p2-compact-branch-certificate-template}
A theorem-level certificate discharging the retained P1/P2 compact-branch
interfaces, rather than assuming them as in
Assumption~\ref{ass:regularity}, would be a fixed same-branch tuple
\begin{equation}
\label{eq:p12-certificate-tuple}
  \mathcal T_{\rm P12}
  =(K,\chi,\mathcal L,\mathcal G_h,\mathcal E_h,\mathcal C_h,
    \mathcal P_h,M_Z,L_Z,r_Z,M_E,r_E,C_s,\Gamma_s,C_{\mathcal R},h_0).
\end{equation}
chosen before any finite trajectory, rank probe, residual log, or observed
order is read.  Equation~\eqref{eq:p12-certificate-tuple} fixes the P1/P2
certificate object before any diagnostic is inspected.  Here \(K\) is the
compact reduced-chart tube, \(\chi\) is
the chart, \(\mathcal L\) is the smooth FullVA lift, \(\mathcal G_h\) is the
Gauss-predictor branch, \(\mathcal E_h\) is Lie endpoint reconstruction,
\(\mathcal C_h\) is the local endpoint-closure map, and
\(\mathcal P_h=\mathcal C_h\circ\mathcal E_h\) is the reported endpoint map.
The tuple would have to certify, uniformly for \(0<h\le h_0\), all retained
states in \(K\), and all accepted branch-selected transitions, the following
facts:
\begin{enumerate}
  \item a \(C^7\) reduced constrained vector field and \(C^7\) FullVA lift on
  \(K\), with the lifted Gauss stages lying in the chart domain;
  \item a unique local stage root on the Gauss-predictor branch, a uniform
  inverse bound \(M_Z\) for \(D_ZF_{A,h}\), a quadratic-remainder bound
  \(L_Z\), and a branch radius \(r_Z\);
  \item a closed endpoint-anchor family, endpoint-ball radius \(r_E\),
  raw-defect bound, and local right inverse for endpoint closure with
  uniform bound \(M_E\);
  \item a strong local residual inverse on the same branch for the
  inexact-Newton perturbation, compatible with the proof residual norm used
  in Lemma~\ref{lem:inexact-newton};
  \item a one-step stability estimate with factor \(1+C_s h\), a
  tube-retention/first-exit bootstrap over \(N_hh\le T\), and the resulting
  fixed-time transfer constant \(\Gamma_s\);
  \item a reporting-map norm-equivalence constant \(C_{\mathcal R}\) for the
  displayed position--velocity grid norm.
\end{enumerate}
Such a certificate would promote P1/P2 from retained compact-branch
interfaces to a constructed compact-branch theorem for the fixed FullVA
object.  No such certificate is proved in this manuscript.  Finite smooth
branch runs, rank probes, endpoint checks, ordinary Lipschitz estimates,
and mechanism-coverage rows may diagnose compatibility with
\(\mathcal T_{\rm P12}\), but they do not instantiate it unless the tube,
branch, chart, inverses, endpoint anchors, stability constants, and reporting
norm were fixed before those diagnostics were read.
\end{definition}

\noindent\textbf{P1/P2 compact-tube interpretation.}
The displayed separation is
not an assumption discharge. It separates smooth-lift
regularity, compact-tube inverse, endpoint-ball/right-inverse, strong local
residual inverse, and tube-retention stability from finite \(h\)-sweeps, rank
probes, endpoint checks, local mechanism rows, ordinary Lipschitz estimates,
residual/reaction tables, source-policy rows, and reproducibility package evidence. It
does not prove P1/P2, a solver-policy theorem, or a P7 residual-to-error
transfer, and it does not supply source-policy/full-\tfe{} readiness or
reproducibility package readiness evidence.
The load-bearing P1/P2 data are not pointwise diagnostics. The theorem
requires a single compact tube, an \(h\)-independent radius and constants,
and a branch selection map fixed before the asymptotic limit: every accepted
predictor, stage, endpoint-correction, and reported reduced-chart state must
stay in that tube; the stage residual must be written in the same row/column
chart; the stage Jacobian and endpoint-closure Jacobian must have uniform
right inverses on the tube; and the one-step map must keep the accepted
trajectory inside the tube under the stated bootstrap. Finite rank probes,
endpoint checks, ordinary Lipschitz estimates, and observed convergence rows
sample this structure on visited states only. They do not establish tube
retention, a uniform lower margin for all nearby accepted states, or
branch-jump exclusion; the constants used in the endpoint, inexact-Newton,
and local-to-global lemmas therefore cannot be fitted from observed
convergence rows.
The \mbox{uniform right inverses on the tube} are retained hypotheses, not
finite-run conclusions.
Equivalently, the manuscript retains P1/P2 as Assumption~\ref{ass:regularity}
interfaces rather than proving the certificate in
Definition~\ref{def:p1p2-compact-branch-certificate-template}.

PC1/P3 are therefore read as implementation-path checks: the implemented
formula-row checks and the AD-expanded implementation-path certificate identify the
direct-route residual/Jacobian path. They do not close the primitive dynamic
symbolic oracle, the primitive/Taylor symbolic-defect route, source-paper
residual reproduction, or a discharge of the retained compact-tube
hypotheses.
Thus a still-open primitive dynamic symbolic oracle is not a missing theorem
input for the accepted route; it is only the alternative primitive/Taylor lane
that would have to be completed before replacing the direct residual
certificate.

\noindent\textbf{P3/P4 implementation-defect boundary.}
P3 and P4 remain retained
theorem-interface inputs rather than standalone discharged interfaces.
Formula-row checks and AD-expanded runtime binding identify
the implemented direct-route residual/Jacobian path; the AD-expanded
implementation-path certificate is closed for that active route and is
consumed by the same-branch direct residual bridge.
They do not prove the primitive dynamic symbolic oracle or the
primitive/Taylor symbolic-defect route, and they are not source-paper residual
reproduction. They also do not claim primitive/Taylor symbolic-defect route
closure. The 96-row non-dynamic certificate covers only kinematic and
lower-pair rows; it does not by itself prove the full 132-row stage residual,
which uses the separate 36-row D5 direct-substitution certificate for
Newton--Euler rows. The full 132-row stage residual required for the order
argument is supplied only after this dynamic-row handoff is combined with the
non-dynamic certificate through Lemma~\ref{lem:full-132-row-residual-bridge}.
The separate 36-row D5 direct-substitution phrase refers only to this
dynamic-row handoff, not to P3/P4 closure.
Within the theorem, P3 is used only as a fixed map-identity interface: before
the local perturbation argument, the mathematical row formulas, row ordering,
row scaling, and AD-differentiated implementation path are fixed as one
residual map. The proof then evaluates this fixed map on the lifted Gauss
stage. Finite runtime probes and AD checks identify that implementation path;
they are not used backward to infer the \(h^7\) defect from reported
trajectories or Newton residuals.

\begin{table}[htbp]
\centering
\caption{P3/P4 implementation-defect interface table for
Theorem~\ref{thm:g6fullva-order}.}
\label{tab:p3p4-implementation-defect-obligation-ledger}
\scriptsize
\setlength{\tabcolsep}{2pt}
\renewcommand{\arraystretch}{0.9}
\begin{tabular}{P{0.18\linewidth}P{0.26\linewidth}P{0.25\linewidth}P{0.23\linewidth}}
\toprule
Retained object & Consumed by & Current evidence & Non-discharge \\
\midrule
P3 runtime formula/AD binding &
The implemented direct-route residual/Jacobian path. &
Formula-row checks and AD-expanded runtime binding. &
This does not prove the primitive dynamic symbolic oracle or primitive/Taylor
symbolic-defect route. \\
Excluded full-\tfe{} source-paper residual identity &
The stronger full-\tfe{} replacement and source-policy route only. &
The current manuscript records source-policy rows separately. &
This is not source-paper residual reproduction or full-\tfe{} readiness. \\
P4 96-row non-dynamic certificate &
The non-dynamic part of the accepted local defect argument. &
The kinematic and lower-pair rows are certified on the smooth lift. &
This does not by itself prove the full 132-row stage residual. \\
P4 handoff to dynamic rows &
The direct residual route for PC2 together with P5. &
The separate 36-row D5 direct-substitution certificate supplies the
Newton--Euler rows. &
The handoff does not convert P3/P4 alone into a standalone 132-row
direct-route certificate. \\
P3/P4 implementation-domain interface &
The conditional order theorem's implementation-defect premise. &
Dependency tables map the implemented runtime path and non-dynamic rows. &
It does not prove P3/P4 as standalone primitive/source-policy closures,
a solver-policy theorem, a P7 residual-to-error boundary,
source-paper comparison, full-\tfe{} replacement, or external
reproducibility package claims. \\
\bottomrule
\end{tabular}
\end{table}

\noindent\textbf{P3/P4 implementation-defect boundary.}
This is not a standalone
implementation-defect proof. It separates formula-row checks, AD-expanded runtime binding, the
96-row non-dynamic certificate, and the separate P5 direct Newton--Euler
certificate from primitive dynamic symbolic oracle closure,
primitive/Taylor symbolic-defect route closure, source-paper residual
reproduction, full 132-row stage residual promotion by P3/P4 alone,
source-paper comparison, full-\tfe{} replacement, and external
reproducibility package evidence. It
does not close P3/P4 as standalone primitive/source-policy closures or promote
P3/P4 alone into an independent direct-route theorem input.

\noindent\textbf{P3/P4 evidence convention.}
Their theorem use is through a specified
same-branch interface. P3 is backed by the implemented formula-row and
AD-expanded runtime binding, and P4 is backed by the 96-row non-dynamic
certificate. Their retained status means
that these checks are consumed through the same-branch direct residual bridge
and still require the separate P5 dynamic-row handoff; they are not, by
themselves, a primitive dynamic symbolic-oracle closure, a source-paper
residual replacement, or a full 132-row certificate independent of P5.

\noindent\textbf{P5 direct-route boundary: P5.}
P5 is proved only as a stage-local
Newton--Euler residual identity on the smooth Gauss lift. It
supplies the 36 dynamic rows required by PC2 through the D5 direct-substitution
certificate; it does not assert residual control at the inexact Newton
iterate, finite-run trajectories, closed-loop mechanism residual tables,
source-policy rows, or residual-to-error promotion.

\noindent\textbf{P5 inverse-boundary.}
The direct Newton--Euler zero rows close
the dynamic block of the residual vector on \(Z_G\), but they do not provide a
stage-Jacobian inverse, contraction radius, branch-selection theorem,
endpoint right inverse, stability factor, or \(\eta_h\) solver hypothesis. Those
inputs enter through P1/P2/P6 and Lemmas~\ref{lem:stage-residual-defect},
\ref{lem:endpoint-closure}, \ref{lem:inexact-newton}, and
\ref{lem:local-global-transfer}.

\begin{table}[htbp]
\centering
\caption{P5 direct-route satisfaction for
Theorem~\ref{thm:g6fullva-order}.}
\label{tab:p5-direct-route-satisfaction-ledger}
\scriptsize
\setlength{\tabcolsep}{2pt}
\renewcommand{\arraystretch}{0.9}
\begin{tabular}{P{0.18\linewidth}P{0.26\linewidth}P{0.25\linewidth}P{0.23\linewidth}}
\toprule
Satisfied object & Consumed by & Current mathematical input & Non-promotion rule \\
\midrule
P5 stage-local dynamic rows &
PC2 and Lemma~\ref{lem:stage-residual-defect}. &
The D5 direct-substitution certificate evaluates the 36 Newton--Euler rows on
the smooth Gauss lift \(Z_G\). &
This is not residual control at the inexact Newton iterate. \\
P5 smooth-lift residual identity &
The direct residual route only. &
The row-target map and dynamic residual-identity table bind the accepted
Newton--Euler rows. &
It is not a finite-run trajectory statement or a closed-loop residual table. \\
P5 handoff from P3/P4 &
The accepted 132-row local-defect route. &
The 96-row non-dynamic certificate plus the 36-row P5 certificate cover the
implemented direct route. &
This handoff does not close P3/P4 as primitive dynamic symbolic oracle
evidence. \\
P5 scope/exclusion rule &
The theorem-interface table. &
P5 is the only stage-local dynamic identity proved by direct substitution. &
It does not prove P1/P2/P3 theorem interfaces, the retained P4 binding
interface, the P6 solver-scale interface, the P7 residual-to-error boundary
interface, source-paper comparison, full-\tfe{} replacement, or external
reproducibility package claims. \\
\bottomrule
\end{tabular}
\end{table}

P5 direct-route satisfaction rule: the displayed evidence is stage-local, not
a global residual or trajectory certificate. It separates the D5
direct-substitution certificate on the smooth
Gauss lift, the 36 Newton--Euler row identity, and the P3/P4 handoff from
inexact-Newton residual control, finite-run trajectories, closed-loop
mechanism residual tables, source-policy rows, residual-to-error promotion,
primitive dynamic symbolic oracle closure, and external reproducibility package
evidence. It
\mbox{does not promote P5 beyond the direct stage-local} Newton--Euler route.

\begin{table}[htbp]
\centering
\caption{Proof causality for Theorem~\ref{thm:g6fullva-order}.}
\label{tab:proof-causality-ledger}
\scriptsize
\setlength{\tabcolsep}{2pt}
\renewcommand{\arraystretch}{0.92}
\begin{tabular}{P{0.16\linewidth}P{0.30\linewidth}P{0.28\linewidth}P{0.18\linewidth}}
\toprule
Mathematical input & Consumed before & Allowed source & Explicit non-source \\
\midrule
P1--P2 compact tube &
The local stage-root, endpoint-closure, inexact Newton, and
local-to-global estimates. &
Smooth lift, compact-tube inverse, endpoint-ball/right-inverse, strong
residual inverse, and tube-retention stability hypotheses. &
Finite \(h\)-sweeps, rank probes, endpoint checks, or ordinary Lipschitz
diagnostics. \\
Direct stage residual &
The stage-root perturbation bound and the four-term local defect sum. &
The 96-row non-dynamic certificate plus the 36-row D5 direct-substitution
identity evaluated on \(Z_G\). &
Computed trajectory error, inexact Newton iterates, residual/reaction tables,
or source-policy rows. \\
P6 solver scale &
The \(C_Nc_\eta h^7\) inexact-Newton endpoint contribution. &
The compact-tube branch-solver hypothesis
\(\eta_h^{\rm tube}\le c_\eta h^7\). &
Fixed production tolerances or finite scaled-tolerance probes as theorem
evidence. \\
P7 and boundary records &
The theorem-use and reproducibility-boundary exclusions. &
Residual-to-error obligations, source-policy records, and reproducibility package
checks as downstream boundary records. &
Inputs to PC2, a theorem-level solver-policy theorem, or source-policy/full-\tfe{}
readiness. \\
\bottomrule
\end{tabular}
\end{table}

Proof-causality rule: each mathematical input is consumed
before the theorem conclusion that depends on it. The D5 direct-substitution
certificate is evaluated on the lifted Gauss stage \(Z_G\) before the local
perturbation argument; the local defect bound is established before the
discrete Gronwall transfer; and the position--velocity reporting map is used
only after the reduced-chart grid estimate. Finite \(h\)-sweeps, inexact
Newton iterates, residual/reaction tables, source-policy rows, and local
reproducibility records are downstream diagnostics or reproducibility-boundary
evidence. They do not
feed PC2, P6, or P7, and they do not prove a theorem-level
solver-policy theorem, a separate residual-to-error transfer theorem,
source-policy readiness, or full-TFE readiness.

\begin{lemma}[Endpoint-closure perturbation]
\label{lem:endpoint-closure}
Let \(z_u\) be the same-branch unclosed endpoint to which the accepted
endpoint-closure map is applied. In the theorem application
\(z_u=\Psi_h^A=\mathcal E_h(Z_A)\), the endpoint reconstructed from the exact
accepted local stage root before endpoint closure; it is not the Gauss
predictor endpoint \(\Psi_h^G\). Let
$E(z)=0$ denote the square endpoint closure subsystem actually corrected by
the accepted endpoint-closure map--the velocity-level/KKT closure equations in
the reported implementation. This \(E\) is not the set of every raw endpoint
diagnostic reported in Table~\ref{tab:endpoint-closure-scaling}; in
particular, the raw position residual in that table is a separate monitor and
not the local \(E(z_u)\) input below. For each retained transition \((y,h)\),
let \(z_\ast=z_\ast(y,h)\) be the same-branch closed endpoint anchor with
\(E(z_\ast)=0\). Assume that $E$ is $C^2$ on a ball
$B_{r_E}(z_\ast)$, that
$z_u=z_u(y,h)\in B_{r_E/2}(z_\ast)$ for all sufficiently small $h$, and that
$D E(z_\ast)$ has a right inverse $B_\ast=B_\ast(y,h)$ with
$\|B_\ast\|\le M_{E,\mathrm{ri}}$. The constants
\(r_E\), \(M_{E,\mathrm{ri}}\), \(L_E\), \(C_{E,\mathrm{raw}}\), and the
second-derivative bound used below are uniform over retained \(y\in K\) and
the admissible small-step window.
\noindent\textbf{Endpoint-time binding.}
If the endpoint-closure rows are nonautonomous because prescribed constraints,
loads, or drive functions enter the closure equations, then \(E\),
\(\mathcal C_h\), \(z_u\), and \(z_\ast\) in this lemma are fixed-transition
abbreviations for \(E_{h,t_n}\), \(\mathcal C_{h,t_n}\),
\(z_u(y,t_n,h)\), and \(z_\ast(y,t_n,h)\).  All endpoint data are evaluated
at the endpoint time \(t_n+h\) and with the stage history of that same
transition.  The right-inverse and derivative constants are uniform over the
retained time interval, but no row in \(E\) may be evaluated at a neighboring
transition time, a resampled output time, or an unreported terminal \(T\) and
then used in this endpoint perturbation bound.
Assume also the derivative perturbation and smallness conditions in
Eq.~\eqref{eq:endpoint-derivative-smallness}:
\begin{equation}
\label{eq:endpoint-derivative-smallness}
  \|D E(z)-D E(z_\ast)\|\le L_E\|z-z_\ast\|,
  \qquad M_{E,\mathrm{ri}} L_E r_E\le \frac12 ,
\end{equation}
and that the unprojected endpoint closure defect has the explicit local
constant in Eq.~\eqref{eq:endpoint-raw-defect-bound}:
\begin{equation}
\label{eq:endpoint-raw-defect-bound}
  \|E(z_u)\|\le C_{E,\mathrm{raw}}h^7 .
\end{equation}
This raw endpoint defect is a same-object local statement: \(z_u\) is the
unclosed endpoint on the same accepted branch and is evaluated by the same
endpoint functional \(E\) used by the accepted closure map. In the theorem's
four-map chain this is precisely the \(\Psi_h^E-\Psi_h^A\) comparison; it is
a retained P2 endpoint-closure hypothesis for this lemma and
Theorem~\ref{thm:g6fullva-order}, not a post-Newton projection log, not a
terminal-time correction, and not a constraint/reaction residual table.
Endpoint diagnostics may record this scale on reported runs, but they do not
discharge the compact-tube endpoint raw-defect/right-inverse hypothesis.
The endpoint-anchor family, uniform ball radius, right inverse, derivative
constants, and endpoint norm are fixed on the compact proof tube before any reported
endpoint residual, projection diagnostic, step-size sweep, or per-step
projection log is inspected.
\noindent\textbf{Endpoint diagnostic-scope checkpoint.}
The order-fit row in Table~\ref{tab:endpoint-closure-scaling} is not used to
infer the hypothesis \(\|E(z_u)\|\le C_{E,\mathrm{raw}}h^7\).  For the
theorem invocation, \(E\), the endpoint norm, the endpoint-anchor family
\(z_\ast(y,h)\), \(r_E\), \(B_\ast(y,h)\), \(M_{E,\mathrm{ri}}\), and
\(C_{E,\mathrm{raw}}\) are fixed uniform compact-tube data before that table
is read.  Thus the raw position monitor
may have a different fitted slope without contradicting the lemma, because it
is not the local closure functional \(E\); conversely, a small raw position
monitor would not discharge the velocity-level/KKT defect bound, the endpoint
ball, or the right-inverse condition.  The table is a consistency diagnostic
for the reported closure subsystem, not an empirical source of the endpoint
perturbation hypothesis.
If the accepted endpoint correction $\Delta z_E$ is obtained by first solving
the minimum-norm linearized closure equation in
Eq.~\eqref{eq:endpoint-linearized-closure}
\begin{equation}
\label{eq:endpoint-linearized-closure}
  D E(z_u)\,\delta z_1=-E(z_u),
\end{equation}
and then, if needed, refined by the same local Newton/right-inverse closure
solve in Eq.~\eqref{eq:endpoint-newton-refinement}
\begin{equation}
\label{eq:endpoint-newton-refinement}
  \begin{aligned}
  z_1&=z_u+\delta z_1, & r_j&=E(z_j),\\
  D E(z_j)\,\delta z_{j+1}&=-r_j, &
  z_{j+1}&=z_j+\delta z_{j+1},
  \end{aligned}
\end{equation}
with each later correction chosen no larger than a local right-inverse
particular solution while the iterates remain in \(B_{r_E}(z_\ast)\), then,
after possibly shrinking the admissible step-size window, the correction obeys
Eq.~\eqref{eq:endpoint-correction-bound}:
\begin{equation}
\label{eq:endpoint-correction-bound}
  \|\Delta z_E\|\le C_E h^7,\qquad
  C_E=4M_{E,\mathrm{ri}}C_{E,\mathrm{raw}} .
\end{equation}
In the four-map local-defect chain this endpoint correction is the map
difference in Eq.~\eqref{eq:endpoint-map-correction-identity}:
\begin{equation}
\label{eq:endpoint-map-correction-identity}
\begin{aligned}
  \Psi_h^A(y)&=z_u,\\
  \Psi_h^E(y)&=\mathcal C_h(z_u)=z_u+\Delta z_E,\\
  \|\Psi_h^E(y)-\Psi_h^A(y)\|_{\rm end}&\le C_Eh^7 .
\end{aligned}
\end{equation}
The endpoint norm \(\|\cdot\|_{\rm end}\) in
Eq.~\eqref{eq:endpoint-map-correction-identity} is the fixed endpoint-output
component of the proof norm used in the local-defect sum; any
finite-dimensional conversion between the closure-solve coordinates and the
reduced-chart endpoint norm is fixed on the compact tube and absorbed into
\(C_E\) before Theorem~\ref{thm:g6fullva-order} chooses its constants.  Thus
the corrected endpoint \(z_u+\Delta z_E\) differs from the same-branch
unclosed endpoint by only this explicit local \(O(h^7)\) perturbation.
\end{lemma}

\begin{proof}
Let $A_\ast=D E(z_\ast)$ and $A_u=D E(z_u)$.  The assumed perturbation bound
gives $\|(A_u-A_\ast)B_\ast\|\le 1/2$.  Hence
$A_uB_\ast=I+(A_u-A_\ast)B_\ast$ is invertible, and the induced right inverse is
Eq.~\eqref{eq:endpoint-right-inverse}:
\begin{equation}
\label{eq:endpoint-right-inverse}
  B_u=B_\ast(A_uB_\ast)^{-1}
\end{equation}
is a right inverse for $A_u$ with $\|B_u\|\le 2M_{E,\mathrm{ri}}$.  The minimum-norm solution
of the linearized closure equation is no larger than the particular solution
$-B_uE(z_u)$, so Eq.~\eqref{eq:endpoint-first-correction-bound} gives
\begin{equation}
\label{eq:endpoint-first-correction-bound}
  \|\delta z_1\|\le 2M_{E,\mathrm{ri}}\|E(z_u)\|
  \le 2M_{E,\mathrm{ri}}C_{E,\mathrm{raw}}h^7 .
\end{equation}
Since \(z_u\in B_{r_E/2}(z_\ast)\), choosing \(h\) small enough gives
\(2\|\delta z_1\|<r_E/2\), and hence
\(z_u+\delta z_1\in B_{r_E}(z_\ast)\).
This endpoint-ball margin keeps the linear correction and the following
local Newton refinement inside the same ball on which the derivative and
right-inverse perturbation bounds are valid.
The same endpoint-ball constants and endpoint norm are used for every
retained-tube state. The admissible \(h_0\) may shrink after these constants
are fixed, but the proof does not choose the endpoint-anchor family, ball radius, or right inverse
from the computed endpoint correction or from a per-step projection log.
Let \(C_{E,2}=\sup_{B_{r_E}(z_\ast)}\|D^2E\|\).  Since
\(A_u\delta z_1=-E(z_u)\), Taylor's formula gives the explicit post-linear
closure residual in Eq.~\eqref{eq:endpoint-postlinear-residual}:
\begin{equation}
\label{eq:endpoint-postlinear-residual}
  r_1=E(z_u+\delta z_1),\qquad
  \|r_1\|\le \frac12 C_{E,2}\|\delta z_1\|^2
  \le 2C_{E,2}M_{E,\mathrm{ri}}^2C_{E,\mathrm{raw}}^2h^{14}.
\end{equation}
For every later iterate that remains in the ball, the same Neumann-series
right-inverse argument gives a local right inverse of \(D E(z_j)\) with norm
at most \(\overline M_E=2M_{E,\mathrm{ri}}\). Hence the chosen correction
satisfies Eq.~\eqref{eq:endpoint-refinement-recurrence}:
\begin{equation}
\label{eq:endpoint-refinement-recurrence}
  \|\delta z_{j+1}\|\le \overline M_E\|r_j\|,
  \qquad
  \|r_{j+1}\|\le \frac12 C_{E,2}\overline M_E^2\|r_j\|^2 .
\end{equation}
After shrinking the step-size window, \(r_1\) is small enough that the
residuals decrease at least geometrically, while
\(\overline M_E C_{E,2}\|\delta z_1\|\le 1\).  Consequently the total
refinement tail is bounded by Eq.~\eqref{eq:endpoint-refinement-tail-bound}:
\begin{equation}
\label{eq:endpoint-refinement-tail-bound}
  \sum_{j\ge 1}\|\delta z_{j+1}\|
  \le 2\overline M_E\|r_1\|
  \le \overline M_E C_{E,2}\|\delta z_1\|^2
  \le \|\delta z_1\|.
\end{equation}
The same bound also applies to any finite accepted refinement prefix.  Thus
the cumulative endpoint displacement from \(z_u\) is at most
\(2\|\delta z_1\|<r_E/2\), so all iterates stay in \(B_{r_E}(z_\ast)\) by the
endpoint-ball margin, and Eq.~\eqref{eq:endpoint-final-displacement-bound}
closes the endpoint perturbation estimate:
\begin{equation}
\label{eq:endpoint-final-displacement-bound}
  \|\Delta z_E\|\le 2\|\delta z_1\|
  \le 4M_{E,\mathrm{ri}}C_{E,\mathrm{raw}}h^7=C_Eh^7 .
\end{equation}
Thus endpoint closure contributes at most the explicit \(C_Eh^7\) local
perturbation to the one-step map.
\end{proof}

Lemma~\ref{lem:endpoint-closure} does not by itself prove the full
Assumption~\ref{ass:regularity}. It isolates the endpoint-closure part of the
perturbation budget: once the raw endpoint closure defect has the same local
scale as the Gauss truncation defect and the closure Jacobian remains
right-invertible under the stated local splitting bound, the endpoint-closure
correction cannot reduce the asymptotic order. Pointwise right-invertibility
at one endpoint is not used as a shortcut; the closed endpoint-anchor family
\(E(z_\ast(y,h))=0\), the uniform endpoint-ball margin, and the compact-tube
perturbation bound are all part of the condition. This endpoint-closure lemma by itself does not address the
stage-row residual obligation. The companion direct-substitution certificate
verifies the accepted stage residual on the lifted branch; the separate
symbolic/primitive route and any residual-to-error transfer theorem remain
separate non-active/open proof routes rather than endpoint-closure
consequences or external proof-completion claims.
Endpoint-refinement boundary: the refinement in
Lemma~\ref{lem:endpoint-closure} is a local closure solve inside the same
endpoint ball. It is not an extra terminal solve, not a global projection
proof, and not a fitted KKT correction norm; its only theorem-level use is
the \(O(h^7)\) endpoint-increment bound under the stated local right-inverse
and endpoint-ball assumptions. This is a local right-inverse perturbation
bound for endpoint closure only. The refinement is branch-local and
endpoint-norm fixed; it cannot define a new terminal output time or calibrate
an \(h\)-dependent projection constant from finite runs.
Endpoint-output nonfeedback rule: the accepted endpoint closure
\(\mathcal C_h\) is applied after the stage residual solve and Lie endpoint
reconstruction.  Its correction changes only the endpoint-output value
\(\mathcal P_h(Z)=\mathcal C_h(\mathcal E_h(Z))\) used in the one-step map;
it does not redefine the stage vector \(Z\), the lifted Gauss base point
\(Z_G\), the exact accepted stage root \(Z_A\), the implemented residual
\(F_{A,h}\), or the D5 Newton--Euler row identity.  Consequently endpoint
closure cannot repair, re-center, or weaken the \(132\)-row residual bridge:
the bridge is evaluated before endpoint closure on the fixed stage object,
and any method that fed the endpoint correction back into the stage equations
would be a different residual map requiring a separate proof.
Endpoint-time nonfeedback rule: for a nonautonomous transition the closure map
in this paragraph is \(\mathcal C_{h,t_n}\), evaluated at the endpoint time
\(t_n+h\) for the same transition that produced \(Z_A\) or
\(\widetilde Z_A\).  It is not a dense-output interpolant, not a terminal-time
postprocessor for an unreported \(T\), and not a way to change the time slice
of the residual bridge or the subsequent Gronwall recurrence.

\begin{lemma}[FullVA stage-row lift consistency]
\label{lem:fullva-stage-lift}
Assume the smooth reduced chart in Assumption~\ref{ass:regularity} admits a
smooth FullVA lift, written in Eq.~\eqref{eq:fullva-lift-map},
\begin{equation}
\label{eq:fullva-lift-map}
  \mathcal L(y,t)=(q,v,a,\lambda)
\end{equation}
for which the kinematic identities, Newton--Euler balance equations, and the
position-, velocity-, and acceleration-level lower-pair constraints in
Eqs.~\eqref{eq:trans-collocation}--\eqref{eq:fullva} hold pointwise on the
trajectory tube. Let $Z_G$ be the lifted stage tuple obtained from the
three-stage Gauss collocation solution of the reduced equation
$\dot y=f(y,t)$, using the same Gauss abscissae, transition input, local Lie
chart, and endpoint-output convention that will be used in the implemented
residual bridge. Then the mathematical accepted stage residual
$R_{\mathrm{G6FVA}}$ in Eq.~\eqref{eq:residual}, evaluated in exact
arithmetic at $Z_G$, satisfies the ideal row identities exactly in all six
row families. Equivalently, its mathematical residual defect is zero, and in
particular it has the \(O(h^7)\) scale required by
Lemma~\ref{lem:stage-residual-defect} before implementation binding is
considered.
This statement concerns the ideal target-equation residual only; it is not a
statement about the row-scaled implemented residual \(F_{A,h}\).  The
implemented \(132\)-row bridge consumed by the theorem is supplied later,
after P3/P4/P5 bind row order, row scaling, the AD/formula path, and the
direct Newton--Euler substitution.
Thus the ideal zero statement is not used as a shortcut around the implemented
PC2 discharge: P3 fixes the formula path and row scaling, P4 supplies the
96-row non-dynamic bound, and P5 supplies the 36 dynamic zero rows consumed by
Lemma~\ref{lem:full-132-row-residual-bridge}.
The object \(Z_G\) in this lemma is the same lifted Gauss object consumed by
that bridge; it is not replaced by an accepted Newton root, an inexact Newton
iterate, or an endpoint-closed postprocessed state.
\end{lemma}

\begin{proof}
The reduced Gauss stage solution satisfies the collocation equations for
$\dot y=f(y,t)$. Applying the smooth lift $\mathcal L$ gives stage positions,
orientations, velocities, accelerations, and multipliers on the same smooth
constraint manifold. The translational and rotational position rows
\eqref{eq:trans-collocation}--\eqref{eq:rot-collocation} are exactly the
integral form of the lifted kinematic identities
$\dot r=v$ and $\dot Q=Q\widehat\omega$ written in the local Lie chart; the
velocity rows \eqref{eq:v-collocation}--\eqref{eq:w-collocation} are the
corresponding integral form of $\dot v=a$ and $\dot\omega=\alpha$. The
Newton--Euler rows \eqref{eq:trans-balance}--\eqref{eq:rot-balance} vanish
because the lift supplies the stage accelerations and multipliers satisfying
force and moment balance. Finally, the FullVA lower-pair rows
\eqref{eq:fullva} vanish because the same lift lies on the position
constraint manifold and its first two time derivatives. Hence each row
family of $R_{\mathrm{G6FVA}}$ is the lifted reduced Gauss stage equation or
a pointwise algebraic consequence of the lifted constrained dynamics. At this
ideal mathematical level the residual defect is therefore zero, which is
stronger than \(O(h^7)\); the subsequent implemented bridge is still needed
to bind this row identity to the residual path actually differentiated and
solved.
\end{proof}

Lemma~\ref{lem:fullva-stage-lift} proves only the ideal mathematical lift
identity for the stage-row equations.  It does not itself discharge the
implemented PC2 residual bridge; that discharge is the later \(132\)-row
bridge under P3/P4/P5.  This zero statement is not the same object as the
endpoint truncation estimate
\(\|\Psi_h^G-\varphi_h\|\le C_Gh^7\): the former says that the lifted reduced
Gauss stage satisfies the fixed stage equations, while the latter compares the
endpoint reconstructed from that Gauss stage with the exact reduced flow.  The
two estimates are therefore consumed in different terms of the four-map local
defect decomposition.  The remaining implementation question is row identity:
the residual appearing in Eq.~\eqref{eq:residual} must be the same residual
differentiated and solved in the implemented residual path, with the same
132-row ordering, scaling, and chart conventions. The source-identity and
row-layout records fix this correspondence. With that correspondence fixed,
the proof certificates supply the accepted direct-route residual/Jacobian
binding used by the bridge: the non-dynamic certificate covers the 96
collocation and lower-pair rows on the smooth FullVA lift, the D5
direct-substitution certificate covers the 36 Newton--Euler weak-balance rows
on the lifted Gauss stage, and the AD-expanded implementation-path
derivative-cell certificate records the derivative identities used by that
binding for those 36 dynamic rows. These certificates are used only to bind the
implemented residual to the mathematical row identities in this proof; they do
not convert the method into a source-paper temporal finite-element residual
replacement.
Mathematical-lift versus implemented-residual split: the zero-defect statement
in Lemma~\ref{lem:fullva-stage-lift} is an ideal row identity after the
smooth FullVA lift has been defined. It is not, by itself, the theorem's
implemented residual-defect certificate. The theorem consumes the implemented
residual only after P3 fixes the formula/AD path and row ordering, P4 supplies
the 96-row non-dynamic bound, and P5 supplies the 36-row D5 direct-substitution
identity. Thus the ideal lift does not bypass P3/P4/P5, does not choose the
Newton branch, and does not discharge the compact-tube inverse or solver-scale
interfaces.
Ideal lift is not implementation discharge: the ideal zero statement is a
target-equation identity for the smooth lift, while the implemented
stage-residual input is the later P3/P4/P5 bridge on the code-facing
\(132\)-row residual path.
There is no contradiction between the ideal zero-defect lift and the later
\(O(h^7)\) implemented bridge.  The theorem deliberately consumes the
code-facing implemented \(O(h^7)\) bridge statement because the residual used
by the code must first be tied to
one formula path, row ordering, and row scaling.  The zero mathematical lift
identifies the target equations; the P3/P4/P5 bridge proves that the
implemented residual actually used in Newton has the required residual value
on \(Z_G\).  This bridge is also the only place where implementation-specific
row binding enters the local-defect proof; it is upstream of the endpoint
closure, inexact Newton, and global error recursion, and downstream of the
reduced Gauss truncation theorem.
Thus the D5 certificate is a stage-local identity on \(Z_G\). It is not evidence about residuals evaluated at the inexact Newton iterate
\(\widetilde Z_A\), finite-run diagnostic trajectories, or closed-loop
mechanism residual tables; those objects enter only through
Lemma~\ref{lem:inexact-newton}, numerical evidence, or the still-open residual-to-error boundary, respectively.
Equivalently, the static implementation certificate and source-identity
record identify one residual/Jacobian implementation path, while the
row-layout records fix the 132-row ordering used in the proof.

For the dynamic rows, the active proof route is the direct-substitution
certificate, with five mathematical sub-obligations feeding the residual
identity. The older symbolic-defect certificate remains a provenance record
for the primitive/Taylor route.  In the direct route, each of the 36 dynamic
rows is expanded into its body-specific proximal/distal force and torque
terms. The multiplier contribution is fixed by the D3 virtual-work identity,
including the six row-expanded lower-pair multiplier identities in the
implemented site ordering. The force and torque terms are evaluated on the
same smooth $C^7$ compact proof tube, the D1/D2 balance-identity certificate proves
the translational and rotational Newton--Euler balance identities, and the D6
row-ordering/scaling/AD binding certificate fixes the implemented row ordering and
scaling.  Substituting the smooth Gauss lift \(Z_G\) into the accepted dynamic
residual therefore leaves only identity terms and a lifted-stage remainder;
the direct-substitution step reduces that remainder row-locally to zero.
Hence the 36 dynamic rows satisfy the \(O(h^7)\) stage-residual condition
without invoking the primitive/Taylor lift route.

\begin{table}[htbp]
\centering
\caption{Direct-route Newton--Euler row obligations used by the P5 stage-residual bridge; the primitive symbolic-defect route is tracked separately.}
\label{tab:newton-euler-obligations}
\scriptsize
\setlength{\tabcolsep}{3pt}
\renewcommand{\arraystretch}{0.90}
\begin{tabular}{P{0.10\linewidth}P{0.25\linewidth}P{0.42\linewidth}P{0.15\linewidth}}
\toprule
ID & Row scope & Required mathematical statement & Proof state \\
\midrule
D1 & Translational balance rows &
After row binding and role anchoring, substituting the lifted Gauss stage
\(Z_G\) gives zero residual for the translational force-balance rows, hence
the required \(O(h^7)\) bound. & D1 balance identity closed; D5 direct residual certificate closed. \\
D2 & Rotational balance rows &
After row binding and role anchoring, substituting the lifted Gauss stage
\(Z_G\) gives zero residual for the body-frame angular balance, Lie-retraction
increment, and inertia rows, hence the required \(O(h^7)\) bound. &
D2 balance identity closed; D5 direct residual certificate closed. \\
D3 & Multiplier wrench rows &
The lower-pair multiplier reactions used in the position, velocity, and
acceleration rows induce the same generalized wrench as the dynamic rows. &
D3 row-expanded lower-pair multiplier identity closed; used in the D5 direct
residual certificate. \\
D4 & Force/friction smoothness &
All applied force and friction laws used by the accepted smooth branch are
$C^7$ on the trajectory tube and share the same FullVA lift. &
Smooth $C^7$ force/friction lift closed on the accepted compact proof tube;
sharp/small-stribeck branch remains diagnostic. \\
D5 & Dynamic-row defect rate &
Evaluating the assembled dynamic rows on the lifted Gauss stage gives a
uniform $O(h^7)$ residual bound. & Closed by direct-substitution residual
certificate. \\
D6 & Implemented row ordering &
The algebraic row order, scaling, and automatic-differentiation blocks in the
implemented residual path are symbolically equivalent to D1--D5. &
Row ordering, fixed residual row-scaling (the D6 unweighted scaling
convention), and accepted AD binding closed;
D5 direct residual certificate closed. \\
\bottomrule
\end{tabular}
\end{table}

\begin{table}[htbp]
\centering
\caption{Row-level target map for the Newton--Euler weak-balance rows.}
\label{tab:newton-euler-row-target-map}
\scriptsize
\setlength{\tabcolsep}{2pt}
\renewcommand{\arraystretch}{0.92}
\begin{tabular}{P{0.08\linewidth}P{0.12\linewidth}P{0.28\linewidth}P{0.28\linewidth}P{0.11\linewidth}}
\toprule
Stage $s$ & Global rows & Translational targets & Rotational targets & Residual role \\
\midrule
0 & 24--35 &
Rows 24--26 and 30--32: bodies 0--1 and components $x,y,z$
satisfy linear-momentum balance. &
Rows 27--29 and 33--35: bodies 0--1 and components $x,y,z$
satisfy angular-momentum balance. &
D1/D2 balance identities give zero dynamic residual on \(Z_G\). \\
1 & 68--79 &
Rows 68--70 and 74--76: bodies 0--1 and components $x,y,z$
satisfy linear-momentum balance. &
Rows 71--73 and 77--79: bodies 0--1 and components $x,y,z$
satisfy angular-momentum balance. &
D1/D2 balance identities give zero dynamic residual on \(Z_G\). \\
2 & 112--123 &
Rows 112--114 and 118--120: bodies 0--1 and components $x,y,z$
satisfy linear-momentum balance. &
Rows 115--117 and 121--123: bodies 0--1 and components $x,y,z$
satisfy angular-momentum balance. &
D1/D2 balance identities give zero dynamic residual on \(Z_G\). \\
\bottomrule
\end{tabular}
\end{table}

Table~\ref{tab:newton-euler-row-target-map} is the compact form of the
row-level target inventory used by the direct certificate.  Target inventory
complete; the active D5 direct-substitution certificate closes the dynamic
defect on the smooth Gauss lift.
For stage $s\in\{0,1,2\}$ and local dynamic-row
offset $\rho\in\{0,\ldots,11\}$, the implemented global row index is
Eq.~\eqref{eq:dynamic-row-index-map}:
\begin{equation}
\label{eq:dynamic-row-index-map}
  i_{\mathrm{row}}(s,\rho)=44s+24+\rho .
\end{equation}
Offsets 0--2 and 6--8 are the two-body translational balances; offsets
3--5 and 9--11 are the corresponding body-frame rotational balances. This
target map indexes the direct D5 certificate and the closed D1/D2/D3/D4/D6
sub-obligations for the 36 implemented dynamic rows.  The map alone is not
the proof of the defect bound; the proof is the row-local substitution of
\(Z_G\) into the accepted residual, recorded by the D5 direct-substitution
certificate.

The row map separates dynamic-row correspondence into two mathematical checks.
First, each target row is mapped to the corresponding translational or
rotational template, so the force, moment, multiplier, and inertia terms occur
in the intended balance equation. Second, after role-anchoring the proximal
and distal joint forces, multiplier torques, inertia terms, and external
loads, the scalar Newton--Euler templates and the body-specific symbolic
expansions simplify to zero difference. In compressed notation the checked
scalar templates are Eq.~\eqref{eq:ne-template-audit}:
\begin{equation}
\begin{aligned}
T^{\mathrm{tr}}_0 &= ma-mg-f_e-(F_0-F_1), &
T^{\mathrm{tr}}_1 &= ma-mg-f_e-F_1,\\
T^{\mathrm{rot}}_0 &= I-\tau_e-(P+D+A-B), &
T^{\mathrm{rot}}_1 &= I-\tau_e-(P+A),
\end{aligned}
\label{eq:ne-template-audit}
\end{equation}
where $I$ abbreviates the body-frame inertial term, $P,D$ are proximal and
distal moment-arm torques, and $A,B$ are the corresponding axis-torque
contributions. The runtime scalar template and the body-specific symbolic
template give identical expressions after this role anchoring for every
stage, body, and Cartesian component. This closes the template-level C2
subcheck and, together with the D6 row-ordering/scaling/AD certificate, closes the
implemented row-order and unweighted residual-scaling equivalence for every
dynamic row. The D1/D2 balance-identity certificate combines this template
equivalence with the D3/D4/D6 closure evidence and closes the implemented
Newton--Euler balance identities for all 36 rows.  This supplies the
\(E_{\mathrm{bal}}\) term in the direct-substitution proof.  The D5
direct-substitution certificate combines it with D3, D4, and D6 and certifies
zero dynamic residual on \(Z_G\) for every dynamic row; the older symbolic
defect certificate is retained only as row-level correspondence for the
primitive/Taylor route.
Representative row substitution: for any accepted stage \(s\), body \(b\),
and Cartesian or body-frame component \(c\), the direct D5 certificate reduces
one translational row and one rotational row on the lifted Gauss stage to
Eq.~\eqref{eq:d5-representative-row-substitution}:
\begin{equation}
\label{eq:d5-representative-row-substitution}
\begin{aligned}
r^{\rm tr}_{sbc}(Z_G)
&=
\bigl[
m_b a^G_{s,b}
-f^{\rm ext}_{s,b}(Z_G)
-G_{s,b}^{T}\lambda^G_s
\bigr]_c
=0,\\
r^{\rm rot}_{sbc}(Z_G)
&=
\bigl[
J_b\alpha^G_{s,b}
+\omega^G_{s,b}\times J_b\omega^G_{s,b}
-\tau^{\rm ext}_{s,b}(Z_G)
-H_{s,b}^{T}\lambda^G_s
\bigr]_c
=0 .
\end{aligned}
\end{equation}
The displayed zero is not an unexplained cancellation. The D1/D2
balance identities used in this substitution are the componentwise
Newton--Euler equations on the same lifted stage, shown in
Eq.~\eqref{eq:d5-balance-identities-on-lift}:
\begin{equation}
\label{eq:d5-balance-identities-on-lift}
\begin{aligned}
m_b a^G_{s,b}
&=f^{\rm ext}_{s,b}(Z_G)+G_{s,b}^{T}\lambda^G_s,\\
J_b\alpha^G_{s,b}
+\omega^G_{s,b}\times J_b\omega^G_{s,b}
&=\tau^{\rm ext}_{s,b}(Z_G)+H_{s,b}^{T}\lambda^G_s .
\end{aligned}
\end{equation}
Substitution of these two identities into the two residual definitions above
gives \(r^{\rm tr}_{sbc}(Z_G)=0\) and
\(r^{\rm rot}_{sbc}(Z_G)=0\) for each component.  D3 fixes the signs in
\(G_{s,b}^{T}\lambda^G_s\) and \(H_{s,b}^{T}\lambda^G_s\) from the
lower-pair virtual-work identities, while D6 states that the implemented
unweighted scalar rows are exactly these component residuals after the
accepted row ordering and scaling are applied.
Here \(G_{s,b}^{T}\lambda^G_s\) and
\(H_{s,b}^{T}\lambda^G_s\) denote the lower-pair multiplier force and torque
components bound to the FullVA constraint rows by D3; D1/D2 provide the
linear- and angular-momentum balance identities, D4 fixes the smooth
force/torque evaluation on the same lift, and D6 binds the displayed
component equations to the implemented row ordering and scaling. This display
is the representative direct substitution used for all 36 mapped dynamic
rows; it is not a primitive/Taylor subterm proof and does not promote the
separate primitive route.
The finite implementation check also links the 36 row slots to the full
formula-row map: the independent formula vector covers all 132 residual rows,
includes the \texttt{newton\_euler\_weak\_balance} family, and its AD Jacobian
matches the accepted residual-Jacobian path on three finite probes. This
finite row-slice check is retained only as implementation consistency
evidence; the proof of the row identity is the symbolic template reduction
above and the direct substitution of \(Z_G\).
In the supporting source package, this is the row-template instantiation and
template-level algebraic equivalence subcheck layer. It provides runtime
row-slice evidence only, not a symbolic identity proof.

\begin{table}[htbp]
\centering
\caption{Dynamic-row residual-identity table for the Newton--Euler rows.}
\label{tab:dynamic-proof-closure-matrix}
\scriptsize
\setlength{\tabcolsep}{3pt}
\renewcommand{\arraystretch}{0.92}
\begin{tabular}{P{0.12\linewidth}P{0.34\linewidth}P{0.38\linewidth}}
\toprule
Item & Verified input & Proof role and remaining non-use \\
\midrule
D1--D2 &
The 36 dynamic rows are present in the row-family-major residual; the
formula-row oracle matches the runtime residual and AD Jacobian on three
finite probes; and the D1/D2 balance-identity certificate proves the implemented
translational and rotational Newton--Euler identities. &
No remaining D1--D2-specific identity test; use the closed balance identities
inside the D5 lifted-stage $O(h^7)$ dynamic-defect proof. \\
D3 &
The lower-pair FullVA rows and the dynamic rows use the same multiplier
variables in the accepted source path, and the package records a 36-row
virtual-work sign-skeleton certificate covering the proximal/distal force and axis
torque transfer sites plus two template triple-product identities for point
and axis multiplier terms plus six row-expanded lower-pair multiplier
identities in the implemented site ordering. &
No remaining D3-specific closure test; the remaining dynamic proof work is
the D5 direct-substitution residual certificate. \\
D4 &
The smooth benchmark branch records bounded force, torque, prescribed-motion,
and friction inputs; the Brown--McPhee smooth-branch source structure,
regularized normal-load map, stage-local force lift, and compact proof-tube
$C^7$ derivative bound have been checked.
The sharp branch is reported separately as a limitation. &
No remaining D4-specific smoothness test; use the closed D4 certificate in
the D5 direct-substitution residual certificate. \\
D5 &
The 96 non-dynamic rows have a proof-scope certificate, and the 36 dynamic
Newton--Euler rows have a direct-substitution certificate on the smooth Gauss
lift \(Z_G\).  After substituting \(Z_G\), each dynamic residual reduces to the
closed D1/D2 balance identities, the D3 multiplier-wrench identity, the D4
smooth force lift, and the D6 row-ordering/scaling/AD binding. &
The direct stage-residual input is supplied by zero row-local dynamic
residual on \(Z_G\).  This route does
not use a \(P_{\mathrm{state}}\) actual PS3 lift, \(P_{\mathrm{acc}}\),
\(P_{\lambda}\), finite probes, or residual-to-error promotion. \\
D6 &
The source-path certificate identifies the residual/Jacobian path and the 132-row
layout; the D6 certificate independently proves the dynamic row ordering, unweighted
residual scaling, and accepted AD binding for all 36 Newton--Euler rows. &
No remaining D6-specific layout/scaling/AD test; use the closed D6 certificate in
the D5 direct-substitution residual certificate. \\
\bottomrule
\end{tabular}
\end{table}

Table~\ref{tab:dynamic-proof-closure-matrix} records the inputs to the dynamic
row proof.  The PC2 stage-residual condition is discharged by the D5
direct-substitution certificate:
for each Newton--Euler component row the accepted dynamic residual evaluated
on the smooth Gauss lift \(Z_G\) reduces to Eq.~\eqref{eq:d5-dynamic-residual-decomposition},
\begin{equation}
\label{eq:d5-dynamic-residual-decomposition}
R_{\mathrm{dyn}}(Z_G)
=E_{\mathrm{bal}}+E_{\lambda}+E_{\mathrm{smooth}}
 +E_{\mathrm{bind}}+E_{\mathrm{lift}},
\end{equation}
where $E_{\mathrm{bal}}=0$ is supplied by the closed D1/D2
balance-identity certificate, $E_{\lambda}=0$ by the D3 virtual-work multiplier
identity, $E_{\mathrm{smooth}}=0$ on the smooth lifted force/torque inputs by
the D4 force-lift certificate, and $E_{\mathrm{bind}}=0$ by the D6 row-ordering,
unweighted-scaling, and AD-binding certificate.  The remaining
\(E_{\mathrm{lift}}=0\) clause is the P5 direct-substitution identity: after
the smooth FullVA Gauss lift \(Z_G\) is inserted into the implemented dynamic
rows, the lifted acceleration, multiplier-wrench, force, and gyroscopic terms
match the continuous Newton--Euler balance row by row.  Hence all 36 dynamic
rows have zero residual on \(Z_G\), which is stronger than the \(O(h^7)\)
residual required by Lemma~\ref{lem:stage-residual-defect}.  This is an
identity in the 132-row implemented residual, not a cancellation argument over
the separate 162 primitive Taylor subterms.  The proof is non-circular: it
does not use a \(P_{\mathrm{state}}\) actual PS3 lift, \(P_{\mathrm{acc}}\),
\(P_{\lambda}\), finite numerical probes, or residual-to-error promotion.

\begin{lemma}[Same-object instantiation for the direct D5 lift]
\label{lem:d5-same-object-lift-instantiation}
For each stage \(s\), body \(b\), component \(c\), and translational or
rotational row type, the \(E_{\mathrm{lift}}\) term in
Eq.~\eqref{eq:d5-dynamic-residual-decomposition} vanishes by same-object
substitution on the accepted direct route.  It is not an additional Taylor
estimate, not a primitive-route assumption, and not inferred from finite-grid
or residual-to-error evidence.
\end{lemma}

\begin{proof}
Fix the accepted step data, the single lifted object
\begin{equation}
  Z_G=(Q^G,V^G,A^G,\Lambda^G),
  \label{eq:d5-same-object-lift}
\end{equation}
and a single implemented scalar row index in the 132-row residual.  The row in
the residual is then evaluated at the object in
Eq.~\eqref{eq:d5-same-object-lift}, with the same stage time, the same FullVA acceleration entry,
the same lower-pair multiplier wrench, the same smooth force/torque lift, and
the same unweighted row ordering and scaling used by the executable residual.
D1--D2 give the Newton--Euler balance identity on this very lift; D3 identifies
the lower-pair multiplier wrench in the same scalar row; D4 identifies the
force and torque maps on the same smooth branch; and D6 identifies the row
ordering, scaling, and AD-bound formula row with the implemented residual
component.  After those same-object identifications, the implemented row is
literally the left-hand side of the continuous Newton--Euler balance displayed
in Eq.~\eqref{eq:d5-balance-identities-on-lift}, so the lift mismatch term
vanishes.  If any ingredient were taken from a different lift, row scaling,
time slice, multiplier convention, or force evaluation, this same-object
identity would fail and the argument would fall back to the separate
primitive/Taylor route rather than to the accepted direct PC2 proof.
\end{proof}

Table~\ref{tab:d5-taylor-certificate-target} is retained as a separate
primitive-route specification for a more granular primitive/Taylor
verification route.  The theorem-use rule is deliberately one-way and is
summarized as follows: the table records a conditional D5 primitive/Taylor
certificate schema, while the active direct residual bridge closes PC2
independently.  Thus the primitive/Taylor route remains open and is non-active in the theorem:
its unproved inputs define a separate dynamic-row replacement-certificate
route, not an additional premise of the direct residual-bridge theorem.
Equivalently, the primitive/Taylor route is non-active in the theorem and its
obligations are primitive-route inputs only; this route does not discharge the
PC2 stage-residual condition.  Its conditional coverage is nevertheless more
specific than an unstructured separate-route open gap: all \(162/162\)
Newton--Euler primitive subterms have assigned Taylor or mean-value reduction
templates and primitive dependencies.  This is template coverage only;
the certified count of actual primitive Taylor bounds is \(0/162\), and only the
compact-tube constant primitive is discharged.  The remaining non-active
primitive-route inputs are the five open same-object lift/reduction estimates
\(P_{\mathrm{state}},P_{\mathrm{acc}},P_{\lambda},
P_{\mathrm{geom}},P_{\mathrm{gyro}}\) needed to turn those templates into
actual \(O(h^7)\) subterm bounds; those bounds are not part of the accepted
theorem route.

\begin{table}[htbp]
\centering
\caption{D5 primitive/Taylor route specification for the
Newton--Euler rows; not invoked as a premise of
Theorem~\ref{thm:g6fullva-order}.}
\label{tab:d5-taylor-certificate-target}
\scriptsize
\setlength{\tabcolsep}{3pt}
\renewcommand{\arraystretch}{0.92}
\begin{tabular}{P{0.20\linewidth}P{0.48\linewidth}P{0.20\linewidth}}
\toprule
Target & Primitive-route target statement & Proof state \\
\midrule
Stage-lift substitution &
Construct the smooth FullVA Gauss lift \(Z_G\) in the accepted Lie chart and
substitute the lifted stage pose, velocity, acceleration, and multiplier
variables into the implemented dynamic-row expressions. &
Separate primitive-route target only; not invoked by the theorem.  The direct
residual bridge supplies the PC2 input. \\
Translational balance rows &
For each body, stage, and Cartesian component, reduce the expanded row
\(m_b\widehat a_{i,b}-f_b(\widehat q_i,\widehat v_i,t_i)
-G_{i,b}^{T}\widehat\lambda_i\)
to the continuous Newton balance on the lift plus a Taylor remainder bounded
by \(C h^7\). &
Separate primitive-route target only; not invoked by the theorem.  The direct
residual bridge supplies the PC2 input. \\
Rotational balance rows &
For each body, stage, and body-frame component, reduce
\(J_b\widehat\alpha_{i,b}
+\widehat\omega_{i,b}\times J_b\widehat\omega_{i,b}
-\tau_b(\widehat q_i,\widehat v_i,t_i)
-H_{i,b}^{T}\widehat\lambda_i\)
to the lifted angular-momentum balance plus a Taylor remainder bounded by
\(C h^7\). &
Separate primitive-route target only; not invoked by the theorem.  The direct
residual bridge supplies the PC2 input. \\
Closed support certificates &
Use D3 multiplier-wrench consistency, D4 smooth force lift, and D6 row
binding as zero or bounded inputs in the above reductions without changing
the row scaling. &
D3/D4/D6 support closed only; \(P_{\mathrm{state}}\),
\(P_{\mathrm{acc}}\), \(P_{\lambda}\), \(P_{\mathrm{geom}}\), and
\(P_{\mathrm{gyro}}\) remain open; no primitive-route Taylor bounds are
certified. \\
Uniformity &
Show the same constant family is bounded on the compact smooth proof tube,
uniformly for all 36 dynamic component rows and all sufficiently small \(h\). &
Separate primitive-route target only; not invoked by the theorem.  The direct
residual bridge supplies the PC2 input. \\
\bottomrule
\end{tabular}
\end{table}

\paragraph{Reference Taylor boundary for D5}
The Taylor local-error argument in the Lie-group BDF reference
\cite{wieloch2021bdf} serves here only as a proof-organization model: fix the
Lie-group reconstruction, insert the smooth exact solution into the already
defined discrete equations, estimate local defects, and only then propagate
the constrained error recursion.  It does not by itself supply any primitive
FullVA Taylor bound for the present Gauss-stage residual.  In particular, the
reference-paper insertion step is typed by the BLieDF multistep equations
themselves; its analogue here is not a free Taylor expansion of separate
physical subterms, but the same-object evaluation
\(F_{A,h}(Z_G;y)\) after the FullVA stage vector, row ordering, scaling, and
proof norm have been fixed.  A primitive D5 subterm estimate can enter this
analogue only after it has been transported to that same fixed residual map
and after its root lift inputs are proved independently on the same compact
tube.
By contrast, the
BLieDF acceleration and multiplier estimates are tied to a multistep
increment representation, hidden-constraint differences, and a coupled BDF
recursion, whereas the separate D5 primitive route here would need a
non-circular unweighted bound on the acceleration components that enter the
implemented Newton--Euler FullVA rows after the lower-pair constraints, Lie
chart, row ordering, and residual scaling have already been fixed.  Thus the
reference proof architecture serves only to order the local obligations for
\(P_{\mathrm{state}}\), \(P_{\mathrm{acc}}\), \(P_{\lambda}\),
\(P_{\mathrm{geom}}\), and \(P_{\mathrm{gyro}}\). It does not certify, import,
or bound any of the 162 D5 primitive Taylor subterms unless the required
\(P_{\mathrm{state}}\) and \(P_{\mathrm{acc}}\) lift inputs are proved inside
this stage-residual setting.  No Taylor term is therefore certified as proved
from the reference paper or from proof-architecture similarity.

The manuscript statement is deliberately narrower than the full
primitive/Taylor route. The
Newton--Euler symbolic-defect certificate is used here as a provenance
record for the row-expansion and runtime-equivalence checks and as
the 36-row inventory feeding Lemma~\ref{lem:stage-residual-defect}. It is not
the active dynamic-row closure route.  For the direct residual-bridge route
used by PC2,
the dynamic-row residual input is supplied by the D5 direct-substitution
certificate: every expanded Newton--Euler component row has zero residual
after substituting the smooth Gauss lift \(Z_G\). The symbolic-certificate C3
primitive/Taylor defect-bound lane remains open and is not the route consumed
by the PC2 condition here. The reproducibility package includes both this direct
certificate and the older row-local D5 primitive-route record.  The latter lists
the closed D1/D2/D3/D4/D6 inputs and the one open lifted-stage dynamic term
for the alternative primitive/Taylor route; it is not the route that discharges
the PC2 stage-residual condition.
The D5 primitive inventory, recorded as a D5 term-budget record, further
decomposes that open lifted-stage term into \(162\) Taylor subterms: four
translational subterms per translational row and five rotational subterms per
rotational row.  The decomposition lists the available regularity and
row-binding inputs, and it supplies conditional reduction rules for all
\(162/162\) subterms.  This is a conditional finite implication map rather than
an active theorem premise: if the primitive lift inputs below are separately
proved uniformly on the same tube, the table identifies which
Taylor/mean-value estimate closes each subterm.  At present, however, this is
template coverage only: the certified count of actual primitive Taylor bounds is \(0/162\)
on that route, and each subterm still needs an independent uniform \(O(h^7)\)
Taylor bound before the primitive/Taylor route can close.
The obstruction is coupled, not cosmetic: all \(162\) conditional templates
are anti-circular reduction templates, but none is a theorem input until its
blocking primitive is proved; in particular the \(36\) linear acceleration
subterms require unweighted acceleration \(O(h^7)\), so the available
\(h\)-weighted acceleration control cannot promote those terms.
A companion primitive-bound reduction groups the 162 subterms into six primitive obligations:
stage state, acceleration, multiplier, dynamic-row geometry, gyroscopic
bilinear, and compact-tube uniformity bounds.  This is an implication map,
not a proof of the primitive bounds: all 162 subterms have a recorded reduction
rule, but it does not certify the lift-rate primitive bounds.  The
compact-tube constants primitive is discharged below, while no induced Taylor bounds are
certified on the primitive/Taylor route.  Only the compact-tube constant
obligation \(P_{\mathrm{tube}}\) is discharged below; the five lift and
bilinear primitives remain open.
The primitive-obligation closure plan is therefore explicit: prove
\(P_{\mathrm{tube}}\), \(P_{\mathrm{state}}\), \(P_{\mathrm{acc}}\),
\(P_{\lambda}\), \(P_{\mathrm{geom}}\), and \(P_{\mathrm{gyro}}\) uniformly on
the same smooth proof tube.  The first supplies compact-tube constants; the
next three bound stage state, acceleration, and multiplier lifts; the last two
bound dynamic-row geometry and the gyroscopic bilinear term.  The compact-tube
constant obligation \(P_{\mathrm{tube}}\) is closed by
Lemma~\ref{lem:d5-compact-tube-constants}; the five lift and bilinear
obligations remain open.
Primitive dependency graph for the separate Taylor route: \(P_{\mathrm{state}}\)
and \(P_{\mathrm{acc}}\) are root lift obligations.  The former still lacks
the actual PS3 state-lift instantiation, and the latter has only the
\(h\)-weighted acceleration control, not the unweighted \(O(h^7)\)
acceleration lift needed by the primitive route.  The multiplier obligation
\(P_{\lambda}\) has its rate-propagation mechanism available only as a
conditional implication from those state and acceleration inputs.  The geometry obligation
\(P_{\mathrm{geom}}\) is downstream of \(P_{\mathrm{state}}\) and
\(P_{\lambda}\), and the gyroscopic obligation \(P_{\mathrm{gyro}}\) is
downstream of the angular-velocity component of \(P_{\mathrm{state}}\).
Thus the two row-level reductions above are conditional sufficiency statements inside
the Taylor route, and the admissible dependency order is
Eq.~\eqref{eq:d5-primitive-dependency-graph},
\begin{equation}
\label{eq:d5-primitive-dependency-graph}
\begin{aligned}
  P_{\mathrm{state}},\ P_{\mathrm{acc}}
  &\Longrightarrow P_{\lambda},\\
  P_{\mathrm{state}},\ P_{\lambda}
  &\Longrightarrow P_{\mathrm{geom}},\\
  P_{\mathrm{state}}
  &\Longrightarrow P_{\mathrm{gyro}} .
\end{aligned}
\end{equation}
This dependency graph is used together with the already closed
\(P_{\mathrm{tube}}\).
The current primitive-route conditions are therefore:
\begin{table}[htbp]
\centering
\small
\setlength{\tabcolsep}{2pt}
\caption{D5 primitive/Taylor retained-interface map; separate primitive-route
specification, not a
premise of Theorem~\ref{thm:g6fullva-order}.}
\label{tab:d5-primitive-blocker-ledger}
\begin{tabular}{@{}P{0.13\linewidth}P{0.14\linewidth}P{0.11\linewidth}P{0.28\linewidth}P{0.24\linewidth}@{}}
\toprule
Primitive & Primitive-route condition & Affected Taylor rows & Required primitive input & Scope limitation \\
\midrule
\(P_{\mathrm{tube}}\) & closed & 162 template constants &
Compact proof-tube constants only. &
Does not by itself prove any primitive lift rate or term-level Taylor bound. \\
\(P_{\mathrm{state}}\) & root not established & 126 &
Non-circular actual PS3 state lift from the non-dynamic rows. &
The direct full-residual route corollary does not close the primitive route. \\
\(P_{\mathrm{acc}}\) & root not established & 36 &
Unweighted \(O(h^7)\) acceleration lift; current
\(h\|\delta A\|=O(h^7)\) gives only unweighted \(O(h^6)\). &
Cannot be read from dynamic balance or from the \(h\)-weighted control. \\
\(P_{\lambda}\) & conditional on roots & 72 &
Actual \(P_{\mathrm{state}}\) and \(P_{\mathrm{acc}}\) inputs before PL4
propagates multiplier rates. &
Closed PL4 propagation is conditional, not an actual multiplier lift. \\
\(P_{\mathrm{geom}}\) & conditional downstream & 36 &
Actual \(P_{\mathrm{state}}\) and \(P_{\lambda}\) inputs. &
Geometry reduction alone is not an actual term closure. \\
\(P_{\mathrm{gyro}}\) & conditional downstream & 18 &
Angular-rate part of the actual \(P_{\mathrm{state}}\) lift. &
Bilinear reduction alone is not an actual term closure. \\
\bottomrule
\end{tabular}
\end{table}
This table is a primitive-route obstruction map: no actual D5 primitive
Taylor bound from this route is invoked by the theorem, and the separate
primitive-route implication remains a separate sufficient route for the
dynamic-row PC2 input.
Here ``conditional'' means that the local reduction lemma is closed under
the listed inputs, while the actual primitive \(O(h^7)\) bound is not
established on the accepted theorem branch.
The two root targets are precise.  To close \(P_{\mathrm{state}}\) on this
route one must prove a non-circular actual PS3 statement: using only the
non-dynamic 96-row defect, the compact-tube PS2 inverse, and an acceleration
input independent of the state perturbation being estimated, obtain
\(\|\delta S\|\le C_Sh^7\) uniformly on the tube.  To close
\(P_{\mathrm{acc}}\), one must prove the unweighted estimate
\(\|\delta A\|\le C_Ah^7\) in the acceleration variables that enter the
Newton--Euler rows.  The currently recorded velocity-collocation algebra gives
only \(h\|\delta A\|=O(h^7)\), equivalently \(\|\delta A\|=O(h^6)\), so a
Taylor-route proof would need either \(O(h^8)\) velocity/velocity-row inputs,
an independent non-dynamic inverse controlling acceleration directly, or an
explicit \(P_{\mathrm{acc}}\) assumption.  Once these two root targets are
proved, PL4 supplies \(P_{\lambda}\), and the row-level geometry and
gyroscopic corollaries below supply \(P_{\mathrm{geom}}\) and
\(P_{\mathrm{gyro}}\).  Without those inputs, the primitive/Taylor route remains a
separate sufficient-route specification for the dynamic-block residual value,
not the route that discharges PC2 in the theorem.

\begin{lemma}[Compact proof-tube constants for D5]
\label{lem:d5-compact-tube-constants}
Under Assumption~\ref{ass:regularity}, the accepted smooth branch has a
compact proof tube \(K\) and a step threshold \(h_0>0\) such that, for
\(0<h<h_0\), the exact reduced trajectory, the lifted Gauss stages, and the
local accepted Newton stage variables remain in \(K\).  On this tube all
force, torque, friction, constraint-geometry, and gyroscopic maps that enter
the 36 Newton--Euler dynamic rows have uniformly bounded derivatives through
the order needed by the D5 Taylor reductions.  Hence the primitive obligation
\(P_{\mathrm{tube}}\) holds.
\end{lemma}

\begin{proof}
The image of the exact smooth reduced trajectory on the fixed interval is
compact.  Assumption~\ref{ass:regularity} gives a smooth \(C^7\) reduced
vector field, a smooth FullVA lift on the same tube, and the
Gauss-predictor inverse and derivative bounds consumed by
Lemma~\ref{lem:stage-residual-defect}.  That lemma, not an isolated-root
premise in this compact-tube argument, supplies the local accepted stage root
\(Z_A(h)\) for sufficiently small \(h\) on the selected branch.  To make the
tube selection explicit, take the compact image of the exact trajectory and
its smooth Gauss-stage lift for \(0\le c_i\le 1\), choose a finite chart cover
whose closures stay inside the regularity domain, and let \(d_K>0\) be the
minimum distance from those closures to the chart, rank, and smooth-force
boundaries.  The Gauss lift is continuous in \(h\), while the accepted-stage
bound from Lemma~\ref{lem:stage-residual-defect} gives
\(\|Z_A-Z_G\|\le C_Zh^7\); hence \(h_0\) is chosen so that both the
Gauss-stage displacement from the exact compact image and \(C_Zh^7\) are
less than \(d_K/3\).  The closed \(d_K/3\)-neighborhood inside the finite
cover is then one compact tube \(K\) containing the exact trajectory, the
reduced Gauss stage reconstruction, and the lifted FullVA stage variables for
every \(0<h<h_0\) on the selected branch.

The D4 smooth-branch certificate supplies \(C^7\) force and friction maps on
this fixed tube, and the constraint geometry, moment-arm maps, and the
bilinear gyroscopic map are smooth compositions of the same lifted pose and
velocity variables.  Since the list of row maps and D5 primitive subterms is
finite, the maxima of the required derivative suprema over \(K\) are finite
and row-independent.  This proves only the uniform-constant primitive
\(P_{\mathrm{tube}}\); it does not prove the state, acceleration,
multiplier, geometry, or gyroscopic lift rates.
\end{proof}

\begin{lemma}[Non-dynamic stage-map definition for \(P_{\mathrm{state}}\)]
\label{lem:d5-p-state-map-definition}
Let Eq.~\eqref{eq:d5-p-state-stage-blocks} collect the accepted stage pose and
velocity block and the auxiliary acceleration block for all bodies in the
accepted Lie chart:
\begin{equation}
\label{eq:d5-p-state-stage-blocks}
  S=(r_i,\eta_i,v_i,\omega_i)_{i=1}^3,\qquad
  A=(a_i,\alpha_i)_{i=1}^3
\end{equation}
Define \(\mathcal N_h^{\mathrm{nd}}(S,A)\) as the ordered vector of the five
non-dynamic row families in Eq.~\eqref{eq:d5-p-state-nondynamic-map},
\begin{equation}
\label{eq:d5-p-state-nondynamic-map}
\begin{aligned}
&r_i-r_n-h\sum_j A_{ij}v_j,\qquad
  \eta_i-h\sum_j A_{ij}J_r^{-1}(\eta_j)\omega_j,\\
&v_i-v_n-h\sum_j A_{ij}a_j,\qquad
  \omega_i-\omega_n-h\sum_j A_{ij}\alpha_j,\\
&\Phi(q_i,t_i)-g_i^{(0)},\quad
  \Phi_q(q_i,t_i)v_i-g_i^{(1)},\quad
  \Phi_q(q_i,t_i)a_i+\dot\Phi_q(q_i,v_i,t_i)v_i-g_i^{(2)} .
\end{aligned}
\end{equation}
This map has \(96\) scalar rows in the cylindrical-chain stage layout: \(18\)
translational position-collocation rows, \(18\) rotational Lie-collocation
rows, \(18\) translational velocity-collocation rows, \(18\) angular
velocity-collocation rows, and \(24\) lower-pair FullVA rows.  It is exactly
the non-dynamic map used by the 96-row certificate. This closes only the PS1
map-definition subproof for \(P_{\mathrm{state}}\). It does not assert a
bounded inverse, an inf-sup estimate, an \(O(h^7)\) state-lift rate, or PC2.
\end{lemma}

\begin{proof}
The row families listed above are the non-dynamic rows in
Eqs.~\eqref{eq:trans-collocation}, \eqref{eq:rot-collocation},
\eqref{eq:v-collocation}, \eqref{eq:w-collocation}, and
\eqref{eq:fullva}.  Their row counts are the same counts recorded in the
kinematic row-defect certificate: four \(18\)-row collocation families and one
\(24\)-row lower-pair FullVA family. The multiplier block is absent from
\(\mathcal N_h^{\mathrm{nd}}\), and the Newton--Euler balance family is
excluded. Therefore the map interface needed by PS1 is fixed in the accepted
Lie chart and is bound to the same \(96\) certified rows. No estimate for
\((S,A)\) follows from this definition alone; a separate local inverse or
inf-sup proof is still needed for PS2 before \(P_{\mathrm{state}}\) can be
closed.
\end{proof}

The next primitive obligation, \(P_{\mathrm{state}}\), is deliberately kept
open.  The 96-row non-dynamic certificate proves that the kinematic and
lower-pair residual rows vanish on the smooth FullVA lift, but a residual-row
identity is not by itself a variable-lift estimate.  The aggregate PS2 inverse
needed by this route is proved below as a compact-tube estimate, but closing
\(P_{\mathrm{state}}\) still requires an actual PS3 instantiation: the
non-dynamic residual defect and the \(h\)-weighted acceleration input must be
combined to prove that the stage pose, translational velocity, and angular
velocity variables differ from the smooth lifted Gauss variables by
\(O(h^7)\).  This proof must not use
Lemma~\ref{lem:stage-residual-defect}, because that lemma already assumes the
dynamic residual defect that \(P_{\mathrm{state}}\) is meant to help prove.

\begin{lemma}[Anti-circular route for \(P_{\mathrm{state}}\)]
\label{lem:d5-p-state-anticircularity}
The separate PS2--PS3 proof route for \(P_{\mathrm{state}}\) may use only the
96-row non-dynamic residual certificate, the stage-map definition
\(\mathcal N_h^{\mathrm{nd}}\), and compact proof-tube smoothness constants.
The D5 dynamic residual defect is disallowed as an input to
\(P_{\mathrm{state}}\).
% The D5 dynamic residual defect is disallowed as an input to P_state.
Likewise Lemma~\ref{lem:stage-residual-defect} and residual identities alone are
disallowed as inputs for converting row defects into state-variable lift
rates. This closes the PS4 anti-circularity subproof, but it does not close
PS2, PS3, \(P_{\mathrm{state}}\), any induced D5 Taylor bound, or PC2.
\end{lemma}

\begin{lemma}[Weighted PS2 target for \(P_{\mathrm{state}}\)]
\label{lem:d5-p-state-ps2-weighted-target}
The PS2 inverse statement for \(P_{\mathrm{state}}\) is a weighted PS2 target,
not a literal inverse for all state and acceleration variables.  In the
cylindrical-chain chart, the state block
\(S=(r_i,\eta_i,v_i,\omega_i)_{i=1}^3\) has \(72\) scalar coordinates, the
auxiliary acceleration block \(A=(a_i,\alpha_i)_{i=1}^3\) has \(36\), and
\(\mathcal N_h^{\mathrm{nd}}\) has \(96\) residual rows.  Thus the correct
finite-dimensional target is the state-block estimate in
Eq.~\eqref{eq:d5-p-state-weighted-ps2-target},
\begin{equation}
\label{eq:d5-p-state-weighted-ps2-target}
\|\delta S\|\le C_{\mathrm{PS2}}
\bigl(\|\delta \mathcal N_h^{\mathrm{nd}}\|+h\|\delta A\|\bigr),
\end{equation}
with acceleration treated as a parameter in the \(h\)-weighted norm
\(h\|\delta A\|\).  This specification fixes the PS2 target, but it does not
close PS2, does not close PS3, does not prove the uniform constant
\(C_{\mathrm{PS2}}\), and does not close \(P_{\mathrm{state}}\) or PC2.
\end{lemma}

\begin{lemma}[Kinematic-block certificate for the weighted PS2 route]
\label{lem:d5-p-state-ps2-kinematic-block}
For one scalar position--velocity channel, the three-stage Gauss kinematic
subblock has the linearized form in
Eq.~\eqref{eq:d5-p-state-gauss-kinematic-linearization},
\begin{equation}
\label{eq:d5-p-state-gauss-kinematic-linearization}
  \delta R_q=\delta q-hA_G\delta v,\qquad
  \delta R_v=\delta v-hA_G\delta a ,
\end{equation}
where \(A_G\) is the fixed three-stage Gauss matrix.  Therefore
Eq.~\eqref{eq:d5-p-state-gauss-kinematic-inverse} gives the corresponding
block inverse:
\begin{equation}
\label{eq:d5-p-state-gauss-kinematic-inverse}
\begin{bmatrix}\delta q\\ \delta v\end{bmatrix}
=
\begin{bmatrix}
I & hA_G & hA_G^2\\
0 & I & A_G
\end{bmatrix}
\begin{bmatrix}
\delta R_q\\ \delta R_v\\ h\delta a
\end{bmatrix}.
\end{equation}
For \(0<h\le 0.04\) this matrix has a bounded Euclidean operator norm
depending only on \(A_G\) and the step threshold.  Kronecker lifting gives the
same bound for all translational channels and for the angular-velocity
collocation channels after the accepted chart norm equivalence is applied.
Thus the 72 kinematic rows already provide a concrete analytic subblock for
the weighted PS2 estimate; the 24 lower-pair rows are surplus for this
subblock.  This certificate records the algebraic constant needed by the PS2
route, but the full nonlinear PS2 proof still has to bind the implemented
rotational nonlinear row and the 72-to-96 residual-row injection to this
template.  Consequently it does not by itself close PS2, PS3,
\(P_{\mathrm{state}}\), any induced Taylor bound, or PC2.
\end{lemma}

\begin{lemma}[SO(3) chart binding for the weighted PS2 route]
\label{lem:d5-p-state-ps2-lie-chart-binding}
On the compact proof tube, restrict the accepted rotational chart to
\(\|\eta\|\le \rho\) with \(\rho=0.5<\pi\).  The SO(3) right Jacobian and its
inverse obey the uniform bounds in
Eq.~\eqref{eq:d5-p-state-lie-chart-bounds},
\begin{equation}
\label{eq:d5-p-state-lie-chart-bounds}
\|J_r(\eta)\|_2\le 1,\qquad
\|J_r(\eta)^{-1}\|_2\le
\frac{\rho}{2\sin(\rho/2)}=1.010493\ldots .
\end{equation}
Thus chart-coordinate perturbations and body-angular-velocity perturbations
may be measured in equivalent Euclidean norms with an \(h\)-independent
constant on the PS2 proof tube.  This records the pure Lie-chart
norm-equivalence input needed by the kinematic-block certificate.  It does not
supply the nonlinear mean-value estimate for
\(\eta_i-h\sum_j A_{ij}J_r^{-1}(\eta_j)\omega_j\), does not prove the
72-to-96 row injection/scaling estimate, and does not close PS2, PS3,
\(P_{\mathrm{state}}\), any induced Taylor bound, or PC2.
\end{lemma}

\begin{lemma}[72-to-96 row injection for the weighted PS2 route]
\label{lem:d5-p-state-ps2-row-injection}
Let \(\Pi_{\mathrm{kin}}\) select the four kinematic row families from the
accepted \(96\)-row non-dynamic residual \(\mathcal N_h^{\mathrm{nd}}\):
translational position collocation, rotational Lie collocation, translational
velocity collocation, and angular-velocity collocation.  These four families
have \(72\) scalar rows, and the remaining \(24\) scalar rows are the lower-pair
FullVA rows.  With the unweighted residual ordering used in the implemented
residual and in the 96-row certificate, the injection bound is
Eq.~\eqref{eq:d5-p-state-row-injection-bound}:
\begin{equation}
\label{eq:d5-p-state-row-injection-bound}
\|\Pi_{\mathrm{kin}}\|_2=1,\qquad
\|\Pi_{\mathrm{kin}}\delta\mathcal N_h^{\mathrm{nd}}\|
\le \|\delta\mathcal N_h^{\mathrm{nd}}\| .
\end{equation}
Consequently the unweighted non-dynamic residual norm dominates the kinematic
residual norm used by the PS2 kinematic block.  This closes only the
72-to-96 row injection/scaling subclaim.  It does not supply the nonlinear
mean-value estimate for the rotational row, does not close PS2, PS3,
\(P_{\mathrm{state}}\), any induced Taylor bound, or PC2.
\end{lemma}

\begin{lemma}[Nonlinear rotational-row binding for the weighted PS2 route]
\label{lem:d5-p-state-ps2-nonlinear-binding}
Let the rotational chart velocity map be
\begin{equation}
\label{eq:d5-p-state-nonlinear-f-map}
F(\eta,\omega)=J_r(\eta)^{-1}\omega .
\end{equation}
On the accepted compact proof tube, the chart radius from
Lemma~\ref{lem:d5-p-state-ps2-lie-chart-binding} keeps \(J_r^{-1}\) smooth,
and Lemma~\ref{lem:d5-compact-tube-constants} supplies for the map in
Eq.~\eqref{eq:d5-p-state-nonlinear-f-map} a finite Lipschitz
constant \(L_F\) for \(F\).  Hence, for two stage states in the tube,
\begin{equation}
\label{eq:d5-p-state-nonlinear-lipschitz}
\|\Delta F\|\le L_F\bigl(\|\Delta\eta\|+\|\Delta\omega\|\bigr).
\end{equation}
For the implemented rotational row
\begin{equation}
\label{eq:d5-p-state-rot-row}
R_{\eta,i}=\eta_i-h\sum_j A_{ij}J_r(\eta_j)^{-1}\omega_j ,
\end{equation}
the nonlinear part therefore enters as an \(h\|A_G\|_2L_F\) perturbation of
the triangular Gauss kinematic block by Eqs.~\eqref{eq:d5-p-state-nonlinear-lipschitz}
and~\eqref{eq:d5-p-state-rot-row}.  After shrinking the asymptotic step
threshold, if necessary, so that \(h\|A_G\|_2L_F\le 1/2\), this perturbation is
absorbed on the left and gives the compact-tube mean-value estimate
\begin{equation}
\label{eq:d5-p-state-rot-bound}
\|\Delta\eta\|+\|\Delta\omega\|
\le C_{\mathrm{rot}}
\bigl(\|\Delta R_\eta\|+\|\Delta R_\omega\|+h\|\Delta\alpha\|\bigr).
\end{equation}
Equation~\eqref{eq:d5-p-state-rot-bound} is the small-step absorption that
binds the implemented rotational Lie row to the
weighted PS2 kinematic template without using the D5 dynamic residual theorem
or Lemma~\ref{lem:stage-residual-defect}.  It closes the rotational nonlinear
mean-value subclaim.  It does not by itself promote the aggregate PS2 inverse,
does not close PS3, \(P_{\mathrm{state}}\), any induced Taylor bound, or PC2.
\end{lemma}

\begin{lemma}[Aggregate weighted PS2 promotion for \(P_{\mathrm{state}}\)]
\label{lem:d5-p-state-ps2-aggregate-promotion}
There are constants \(h_0>0\) and \(C_{\mathrm{PS2}}>0\), depending only on
the compact proof tube and on the fixed Gauss matrix, such that for
\(0<h<h_0\) every pair of accepted non-dynamic stage states in the tube
satisfies the state-block estimate
\begin{equation}
\label{eq:d5-p-state-ps2-estimate}
  \|\delta S\|\le C_{\mathrm{PS2}}
  \bigl(\|\delta\mathcal N_h^{\mathrm{nd}}\|+h\|\delta A\|\bigr).
\end{equation}
Thus Eq.~\eqref{eq:d5-p-state-ps2-estimate} promotes the component certificates to a uniform weighted PS2 inverse for
the non-dynamic stage map.  The finite linearization probe is diagnostic only
and is not used in the theorem proof.  This lemma does not instantiate PS3,
does not prove an \(O(h^7)\) state lift, does not certify any D5 Taylor term
bound, and does not close \(P_{\mathrm{state}}\) or PC2.
\end{lemma}

\begin{proof}
Write the selected kinematic residual as
\(\Pi_{\mathrm{kin}}\mathcal N_h^{\mathrm{nd}}\), where
\(\Pi_{\mathrm{kin}}\) selects the four 18-row kinematic families.  The
72-to-96 selector has norm one by
Lemma~\ref{lem:d5-p-state-ps2-row-injection}, hence
\begin{equation}
\label{eq:d5-p-state-kin-selector-bound}
  \|\Pi_{\mathrm{kin}}\delta\mathcal N_h^{\mathrm{nd}}\|
  \le \|\delta\mathcal N_h^{\mathrm{nd}}\|.
\end{equation}
The lower-pair rows are therefore harmless surplus rows for the estimate in
Eq.~\eqref{eq:d5-p-state-kin-selector-bound}.

For the translational position--velocity channels, and for the
angular-velocity collocation channels after the chart norm equivalence of
Lemma~\ref{lem:d5-p-state-ps2-lie-chart-binding}, the linearized row map is
the fixed three-stage triangular Gauss block of
Lemma~\ref{lem:d5-p-state-ps2-kinematic-block}.  Its inverse is explicit and
has an \(h\)-independent bound for \(0<h\le h_0\).  Consequently those
channels satisfy
\begin{equation}
\label{eq:d5-p-state-trans-ang-bound}
  \|\delta r\|+\|\delta v\|+\|\delta\omega\|
  \le C_0\bigl(\|\delta R_r\|+\|\delta R_v\|
      +\|\delta R_\omega\|+h\|\delta a\|+h\|\delta\alpha\|\bigr)
\end{equation}
with \(C_0\) depending only on \(A_G\), the chart-equivalence constant, and
the chosen step threshold.  Eq.~\eqref{eq:d5-p-state-trans-ang-bound} is the
linear kinematic part that remains before the nonlinear rotational row is
absorbed.

It remains to bind the nonlinear rotational Lie row.  Let
\(F(\eta,\omega)=J_r(\eta)^{-1}\omega\).  Since
\(\|\eta\|\le\rho<\pi\) on the compact proof tube, \(F\) is smooth there.
For two stage vectors in the tube, Taylor's formula with integral remainder,
equivalently the finite-dimensional mean-value theorem, gives
\begin{align}
\label{eq:d5-p-state-f-taylor-remainder-identity}
  F(\eta+\delta\eta,\omega+\delta\omega)-F(\eta,\omega)
  =DF(\eta,\omega)[\delta\eta,\delta\omega]
   +\mathcal R_F,
\end{align}
where
\begin{equation}
\label{eq:d5-p-state-f-remainder-bound}
  \|\mathcal R_F\|
  \le C_F\bigl(\|\delta\eta\|+\|\delta\omega\|\bigr)^2
\end{equation}
Equation~\eqref{eq:d5-p-state-f-remainder-bound} is the quadratic remainder
bound used only inside the rotational-row mean-value argument.
for a tube-uniform second-derivative bound \(C_F\).  The linear mean-value
form also gives the Lipschitz estimate
\begin{equation}
\label{eq:d5-p-state-f-lipschitz-proof}
  \|\Delta F\|\le L_F(\|\Delta\eta\|+\|\Delta\omega\|).
\end{equation}
Thus the difference of the implemented rotational rows can be written as the
triangular Gauss rotational block plus the perturbation
\(-hA_G\Delta F\), with the Taylor and Lipschitz controls in
Eqs.~\eqref{eq:d5-p-state-f-taylor-remainder-identity}--\eqref{eq:d5-p-state-f-lipschitz-proof}.
Shrinking \(h_0\), and if necessary the fixed local
proof-tube radius, makes
\begin{equation}
\label{eq:d5-p-state-small-step-absorption}
  h\|A_G\|_2L_F\le \frac12
\end{equation}
on the tube.  The perturbation is then absorbed into the left-hand side of
the triangular inverse estimate, yielding
\begin{equation}
\label{eq:d5-p-state-rot-absorbed-bound}
  \|\delta\eta\|+\|\delta\omega\|
  \le C_1\bigl(\|\delta R_\eta\|+\|\delta R_\omega\|
       +h\|\delta\alpha\|\bigr),
\end{equation}
with \(C_1\) independent of \(h\).  This is a genuine Taylor/mean-value proof
on the compact tube: Eq.~\eqref{eq:d5-p-state-small-step-absorption} is the
small-step hypothesis and Eq.~\eqref{eq:d5-p-state-rot-absorbed-bound} is the
absorbed rotational estimate, not a finite-rank numerical probe.

Combining the translational, angular-velocity, and rotational estimates and
taking \(C_{\mathrm{PS2}}\) to dominate the finitely many constants gives
\begin{equation}
\label{eq:d5-p-state-ps2-combined-bound}
  \|\delta S\|\le C_{\mathrm{PS2}}
  \bigl(\|\Pi_{\mathrm{kin}}\delta\mathcal N_h^{\mathrm{nd}}\|
      +h\|\delta A\|\bigr)
  \le C_{\mathrm{PS2}}
  \bigl(\|\delta\mathcal N_h^{\mathrm{nd}}\|+h\|\delta A\|\bigr).
\end{equation}
The proof uses only the non-dynamic map, compact-tube smoothness, the fixed
Gauss block, and the row selector in Eq.~\eqref{eq:d5-p-state-ps2-combined-bound}.
It does not use the D5 dynamic residual
defect or Lemma~\ref{lem:stage-residual-defect}.  Therefore it closes the
uniform weighted PS2 inverse and no stronger claim.
\end{proof}

\begin{lemma}[Finite PS2 linearization diagnostic for \(P_{\mathrm{state}}\)]
\label{lem:d5-p-state-ps2-linearization-probe}
At the solved one-step accepted smooth cylindrical-chain stages for
\(h\in\{0.04,0.02,0.01\}\), the implemented residual Jacobian has the finite
weighted PS2 operator
\begin{equation}
\label{eq:d5-p-state-finite-ps2-operator}
M_h=\begin{bmatrix}
D_{S,A}\mathcal N_h^{\mathrm{nd}}\\[2pt]
0\quad hI_A
\end{bmatrix},
\qquad
(\delta S,\delta A)\mapsto
\bigl(D\mathcal N_h^{\mathrm{nd}}[\delta S,\delta A],h\delta A\bigr),
\end{equation}
where \(S=(r,\eta,v,\omega)\) has \(72\) coordinates, \(A=(a,\alpha)\)
has \(36\) coordinates, and \(\mathcal N_h^{\mathrm{nd}}\) has \(96\)
non-dynamic rows.  The finite operator in Eq.~\eqref{eq:d5-p-state-finite-ps2-operator}
records full column rank \(108/108\)
for all three probes, a minimum singular value \(9.999780\times 10^{-3}\), and
a maximum finite state-projection constant \(2.019470\).  This is a finite
solved-stage rank diagnostic for the PS2 target; it is not a
uniform compact-tube inverse or inf-sup proof and does not close PS2, PS3,
\(P_{\mathrm{state}}\), any D5 Taylor bound, or PC2.
\end{lemma}

\begin{lemma}[Conditional PS3 conversion for \(P_{\mathrm{state}}\)]
\label{lem:d5-p-state-ps3-conditional-conversion}
Assume the weighted PS2 target above holds with a uniform constant on the
compact proof tube.  If the 96-row non-dynamic residual defect satisfies
\(\|\delta\mathcal N_h^{\mathrm{nd}}\|=O(h^7)\) and the acceleration input
satisfies \(\|\delta A\|=O(h^6)\), then the state-block perturbation satisfies
\(\|\delta S\|=O(h^7)\).
\end{lemma}

\begin{proof}
The weighted PS2 estimate gives
\begin{equation}
\label{eq:d5-p-state-ps3-conversion-bound}
\|\delta S\|\le C_{\mathrm{PS2}}
\bigl(O(h^7)+h\,O(h^6)\bigr)=O(h^7).
\end{equation}
Equation~\eqref{eq:d5-p-state-ps3-conversion-bound} records the conditional
PS3 conversion and explains why the
\(O(h^6)\) acceleration-rate obstruction is compatible with the
\(h\|\delta A\|\) PS2 target.  It does not by itself instantiate the actual
PS3 lift, because the required residual and acceleration-rate hypotheses still
have to be supplied in the lift table; it does not certify any induced Taylor
bound, and does not close \(P_{\mathrm{state}}\) or PC2.
\end{proof}

\begin{lemma}[Direct-route state-block corollary; not primitive \(P_{\mathrm{state}}\) closure]
\label{lem:d5-p-state-ps3-full-residual-route}
Assume the aggregate PS2 estimate of
Lemma~\ref{lem:d5-p-state-ps2-aggregate-promotion}, the 96-row
non-dynamic certificate, the D5 direct-substitution certificate for the
36 Newton--Euler rows, and the uniform inverse hypothesis in
Lemma~\ref{lem:stage-residual-defect}.  Then the full 132-row residual
route supplies a direct-route state-block corollary only, not a primitive
PS3 input for proving \(P_{\mathrm{state}}\):
\begin{equation}
\label{eq:d5-p-state-direct-corollary-state-acc}
  \|\delta S\|=O(h^7),
  \qquad
  h\|\delta A\|=O(h^7).
\end{equation}
The route does not use velocity collocation, does not use
\(P_{\mathrm{state}}\), \(P_{\mathrm{acc}}\), or \(P_{\lambda}\) as an
input, and Eq.~\eqref{eq:d5-p-state-direct-corollary-state-acc} does not
promote a finite rank probe to a compact-tube theorem.
It closes only the direct-route PS3 corollary; it does not certify the
primitive/Taylor route or any of the 162 induced D5 Taylor subterm bounds.
\end{lemma}

\begin{proof}
Let \(Z_G\) be the smooth FullVA lift of the reduced Gauss stage and
\(Z_A\) the accepted FullVA stage zero.  The non-dynamic certificate gives
\begin{equation}
\label{eq:d5-p-state-nondyn-residual-on-lift}
  \|\mathcal N_h^{\mathrm{nd}}(Z_G)\|=O(h^7)
\end{equation}
for the 96 non-dynamic rows.  Equation~\eqref{eq:d5-p-state-nondyn-residual-on-lift}
is the 96-row non-dynamic input to the direct residual bridge.
For each Newton--Euler component row, the D5
direct-substitution certificate expands the implemented row into the closed
D1/D2 balance identity, the D3 multiplier-wrench identity, the D4 smooth
force/torque input, the D6 row binding, and the lifted-stage balance term.
After substituting \(Z_G\), the direct-route residual decomposition has zero
residual remainder in each of the 36 dynamic rows. This supplies
\(R_{\rm dyn}(Z_G)=0\), hence the \(O(h^7)\) input to
Lemma~\ref{lem:stage-residual-defect}; it is not a certification of the
primitive 162-subterm Taylor route. In that separate verification route no
induced Taylor bounds are certified, the 162 subterms remain under open
primitive interfaces, and the 36 \(h\)-weighted acceleration primitive terms
remain insufficient.
Combining the finite row families gives
\begin{equation}
\label{eq:d5-p-state-full-residual-on-lift}
  \|F_{A,h}(Z_G;y)\|=O(h^7)
\end{equation}
for the full 132-row residual.

Lemma~\ref{lem:stage-residual-defect} then gives
\(\|Z_A-Z_G\|=O(h^7)\) in the weighted local stage norm.  In particular the
state block and weighted acceleration block obey
\begin{equation}
\label{eq:d5-p-state-stage-norm-bound}
  \|\delta S\| + h\|\delta A\| \le C\|Z_A-Z_G\|=O(h^7).
\end{equation}
Equivalently, the h-weighted acceleration input required by the aggregate
PS2 estimate is supplied by the full residual perturbation itself through
Eqs.~\eqref{eq:d5-p-state-full-residual-on-lift} and
\eqref{eq:d5-p-state-stage-norm-bound}.  The
argument never rewrites velocity collocation as
\begin{equation}
\label{eq:d5-p-state-velocity-collocation-forbidden}
  \delta A=h^{-1}(A_G^{-1}\otimes I)(\delta V-\rho_V),
\end{equation}
so it does not assume the \(P_{\mathrm{state}}\) velocity perturbation that it
is trying to control.  Avoiding Eq.~\eqref{eq:d5-p-state-velocity-collocation-forbidden}
is why the route is non-circular.

Finally, applying the aggregate PS2 estimate,
\begin{equation}
\label{eq:d5-p-state-final-ps2-application}
  \|\delta S\|\le C_{\mathrm{PS2}}
  \bigl(\|\delta\mathcal N_h^{\mathrm{nd}}\|+h\|\delta A\|\bigr),
\end{equation}
Equation~\eqref{eq:d5-p-state-final-ps2-application} gives
\(\|\delta S\|=O(h^7)\).  This direct
residual-bridge/Kantorovich perturbation route is a direct-route corollary of
the full residual certificate.  It is not recycled as an independent proof of
the primitive/Taylor route, because that route
must derive \(P_{\mathrm{state}}\), \(P_{\mathrm{acc}}\), and \(P_{\lambda}\)
before using them to certify the D5 Taylor subterms.
\end{proof}

\begin{lemma}[No direct-route feedback into primitive \(P_{\mathrm{state}}\)]
\label{lem:d5-p-state-direct-feedback-nonclosure}
The direct-route estimate of
Lemma~\ref{lem:d5-p-state-ps3-full-residual-route} cannot be used as the
actual \(P_{\mathrm{state}}\) primitive in the separate D5 primitive/Taylor
route.  In that route the \(P_{\mathrm{state}}\) input must be obtained before
the 162 primitive Newton--Euler Taylor subterms are certified, using only the
accepted non-dynamic rows, the compact-tube PS2 inverse data, and an
acceleration input independent of the state perturbation being estimated.
The direct-route estimate instead consumes the full 132-row residual bridge,
including the D5 direct-substitution certificate and
Lemma~\ref{lem:stage-residual-defect}.  It is therefore a theorem-route
corollary, not an admissible primitive-route input.
\end{lemma}

\begin{proof}
The separate primitive/Taylor route has a fixed dependency order:
first prove the root lift primitives \(P_{\mathrm{state}}\) and
\(P_{\mathrm{acc}}\) on the accepted compact branch, then propagate
\(P_{\lambda}\), and only then bound the primitive D5 subterms.  The
\(P_{\mathrm{state}}\) slot in that order is required to be independent of the
D5 dynamic residual defect and of the full-map perturbation lemma whose input
it would later help certify.

Lemma~\ref{lem:d5-p-state-ps3-full-residual-route} has a different
causality.  It first uses the accepted direct PC2 bridge: the 96 non-dynamic
row certificate, the 36 dynamic rows supplied as a same-branch base-point
zero block at the smooth FullVA lift, and the full-map stage perturbation
result.  Only after that full residual object is fixed does it read off the state-block and
\(h\)-weighted acceleration-block estimates.  Feeding this output back as the
primitive \(P_{\mathrm{state}}\) input would reverse the proof order and would
make the primitive Taylor route depend on the direct route it is meant to be
independent from.  Hence the direct-route estimate may be cited only as a
direct-route corollary for the theorem proof; it leaves actual PS3,
\(P_{\mathrm{state}}\), the \(126\) \(P_{\mathrm{state}}\)-dependent
primitive subterms, and the \(0/162\) primitive Taylor inventory unchanged.
\end{proof}

\begin{lemma}[Acceleration-map definition for \(P_{\mathrm{acc}}\)]
\label{lem:d5-p-acc-map-definition}
Let the acceleration block be
\begin{equation}
\label{eq:d5-p-acc-stage-acc-block}
  A=(a_i,\alpha_i)_{i=1}^3
\end{equation}
collect the accepted translational and angular stage accelerations for all
bodies in the same accepted stage ordering as the FullVA residual.  Define
\(\mathcal A_h^{\mathrm{acc}}(A)\) as the ordered vector of acceleration-lift
Taylor-row interfaces
\begin{equation}
\label{eq:d5-p-acc-taylor-row-interfaces}
  m_b(\widehat a_{b,i}-a_b(t_i)),\qquad
  J_b(\widehat\alpha_{b,i}-\alpha_b(t_i))
\end{equation}
for every stage \(i\), body \(b\), and component. Equivalently, the
acceleration block is Eq.~\eqref{eq:d5-p-acc-stage-acc-block}.  The
interfaces in Eq.~\eqref{eq:d5-p-acc-taylor-row-interfaces} give 36 scalar rows:
18 translational acceleration-lift rows and 18 angular acceleration-lift
rows. It is exactly the acceleration variable interface used by the D5
term-budget rows \(T_{\mathrm{acceleration\_lift}}\) and
\(R_{\mathrm{angular\_acceleration\_lift}}\). This closes only the PA1
map-definition subproof for \(P_{\mathrm{acc}}\). It does not prove an
acceleration lift rate, a velocity-collocation implication, a dynamic-row
defect bound, or PC2.
\end{lemma}

\begin{proof}
The accepted residual has three stages, two moving bodies, and three
Cartesian components per translational acceleration block, giving
\(3\times2\times3=18\) translational acceleration-lift interfaces.  The
angular acceleration block has the same stage-body-component count, giving
18 angular acceleration-lift interfaces.  The D5 Taylor-term inventory names
these two families as \(T_{\mathrm{acceleration\_lift}}\) and
\(R_{\mathrm{angular\_acceleration\_lift}}\), so the definition above fixes
the row ordering and variable interface for the 36 \(P_{\mathrm{acc}}\)
term rows.  No rate estimate follows from this definition alone: closing
\(P_{\mathrm{acc}}\) still requires a non-circular proof that the accepted
acceleration variables differ from the smooth lifted accelerations by
\(O(h^7)\), without reading that estimate from the Newton--Euler dynamic
balance.
\end{proof}

\begin{lemma}[Acceleration-row binding for \(P_{\mathrm{acc}}\)]
\label{lem:d5-p-acc-row-binding}
Under the accepted residual ordering verified in D6, the 36 components of
\(\mathcal A_h^{\mathrm{acc}}(A)\) occupy global rows 24--35, 68--79, and
112--123.  Within each stage, the row order is body-0 translational,
body-0 rotational, body-1 translational, and body-1 rotational, with three
Cartesian components per block.  The translational components are exactly the
D5 term-budget rows \(T_{\mathrm{acceleration\_lift}}\), and the rotational
components are exactly the rows
\(R_{\mathrm{angular\_acceleration\_lift}}\). This closes only the PA4
row-ordering/Taylor-term binding subproof for \(P_{\mathrm{acc}}\). It does
not prove an acceleration lift rate, a velocity-collocation implication, a
Taylor bound, or PC2.
\end{lemma}

\begin{proof}
D6 fixes the dynamic-row layout as a 44-row stage block with dynamic offset
24 and dynamic width 12.  Therefore the dynamic rows in stages
\(s=0,1,2\) are \(44s+24,\ldots,44s+35\), namely global rows 24--35, 68--79,
and 112--123.  Its per-stage order lists the translational source rows for
body 0, the rotational source rows for body 0, the translational source rows
for body 1, and the rotational source rows for body 1.  The D5 Taylor-term
budget assigns the acceleration-lift translational terms to the two
translational Newton-balance blocks and the angular acceleration-lift terms
to the two rotational Euler-balance blocks, with the same stage, body, and
component indices.  Hence every acceleration-lift Taylor row has the same
global row, balance block, body, component, and runtime source component in
the D5 budget and in the D6 accepted residual ordering. This establishes only
row and term binding.  It gives no estimate for
\(\widehat a-a(t_i)\) or \(\widehat\alpha-\alpha(t_i)\).  Thus PA2 remains
open; PA3 is the separate non-dynamic input-independence subproof closed in
Lemma~\ref{lem:d5-p-acc-independence}, not an acceleration-rate estimate.
\end{proof}

\begin{lemma}[Non-dynamic input boundary for \(P_{\mathrm{acc}}\)]
\label{lem:d5-p-acc-independence}
The proof route for \(P_{\mathrm{acc}}\) can be restricted to the
non-dynamic FullVA row families
\begin{equation}
\label{eq:d5-p-acc-nondynamic-row-families}
\begin{gathered}
\texttt{translational\_velocity\_weak\_defect},\\
\texttt{angular\_velocity\_weak\_defect},\\
\texttt{lower\_pair\_index3\_weak\_constraints},
\end{gathered}
\end{equation}
together with the 96-row non-dynamic certificate.  Newton--Euler
weak-balance rows are not inputs to this route; Eq.~\eqref{eq:d5-p-acc-nondynamic-row-families}
is the allowed row-family list.  This closes only the PA3
non-circularity subproof for \(P_{\mathrm{acc}}\). It does not prove an
acceleration lift rate, a velocity-collocation implication, a Taylor bound,
or the primitive-and-Taylor route to PC2; PA2 remains open.
\end{lemma}

\begin{proof}
The accepted stage block has the certified row partition
\begin{equation}
\label{eq:d5-p-acc-stage-row-partition}
0{:}6,\quad 6{:}12,\quad 12{:}18,\quad 18{:}24,\quad
24{:}36,\quad 36{:}44,
\end{equation}
with the Newton--Euler weak-balance family occupying only the \(24{:}36\)
subblock.  The rows used to formulate the lower-pair velocity and
acceleration-level FullVA relations occupy the non-dynamic families at
offsets \(12\), \(18\), and \(36\) in Eq.~\eqref{eq:d5-p-acc-stage-row-partition}.
These three families are included in
the 96-row kinematic row-defect certificate, while the Newton--Euler
weak-balance family is explicitly excluded from that certificate.  Thus a
later PA2 proof may use the certified non-dynamic FullVA relations to derive
acceleration lift rates, but it cannot read those rates from the dynamic
balance equations whose D5 defect is under proof.  The lemma therefore
establishes only the non-circular input boundary.
\end{proof}

\begin{lemma}[Velocity-collocation obstruction for \(P_{\mathrm{acc}}\)]
\label{lem:d5-p-acc-pa2-obstruction}
The velocity-collocation rows alone are insufficient to prove the unweighted
\(O(h^7)\) acceleration-lift primitive \(P_{\mathrm{acc}}\).  Let
\(\delta V\) denote the stacked translational and angular stage-velocity
difference between the accepted stage variables and the smooth lifted Gauss
stages, let \(\delta A\) denote the corresponding stacked acceleration
difference, and let \(\rho_V\) be the velocity-collocation row perturbation.
For the three-stage Gauss matrix \(A_G\),
\begin{equation}
\label{eq:d5-p-acc-velocity-collocation-loss}
  \delta V-h(A_G\otimes I)\delta A=\rho_V,\qquad
  \delta A=h^{-1}(A_G^{-1}\otimes I)(\delta V-\rho_V).
\end{equation}
Thus an \(O(h^7)\) bound on \(\delta V\) and \(\rho_V\) gives only an
unweighted \(O(h^6)\) acceleration bound through
Eq.~\eqref{eq:d5-p-acc-velocity-collocation-loss}.  PA2
therefore remains open unless one proves a stronger \(O(h^8)\)
stage-velocity/velocity-row estimate, an independent non-dynamic inverse that
controls acceleration directly, or an explicit \(P_{\mathrm{acc}}\)
regularity assumption.
\end{lemma}

\begin{proof}
The three-stage Gauss matrix is nonsingular, with entries fixed
independently of \(h\).  Subtracting the lifted velocity-collocation equation
from the accepted velocity-collocation equation gives the first displayed
identity.  Multiplication by \(A_G^{-1}\otimes I\) gives the second identity
and introduces the factor \(h^{-1}\).  Since \(\|A_G^{-1}\|\) is a fixed
constant, the estimate obtained from velocity-collocation alone is
\(\|\delta A\|\le C h^{-1}(\|\delta V\|+\|\rho_V\|)\).  With only
\(\delta V,\rho_V=O(h^7)\), this is \(O(h^6)\).  The current D5 Taylor budget
uses the unweighted dynamic-row acceleration terms \(m_b\delta a\) and
\(J_b\delta\alpha\), so a bound only on \(h\delta A\) is not enough to close
the existing \(P_{\mathrm{acc}}\) primitive.  This proves the obstruction and
does not close PA2, \(P_{\mathrm{acc}}\), or PC2.
\end{proof}

\noindent\textbf{Reference-proof boundary for \(P_{\mathrm{acc}}\).}
This obstruction is exactly why the constrained-BDF proof of
Wieloch and Arnold is invoked only as proof-order discipline.  Their analysis
also exposes step-size-scaled derivative and constraint-difference quantities
inside a coupled error recursion; it does not provide a FullVA stage estimate
for the unweighted acceleration variables appearing in the present
Newton--Euler rows.  Importing that reference proof at the primitive level
would therefore skip the \(h^{-1}\) loss displayed above.  The admissible
uses of the reference remain the structural ones stated in
Table~\ref{tab:reference-proof-order-correspondence}: fix the Lie-group
discrete equations first, prove the local defect for that fixed object, and
then propagate the accepted-branch error.  A primitive Taylor proof route
would still require an independent unweighted \(P_{\mathrm{acc}}\) proof, an
\(O(h^8)\) velocity-stage estimate before inversion, a direct non-dynamic
acceleration inverse, or an explicit extra theorem assumption.

\paragraph{PA2 weighted-inverse diagnostic}
The finite weighted-inverse diagnostic evaluates the operator
\((\delta S,\delta A)\mapsto(DN_h^{\mathrm{nd}}[\delta S,\delta A],
h\delta A)\) on the same solved stages used by the PS2 probe.  It records
bounded finite control of \(h\delta A\), separating the weighted
acceleration diagnostic from the unweighted acceleration estimate required by
the D5 Taylor budget.  This does not prove a uniform unweighted acceleration
lift.  PA2 remains open, \(P_{\mathrm{acc}}\) remains open, and the
primitive-and-Taylor route to PC2 remains open.  The finite weighted-inverse
constants are reference-boundary diagnostics only: they may guide a separate
PA2 proof, but they are not theorem inputs, do not certify an unweighted
acceleration lift, and do not change the non-active primitive-route
Taylor-bound status.

\begin{lemma}[Lower-pair acceleration rows do not close \(P_{\mathrm{acc}}\)]
\label{lem:d5-p-acc-lower-pair-nonclosure}
The lower-pair acceleration-level rows, even together with the finite
weighted-inverse PA2 diagnostic, do not by themselves imply the unweighted
full-body acceleration primitive \(P_{\mathrm{acc}}\).  To use these
non-dynamic rows as a PA2 closure one would need an \(h\)-uniform right
inverse, or an equivalent compact-tube inf-sup estimate, from the accepted
non-dynamic acceleration-row residuals to every translational and angular
stage-acceleration component in the D5 acceleration inventory.  No such
operator is constructed in this manuscript.  Consequently these rows cannot
serve as a hidden proof of PA2, \(P_{\mathrm{acc}}\), or PC2.
\end{lemma}

\begin{proof}
The \(P_{\mathrm{acc}}\) primitive used by the D5 Taylor budget is an
unweighted estimate for the full stacked translational and angular
accelerations appearing in the Newton--Euler rows.  The lower-pair
acceleration relations constrain constraint-direction combinations of these
accelerations, with coefficients determined by the lower-pair Jacobians and
with additional terms depending on already chosen position and velocity
variables.  They are therefore row relations for the mechanism constraints;
they are not, without an additional rank and right-inverse certificate, a
reconstruction formula for all body acceleration components in the D5 dynamic
row inventory.

The finite PA2 diagnostic adds bounded control of the weighted component
\(h\delta A\).  This is compatible with the velocity-collocation obstruction
of Lemma~\ref{lem:d5-p-acc-pa2-obstruction}: after removing the factor \(h\),
the available estimate loses one power of \(h\) unless a stronger velocity
estimate or an independent unweighted acceleration inverse is proved.  Hence
the current manuscript has no implication of the form
\begin{equation}
\label{eq:d5-p-acc-missing-unweighted-inverse}
  \|\delta A\|
  \le C\bigl(\|N_h^{\mathrm{acc}}\|+
             \hbox{already proved state and velocity errors}\bigr)
\end{equation}
Equation~\eqref{eq:d5-p-acc-missing-unweighted-inverse} records the
unavailable unweighted implication: it would need \(C\) independent of \(h\)
and range covering all D5 acceleration components.  The lower-pair rows and
the weighted diagnostic remain useful
diagnostics and possible ingredients for a separate PA2 proof, but they do not
close the unweighted \(P_{\mathrm{acc}}\) assumption used below.
\end{proof}

\begin{lemma}[No PS2--PA2 bootstrap closure]
\label{lem:d5-p-state-pacc-bootstrap-nonclosure}
The weighted PS2 estimate for \(P_{\mathrm{state}}\) and the
velocity-collocation estimate for \(P_{\mathrm{acc}}\) cannot close each
other by a small-step bootstrap.  Let
\(x_h=\|\delta S\|\), \(a_h=\|\delta A\|\), and let
\(r_h\) collect the \(O(h^7)\) non-dynamic row defects and velocity-row
defects on the same compact branch.  The available estimates have only the
form in Eq.~\eqref{eq:d5-ps2-pa2-coupled-available-bounds}:
\begin{equation}
\label{eq:d5-ps2-pa2-coupled-available-bounds}
  x_h\le C_{\mathrm{PS2}}(r_h+h a_h),
  \qquad
  a_h\le C_A h^{-1}(x_h+r_h).
\end{equation}
Substituting the second bound in the first gives
Eq.~\eqref{eq:d5-ps2-pa2-feedback-bound}:
\begin{equation}
\label{eq:d5-ps2-pa2-feedback-bound}
  x_h\le C_{\mathrm{PS2}}(1+C_A)r_h
       +C_{\mathrm{PS2}}C_A x_h .
\end{equation}
The feedback coefficient in Eq.~\eqref{eq:d5-ps2-pa2-feedback-bound} is
independent of \(h\), so shrinking the step size does not make this a
contraction.  Moreover, even if \(x_h=O(h^7)\) were supplied by some other
non-circular argument, the second estimate in
Eq.~\eqref{eq:d5-ps2-pa2-coupled-available-bounds} would give only
\(a_h=O(h^6)\) from \(r_h=O(h^7)\).  Thus the current PS2 and PA2 ingredients
can at most explain the weighted input \(h a_h=O(h^7)\); they do not prove
the unweighted \(P_{\mathrm{acc}}\) primitive, the actual primitive
\(P_{\mathrm{state}}\) lift, or any D5 Taylor subterm bound.
\end{lemma}

\begin{proof}
Equation~\eqref{eq:d5-p-state-ps2-estimate} gives the first inequality in
Eq.~\eqref{eq:d5-ps2-pa2-coupled-available-bounds} after the non-dynamic
defect terms are collected into \(r_h\).  The velocity-collocation loss in
Eq.~\eqref{eq:d5-p-acc-velocity-collocation-loss} gives the second inequality
with a constant \(C_A\) depending only on the fixed Gauss matrix and the
chosen component norm.  Composing the two estimates eliminates the small
factor \(h\): the \(h\) multiplying \(a_h\) in PS2 is cancelled by the
\(h^{-1}\) in the PA2 velocity-collocation inverse.  Therefore the available
closed estimates do not contain a small parameter that can absorb the
\(x_h\)-term in Eq.~\eqref{eq:d5-ps2-pa2-feedback-bound}.  A valid closure
would need an additional proof, such as an \(O(h^8)\) velocity-stage
estimate, an independent unweighted acceleration inverse, or an explicit
retained \(P_{\mathrm{acc}}\) assumption.  No such input is proved here, so
the primitive/Taylor route remains non-closing.
\end{proof}

\begin{corollary}[Row-level acceleration Taylor bound under
\(P_{\mathrm{acc}}\)]
\label{cor:d5-p-acc-row-bound-under-pacc}
Let the acceleration-row set be
\begin{equation}
\label{eq:d5-p-acc-row-set}
\begin{aligned}
\mathcal R_{\mathrm{acc}}=\{&
24,25,26,27,28,29,30,31,32,33,34,35,\\
&68,69,70,71,72,73,74,75,76,77,78,79,\\
&112,113,114,115,116,117,118,119,120,121,122,123\}
\end{aligned}
\end{equation}
be the global-row set for the translational and angular acceleration
Newton--Euler Taylor components.  Assume the acceleration part of
\(P_{\mathrm{acc}}\) gives
\begin{equation}
\label{eq:d5-p-acc-assumed-lift-rate}
  \|\widehat a_{s,b}-a_b(t_s)\|\le C_a h^7,\qquad
  \|\widehat\alpha_{s,b}-\alpha_b(t_s)\|\le C_\alpha h^7
\end{equation}
uniformly for every Gauss stage \(s\) and body \(b\).  If
\begin{equation}
\label{eq:d5-p-acc-body-constants}
  m_{\max}=\max_b m_b,\qquad J_{\max}=\max_b\|J_b\|_2,
\end{equation}
Equations~\eqref{eq:d5-p-acc-assumed-lift-rate} and
\eqref{eq:d5-p-acc-body-constants} are the conditional lift-rate and finite
body-constant inputs for this row-level corollary.
then every scalar row \(r\in\mathcal R_{\mathrm{acc}}\) satisfies
\begin{equation}
\label{eq:d5-p-acc-trans-row-bound}
  |e_r^T m_b(\widehat a_{s,b}-a_b(t_s))|
  \le m_{\max}C_a h^7
\end{equation}
on translational rows and
\begin{equation}
\label{eq:d5-p-acc-rot-row-bound}
  |e_r^T J_b(\widehat\alpha_{s,b}-\alpha_b(t_s))|
  \le J_{\max}C_\alpha h^7
\end{equation}
on rotational rows.  Thus Eqs.~\eqref{eq:d5-p-acc-row-set}--\eqref{eq:d5-p-acc-rot-row-bound}
show that the \(36\) acceleration Taylor rows satisfy the conditional T3 target
\(O(h^7)\) row bound under \(P_{\mathrm{acc}}\) with zero Taylor
remainder.  This is a conditional row-level consequence only: it does not
prove the unweighted acceleration lift, does not close PA2 or
\(P_{\mathrm{acc}}\), and does not change the PC2 status.
\end{corollary}

\begin{proof}
The acceleration contributions in the expanded Newton--Euler rows are the
linear maps \(A\mapsto m_b A\) and \(A\mapsto J_b A\).  Hence subtracting the
lifted and exact acceleration terms gives exactly
\begin{equation}
\label{eq:d5-p-acc-linear-lift-terms}
  m_b(\widehat a_{s,b}-a_b(t_s)),\qquad
  J_b(\widehat\alpha_{s,b}-\alpha_b(t_s)),
\end{equation}
Equation~\eqref{eq:d5-p-acc-linear-lift-terms} is the exact linear
acceleration contribution after subtracting the lifted and exact rows.
with no quadratic or higher Taylor remainder.  Projection onto a scalar
component can only decrease the Euclidean norm, and the finite body set gives
the constants \(m_{\max}\) and \(J_{\max}\).  The assumed \(P_{\mathrm{acc}}\)
rates therefore imply the \(O(h^7)\) row bounds in
Eqs.~\eqref{eq:d5-p-acc-trans-row-bound} and~\eqref{eq:d5-p-acc-rot-row-bound}
for exactly the
listed translational and rotational acceleration rows.  The proof consumes
the \(P_{\mathrm{acc}}\) lift as an input and does not derive that lift from
velocity collocation, the weighted PA2 diagnostic, or the dynamic residual.
\end{proof}

\begin{lemma}[KKT multiplier-column interface for \(P_{\lambda}\)]
\label{lem:d5-p-lambda-interface}
Let the accepted three-stage lower-pair multiplier array be
\begin{equation}
\label{eq:d5-p-lambda-stage-array}
\Lambda_h=\{\lambda_{i,j}\}_{i=1,\ldots,3;\,j=0,1},\qquad
\lambda_{i,j}\in\mathbb R^4,
\end{equation}
Equation~\eqref{eq:d5-p-lambda-stage-array} denotes the lower-pair multiplier variables in the accepted three-stage
residual.  In each 44-variable stage block, the two 18-variable body blocks
are followed by \(2\times4=8\) lower-pair multiplier variables; hence the
full Newton vector has 24 multiplier columns.  The first two components of
\(\lambda_{i,j}\) generate the normal force
\(\lambda_{i,j,0}b_{j,0}+\lambda_{i,j,1}b_{j,1}\) in the joint basis, and
the last two components generate the two axis/twist torque multiplier
directions used in the rotational balance.

The D5 multiplier interface therefore exposes the direct wrench terms
\begin{equation}
\label{eq:d5-p-lambda-wrench-interface}
G_b(q_i)^{T}\lambda_i,\qquad H_b(q_i)^{T}\lambda_i
\end{equation}
in the translational and rotational Newton--Euler rows.  The direct
multiplier-wrench Taylor rows exposed by Eq.~\eqref{eq:d5-p-lambda-wrench-interface}
are the 36 rows named
\(T_{\mathrm{multiplier\_force\_lift}}\) and
\(R_{\mathrm{multiplier\_torque\_lift}}\), while the primitive-reduction
table records 72 total D5 Taylor rows depending on \(P_{\lambda}\).
The D3 virtual-work certificate supplies the row-expanded wrench-consistency
identity for the same 36 dynamic rows.  This closes only the PL1
multiplier-interface subproof for \(P_{\lambda}\). It does not prove a
multiplier lift rate, does not prove a uniform inf-sup bound, does not
certify any Taylor bound. It also does not close the primitive-and-Taylor
route to PC2.
\end{lemma}

\begin{proof}
The runtime unpacking is
\begin{equation}
\label{eq:d5-p-lambda-runtime-unpack}
\lambda_{i,j}=x_{44i+36+4j:\,44i+36+4j+4}
\end{equation}
The slice in Eq.~\eqref{eq:d5-p-lambda-runtime-unpack} appears after the body variables in stage \(i\), which is the same column group
identified as \(\texttt{lower\_pair\_lambda}\).  The translational balance uses
the first two components only through the joint-basis normal force, and the
rotational balance uses the last two components through the corresponding
axis/twist torque moment directions at the proximal and distal lower pairs.
Thus the KKT multiplier columns that can affect the Newton--Euler weak rows
are explicitly identified before any lift estimate is invoked.  The D3
wrench identity checks that these columns produce the expected generalized
wrench signs row by row, but that identity is only an algebraic interface
check.  The PL2 proof below supplies the uniform constrained linearized
estimate, and the PL4 proof below supplies only the conditional propagation
from state and acceleration lifts to
\(\lambda_{i,j}-\lambda(t_i)=O(h^7)\); the required input lifts themselves
remain open.
\end{proof}

\begin{lemma}[Non-circular D3 use for \(P_{\lambda}\)]
\label{lem:d5-p-lambda-d3-noncircularity}
The row-expanded D3 virtual-work identity may be used in the
\(P_{\lambda}\) proof chain only as an algebraic mapping from lower-pair
multiplier components to the generalized force and torque entries in the 36
Newton--Euler weak-balance rows.  This closes PL3: the D3 wrench-consistency
certificate is used non-circularly.  It does not prove a multiplier lift rate, does
not prove a uniform inf-sup bound, and does not certify a Taylor bound. It
also does not close the primitive-and-Taylor route to PC2.
\end{lemma}

\begin{proof}
The D3 certificate expands the virtual work of the lower-pair reaction forces and
torques row by row and checks the same translational and rotational signs as
the accepted residual.  That check is purely algebraic in the runtime row
formulas, joint bases, moment arms, and multiplier slots.  It neither compares
the implemented stage variables with the lifted exact Gauss stages nor invokes
the D5 dynamic residual defect rate.  Therefore it can identify the correct
wrench map for a separate \(P_{\lambda}\) estimate without assuming the estimate
itself.  The PL2 compact-tube inf-sup bound is supplied by the symbolic
geometric proof below.  The remaining rate boundary is that PL4 can only be
instantiated after independent state and acceleration lift inputs are
available.
\end{proof}

The finite diagnostic probe
\(\texttt{D5\_P\_LAMBDA\_INF\_SUP\_PROBE}\) evaluates the accepted
dynamic-row multiplier-column block
\begin{equation}
\label{eq:d5-p-lambda-inf-sup-probe-operator}
  \delta\lambda\mapsto D_{\lambda}R_{\mathrm{dyn}}(Z_G)\delta\lambda
\end{equation}
Equation~\eqref{eq:d5-p-lambda-inf-sup-probe-operator} is evaluated on solved one-step stages at \(h=0.04,0.02,0.01\).  The measured
\(36\times24\) block has rank \(24\) in all three probes, with minimum
singular value about \(6.14\times10^{-1}\) and maximum condition number about
\(2.77\).  This is useful finite evidence that the PL2 target is the right
linearized multiplier interface.  It is not a uniform compact-tube inf-sup
proof, it does not prove a multiplier lift rate, and it does not close
\(P_{\lambda}\) or PC2.

\paragraph{PL2 geometric-margin diagnostic}
The finite geometric-margin diagnostic separates the PL2 operator into the
stage-local dynamic--lambda block, the
translational normal-force subblock, the rotational axis-torque subblock, and
their direct-sum diagnostic.  On the recorded solved stages at
\(h=0.04,0.02,0.01\), these finite blocks are full column rank and have
positive smallest singular values.  These finite margins identify the
geometric proof target for the multiplier columns, but they are not a uniform
compact-tube inf-sup proof.  The symbolic proof below gives that compact-tube
lower bound and closes PL2 as a uniform inf-sup subproof.  In the
primitive-and-Taylor route, the \(P_{\lambda}\) route and the PC2 route both
remain open because the PL4
propagation below is conditional on the still-open state and acceleration
lift inputs.

\begin{lemma}[PL2 multiplier-column symbolic structure]
\label{lem:d5-p-lambda-pl2-symbolic-structure}
For one lower pair, let \(a\) be the cylindrical axis, let
\(B=[b_0\ b_1]\) be the two normal basis columns with \(B^T B=I\) and
\(B^T a=0\), and write the normal multiplier as
\(\lambda_N\in\mathbb R^2\).  The implemented normal-force and
smooth-friction subblock has the form
\begin{equation}
\label{eq:d5-p-lambda-normal-force-map}
  F_N(\lambda_N)=B\lambda_N+\varphi(v,\|B\lambda_N\|)a ,
\end{equation}
where \(\varphi\) is the smooth Brown--McPhee scalar on the compact proof
tube.  Differentiating Eq.~\eqref{eq:d5-p-lambda-normal-force-map} gives
\begin{equation}
\label{eq:d5-p-lambda-normal-force-gram}
\begin{aligned}
  D_{\lambda_N}F_N=B+a c^T,\qquad
  (B+a c^T)^T(B+a c^T)=I+c c^T,
\end{aligned}
\end{equation}
Equation~\eqref{eq:d5-p-lambda-normal-force-gram} shows that the
smooth-friction normal-load derivative adds only a nonnegative Gram
correction.
for a bounded row vector \(c\).  Therefore each joint normal-force column
map has smallest singular value at least one; the smooth-friction
normal-load derivative cannot reduce this margin.  The two-joint
translational action--reaction block consequently has the symbolic bound
\begin{equation}
\label{eq:d5-p-lambda-trans-symbolic-bound}
  \sigma_{\min}\ge {(\sqrt5-1)}/{2}.
\end{equation}
For an axis-torque block
\begin{equation}
\label{eq:d5-p-lambda-axis-torque-map}
  C(a,R,B)\eta=\eta_0\,a\times R^T b_0+\eta_1\,a\times R^T b_1,
  \qquad n_B=b_0\times b_1 ,
\end{equation}
one has
\begin{equation}
\label{eq:d5-p-lambda-axis-torque-gram}
\begin{aligned}
  C^T C=I-c_a c_a^T,\qquad (c_a)_i=(Ra)\cdot b_i,\qquad
  \sigma_{\min}(C)=|(Ra)\cdot n_B|.
\end{aligned}
\end{equation}
Thus PL2 is reduced to an axis-plane margin:
\(|(Ra)\cdot n_B|\ge\chi_{\mathrm{axis}}>0\) on the compact tube for the
axis-torque columns by Eqs.~\eqref{eq:d5-p-lambda-axis-torque-map}
and~\eqref{eq:d5-p-lambda-axis-torque-gram}.  With that margin and the compact-tube upper bounds for
distal torque and normal-to-rotational couplings, the recorded block
triangular estimate gives a full dynamic-lambda lower bound.  The
normal-force and smooth-friction part is closed by this symbolic argument;
the remaining input for PL2 is the compact-tube axis-plane margin proved in
Lemma~\ref{lem:d5-p-lambda-pl2-axis-plane-margin}.
\end{lemma}

\begin{proof}
The first two identities use only the orthogonal decomposition
\(\mathbb R^3=\operatorname{span}\{b_0,b_1\}\oplus\operatorname{span}\{a\}\).
Since \(B^T a=0\), the rank-one friction derivative contributes
\(a c^T\) in a direction orthogonal to the normal basis, and the Gram matrix
is \(I+c c^T\).  Its eigenvalues are at least one.  For the two lower pairs,
write \(u=A_0x\) and \(v=A_1y\), where each \(A_j\) is such a normal-force
map.  Then
\begin{equation}
\label{eq:d5-p-lambda-two-joint-trans-bound}
  \|-u+v\|^2+\|v\|^2
  \ge \frac{3-\sqrt5}{2}\,(\|u\|^2+\|v\|^2)
  \ge \frac{3-\sqrt5}{2}\,(\|x\|^2+\|y\|^2),
\end{equation}
Equation~\eqref{eq:d5-p-lambda-two-joint-trans-bound} gives the
translational lower bound in
Eq.~\eqref{eq:d5-p-lambda-trans-symbolic-bound}.

For the axis-torque identity, use
\((a\times q_i)\cdot(a\times q_j)=q_i\cdot q_j-(a\cdot q_i)(a\cdot q_j)\)
with \(q_i=R^T b_i\).  Because \(R\) is orthogonal and \(b_0,b_1\) are
orthonormal, \(q_i\cdot q_j=\delta_{ij}\), and
\(a\cdot q_i=(Ra)\cdot b_i\).  The two nonzero singular values are therefore
the square roots of the eigenvalues of \(I-c_a c_a^T\), whose smaller value
is \(|(Ra)\cdot n_B|\).  The remaining full-block statement is the standard
left-inverse estimate for a block triangular operator: the translational
normal block controls the normal multipliers, the rotational block controls
the axis multipliers once the axis-plane margin is available, and the compact
tube bounds the off-diagonal coupling.  This lemma records the strict
symbolic reduction; the next lemma supplies the required compact-tube margin.
\end{proof}

\begin{lemma}[PL2 axis-plane margin on the accepted compact tube]
\label{lem:d5-p-lambda-pl2-axis-plane-margin}
For each proximal axis-torque pair in the accepted lower-pair chart, let
\begin{equation}
\label{eq:d5-p-lambda-axis-plane-margin-function}
  m_j(z)= |(R_j(z)a_j)\cdot n_j|,
  \qquad n_j=b_{j,0}\times b_{j,1},
\end{equation}
where \(b_{j,0},b_{j,1}\) are the fixed normal basis vectors used by the
runtime lower-pair multiplier columns.  Under
Assumption~\ref{ass:regularity}, after shrinking the compact proof tube
inside the same accepted chart if necessary, there is a constant
\begin{equation}
\label{eq:d5-p-lambda-axis-plane-margin-constant}
\chi_{\mathrm{axis}}
  :=\min_{z\in K}\min_j m_j(z)>0 .
\end{equation}
Thus Eqs.~\eqref{eq:d5-p-lambda-axis-plane-margin-function}
and~\eqref{eq:d5-p-lambda-axis-plane-margin-constant} give the proximal axis-torque blocks a uniform compact-tube lower
singular-value bound.  Finite probes are not used to establish the margin.
\end{lemma}

\begin{proof}
The previous lemma gives the exact identity
\begin{equation}
\label{eq:d5-p-lambda-axis-plane-singular-identity}
  \sigma_{\min} C_j(z)=m_j(z)
\end{equation}
Equation~\eqref{eq:d5-p-lambda-axis-plane-singular-identity} identifies the
axis-plane margin with the singular value of the proximal torque block.
for each proximal two-column axis-torque block.  If \(m_j(z)=0\), then the
two axis-torque multiplier columns are linearly dependent and the
axis-torque coordinate part of the lower-pair chart loses rank at \(z\).
This is precisely the chart transversality failure excluded by
Assumption~\ref{ass:regularity}: the active lower-pair constraint Jacobian has
constant full row rank and the accepted branch is represented in a smooth
lower-pair chart.  Hence \(m_j(z)>0\) on the accepted smooth branch for every
finite proximal axis-torque pair \(j\).

Each \(m_j\) is continuous because \(R_j(z)\) is smooth in the accepted chart
and the basis vectors are fixed.  Lemma~\ref{lem:d5-compact-tube-constants}
provides a compact tube \(K\) around the accepted branch; shrinking this tube
inside the same nondegenerate chart preserves the branch and keeps
\(m_j>0\) on \(K\).  The minimum over the finite family of continuous
positive functions \(m_j\) on compact \(K\) is therefore positive.  This
compactness argument uses the symbolic singular-value identity and chart
regularity only; the finite singular-value probes are diagnostic and are not
used as proof.
\end{proof}

\begin{lemma}[PL2 compact-tube reduction]
\label{lem:d5-p-lambda-pl2-compact-reduction}
Let \(K\) be the compact smooth proof tube and let
\(B(z,h)=D_{\lambda}R_{\mathrm{dyn}}(z,h)\) denote the accepted
dynamic-row multiplier-column map restricted to one stage.  Let
\begin{equation}
\label{eq:d5-p-lambda-pl2-gamma-tr}
\begin{aligned}
  \gamma_T={(\sqrt5-1)}/{2},\qquad
  \gamma_R=\left(2\chi_{\mathrm{axis}}^{-2}
     +\beta_D^2\chi_{\mathrm{axis}}^{-4}\right)^{-1/2},
\end{aligned}
\end{equation}
where \(\beta_D<\infty\) is the compact-tube bound for the distal
axis-torque coupling.  If \(\beta_N<\infty\) is the compact-tube bound for
normal-to-rotational coupling, define
\begin{equation}
\label{eq:d5-p-lambda-pl2-gamma-full}
  \gamma_{\mathrm{PL2}}
  =\left(\gamma_T^{-2}+\gamma_R^{-2}
     +\beta_N^2\gamma_T^{-2}\gamma_R^{-2}\right)^{-1/2}.
\end{equation}
Equation~\eqref{eq:d5-p-lambda-pl2-gamma-full} combines the translational,
rotational, and normal-to-rotational coupling margins into the PL2 constant.
Then \(\gamma_{\mathrm{PL2}}>0\), and PL2 has a uniform inf-sup constant on
the tube:
\begin{equation}
\label{eq:d5-p-lambda-pl2-uniform-infsup}
  \|B(z,h)\delta\lambda\|
  \ge \gamma_{\mathrm{PL2}}\|\delta\lambda\|,
  \qquad (z,h)\in K\times(0,h_0].
\end{equation}
The finite PL2 probes are not a uniform compact-tube inf-sup proof; they only identify the
normal-force, axis-torque, and smooth-friction perturbation maps.  PL2 is
closed as a uniform inf-sup subproof by
Eqs.~\eqref{eq:d5-p-lambda-pl2-gamma-tr}--\eqref{eq:d5-p-lambda-pl2-uniform-infsup}.
This does
not prove the PL4 multiplier-rate propagation, \(P_{\lambda}\), or PC2.
\end{lemma}

\begin{proof}
The preceding symbolic-structure lemma gives the translational normal block
the lower bound \(\gamma_T\).  Lemma~\ref{lem:d5-p-lambda-pl2-axis-plane-margin}
gives the proximal axis-torque lower bound \(\chi_{\mathrm{axis}}\), while
Lemma~\ref{lem:d5-compact-tube-constants} bounds the distal torque coupling
by \(\beta_D\).  The standard two-block left-inverse estimate for that
rotational action-reaction block gives \(\gamma_R\).  Applying the same
block-triangular estimate to the full dynamic-lambda map, with
normal-to-rotational coupling bounded by \(\beta_N\), gives
\(\gamma_{\mathrm{PL2}}\).

All constants in this formula are positive or finite on \(K\), and the finite
number of stages and lower pairs lets us take the minimum over the stage-local
constants.  Hence Eq.~\eqref{eq:d5-p-lambda-pl2-uniform-infsup} holds uniformly for all accepted
stage states in the tube and all sufficiently small \(h\).  The proof uses
only the symbolic identities, chart transversality, and compact-tube
constants; finite probes do not establish the inf-sup constant.
\end{proof}

\begin{lemma}[PL4 conditional multiplier-rate propagation]
\label{lem:d5-p-lambda-pl4-rate-propagation}
Let \(y=(z,a)\) collect the accepted state and acceleration variables in the
lower-pair chart, and let \(F(y,\lambda,h)\) be the accepted lower-pair
dynamic-row map after the PL1 multiplier-column interface and the PL3 D3
wrench mapping have fixed the multiplier coordinates.  On the compact proof
tube, assume the PL2 bound
\begin{equation}
\label{eq:d5-p-lambda-pl4-pl2-assumption}
  \|D_{\lambda}F(y,\lambda,h)\delta\lambda\|
  \ge \gamma_{\mathrm{PL2}}\|\delta\lambda\|
\end{equation}
with \(\gamma_{\mathrm{PL2}}>0\).  Then the local lower-pair multiplier
solution branch is uniformly Lipschitz in \(y\):
\begin{equation}
\label{eq:d5-p-lambda-pl4-lipschitz}
  \|\widehat\lambda-\lambda\| \le C_{\lambda y}
     \bigl(\|\widehat z-z\|+\|\widehat a-a\|\bigr).
\end{equation}
Equations~\eqref{eq:d5-p-lambda-pl4-pl2-assumption} and
\eqref{eq:d5-p-lambda-pl4-lipschitz} state the PL2 input and the resulting
conditional PL4 propagation estimate.
Consequently, if the still-open \(P_{\mathrm{state}}\) and
\(P_{\mathrm{acc}}\) primitives supply
\(\|\widehat z-z\|=O(h^7)\) and
\(\|\widehat a-a\|=O(h^7)\), then
\(\|\widehat\lambda-\lambda\|=O(h^7)\).  This closes only the conditional PL4
propagation step; it does not close \(P_{\lambda}\), certify any D5 Taylor
term, or close PC2.
\end{lemma}

\begin{proof}
Fix two points on the same accepted lower-pair solution branch inside the
compact tube, and write
\(\delta y=\widehat y-y\) and
\(\delta\lambda=\widehat\lambda-\lambda\).  Since both points satisfy the
same lower-pair dynamic-row equation in the accepted chart, the first-order
Taylor formula with integral remainder gives
\begin{equation}
\label{eq:d5-p-lambda-pl4-taylor-identity}
\begin{aligned}
  0
  &=F(\widehat y,\widehat\lambda,h)-F(y,\lambda,h)\\
  &=B(y,\lambda,h)\delta\lambda
    +D_yF(y,\lambda,h)\delta y+\rho,\\
  B&=D_{\lambda}F .
\end{aligned}
\end{equation}
where compact \(C^2\) bounds give
\begin{equation}
\label{eq:d5-p-lambda-pl4-remainder-bound}
\begin{aligned}
  \|\rho\|
  \le C_R\bigl(\|\delta y\|+\|\delta\lambda\|\bigr)^2,
  \qquad
  \|D_yF(y,\lambda,h)\delta y\|\le L_y\|\delta y\|.
\end{aligned}
\end{equation}
PL2 gives a uniform left-inverse estimate for \(B\), hence
\begin{equation}
\label{eq:d5-p-lambda-pl4-preabsorption}
  \gamma_{\mathrm{PL2}}\|\delta\lambda\|
  \le L_y\|\delta y\|
     +C_R\bigl(\|\delta y\|+\|\delta\lambda\|\bigr)^2 .
\end{equation}
Equations~\eqref{eq:d5-p-lambda-pl4-taylor-identity},
\eqref{eq:d5-p-lambda-pl4-remainder-bound}, and
\eqref{eq:d5-p-lambda-pl4-preabsorption} are the Taylor identity,
compact-tube remainder bound, and pre-absorption estimate for PL4.
The absorption is a local-branch condition, not a global multiplier-root
selection theorem.  Fix a PL4 comparison radius \(r_{\lambda}>0\), contained
in the same compact-tube chart and accepted lower-pair branch, such that
\begin{equation}
\label{eq:d5-p-lambda-pl4-local-radius}
  \|\delta y\|+\|\delta\lambda\|\le r_{\lambda},
  \qquad
  C_R\bigl(\|\delta y\|+\|\delta\lambda\|\bigr)
  \le \frac{1}{2}\gamma_{\mathrm{PL2}} .
\end{equation}
Under Eq.~\eqref{eq:d5-p-lambda-pl4-local-radius}, the quadratic term in
Eq.~\eqref{eq:d5-p-lambda-pl4-preabsorption} is absorbed into the left-hand
side, giving
\begin{equation}
\label{eq:d5-p-lambda-pl4-absorbed-bound}
  \|\delta\lambda\|
  \le \frac{2L_y}{\gamma_{\mathrm{PL2}}}\|\delta y\|
  =C_{\lambda y}\|\delta y\|.
\end{equation}
Equation~\eqref{eq:d5-p-lambda-pl4-absorbed-bound} is uniform over the finite
set of stages and lower pairs once both compared points are already known to
remain in the PL4 local branch ball in
Eq.~\eqref{eq:d5-p-lambda-pl4-local-radius}.  Small \(h\) enters only through
the independent state and acceleration input lifts that place the compared
points in that ball; it does not prove global branch selection or exclude
remote multiplier roots.  Inserting the independent input rates
\(\|\delta z\|=O(h^7)\) and \(\|\delta a\|=O(h^7)\) gives the claimed
conditional multiplier rate.  The argument uses the Taylor expansion,
compact-tube constants, and the PL2 inf-sup constant; it does not use
Lemma~\ref{lem:stage-residual-defect}, the D5 dynamic residual defect,
direct substitution, finite probes, or the D3 identity as a rate proof.
\end{proof}

\begin{lemma}[PL4 closure is not \(P_{\lambda}\) closure]
\label{lem:d5-p-lambda-conditional-nonclosure}
Closing the four local PL subproofs does not by itself close the
multiplier-lift primitive \(P_{\lambda}\).  The output of PL4 is the
conditional implication
\begin{equation}
\label{eq:d5-p-lambda-pl4-conditional-output}
  \|\delta\lambda\|
  \le C_{\lambda y}\bigl(\|\delta z\|+\|\delta a\|\bigr),
\end{equation}
on the accepted compact branch and in the fixed FullVA norm.  It becomes an
actual \(O(h^7)\) multiplier estimate only after the state lift
\(\|\delta z\|=O(h^7)\) and the unweighted acceleration lift
\(\|\delta a\|=O(h^7)\) have been proved independently for the same branch.
Since the current primitive status keeps \(P_{\mathrm{state}}\) and
\(P_{\mathrm{acc}}\) open, \(P_{\lambda}\) remains open even though PL4 is
available as a conditional propagation implication.
\end{lemma}

\begin{proof}
Lemma~\ref{lem:d5-p-lambda-pl4-rate-propagation} proves a Lipschitz
dependence of the multiplier coordinates on the already chosen state and
acceleration variables.  This dependence is an implication from input lift
rates to multiplier lift rates; it is not a source of those input lift rates.
Substituting \(O(h^7)\) input estimates into
Eq.~\eqref{eq:d5-p-lambda-pl4-conditional-output} would
indeed give \(\|\delta\lambda\|=O(h^7)\), but that substitution is admissible
only if the two inputs have already been established without using the D5
dynamic residual defect, Lemma~\ref{lem:stage-residual-defect}, the direct
PC2 substitution certificate, finite probes, or the D3 wrench identity as a
rate proof.

Those independent inputs are precisely the still-open
\(P_{\mathrm{state}}\) and \(P_{\mathrm{acc}}\) primitives recorded in the
D5 primitive inventory.  Therefore the PL subproof count \(4/4\) means that the
conditional multiplier propagation mechanism is complete; it does not mean
that the actual multiplier primitive is complete.  The 72 D5 Taylor subterms
that consume \(P_{\lambda}\) remain uncertified until the two root lifts are
closed first.
\end{proof}

\noindent\textbf{Reference-proof boundary for \(P_{\lambda}\).}
The multiplier estimates in the constrained-BDF reference are not imported
as a \(P_{\lambda}\) lift.  In that proof the multiplier recursion is tied to
the BLieDF multistep defect, hidden-constraint coordinates, and its own
start-value assumptions.  The present lower-pair multipliers are instead
stage unknowns in the \(132\)-row FullVA residual, with their columns fixed by
PL1, their compact inf-sup geometry fixed by PL2--PL3, and their rate
controlled only conditionally by PL4 from the state and acceleration lifts.
Thus the reference proof only motivates the ordering discipline recorded in
Table~\ref{tab:reference-proof-order-correspondence}: first fix the discrete
object and multiplier interface, then prove the local defect for that fixed
object, and only then propagate the accepted-branch error.  It does not
provide the required actual \(P_{\lambda}\) rate, does not bypass the open
\(P_{\mathrm{state}}\) and \(P_{\mathrm{acc}}\) inputs, and does not certify
any of the 72 primitive Taylor subterms that depend on \(P_{\lambda}\).

\begin{table}[htbp]
\centering
\scriptsize
\setlength{\tabcolsep}{2pt}
\renewcommand{\arraystretch}{0.88}
\caption{Retained primitive proof interfaces for D5.}
\label{tab:d5-open-primitive-proof-interfaces}
\begin{tabular}{@{}P{0.14\linewidth}P{0.10\linewidth}P{0.42\linewidth}P{0.20\linewidth}@{}}
\toprule
Primitive & Term rows & Missing independent proof & Forbidden shortcut \\
\midrule
\(P_{\mathrm{state}}\) & 126 &
The PS2 weighted inverse and row-injection route are recorded; the required
step is the actual PS3 instantiation that converts the 96 certified
non-dynamic residual rows, together with an acceleration input independent of
the state perturbation being estimated, into \(O(h^7)\) pose and velocity
lift rates. &
Using Lemma~\ref{lem:stage-residual-defect} or the D5 dynamic residual defect. \\
\(P_{\mathrm{acc}}\) & 36 &
Unweighted \(O(h^7)\) translational and angular acceleration lift rates in
the accepted variable ordering, supplied by one of the admissible PA2 routes:
an \(O(h^8)\) velocity/velocity-row estimate, an independent non-dynamic
inverse that controls acceleration directly, or an explicit
\(P_{\mathrm{acc}}\) theorem assumption. &
Reading acceleration rates from dynamic balance. \\
\(P_{\lambda}\) & 72 &
Instantiate the closed conditional PL4 propagation with independent
\(P_{\mathrm{state}}\) and \(P_{\mathrm{acc}}\) lift inputs; the PL
subproof count is complete only in this conditional sense, and the actual
multiplier lift remains open. &
Using the D3 wrench identity as a multiplier-rate proof. \\
\(P_{\mathrm{geom}}\) & 36 &
Apply the closed accepted-chart geometry reduction below using the pose and
multiplier lifts supplied by \(P_{\mathrm{state}}\) and \(P_{\lambda}\). &
Treating the chart reduction as a pose or multiplier lift proof. \\
\(P_{\mathrm{gyro}}\) & 18 &
Apply the closed bilinear reduction below using the angular-velocity lift
supplied by \(P_{\mathrm{state}}\), then bind it to the rotational rows. &
Treating the bilinear estimate as a velocity-lift proof. \\
\bottomrule
\end{tabular}
\end{table}

This table is an obligation map for the separate primitive route: it does not
certify actual D5 Taylor subterm bounds and it identifies the independent
primitive estimates that would be needed before that separate primitive route could
close.
None of the five retained primitive rows is allowed to use
Lemma~\ref{lem:stage-residual-defect} or the D5 dynamic residual rate as an
input. This keeps the proof non-circular:
Table~\ref{tab:d5-open-primitive-proof-interfaces} states exactly which
independent estimates are still needed before the primitive-and-Taylor route
to PC2 can be closed.

\noindent\textbf{Finite dependency certificate for the separate route.}
The \(162\)-subterm statement below is a finite dependency certificate, not a
generic smoothness assertion.  The recorded D5 inventory splits into four
template families in Eq.~\eqref{eq:d5-primitive-family-count}:
\begin{equation}
\label{eq:d5-primitive-family-count}
\begin{aligned}
&36\ \text{acceleration lift}
 +72\ \text{smooth force/torque/friction lift}\\
&\quad
 +36\ \text{multiplier-geometry lift}
 +18\ \text{gyroscopic bilinear lift}
 =162 .
\end{aligned}
\end{equation}
Each family has a fixed primitive-input list before any Taylor estimate is
read: acceleration rows consume \(P_{\mathrm{acc}}\); smooth
force/torque/friction rows consume \(P_{\mathrm{state}}\) and, for the
lambda-dependent friction terms, \(P_{\lambda}\); multiplier-geometry rows
consume \(P_{\mathrm{state}}\) and \(P_{\lambda}\); and gyroscopic rows
consume the angular-velocity part of \(P_{\mathrm{state}}\).  The compact
tube primitive supplies the uniform derivative or bilinear constants for all
four families.  This dependency graph is acyclic in the separate route:
\(P_{\mathrm{state}}\) and \(P_{\mathrm{acc}}\) are root lift inputs,
\(P_{\lambda}\) is propagated only from those roots, and
\(P_{\mathrm{geom}}\) and \(P_{\mathrm{gyro}}\) are downstream reductions.
Consequently the \(162/162\) coverage count means that every primitive
subterm has an assigned template and primitive-input set. It is template
coverage only: the certified count of actual primitive Taylor bounds is \(0/162\); it
does not mean that any of those five open primitive inputs has been proved,
and it is not a projection of the accepted \(132\)-row direct residual
estimate.

\begin{proposition}[Conditional implication order for the separate D5 primitive/Taylor route]
\label{prop:d5-primitive-taylor-conditional-implication}
Assume, in addition to the already closed compact-tube primitive
\(P_{\mathrm{tube}}\), that two independent root-lift estimates are proved on
the same accepted compact tube, as in Eq.~\eqref{eq:d5-root-lift-hypotheses}:
\begin{equation}
\label{eq:d5-root-lift-hypotheses}
  \|\delta S\|\le C_Sh^7,\qquad
  \|\delta A\|\le C_Ah^7 .
\end{equation}
Assume further that these two estimates are obtained without using
Lemma~\ref{lem:stage-residual-defect}, the D5 dynamic residual defect, the D5
direct-substitution certificate, finite numerical probes, or the D3
wrench identity as a rate proof.  Then PL4 gives the conditional multiplier
rate \(P_{\lambda}\); the accepted-chart geometry and gyroscopic reductions
give \(P_{\mathrm{geom}}\) and \(P_{\mathrm{gyro}}\); and the finite sum over
the 162 D5 primitive Taylor subterms yields
Eq.~\eqref{eq:d5-prim-dynamic-block-bound},
\begin{equation}
\label{eq:d5-prim-dynamic-block-bound}
  \|R_{\mathrm{dyn}}^{\mathrm{prim}}(Z_G)\|
  \le C_{\mathrm{prim}}h^7 .
\end{equation}
Thus the separate primitive/Taylor route has the dependency order
in Eq.~\eqref{eq:d5-prim-dependency-order}:
\begin{equation}
\label{eq:d5-prim-dependency-order}
\begin{aligned}
  P_{\mathrm{state}},\ P_{\mathrm{acc}}
  &\Rightarrow P_{\lambda}
  \Rightarrow P_{\mathrm{geom}},\ P_{\mathrm{gyro}}\\
  &\Rightarrow \text{162 D5 Taylor subterms}.
\end{aligned}
\end{equation}
The proposition is only a conditional sufficiency statement.  Its root-lift
prerequisites remain unproved in the current manuscript, it does not establish
actual primitive Taylor bounds, and it does not
make Theorem~\ref{thm:g6fullva-order} depend on the primitive/Taylor route.
\end{proposition}

\begin{proof}
The proof is a finite implication chain inside the separate primitive/Taylor sufficiency route.  The
assumed \(P_{\mathrm{state}}\) and \(P_{\mathrm{acc}}\) estimates are exactly
the not-yet-established root inputs listed in
Table~\ref{tab:d5-primitive-blocker-ledger}.
With PL2 and PL3 already fixed in the accepted lower-pair chart, PL4
propagates those root rates to
\(\|\delta\lambda\|=O(h^7)\), so \(P_{\lambda}\) would follow conditionally.  The
row-level smooth force/torque/friction corollary then bounds the 72 smooth-map
subterms from \(P_{\mathrm{state}}\) and, for friction terms,
\(P_{\lambda}\).  The multiplier-geometry chart reduction bounds the 36 geometry-dependent
subterms from \(P_{\mathrm{state}}\) and \(P_{\lambda}\), and the gyroscopic
bilinear reduction bounds the 18 rotational bilinear subterms from the
angular-rate part of \(P_{\mathrm{state}}\).  The compact-tube primitive
supplies the uniform derivative constants for all smooth maps.  Applying the
recorded Taylor/mean-value template to each of the finite 162 subterms and
taking the finite row-scaled dynamic-block norm gives the displayed bound.
More explicitly, enumerate the primitive Taylor subterms as
\(\tau_i\), \(i=1,\ldots,162\), after the row scaling and component ordering
used by the D5 inventory have been fixed.  Conditional on the retained
primitive inputs, the row corollaries provide constants \(C_i\), independent
of \(h\), with \(\|\tau_i\|\le C_i h^7\) on the accepted compact tube.  The
dynamic-block norm is finite dimensional and the row weights are fixed before
the estimate is read, so a norm-equivalence constant \(C_W\), also independent
of \(h\), gives the aggregate constant
\(C_{\mathrm{prim}}=C_W\sum_{i=1}^{162} C_i\).  This yields
\(\|R_{\mathrm{dyn}}^{\mathrm{prim}}(Z_G)\|\le C_{\mathrm{prim}}h^7\).

No inverse projection from the aggregate \(132\)-row residual is used in this
argument, and no direct-route corollary is promoted into a primitive
estimate.  Since the two root estimates are not currently proved, the
argument records only the conditional implication order that a separate
primitive/Taylor proof for that route must satisfy.
\end{proof}

This proposition records the separate primitive route as a conditional
sufficiency map, not as actual primitive-route closure and not as
an input to Theorem~\ref{thm:g6fullva-order}: actual primitive Taylor-term
closure is not established on the accepted theorem route.  No theorem arrow consumes this
proposition unless the root-lift estimates are separately proved on the fixed
theorem branch before the theorem constants and admissible window are chosen.

\begin{corollary}[Row-level smooth force/torque/friction Taylor bound under
\(P_{\mathrm{state}}\) and \(P_{\lambda}\)]
\label{cor:d5-smooth-force-row-bound-under-lifts}
Let \(\mathcal T_{\mathrm{smooth}}\) be the \(72\) scalar smooth-map D5
subterm slots in Eq.~\eqref{eq:d5-smooth-template-slots}:
\begin{equation}
\label{eq:d5-smooth-template-slots}
\begin{gathered}
T_{\mathrm{external\_force\_lift}},\quad
T_{\mathrm{friction\_force\_lift}},\quad
R_{\mathrm{external\_torque\_lift}},\quad
R_{\mathrm{friction\_torque\_lift}}
\end{gathered}
\end{equation}
over the three Gauss stages, two bodies, and three scalar components.  Write
\(z=(q,v,\omega)\) for the accepted stage state variables used by these maps.
Assume the state part of \(P_{\mathrm{state}}\) gives
Eq.~\eqref{eq:d5-smooth-state-lift},
\begin{equation}
\label{eq:d5-smooth-state-lift}
  \|\widehat z_{s,b}-z_b(t_s)\|\le C_z h^7
\end{equation}
uniformly for every stage \(s\) and body \(b\), and assume the multiplier
part of \(P_{\lambda}\) gives Eq.~\eqref{eq:d5-smooth-lambda-lift}:
\begin{equation}
\label{eq:d5-smooth-lambda-lift}
  \|\widehat\lambda_s-\lambda(t_s)\|\le C_\lambda h^7 .
\end{equation}
On the compact smooth proof tube, let \(L_{\mathrm{ext}}\) bound the
derivatives of the external force and external torque maps with respect to
state, and let \(L_{\mathrm{fric}}\) bound the derivatives of the regularized
friction force and torque maps with respect to state and multiplier
variables.  Then every external force/torque scalar subterm satisfies
Eq.~\eqref{eq:d5-smooth-external-row-bound},
\begin{equation}
\label{eq:d5-smooth-external-row-bound}
  |e_r^T(F_{\mathrm{ext}}(\widehat z_{s,b},t_s)
        -F_{\mathrm{ext}}(z_b(t_s),t_s))|
  \le L_{\mathrm{ext}}C_z h^7 ,
\end{equation}
and every regularized friction force/torque scalar subterm satisfies
Eq.~\eqref{eq:d5-smooth-friction-row-bound}:
\begin{equation}
\label{eq:d5-smooth-friction-row-bound}
  |e_r^T(F_{\mathrm{fric}}(\widehat z_{s,b},\widehat\lambda_s)
        -F_{\mathrm{fric}}(z_b(t_s),\lambda(t_s)))|
  \le L_{\mathrm{fric}}(C_z+C_\lambda)h^7 .
\end{equation}
Thus the \(72\) smooth force, torque, and regularized-friction Taylor
subterms have the required \(O(h^7)\) row-level bounds under
\(P_{\mathrm{state}}\) and \(P_{\lambda}\).  This is a conditional
consequence only: it does not prove either lift, does not turn the D4
direct-route smoothness certificate into a primitive-rate proof, and does not
change the primitive/Taylor PC2 status.  It also does not use
Lemma~\ref{lem:stage-residual-defect}, direct substitution, or any
residual-to-root perturbation estimate as a primitive-bound proof.
\end{corollary}

\begin{proof}
For an external force or torque map, Taylor's formula with integral remainder
on the compact tube gives Eq.~\eqref{eq:d5-smooth-external-taylor-identity}:
\begin{equation}
\label{eq:d5-smooth-external-taylor-identity}
\begin{aligned}
  F_{\mathrm{ext}}(\widehat z,t)-F_{\mathrm{ext}}(z,t)
  &=
  \int_0^1 D_zF_{\mathrm{ext}}(z+\theta(\widehat z-z),t)\\
  &\quad{}\times
  (\widehat z-z)\,d\theta .
\end{aligned}
\end{equation}
The derivative norm is bounded by \(L_{\mathrm{ext}}\), so the
\(P_{\mathrm{state}}\) input gives Eq.~\eqref{eq:d5-smooth-external-row-bound}.
For the regularized friction maps, apply the same formula to the joint
variable \((z,\lambda)\).  The compact-tube derivative bound
\(L_{\mathrm{fric}}\) and the two conditional lift inputs give
Eq.~\eqref{eq:d5-smooth-friction-proof-bound}:
\begin{equation}
\label{eq:d5-smooth-friction-proof-bound}
\begin{aligned}
\|F_{\mathrm{fric}}(\widehat z,\widehat\lambda)
        -F_{\mathrm{fric}}(z,\lambda)\|
&\le L_{\mathrm{fric}}\bigl(\|\widehat z-z\|
      +\|\widehat\lambda-\lambda\|\bigr)\\
&\le L_{\mathrm{fric}}(C_z+C_\lambda)h^7 .
\end{aligned}
\end{equation}
Projection onto any scalar component can only decrease the Euclidean norm,
and the D5 inventory has exactly \(18\) external force, \(18\) friction force,
\(18\) external torque, and \(18\) friction torque scalar subterm slots.  The
result therefore holds uniformly over all \(72\) listed smooth-map subterms.
The proof uses only compact smoothness and the conditional lift inputs; it
does not use Lemma~\ref{lem:stage-residual-defect}, direct substitution, or
finite probes as rate evidence.
\end{proof}

\begin{lemma}[Conditional multiplier-geometry chart reduction]
\label{lem:d5-geom-chart-reduction}
Let \(G_b(q)\) and \(H_b(q)\) denote the translational and rotational
lower-pair multiplier-wrench geometry maps for body \(b\) in the accepted
Lie chart, and define the wrench maps in
Eq.~\eqref{eq:d5-geom-wrench-maps}:
\begin{equation}
\label{eq:d5-geom-wrench-maps}
  W_G(q,\lambda)=G_b(q)^T\lambda,\qquad
  W_H(q,\lambda)=H_b(q)^T\lambda .
\end{equation}
On the compact proof tube from Lemma~\ref{lem:d5-compact-tube-constants},
there are multiplier-geometry constants \(C_{WG},C_{WH}\) such that
Eq.~\eqref{eq:d5-geom-chart-lipschitz-wg} holds,
\begin{equation}
\label{eq:d5-geom-chart-lipschitz-wg}
  \|W_G(\widehat q,\widehat\lambda)-W_G(q,\lambda)\|
  \le C_{WG}\bigl(\|\widehat q-q\|+\|\widehat\lambda-\lambda\|\bigr),
\end{equation}
and the same estimate holds for \(W_H\) with \(C_{WH}\).  Hence the \(36\)
multiplier-geometry Taylor subterms are \(O(h^7)\) once the pose and
multiplier lifts supplied by \(P_{\mathrm{state}}\) and \(P_{\lambda}\) are
\(O(h^7)\).  This records the conditional chart reduction for
\(P_{\mathrm{geom}}\), but it does not close \(P_{\mathrm{geom}}\) or PC2
because those lift estimates remain open.  Here ``conditional chart
reduction'' means only this downstream implication: the Lipschitz map from the open
\(P_{\mathrm{state}}\) and \(P_{\lambda}\) inputs to the \(36\) geometry
rows has no further internal chart obligation.  It is not a primitive
\(P_{\mathrm{geom}}\) proof and certifies no actual D5 primitive Taylor
subterm without those lifts.
\end{lemma}

\begin{proof}
Let \(A(q)\) denote either \(G_b(q)\) or \(H_b(q)\).  Add and subtract
\(A(q)^T\widehat\lambda\) to get
Eq.~\eqref{eq:d5-geom-add-subtract}:
\begin{equation}
\label{eq:d5-geom-add-subtract}
\begin{aligned}
 A(\widehat q)^T\widehat\lambda-A(q)^T\lambda
 &=
 \bigl(A(\widehat q)-A(q)\bigr)^T\widehat\lambda
 + A(q)^T(\widehat\lambda-\lambda).
\end{aligned}
\end{equation}
The compact-tube constant lemma gives uniform bounds for \(A\), \(D A\), and
\(\widehat\lambda\).  The mean-value theorem then yields
Eq.~\eqref{eq:d5-geom-proof-lipschitz}:
\begin{equation}
\label{eq:d5-geom-proof-lipschitz}
\|A(\widehat q)^T\widehat\lambda-A(q)^T\lambda\|
\le C_W\bigl(\|\widehat q-q\|+\|\widehat\lambda-\lambda\|\bigr).
\end{equation}
The term-budget analysis identifies \(18\) translational
\(G_b(q)^T\lambda\) terms and \(18\) rotational \(H_b(q)^T\lambda\) terms.
Applying the same estimate to this finite row set only changes the constant.
The result is conditional on the independent \(P_{\mathrm{state}}\) and
\(P_{\lambda}\) lift estimates and cannot be used to prove either estimate.
\end{proof}

\begin{corollary}[Row-level multiplier-geometry Taylor bound under
\(P_{\mathrm{state}}\) and \(P_{\lambda}\)]
\label{cor:d5-geom-row-bound-under-lifts}
Let Eq.~\eqref{eq:d5-geom-row-set} define the global-row set for the
multiplier-geometry Newton--Euler Taylor components:
\begin{equation}
\label{eq:d5-geom-row-set}
\begin{aligned}
\mathcal R_{\mathrm{geom}}=\{&
24,25,26,27,28,29,30,31,32,33,34,35,\\
&68,69,70,71,72,73,74,75,76,77,78,79,\\
&112,113,114,115,116,117,118,119,120,121,122,123\}
\end{aligned}
\end{equation}
Assume the pose and multiplier parts of
\(P_{\mathrm{state}}\) and \(P_{\lambda}\) give
Eq.~\eqref{eq:d5-geom-lift-assumptions},
\begin{equation}
\label{eq:d5-geom-lift-assumptions}
  \|\widehat q_{s,b}-q_b(t_s)\|\le C_qh^7,\qquad
  \|\widehat\lambda_s-\lambda(t_s)\|\le C_\lambda h^7
\end{equation}
uniformly for every Gauss stage \(s\) and body \(b\).  If
Eq.~\eqref{eq:d5-geom-wrench-lipschitz-assumptions} holds,
\begin{equation}
\label{eq:d5-geom-wrench-lipschitz-assumptions}
  \begin{aligned}
  \|W_G(\widehat q,\widehat\lambda)-W_G(q,\lambda)\|
  &\le C_{WG}\bigl(\|\widehat q-q\|+\|\widehat\lambda-\lambda\|\bigr),\\
  \|W_H(\widehat q,\widehat\lambda)-W_H(q,\lambda)\|
  &\le C_{WH}\bigl(\|\widehat q-q\|+\|\widehat\lambda-\lambda\|\bigr),
  \end{aligned}
\end{equation}
and \(C_{W,\max}=\max(C_{WG},C_{WH})\), then every scalar row
\(r\in\mathcal R_{\mathrm{geom}}\) satisfies
Eq.~\eqref{eq:d5-geom-scalar-row-bound}:
\begin{equation}
\label{eq:d5-geom-scalar-row-bound}
\begin{aligned}
  |e_r^T(W_\star(\widehat q_{s,b},\widehat\lambda_s)
        -W_\star(q_b(t_s),\lambda(t_s)))|
  &\le C_{\mathrm{geom}}h^7,\\
  C_{\mathrm{geom}}&=C_{W,\max}(C_q+C_\lambda),
\end{aligned}
\end{equation}
where \(W_\star\) is \(W_G\) on translational multiplier-force rows and
\(W_H\) on rotational multiplier-torque rows.  Thus all \(36\)
multiplier-geometry Taylor rows satisfy the conditional T3 target \(O(h^7)\)
row bound under the two lift inputs.  This is a conditional row-level consequence only: it
does not prove either lift, does not close \(P_{\mathrm{geom}}\), and does
not change the PC2 status.
\end{corollary}

\begin{proof}
Apply Lemma~\ref{lem:d5-geom-chart-reduction} at each listed stage-body pair.
The assumed \(P_{\mathrm{state}}\) and \(P_{\lambda}\) estimates give
Eq.~\eqref{eq:d5-geom-proof-row-bound}:
\begin{equation}
\label{eq:d5-geom-proof-row-bound}
  \|W_\star(\widehat q_{s,b},\widehat\lambda_s)
        -W_\star(q_b(t_s),\lambda(t_s))\|
  \le C_{W,\max}(C_q+C_\lambda)h^7 .
\end{equation}
Projection onto any scalar component can only decrease the Euclidean norm.
The row inventory contains exactly the listed translational
\(G_b(q)^T\lambda\) rows and rotational \(H_b(q)^T\lambda\) rows for each
body at each Gauss stage, so the estimate holds uniformly over all rows in
\(\mathcal R_{\mathrm{geom}}\).
\end{proof}

\begin{lemma}[Gyroscopic bilinear reduction]
\label{lem:d5-gyro-bilinear-reduction}
Let \(G(\omega)=\omega\times J\omega\), where the body inertia \(J\) is
constant.  On the compact proof tube from
Lemma~\ref{lem:d5-compact-tube-constants}, the gyroscopic contribution is
Lipschitz in the angular velocity: if
\(\|\omega\|,\|\widehat\omega\|\le M\), then
Eq.~\eqref{eq:d5-gyro-lipschitz-bound} holds:
\begin{equation}
\label{eq:d5-gyro-lipschitz-bound}
  \|G(\widehat\omega)-G(\omega)\|
  \le 2M\|J\|\,\|\widehat\omega-\omega\|.
\end{equation}
Consequently, the \(18\) rotational Newton--Euler gyroscopic subterms are
\(O(h^7)\) once the angular-velocity lift supplied by
\(P_{\mathrm{state}}\) is \(O(h^7)\).  This closes the algebraic reduction
for \(P_{\mathrm{gyro}}\), but it does not close \(P_{\mathrm{gyro}}\) or
PC2 because \(P_{\mathrm{state}}\) remains open.  Here ``closes the
algebraic reduction'' means only that the bilinear gyroscopic map has no
remaining algebraic subproof under the open
\(P_{\mathrm{state}}\) angular-velocity input.  It is not a primitive
\(P_{\mathrm{gyro}}\) proof and certifies no actual D5 primitive Taylor
subterm without that lift.
\end{lemma}

\begin{proof}
Using bilinearity of the cross product and the constant inertia tensor gives
Eq.~\eqref{eq:d5-gyro-bilinear-identity}:
\begin{equation}
\label{eq:d5-gyro-bilinear-identity}
\begin{aligned}
G(\widehat\omega)-G(\omega)
  &=(\widehat\omega-\omega)\times J\widehat\omega
   +\omega\times J(\widehat\omega-\omega).
\end{aligned}
\end{equation}
The Euclidean cross-product bound and the compact-tube angular-velocity bound
give Eq.~\eqref{eq:d5-gyro-proof-lipschitz}:
\begin{equation}
\label{eq:d5-gyro-proof-lipschitz}
\begin{aligned}
\|G(\widehat\omega)-G(\omega)\|
&\le \|\widehat\omega-\omega\|\,\|J\|\,\|\widehat\omega\|
   +\|\omega\|\,\|J\|\,\|\widehat\omega-\omega\|\\
&\le 2M\|J\|\,\|\widehat\omega-\omega\|.
\end{aligned}
\end{equation}
The rotational-row inventory contains exactly three gyroscopic components for
each body at each of the three Gauss stages, hence \(18\) rows.  Applying the
same bound to this finite row set only changes the constant.  The result is
therefore conditional on the independent angular-velocity lift estimate from
\(P_{\mathrm{state}}\); it is not a proof of that lift estimate and cannot be
used to close the dynamic residual criterion.
\end{proof}

\begin{corollary}[Row-level gyroscopic Taylor bound under \(P_{\mathrm{state}}\)]
\label{cor:d5-gyro-row-bound-under-pstate}
Let Eq.~\eqref{eq:d5-gyro-row-set} define the global-row set for the
rotational gyroscopic Newton--Euler components:
\begin{equation}
\label{eq:d5-gyro-row-set}
\begin{aligned}
\mathcal R_{\mathrm{gyro}}=\{&
27,28,29,33,34,35,\\
&71,72,73,77,78,79,\\
&115,116,117,121,122,123\}
\end{aligned}
\end{equation}
Assume the angular-velocity part of \(P_{\mathrm{state}}\) gives
Eq.~\eqref{eq:d5-gyro-lift-assumption},
\begin{equation}
\label{eq:d5-gyro-lift-assumption}
  \|\widehat\omega_{s,b}-\omega_b(t_s)\|\le C_\omega h^7
\end{equation}
uniformly for every Gauss stage \(s\) and body \(b\).  If the compact tube
also gives Eq.~\eqref{eq:d5-gyro-compact-inertia-bound},
\begin{equation}
\label{eq:d5-gyro-compact-inertia-bound}
  \|\omega_b(t_s)\|,\|\widehat\omega_{s,b}\|\le M_\omega,
  \qquad
  J_{\max}=\max_b\|J_b\|_2,
\end{equation}
then every scalar row \(r\in\mathcal R_{\mathrm{gyro}}\) satisfies
Eq.~\eqref{eq:d5-gyro-scalar-row-bound}:
\begin{equation}
\label{eq:d5-gyro-scalar-row-bound}
\begin{aligned}
  |e_r^T(G_b(\widehat\omega_{s,b})-G_b(\omega_b(t_s)))|
  &\le C_{\mathrm{gyro}} h^7,\\
  C_{\mathrm{gyro}}&=2M_\omega J_{\max}C_\omega .
\end{aligned}
\end{equation}
Thus the \(18\) gyroscopic Taylor rows satisfy the conditional T3 target
\(O(h^7)\) row bound under \(P_{\mathrm{state}}\).  This is a conditional row-level consequence
only: it does not prove the angular-velocity lift, does not close
\(P_{\mathrm{gyro}}\), and does not change the PC2 status.
\end{corollary}

\begin{proof}
Apply Lemma~\ref{lem:d5-gyro-bilinear-reduction} with
\(\widehat\omega=\widehat\omega_{s,b}\), \(\omega=\omega_b(t_s)\), and
\(\|J_b\|_2\le J_{\max}\).  The assumed \(P_{\mathrm{state}}\) angular-rate
estimate gives Eq.~\eqref{eq:d5-gyro-proof-row-bound}:
\begin{equation}
\label{eq:d5-gyro-proof-row-bound}
\|G_b(\widehat\omega_{s,b})-G_b(\omega_b(t_s))\|
\le 2M_\omega J_{\max} C_\omega h^7 .
\end{equation}
Projection onto any scalar component can only decrease the Euclidean norm.
The row inventory contains exactly the listed three rotational
components for each body at each Gauss stage, so the estimate holds uniformly
over all rows in \(\mathcal R_{\mathrm{gyro}}\).
\end{proof}

\paragraph{Primitive/Taylor nonclosure invariant}
The row-level acceleration, smooth-force, geometry, and gyroscopic corollaries above strengthen the
conditional Taylor-reduction record, but they do not certify actual
primitive-route Taylor bounds and do not change the primitive route status.
They are downstream conditional maps, not roots of the primitive dependency
graph.
They are conditional consequences of the still-open lift estimates
\(P_{\mathrm{state}}\), \(P_{\mathrm{acc}}\), \(P_{\lambda}\),
\(P_{\mathrm{geom}}\), and \(P_{\mathrm{gyro}}\).  Hence the primitive
route still has zero actual primitive-route Taylor subterm bounds out of the 162
D5 subterms until those five primitive obligations are proved uniformly on
the same compact proof tube.  The PC2 stage-residual condition used in
Theorem~\ref{thm:g6fullva-order} remains discharged by the direct
residual-bridge route, which consumes the D5 direct-substitution certificate rather than the
open primitive/Taylor route.

\paragraph{Primitive-route one-way implication discipline}
The separate primitive implication has the one-way logical form in
Eq.~\eqref{eq:d5-primitive-one-way-discipline}:
\begin{equation}
\label{eq:d5-primitive-one-way-discipline}
\begin{aligned}
&P_{\mathrm{state}},P_{\mathrm{acc}},P_{\lambda},
  P_{\mathrm{geom}},P_{\mathrm{gyro}},P_{\mathrm{tube}}\\
&\quad\Longrightarrow
  \{\text{162 D5 primitive Taylor subterms are }O(h^7)\}\\
&\quad\Longrightarrow R_{\mathrm{dyn}}(Z_G)=O(h^7).
\end{aligned}
\end{equation}
No converse implication is used or available in this manuscript.  In
particular, the active direct route proves the dynamic-row residual value at
\(Z_G\) by direct Newton--Euler substitution and then uses the full
row-scaled residual inverse to obtain a weighted stage-root perturbation.
That output cannot be reinserted as a proof of the primitive inputs.  The
weighted estimate \(h\|\delta A\|=O(h^7)\) does not imply the unweighted
\(P_{\mathrm{acc}}\) estimate \(\|\delta A\|=O(h^7)\), and
Lemma~\ref{lem:d5-p-state-ps3-full-residual-route} is a direct-route
state-block corollary, not a primitive \(P_{\mathrm{state}}\) closure.
Therefore the conditional row-level reductions below may certify the finite
implication from primitive inputs to Taylor subterms, but they do not turn
the already proved direct residual bridge into a primitive 162-subterm
proof.  Actual primitive-route Taylor subterm bounds are not established until
the five open primitive obligations are proved
independently on the same compact proof tube.
The current primitive-status accounting is therefore positive but
non-closing: \(P_{\mathrm{tube}}\) supplies the common compact-tube
constant slot, the finite reduction templates cover all \(162\) recorded D5
subterms, and the remaining proof burden is exactly the five same-object
lift/bilinear inputs
\(P_{\mathrm{state}},P_{\mathrm{acc}},P_{\lambda},
P_{\mathrm{geom}},P_{\mathrm{gyro}}\).
Since none of those five inputs is proved on the accepted theorem tuple, the
recorded template coverage gives no actual primitive Taylor bound and no
primitive-route PC2 certificate.  This is a strength statement, not a
diagnostic downgrade: the route has a typed finite reduction once its root
primitive estimates are supplied, while the accepted theorem uses the direct
\(96+36\) residual bridge already available for the same \(132\)-row object.

\noindent\textbf{Taylor-projection checkpoint.}
The \(132\)-row residual estimate is an aggregate estimate in the fixed block
norm of the row-scaled residual map.  It is not a projection certificate for
the 162 primitive scalar subterms.  To promote an aggregate residual estimate
into primitive-route closure one would first need, before invoking the
theorem-level local defect, a fixed primitive decomposition operator
\(\Pi_{\rm prim}\), an \(h\)-independent primitive projection/right-inverse
or norm-equivalence bound on the compact tube, and independent \(O(h^7)\)
bounds for all primitive lift maps feeding \(\Pi_{\rm prim}\).  No such
primitive projection/right-inverse is used in
Theorem~\ref{thm:g6fullva-order}.  Hence the aggregate residual bridge can
close the T2 residual-defect input while the primitive-route status
remains separate from theorem use; this is a one-way proof route, not a hidden projection from
theorem order back to primitive Taylor subterms.

\noindent\textbf{Primitive-route non-substitution contract.}
If a separate T3 primitive/Taylor route is established, its first use
in the present proof architecture would be limited to replacing the
dynamic-row residual-value certificate inside the full \(132\)-row bridge. It
would not by itself replace the T2 full-map inverse/Kantorovich step, the
endpoint-closure perturbation, the retained P6 solver-scale condition, the
local-to-global stability transfer, or the P7 residual-to-error theorem. Those
objects are separate proof slots with their own fixed branch, norm, constants,
and scope rules. Conversely, the current direct bridge closes the
theorem-level residual value without proving a primitive projection/right-inverse
or any of the \(162\) primitive subterm bounds. Thus T3 closure would be an
alternative certificate for one input to the existing theorem chain, not a
complete order theorem and not a retroactive proof of the current primitive
status.

\noindent\textbf{Reference-proof boundary for the primitive Taylor route.}
The Wieloch--Arnold constrained-BDF proof motivates the proof-order discipline applied
above, but it supplies no primitive FullVA lift estimate.  Its local
truncation proof is tied to its BLieDF discrete equations, hidden-constraint
coordinates, start-value assumptions, and multiplier recursion; it is not an
estimate for \(P_{\mathrm{state}}\), \(P_{\mathrm{acc}}\),
\(P_{\lambda}\), \(P_{\mathrm{geom}}\), \(P_{\mathrm{gyro}}\), or for a
primitive projection \(\Pi_{\rm prim}\) on the \(162\) D5 subterms.  The
criterion below is therefore not a borrowed Taylor local-error theorem.  It
is a FullVA-specific finite implication: after the five primitive inputs are
proved in the accepted stage variables and on the same compact proof tube,
the listed Taylor templates imply the dynamic-row primitive defect.  Until
those inputs are proved independently, the reference paper cannot close the
primitive/Taylor status map or supply PC2.

\noindent\textbf{Reference-use admissibility test for T3.}
Any reference-derived estimate proposed for the separate T3
primitive/Taylor route must pass a same-object test before it can enter the
primitive inventory: it must be stated for the accepted FullVA stage variables, row-scaled
residual norm, lower-pair chart, and compact proof tube used by
Theorem~\ref{thm:g6fullva-order}; it must prove the root lifts
\(P_{\mathrm{state}}\) and \(P_{\mathrm{acc}}\) without using
Lemma~\ref{lem:stage-residual-defect}, the direct D5 residual value, observed
grid orders, or finite residual logs; it must propagate
\(P_{\lambda}\), \(P_{\mathrm{geom}}\), and \(P_{\mathrm{gyro}}\) only
forward from those root lifts; and it must deliver row-level \(O(h^7)\)
bounds in the fixed D5 primitive inventory before any finite summation is
used.  The constrained-BDF estimates do not satisfy this test as direct
inputs because they are estimates for BLieDF hidden constraints, multistep
increments, and multiplier recursion, not for the present FullVA stage-root
variables or the unweighted acceleration primitive.  They therefore remain
ordering discipline only and do not change the non-active primitive-route
Taylor status.

\begin{lemma}[Same-object Taylor transport criterion]
\label{lem:same-object-taylor-transport-gate}
A reference-style Taylor or multiplier estimate from another Lie-group DAE
scheme can be used in the separate T3 primitive route only after it is
transported to the same mathematical object used here.  Concretely, before
such an estimate may enter the D5 primitive inventory, one must construct
\(h\)-uniformly bounded identifications from the reference perturbation
variables to the accepted FullVA stage variables, from the reference defect
components to the fixed D5 primitive inventory, and from the reference norm
to the row-scaled proof norm of \(F_{A,h}\), all on the same compact
Gauss-predictor branch and before any finite-grid or residual-log evidence is
used.  The transported estimate must then prove the FullVA root lift inputs
\(P_{\mathrm{state}}\) and \(P_{\mathrm{acc}}\) in those variables; only after
those two inputs are independently available may the present \(P_{\lambda}\),
\(P_{\mathrm{geom}}\), and \(P_{\mathrm{gyro}}\) propagation/reduction steps
be invoked.  Without such a same-object transport, the reference estimate is
not a T3 input, cannot certify any of the \(162\) primitive D5 Taylor
subterms, and cannot replace Lemma~\ref{lem:full-132-row-residual-bridge}.
\end{lemma}

\begin{proof}
The T3 route is typed by the objects listed in its inventory: the lifted Gauss
stage \(Z_G\), the accepted FullVA stage vector, the fixed lower-pair chart,
the row-scaled residual norm, and the compact proof tube used by
Theorem~\ref{thm:g6fullva-order}.  A Taylor estimate proved for a different
scheme controls a different perturbation vector and a different defect map
unless a uniformly bounded identification maps those quantities into this
typed FullVA inventory.  Because finite-dimensional norm equivalence is
available only after the spaces and maps are fixed, the transport constants
must be chosen before the estimate is consumed and cannot be fitted from
observed orders, residual histories, or source-policy rows.

After such a transport, the primitive dependency graph still has to be
respected.  The state and unweighted acceleration lifts are the root inputs
for the separate route; the multiplier, geometry, and gyroscopic reductions
are downstream same-branch consequences in the present FullVA variables.
Therefore a BLieDF hidden-constraint, multiplier-history, or multistep
recursion estimate does not become a FullVA primitive bound merely because it
has the same formal power of \(h\).  In this manuscript no same-object
transport map from the constrained-BDF reference variables to the D5 FullVA
primitive inventory is constructed.  The reference proof therefore remains an
ordering discipline only, while the primitive-route status remains at
no actual primitive D5 Taylor subterm bounds for a T3 replacement certificate
in the accepted theorem route.
\end{proof}

\noindent\textbf{Conditional primitive Taylor implication criterion (non-closure).}
The next lemma is the exact one-way criterion for the separate primitive
route.  It should be read as a finite proof implication, not as an informal
route summary: if the five open primitive lift/bilinear estimates are supplied on the
same compact tube, then the already listed Taylor templates close all
\(162\) primitive D5 subterms and hence the dynamic-row residual value.
The status recorded before this lemma is deliberately asymmetric: one common
tube primitive is already available, all \(162\) conditional template rows are
typed, but the certified count of actual primitive Taylor bounds is \(0/162\)
because the
five root lift/bilinear estimates remain open.
This is only a conditional 36-row dynamic statement: it is not yet a
theorem-facing \(132\)-row residual-value certificate, because the latter
would still have to assemble the dynamic bound with the same 96-row
non-dynamic certificate, row ordering, row scaling, block norm, branch, and
endpoint convention used by the theorem.
Equally important, the criterion states the open premises explicitly, so
the current primitive non-invocation boundary is preserved until those
premises are proved without using the direct full-residual bridge as an input.

\begin{lemma}[Conditional primitive-obligation implication for D5 (non-closure)]
\label{lem:d5-primitive-obligation-implication}
This separate primitive-route lemma is a conditional finite implication, not
a closure certificate, and is not invoked in
Theorem~\ref{thm:g6fullva-order}; it records only what would follow if the
five primitive obligations are proved separately.
Assume the five remaining primitive obligations
\(P_{\mathrm{state}},P_{\mathrm{acc}},P_{\lambda},P_{\mathrm{geom}}\), and
\(P_{\mathrm{gyro}}\) hold uniformly on the compact smooth proof tube, and use
\(P_{\mathrm{tube}}\) from Lemma~\ref{lem:d5-compact-tube-constants}.  Then,
only as a conditional implication under those five still-open primitive
hypotheses, every D5 Taylor subterm in the 36 Newton--Euler rows is
\(O(h^7)\), and the lifted-stage dynamic residual satisfies
\(R_{\mathrm{dyn}}(Z_G)=O(h^7)\) for each expanded component row, using the
already closed D1/D2, D3, D4, and D6 inputs.  This implication is not an
accepted primitive-route closure, not a PC2 input, and not invoked by
Theorem~\ref{thm:g6fullva-order}.
All quantities in this conditional implication are evaluated on the same
tuple \((h,y,Z_G,F_{A,h},\|\cdot\|_{\rm proof})\) and in the same D5 row
ordering as the active theorem.  Its output is only a 36-row primitive
dynamic-block predicate; it becomes theorem-facing only after a separately
proved \(132\)-row T3 predicate assembles that block with the same 96-row
non-dynamic certificate.  No arrow from the direct residual bridge, the T2
Taylor/Kantorovich conclusion, or reported-grid errors is used to establish
the five primitive premises.
\end{lemma}

\begin{proof}
Each translational dynamic row has four Taylor subterms and each rotational
dynamic row has five.  The recorded reduction map assigns every subterm to a
finite subset of the primitive obligations.  The four Taylor templates are
explicit.  Linear acceleration terms have no remainder, as shown in
Eq.~\eqref{eq:d5-primitive-linear-acc-terms}:
\begin{equation}
\label{eq:d5-primitive-linear-acc-terms}
 m_b(\widehat a-a),\qquad J_b(\widehat\alpha-\alpha).
\end{equation}
The conditional row-level acceleration version is
Corollary~\ref{cor:d5-p-acc-row-bound-under-pacc}; it consumes
\(P_{\mathrm{acc}}\) and does not prove PA2 or the unweighted acceleration
lift.
For any smooth force, torque, or regularized friction map \(F\), Taylor's
formula with integral remainder gives
Eq.~\eqref{eq:d5-primitive-smooth-taylor-template}:
\begin{equation}
\label{eq:d5-primitive-smooth-taylor-template}
\begin{aligned}
 F(\widehat z)-F(z)
 &=
 \int_0^1 DF(z+\theta(\widehat z-z))(\widehat z-z)\,d\theta ,
\end{aligned}
\end{equation}
so the compact-tube derivative bound gives
\(\|F(\widehat z)-F(z)\|\le C_F\|\widehat z-z\|\).  The conditional
row-level smooth-map version is
Corollary~\ref{cor:d5-smooth-force-row-bound-under-lifts}: it covers exactly
the 72 external-force, friction-force, external-torque, and friction-torque
subterms, consumes only the \(P_{\mathrm{state}}\) and \(P_{\lambda}\) lift
inputs, and does not use the D4 direct-route smoothness certificate as a
primitive rate proof.
For a multiplier geometry
map \(A(q)^T\lambda\), use
Eq.~\eqref{eq:d5-primitive-geom-taylor-template}:
\begin{equation}
\label{eq:d5-primitive-geom-taylor-template}
\begin{aligned}
 A(\widehat q)^T\widehat\lambda-A(q)^T\lambda
 &=
 A(q)^T(\widehat\lambda-\lambda)
 +\left(\int_0^1 DA(q+\theta(\widehat q-q))(\widehat q-q)\,d\theta\right)^T
 \widehat\lambda .
\end{aligned}
\end{equation}
For the gyroscopic term, use Eq.~\eqref{eq:d5-primitive-gyro-template}:
\begin{equation}
\label{eq:d5-primitive-gyro-template}
\begin{aligned}
 \widehat\omega\times J\widehat\omega-\omega\times J\omega
 &=
 (\widehat\omega-\omega)\times J\widehat\omega
 +\omega\times J(\widehat\omega-\omega).
\end{aligned}
\end{equation}
By Lemma~\ref{lem:d5-compact-tube-constants}, the constants in these four
templates are uniformly bounded on the proof tube.  The state-, acceleration-,
multiplier-, geometry-, and gyroscopic-lift obligations then bound their
assigned subterms by \(C h^7\).  Equivalently, for each dynamic component row
\(r\), let \(\mathcal I_r\) be its finite primitive-subterm index set
(\(|\mathcal I_r|=4\) for translational rows and \(|\mathcal I_r|=5\) for
rotational rows).  If
\(\|T_{\ell,h}(Z_G;y)\|\le C_\ell h^7\) for every
\(\ell\in\mathcal I_r\), then the row defect satisfies
Eq.~\eqref{eq:d5-primitive-row-summation-bound}:
\begin{equation}
\label{eq:d5-primitive-row-summation-bound}
\begin{aligned}
  \|R_{r,h}^{\rm prim}(Z_G;y)\|
  &\le \sum_{\ell\in\mathcal I_r} C_\ell h^7\\
  = C_{\mathrm{D5},r}h^7 ,
\end{aligned}
\end{equation}
where \(C_{\mathrm{D5},r}\) is independent of \(h\), the reported grid
length, backend, and finite diagnostic tolerances.  Taking the maximum of the
36 constants \(C_{\mathrm{D5},r}\) gives a uniform row-family constant because
the dynamic-row inventory is finite.  Thus finite primitive-row summation is
the only aggregation step in this separate route; it is not an inverse
projection from the aggregate 132-row residual estimate back to primitive
subterms.  The closed D1/D2 balance identities, D3 multiplier wrench
identity, D4 smooth-force input, and D6 row-binding input remove the non-lift
discrepancy terms, leaving a row-local \(O(h^7)\) dynamic defect.
\end{proof}

Lemma~\ref{lem:d5-primitive-obligation-implication} is conditional on the five
remaining primitive obligations and does not close the primitive-and-Taylor
route to PC2.  It closes only the finite implication from primitive bounds to the
162 D5 subterm bounds; the state, acceleration, multiplier, geometry, and
gyroscopic primitive bounds remain outside the active proof unless they are
separately proved on the same compact branch.
Thus the Taylor layer now has two separate theorem-facing statements: the
active full-map Taylor/Kantorovich lemma used by
Theorem~\ref{thm:g6fullva-order}, and this separate primitive implication
criterion for a stricter separate certificate.  The first is closed only after
two same-object inputs are present together: the direct \(132\)-row residual
bridge supplies the residual value \(R_h=F_{A,h}(Z_G;y)=O(h^7)\), and the
retained compact-tube Taylor/Kantorovich hypotheses supply the inverse,
quadratic-remainder, branch-ball, and endpoint-map constants for that same
row-scaled residual map.  Without those compact-inverse and remainder
interfaces, the direct bridge is only a residual-value certificate, not a
stage-root perturbation theorem.  The second records only a conditional
implication and cannot be spliced into the active theorem as partial evidence
or counted as partial closure of the primitive/Taylor route.

\paragraph{Uniform residual-scaling convention}
In Lemma~\ref{lem:stage-residual-defect} and in the full-residual bridge,
write the proof residual as \(F_{A,h}=S F^{\rm raw}_{A,h}\), where \(S\) is
a fixed finite-dimensional row-scaling map chosen with the chart, row
ordering, and proof norm.  The explicit powers of \(h\) appearing in the
collocation and FullVA row formulas are part of \(F^{\rm raw}_{A,h}\) before
the asymptotic estimates begin; there is no \(h\)-dependent left
preconditioner or row reweighting chosen from \(N_h\), backend behavior,
solver logs, or observed slopes.  All norm-equivalence constants associated
with \(S\) are therefore independent of \(h\).  The uniform inverse,
second-derivative, and ball constants \(M,L,r\) in
Lemma~\ref{lem:stage-residual-defect} are retained P2 compact-tube hypotheses
for this fixed scaled residual, not consequences of the D5 row identity,
finite rank probes, or the separate primitive/Taylor route.

\begin{lemma}[Stage-residual perturbation criterion]
\label{lem:stage-residual-defect}
Let \(K\) be the retained compact proof-tube state set from
Assumption~\ref{ass:regularity}.  For \(y\in K\), let
$F_{G,h}(Z;y)=0$ denote the reduced three-stage Gauss collocation stage
equations in the smooth chart and let $Z_G(y,h)$ be their isolated solution.
Below, \(Z_G\) denotes this \(Z_G(y,h)\) for the fixed transition input.
Let \(F_{A,h}(\cdot;y)\) be the accepted FullVA residual map written in the
same local unknowns after the Lie retraction and constraint chart are fixed.
If the rows are nonautonomous, the transition input time \(t_n\) is fixed
before this lemma is invoked and \(F_{A,h}\) is a fixed-time-slice
abbreviation for \(F_{A,h,t_n}\).  In that case \(Z_G\) denotes the lifted
Gauss stage \(Z_G(y,t_n,h)\), all driven rows are evaluated at the stage
times \(t_n+c_i h\), and the constants \(M,L,r,C_R,M_E\) are uniform over the
retained time interval as well as over \(y\in K\).  The Taylor expansion below
is only in the stage variables at fixed \(t_n\); it does not differentiate in
time, interpolate between time slices, or absorb a terminal-time correction
into \(D_ZF_{A,h}\).
The Gauss stage is first lifted by Lemma~\ref{lem:fullva-stage-lift} into this
same accepted stage-coordinate vector before \(F_{A,h}(Z_G;y)\) is evaluated.
Thus \(F_{G,h}\) is used only to define the reference stage \(Z_G\); the
Taylor expansion and inverse estimate below are taken entirely in the fixed
finite-dimensional domain and the row-scaled \(132\)-row codomain of
\(F_{A,h}\), with the proof norm chosen before the small-\(h\) argument.
The equation \(F_{A,h}(Z;y)=0\) defines the exact accepted stage root when the
local root exists; that local existence is the conclusion of the contraction
argument below, not a premise. Assume that, uniformly for \(y\in K\), on a
ball $B_r(Z_G(y,h))$ contained in the compact proof tube, the Jacobian in
Eq.~\eqref{eq:stage-pert-jacobian}
\begin{equation}
\label{eq:stage-pert-jacobian}
  J_h:=D_ZF_{A,h}(Z_G;y)
\end{equation}
is invertible with $\|J_h^{-1}\|\le M$ uniformly in $h$, and that
$D_Z^2F_{A,h}$ is bounded by $L$ on the same ball, with $M,L,r$ independent
of $h$.  Equivalently, for a preselected step window \(0<h\le \bar h\),
the same chart, row ordering, row scaling, and proof norm give the compact
second-derivative bound in Eq.~\eqref{eq:stage-pert-second-derivative-bound},
\begin{equation}
\label{eq:stage-pert-second-derivative-bound}
  \sup_{0<h\le \bar h}\sup_{y\in K}\sup_{Z\in B_r(Z_G(y,h))}
  \|D_Z^2F_{A,h}(Z;y)\|\le L
\end{equation}
on the retained compact proof tube \(K\).  The explicit \(h\)-dependence of
the Gauss weights and residual row factors is already part of the map
\(F_{A,h}\) before this supremum is taken; \(L\) is not fitted from
reported residual logs and is not supplied by the separate 162-subterm
primitive Taylor route.  Assume also, uniformly for \(y\in K\) and
\(0<h\le\bar h\), that the accepted residual evaluated at the Gauss stage
satisfies Eq.~\eqref{eq:stage-pert-gauss-residual-input},
\begin{equation}
\label{eq:stage-pert-gauss-residual-input}
  \|F_{A,h}(Z_G;y)\| \le C_R h^7 .
\end{equation}
In Theorem~\ref{thm:g6fullva-order}, this hypothesis is discharged by
Lemma~\ref{lem:full-132-row-residual-bridge}; here it is isolated only to
state the perturbation criterion.
After possibly shrinking the admissible step window inside \(\bar h\), choose
\(h_0>0\) so that Eq.~\eqref{eq:stage-pert-small-step} holds:
\begin{equation}
\label{eq:stage-pert-small-step}
  2MC_Rh_0^7\le r,\qquad 2M^2LC_Rh_0^7\le \frac12 .
\end{equation}
Then, for \(0<h\le h_0\), with
the contraction radius defined in Eq.~\eqref{eq:stage-pert-radius},
\begin{equation}
\label{eq:stage-pert-radius}
  \rho_h=2M\|F_{A,h}(Z_G;y)\|,
\end{equation}
there is a unique root \(Z_A\) in the contraction ball
\(Z_G+B_{\rho_h}(0)\subset B_r(Z_G)\).  This root is the local accepted
root on the Newton branch selected from the Gauss predictor, and it satisfies
the explicit stage-residual-to-root estimate in
Eq.~\eqref{eq:stage-pert-root-bound},
\begin{equation}
\label{eq:stage-pert-root-bound}
  \|Z_A-Z_G\|\le C_Zh^7,\qquad C_Z=2MC_R .
\end{equation}
If the endpoint reconstruction satisfies the corresponding uniform derivative
bound in Eq.~\eqref{eq:stage-pert-endpoint-derivative-bound},
\begin{equation}
\label{eq:stage-pert-endpoint-derivative-bound}
  \sup_{0<h\le h_0}\sup_{y\in K}
  \sup_{Z\in B_r(Z_G(y,h))}
  \|D_Z\mathcal E_h(Z)\|\le M_E
\end{equation}
on the same compact tube, then its accepted endpoint differs from the Gauss
endpoint by Eq.~\eqref{eq:stage-pert-endpoint-bound},
\begin{equation}
\label{eq:stage-pert-endpoint-bound}
  \|\mathcal E_h(Z_A)-\mathcal E_h(Z_G)\|
  \le C_Ah^7,\qquad C_A=M_EC_Z .
\end{equation}
This is a local Newton-branch statement, not a global uniqueness claim for
remote nonlinear roots of $F_{A,h}$. The uniform inverse hypothesis is a
compact-tube Gauss-predictor-branch condition; finite rank probes are not a
compact-tube inverse proof and are not used by this perturbation criterion.
Compact-inverse quantifier discipline: the constants \(M,L,r\), the proof norm,
and the Gauss-predictor branch are fixed before the residual defect
\(C_Rh^7\) is applied. The proof may shrink \(h_0\) after these constants are
fixed, but it does not choose \(M,L,r\), the tube, or the branch from the
computed accepted trajectory. The uniform stage-root threshold is chosen once
on the compact proof tube; the same contraction radius applies to every
retained-tube state and is not selected separately for each reported step.
The Taylor expansion in this criterion is the expansion of the full scaled
map \(F_{A,h}\) with respect to the complete stage vector \(Z\).  It therefore
uses the single bridged residual value \(R_h=F_{A,h}(Z_G;y)\) as input; it does
not require, and does not infer, a termwise certification of the 162 primitive
Newton--Euler Taylor subterms.
\noindent\textbf{Forward-only Taylor direction.}
This is the FullVA analogue of the constrained-BDF consistency step in the
reference proof: the exact smooth branch is inserted into a fixed discrete
operator, a local defect is proved for that operator, and only then is the
defect transported through the local inverse and global error recursion.  The
analogy is specifically to the reference paper's Section~6.2
Theorems~2--4 sequence \cite{wieloch2021bdf}: local truncation errors are
formed by inserting the analytic constrained solution into the already
specified BLieDF equations, the configuration/velocity error recursion is then
formed, and only afterward are the hidden-constraint, multiplier, and
convergence estimates used.  Here the inserted object is \(Z_G\) and the fixed equations are the
row-scaled \(F_{A,h}\), so the Taylor/Kantorovich inverse acts only after the
same-tuple residual value has been proved.
This direction is not reversible.  A small endpoint error, an observed sixth-order
slope, or an accepted Newton residual cannot be used to infer
Eq.~\eqref{eq:stage-pert-gauss-residual-input}; that residual value must be
proved before the contraction map is invoked.  Conversely, a future primitive
T3 bound would be only a different proof of the same residual-value input for
the same \(132\)-row map.  It would not replace the compact inverse, the
quadratic-remainder bound, the endpoint map, the retained P6 solver-scale
condition, or the local-to-global transfer.
\noindent\textbf{Parameter-dependent Taylor-family rule.}
The family \(h\mapsto F_{A,h}\) is treated as a parameterized family of
finite-dimensional stage maps.  For each theorem instance, \(h\), \(t_n\),
the row scaling, and the branch chart are fixed first; Taylor's formula is
then applied only in the stage variable \(Z\).  The uniform bound on
\(D_Z^2F_{A,h}\) is a supremum over this predeclared family on
\(0<h\le\bar h\), not a derivative with respect to \(h\), not a comparison
between residual maps at two step sizes, and not a license to retune row
weights or residual formulas with \(h\) after inspecting diagnostics.  Thus
the quadratic remainder \(Q_h\) below belongs to the same frozen residual map
as \(R_h\), \(J_h\), and the endpoint reconstruction.
Single-instance requirement: every object in the displayed hypotheses must
belong to the same theorem instance.  The branch ball, inverse bound, second
derivative bound, row scaling, proof norm, residual value, and endpoint map
are fixed together before the estimate is invoked.  Replacing one of these
objects by a finite rank probe, a residual log in a different norm, a
postprocessed endpoint map, or a root found on another nonlinear branch does
not give a weaker conclusion of this lemma; it changes the lemma's input data
and leaves the perturbation estimate unavailable for that assembled object.
\end{lemma}

\begin{proof}
Write $\Delta=Z-Z_G$ and $R_h=F_{A,h}(Z_G;y)$. Taylor's formula with integral
remainder gives Eq.~\eqref{eq:stage-pert-taylor-expansion}, for
$\|\Delta\|\le r$,
\begin{equation}
\label{eq:stage-pert-taylor-expansion}
  F_{A,h}(Z_G+\Delta;y)
  =R_h+J_h\Delta+Q_h(\Delta),
\end{equation}
where the quadratic remainder is defined in
Eq.~\eqref{eq:stage-pert-remainder},
\begin{equation}
\label{eq:stage-pert-remainder}
  Q_h(\Delta)
  =
  \int_0^1(1-\theta)
  D_Z^2F_{A,h}(Z_G+\theta\Delta;y)[\Delta,\Delta]\,d\theta
\end{equation}
and therefore $\|Q_h(\Delta)\|\le (L/2)\|\Delta\|^2$.  This is the only
Taylor remainder used by the perturbation step: the bound is taken on the
full implemented residual family after the finite row-scaling convention is
fixed, rather than by summing primitive Newton--Euler subterm estimates.
Here \(h\) is a frozen parameter of the map being expanded.  No term in
\(Q_h\) differentiates the residual with respect to \(h\), changes the stage
abscissae, or compares \(F_{A,h}\) with \(F_{A,\hat h}\); all such
step-size variation is outside this Taylor remainder and can only enter
through the predeclared uniform constants.  The zero equation is
equivalent to the fixed-point problem in
Eq.~\eqref{eq:stage-pert-fixed-point-map},
\begin{equation}
\label{eq:stage-pert-fixed-point-map}
  \Delta
  =
  \mathcal T_h(\Delta)
  :=
  -J_h^{-1}R_h-J_h^{-1}Q_h(\Delta).
\end{equation}
Let $\rho_h=2M\|R_h\|$.  The chosen \(h_0\) gives, for \(0<h\le h_0\),
$\rho_h\le r$ and $ML\rho_h\le 1/2$.  If $\|\Delta\|\le\rho_h$, then
Eq.~\eqref{eq:stage-pert-ball-map} shows that the ball maps into itself:
\begin{equation}
\label{eq:stage-pert-ball-map}
  \|\mathcal T_h(\Delta)\|
  \le M\|R_h\|+\frac{ML}{2}\rho_h^2
  \le \rho_h .
\end{equation}
Moreover, using the integral mean-value form of $Q_h(\Delta_1)-Q_h(\Delta_2)$
on the same ball gives the contraction estimate in
Eq.~\eqref{eq:stage-pert-contraction},
\begin{equation}
\label{eq:stage-pert-contraction}
  \|\mathcal T_h(\Delta_1)-\mathcal T_h(\Delta_2)\|
  \le ML\rho_h\|\Delta_1-\Delta_2\|
  \le \frac12\|\Delta_1-\Delta_2\|.
\end{equation}
Thus $\mathcal T_h$ is a contraction on $B_{\rho_h}(0)$ and has a unique
fixed point $\Delta_A$ there.  The corresponding zero
$Z_A=Z_G+\Delta_A$ satisfies Eq.~\eqref{eq:stage-pert-proof-root-bound},
\begin{equation}
\label{eq:stage-pert-proof-root-bound}
  \|Z_A-Z_G\|=\|\Delta_A\|\le\rho_h
  \le 2MC_Rh^7=C_Zh^7.
\end{equation}
If $\mathcal E_h$ is any smooth Lie-retraction endpoint reconstruction, then
Taylor's formula in the stage variables gives
Eq.~\eqref{eq:stage-pert-endpoint-taylor},
\begin{equation}
\label{eq:stage-pert-endpoint-taylor}
  \mathcal E_h(Z_A)-\mathcal E_h(Z_G)
  =
  \int_0^1D_Z\mathcal E_h(Z_G+\theta\Delta_A)\Delta_A\,d\theta ,
\end{equation}
and the displayed compact-tube derivative bound gives
\(\|\mathcal E_h(Z_A)-\mathcal E_h(Z_G)\|\le M_EC_Zh^7=C_Ah^7\).
The contraction ball is the local branch used by the accepted Newton solve;
the argument neither requires nor asserts uniqueness of possible roots outside
the compact-tube neighborhood. Since \(M,L,r,C_R\) are compact-tube constants,
the small-step restriction just used is uniform over the retained proof tube
and over the reported transition indices; it is not chosen from a solved
trajectory, a residual log, or a per-step rank probe.  The same fixed residual
map and norm are used in \(R_h\), \(J_h\), \(Q_h\), the contraction ball, and
the endpoint reconstruction.  This same-object condition is what makes the
Kantorovich step a proof of a branch-selected local map rather than a
post-hoc consistency check assembled from separate diagnostics.
\end{proof}

\noindent\textbf{Kantorovich-radius interpretation.}
The perturbation criterion is used only through the fixed compact-tube data
\((M,L,r)\) and the bridged residual value
\(R_h=F_{A,h}(Z_G;y)\).  With
\(\beta_h=M\|R_h\|\) and \(\rho_h=2\beta_h\), the small-step restrictions
above enforce Eq.~\eqref{eq:stage-pert-kantorovich-radius},
\begin{equation}
\label{eq:stage-pert-kantorovich-radius}
  \rho_h\le r,\qquad ML\rho_h\le \frac12 .
\end{equation}
Thus the nonlinear Taylor remainder is absorbed inside the same ball on which
the inverse and second-derivative constants are valid, and the local root
selected by the Gauss predictor satisfies
\(\|Z_A-Z_G\|\le 2M\|R_h\|\).  Lemma~\ref{lem:full-132-row-residual-bridge}
supplies only the scale \(\|R_h\|\le C_Rh^7\); it does not select \(M,L,r\)
or the branch.  Those constants and the proof norm are fixed before the
direct residual bridge, solver logs, or numerical trajectory are read.

\begin{lemma}[Full 132-row residual-defect certificate bridge]
\label{lem:full-132-row-residual-bridge}
Use the accepted row ordering, finite row-scaling map, and fixed
finite-dimensional proof norm for the implemented \(132\)-row stage residual.
If the 96-row non-dynamic certificate gives
\(\|F_{\rm nd,h}(Z_G;y)\|\le C_{\rm nd}h^7\) on the lifted smooth Gauss
stage and the implemented Newton--Euler rows satisfy the direct
same-branch row identity in Eq.~\eqref{eq:full-bridge-dynamic-zero},
\begin{equation}
\label{eq:full-bridge-dynamic-zero}
  F_{\rm dyn,h}(Z_G;y)
  =\bigl(r^{\rm NE}_{j,h}(Z_G;y)\bigr)_{j=1}^{36}
  =0\in\mathbb R^{36},
\end{equation}
then the full
implemented residual satisfies Eq.~\eqref{eq:full-bridge-full-residual},
\begin{equation}
\label{eq:full-bridge-full-residual}
  \|F_{A,h}(Z_G;y)\|\le C_Rh^7
\end{equation}
with \(C_R\) independent of \(y\in K\), \(h\), the reported grid length,
backend, and finite diagnostic tolerances. The \(Z_G\) appearing in the
96-row certificate, the 36-row Newton--Euler identity, and the full residual
evaluation is one fixed lifted Gauss stage with the same stage abscissae,
transition input, chart, row ordering, row scaling, proof norm, and endpoint
output convention.  The bridge is not re-centered at \(Z_A\), at
\(\widetilde Z_A\), or at an endpoint-refined state. The row-scaling map is fixed before the
asymptotic argument and has \(h\)-independent operator norm on this finite
row space, so it can change only the constant \(C_R\), not the \(h^7\)
power. Any comparison with an unscaled residual norm uses an
\(h\)-independent finite-dimensional norm-equivalence constant; no hidden
\(h\)-dependent row weights are introduced beyond the explicit \(h\)-factors
in the residual formulas. Thus, at the lifted Gauss stage \(Z_G\) for this
fixed row-scaled residual, the certified 132-row residual-value bridge supplies the
stage-residual hypothesis of Lemma~\ref{lem:stage-residual-defect} for the
accepted direct-route formula residual. The bridge discharges only this
\(C_Rh^7\) residual-defect input. At this bridge interface, P3 is used only
to identify the implemented residual/Jacobian formula path and fixed row
scaling, P4 is used only through the proved 96-row non-dynamic row-local
certificate, and P5 supplies the 36 dynamic rows; none of P3/P4 is promoted
into a standalone full-132-row theorem-interface discharge or a primitive
symbolic-oracle closure. The missing primitive dynamic symbolic oracle is not
a missing input to this bridge. It does not discharge the uniform stage inverse,
branch-ball, endpoint/right-inverse, stability/tube-retention, or solver-scale
hypotheses, which remain the separate P1/P2/P6 theorem interfaces consumed by
the surrounding perturbation argument. The D5 direct-substitution certificate
is the row-by-row provenance record for the displayed 36-row identity; the
proof-facing input to this lemma is the identity itself on the lifted Gauss
stage.
\end{lemma}

\begin{proof}
After the accepted row ordering is fixed, write the block vector as
Eq.~\eqref{eq:full-bridge-block-vector}:
\begin{equation}
\label{eq:full-bridge-block-vector}
  F_{A,h}(Z_G;y)=
  \bigl(F_{\rm nd,h}(Z_G;y),F_{\rm dyn,h}(Z_G;y)\bigr),
\end{equation}
where \(F_{\rm nd,h}\in\mathbb R^{96}\) collects the kinematic, collocation,
and lower-pair FullVA rows and \(F_{\rm dyn,h}\in\mathbb R^{36}\) collects the
Newton--Euler rows.  The finite row scaling is already included in these two
blocks; because it is a fixed linear map on a fixed finite-dimensional row
space, its operator norm and inverse-equivalence constants are independent of
\(h\).  The non-dynamic certificate gives
\(\|F_{\rm nd,h}(Z_G;y)\|\le C_{\rm nd}h^7\) in the fixed proof norm after
the finite row-scaling map, while the D5 direct-substitution certificate gives
Eq.~\eqref{eq:full-bridge-proof-dynamic-zero},
\begin{equation}
\label{eq:full-bridge-proof-dynamic-zero}
  F_{\rm dyn,h}(Z_G;y)
  =\bigl(r^{\rm NE}_{j,h}(Z_G;y)\bigr)_{j=1}^{36}=0 .
\end{equation}
For nonautonomous rows this block vector is evaluated on the same
time-indexed residual slice \(F_{A,h,t_n}\).  The 96-row block, the 36-row
block, the Jacobian used by the perturbation criterion, and the endpoint map
therefore share the same transition time and stage abscissae; row families may
not be evaluated at different transition times and then assembled into one
residual bridge.
Finite-dimensional norm equivalence for the
fixed 132-row block then gives Eq.~\eqref{eq:full-bridge-block-bound},
\begin{equation}
\label{eq:full-bridge-block-bound}
\begin{aligned}
  \|F_{A,h}(Z_G;y)\|
  &\le C_{\rm blk}\bigl(\|F_{\rm nd,h}(Z_G;y)\|
       +\|F_{\rm dyn,h}(Z_G;y)\|\bigr) \\
  &\le C_Rh^7,\qquad C_R=C_{\rm blk}C_{\rm nd}.
\end{aligned}
\end{equation}
Thus the \(C_Rh^7\) hypothesis in Lemma~\ref{lem:stage-residual-defect} is a
finite block-norm consequence of the 96-row \(O(h^7)\) certificate plus the
36-row zero dynamic certificate, not an empirical residual fit.
The implication is deliberately one-sided.  It uses the fixed block projection
from the implemented \(132\)-row vector to its \(96+36\) row families and the
finite-dimensional norm equivalence after row scaling; it does not use
cancellation between row families, a small assembled residual to infer smaller
primitive subterms, or a dynamic-row identity to improve the non-dynamic row
constant.  In particular, the zero Newton--Euler block is a base-point
identity in the residual vector, not a hidden estimate of the inverse,
second-derivative, endpoint, or solver-tolerance constants used later.
Direct-route sufficiency boundary: Lemma~\ref{lem:stage-residual-defect}
requires a residual value at the lifted Gauss stage, not a primitive symbolic
expansion of every dynamic Taylor term. Once P3 fixes the implemented residual
path and row scaling, P4 supplies the 96 non-dynamic rows, and P5 supplies the
D5 direct-substitution zero rows, the missing primitive dynamic symbolic
oracle is not a missing input to this bridge. The incomplete primitive
symbolic route remains a separate primitive-route record and is not used to prove or
disprove the displayed \(C_Rh^7\) defect.
This discharge is a pre-perturbation statement on \(Z_G\), not an inference
from the accepted numerical trajectory or from small residuals after Newton
termination; the latter enter only through Lemma~\ref{lem:inexact-newton} and
the retained P6 solver-scale condition. Consequently the bridge is not a
compact-tube inverse proof: it does not choose the Newton branch, prove the
\(J_h^{-1}\) bound in Lemma~\ref{lem:stage-residual-defect}, or close the
endpoint, stability, tube-retention, and \(\eta_h^{\rm tube}\) hypotheses.
Those interfaces remain explicit conditions around the residual-defect
handoff.
\end{proof}

\begin{corollary}[Route-separation invariant for the direct residual bridge]
\label{cor:route-separation-invariant}
For the theorem invocation, define the active residual predicate in
Eq.~\eqref{eq:direct-route-predicate},
\begin{equation}
\label{eq:direct-route-predicate}
  \mathsf R_{\rm dir}(h,y):
  \quad
  \|F_{\rm nd,h}(Z_G;y)\|\le C_{\rm nd}h^7
  \ \hbox{and}\ 
  F_{\rm dyn,h}(Z_G;y)=0
\end{equation}
on the fixed branch, chart, row ordering, row scaling, and proof norm.  The
full-residual bridge proves
\(\mathsf R_{\rm dir}(h,y)\Rightarrow
\|F_{A,h}(Z_G;y)\|\le C_Rh^7\), which is the only residual-defect predicate
consumed by Lemma~\ref{lem:stage-residual-defect} and
Theorem~\ref{thm:g6fullva-order}.  The predicate is single-object: the
non-dynamic rows, dynamic rows, block norm, and endpoint map all refer to the
same \(Z_G\); no row family may be evaluated on a different stage vector and
then spliced into this implication.  A stricter primitive predicate
\(\mathsf R_{\rm T3}^{132}(h,y)\) could replace
\(\mathsf R_{\rm dir}(h,y)\) only after it is proved on the same tuple
\((h,y,Z_G,F_{A,h})\).  Its primitive part would first have to prove the
five state-, acceleration-, multiplier-, geometry-, and gyroscopic-lift
obligations and thereby obtain the 36-row dynamic bound; that dynamic bound
would then have to be assembled with the same 96-row non-dynamic certificate,
row ordering, row scaling, block norm, branch, and endpoint convention to
yield
\begin{equation}
\label{eq:t3-replacement-predicate-bound}
  \mathsf R_{\rm T3}^{132}(h,y)\Longrightarrow
  \|F_{A,h}(Z_G;y)\|\le C_R^{\mathrm{T3}}h^7 .
\end{equation}
No such \(132\)-row T3 predicate, in the sense of
Eq.~\eqref{eq:t3-replacement-predicate-bound}, is a premise in this
manuscript's theorem.
Therefore the present open status of the primitive/Taylor route neither
supplies nor removes a theorem hypothesis: the theorem depends on
\(\mathsf R_{\rm dir}\) plus the retained compact-branch, endpoint,
stability, and solver-scale interfaces, and it does not infer source-policy,
full-\tfe{}, or residual-to-error promotion from either route.
\end{corollary}

\begin{proof}
The statement is a direct consequence of
Lemma~\ref{lem:full-132-row-residual-bridge} and
Lemma~\ref{lem:stage-residual-defect}.  The first lemma supplies the
residual value consumed by the second lemma on the fixed implemented residual
object \(F_{A,h}\), with the row families evaluated at the same lifted Gauss
stage and in the same proof norm.  A primitive proof of the 36 dynamic rows
would still become theorem-facing only after it is assembled with the same
96-row non-dynamic block into a new \(132\)-row residual-value predicate on
that object.  Since neither lemma maps a full residual estimate back to
primitive subterms, reported residual histories, or source-policy outputs,
the two route predicates cannot be spliced and neither route closes P6, P7,
or external-comparison readiness.
\end{proof}

\begin{lemma}[Inexact Newton tolerance scale]
\label{lem:inexact-newton}
Let $Z_A$ be the exact accepted stage solution of $F_{A,h}(Z;y)=0$ and let
$\widetilde Z_A$ be the computed Newton iterate. Put
$J_A=D_ZF_{A,h}(Z_A;y)$. Assume that, on a local ball
$B_{r_A}(Z_A)$ contained in the accepted Newton branch,
Eq.~\eqref{eq:inexact-newton-local-hypotheses} holds:
\begin{equation}
\label{eq:inexact-newton-local-hypotheses}
\begin{aligned}
  \|J_A^{-1}\|&\le M_A,\\
  \|D_ZF_{A,h}(Z_1;y)-D_ZF_{A,h}(Z_2;y)\|
    &\le L_A\|Z_1-Z_2\|,\qquad Z_1,Z_2\in B_{r_A}(Z_A),\\
  M_A L_A r_A&\le \frac12 ,
\end{aligned}
\end{equation}
with constants independent of $h$. Assume also that
$\widetilde Z_A\in B_{r_A}(Z_A)$ and
that the proof-norm residual input satisfies
Eq.~\eqref{eq:inexact-newton-residual-input},
\begin{equation}
\label{eq:inexact-newton-residual-input}
  \|F_{A,h}(\widetilde Z_A;y)\| \le \eta_h ,
\end{equation}
and let
\(\mathcal P_h=\mathcal C_h\circ\mathcal E_h\) denote the full endpoint-output
map used after the stage solve, namely Lie endpoint reconstruction followed
by the accepted local endpoint-closure map. Assume that this full
endpoint-output map satisfies
\(\|D_Z\mathcal P_h\|\le M_N\) on the same ball. Then
Eq.~\eqref{eq:inexact-newton-stage-output-bounds} gives the stage and output
perturbation bounds:
\begin{equation}
\label{eq:inexact-newton-stage-output-bounds}
  \|\widetilde Z_A-Z_A\|\le 2M_A\eta_h,\qquad
  \|\mathcal P_h(\widetilde Z_A)-\mathcal P_h(Z_A)\|
  \le C_N\eta_h,\qquad C_N=2M_AM_N .
\end{equation}
Consequently, if the accepted-branch nonlinear solver hypothesis satisfies
the asymptotic scale in Eq.~\eqref{eq:inexact-newton-solver-scale},
\begin{equation}
\label{eq:inexact-newton-solver-scale}
  \eta_h\le c_\eta h^7 ,
\end{equation}
then the endpoint contribution from inexact Newton is bounded explicitly by
Eq.~\eqref{eq:inexact-newton-final-output-bound},
\begin{equation}
\label{eq:inexact-newton-final-output-bound}
  \|\mathcal P_h(\widetilde Z_A)-\mathcal P_h(Z_A)\|
  \le C_Nc_\eta h^7 .
\end{equation}
Here \(M_N\) includes the local Lipschitz factor of the endpoint-closure map;
the inexact-Newton term is therefore measured between two endpoint-closed
maps, not between an unclosed reconstructed endpoint and a closed theorem
output.
A fixed absolute residual tolerance is a finite-run engineering policy and
must be tightened with \(h\) to support the asymptotic theorem. This is again
a local Newton-branch estimate, not a certificate of global nonlinear solver
convergence.
Conversely, for a transition already known to remain in the accepted Newton
branch ball, the proof-norm stopping rule
\(\|F_{A,h}(\widetilde Z_A;y)\|\le c_\eta h^7\) verifies only the
residual-size subcondition of the retained P6 hypothesis for an already
branch-certified transition. It still does not prove Newton convergence to
the branch, choose the proof norm, or verify the condition on every reported
step.
The residual inequality and the branch-membership hypothesis are separate
load-bearing inputs: the lemma does not infer
\(\widetilde Z_A\in B_{r_A}(Z_A)\) from a small residual alone. Residual
control outside this ball, or on a different branch, does not control the
stage error without the averaged-Jacobian inverse used below.

\noindent\textbf{Branch-ball membership checkpoint.}
The theorem-level use of this lemma consumes two independent solver facts for
every accepted transition: the iterate is certified to remain in
\(B_{r_A}(Z_A)\), and its residual is bounded in the fixed proof norm by the
compact-tube envelope \(\eta_h^{\rm tube}\). The second fact cannot replace
the first. A finite log entry with a small residual, an iteration cap, or a
successful line-search flag is therefore not a theorem input unless it is
already tied to the same Gauss-predictor branch, the same residual norm, and
the same local ball. If this membership certificate is unavailable, the
transition is outside the branch-selected map used in
Theorem~\ref{thm:g6fullva-order}; it is not counted as a larger local
constant or as a lower-order inexact Newton step.

Equivalently, the accepted proof uses a strong local residual inverse: the
averaged Jacobian along the segment from $Z_A$ to $\widetilde Z_A$ is required
to be uniformly invertible by the small-neighborhood perturbation bound.
In Theorem~\ref{thm:g6fullva-order}, the lemma-level residual budget
\(\eta_h\) is applied through the compact-tube upper envelope
\(\eta_h^{\rm tube}\). The single reported-grid maximum is only the
trajectory specialization of that envelope, not a replacement for the
uniform compact-tube solver hypothesis.
For this lemma, the residual norm is the proof residual norm of the accepted
132-row residual after the local chart, row ordering, and finite residual
row-scaling convention, including the D6 unweighted scaling convention where
applicable, are fixed. Logged implementation residuals specialize
\(\eta_h\) only through the accepted implementation-path binding and an
\(h\)-independent norm equivalence on the compact tube. Without that binding
they remain finite diagnostics; they do not introduce a second residual
budget, change the proof norm, or prove P6 as a solver theorem.
The scalar \(c_\eta\), residual norm, accepted branch, and constants
\(M_A,M_N\) are fixed proof data before any Newton log, tolerance sweep, or
reported residual ratio is inspected. A run that satisfies a fixed absolute
tolerance but violates the compact-tube \(h^7\) residual scale is simply
outside this lemma's asymptotic use; post-hoc residual filtering cannot create
an accepted transition for the theorem.
Likewise, a Newton correction-norm test, relative residual test, iteration
cap, line-search success flag, or backend convergence flag is not a substitute
for the displayed proof-norm residual inequality. Such a proxy may be used
only if, before the run and on the same accepted branch, it is accompanied by
a uniform conversion that implies
\(\|F_{A,h}(\widetilde Z_A;y)\|\le c_\eta h^7\) in the fixed 132-row proof
norm. Otherwise it is an implementation diagnostic and does not contribute to
the \(C_N\eta_h\) perturbation bound.
\end{lemma}

\begin{proof}
Set $\delta=\widetilde Z_A-Z_A$. Since $F_{A,h}(Z_A;y)=0$, the mean-value form
of the residual gives Eq.~\eqref{eq:inexact-newton-averaged-jacobian},
\begin{equation}
\label{eq:inexact-newton-averaged-jacobian}
  F_{A,h}(\widetilde Z_A;y)
   =\bar J_h\delta,\qquad
  \bar J_h=\int_0^1D_ZF_{A,h}(Z_A+\theta\delta;y)\,d\theta .
\end{equation}
Write $\bar J_h=J_A+E_h$. The Lipschitz bound and
$\widetilde Z_A\in B_{r_A}(Z_A)$ give
Eq.~\eqref{eq:inexact-newton-jacobian-error},
\begin{equation}
\label{eq:inexact-newton-jacobian-error}
  \|E_h\|
  \le \int_0^1 L_A\theta\|\delta\|\,d\theta
  \le L_A r_A .
\end{equation}
Hence $\|J_A^{-1}E_h\|\le 1/2$. By the Neumann perturbation theorem,
$\bar J_h$ is invertible and $\|\bar J_h^{-1}\|\le 2M_A$. Therefore
Eq.~\eqref{eq:inexact-newton-stage-error-proof} follows:
\begin{equation}
\label{eq:inexact-newton-stage-error-proof}
  \|\widetilde Z_A-Z_A\|
   =\|\delta\|
   \le 2M_A\|F_{A,h}(\widetilde Z_A;y)\|
   \le 2M_A\eta_h .
\end{equation}
The full endpoint-output map has \(\|D_Z\mathcal P_h\|\le M_N\) on the same
neighborhood, hence Eq.~\eqref{eq:inexact-newton-output-error-proof},
\begin{equation}
\label{eq:inexact-newton-output-error-proof}
  \|\mathcal P_h(\widetilde Z_A)-\mathcal P_h(Z_A)\|
  \le M_N\|\widetilde Z_A-Z_A\|
  \le 2M_AM_N\eta_h=C_N\eta_h .
\end{equation}
If \(\eta_h\le c_\eta h^7\), the same estimate gives
\(C_N\eta_h\le C_Nc_\eta h^7\), keeping the nonlinear solve error below the
local truncation scale of a sixth-order method with an explicit constant.
The averaged-Jacobian inverse is the transfer operator from residual norm to
stage-error norm on this accepted local branch. Pointwise Jacobian
invertibility alone, or a residual bound away from this ball, is not used as
a shortcut.
\end{proof}

\begin{lemma}[Reported solver logs are not a P6 proof]
\label{lem:p6-reported-log-nonclosure}
A finite reported-grid solver log, even one satisfying
Eq.~\eqref{eq:p6-reported-log-bound},
\begin{equation}
\label{eq:p6-reported-log-bound}
  \eta_h^{\rm reported}
  :=\max_{0\le n<N_h}\eta_{h,n}
  =\widehat\eta_h^{\rm log}\le c_\eta h^7
\end{equation}
on the recorded accepted trajectory, is not a compact-tube solver envelope
and is not a proof of the P6 solver-policy theorem.  It can specialize an
already fixed P6 theorem instance only after the compact tube, proof residual
norm, row scaling, Gauss-predictor branch, branch-ball membership, and
constant \(c_\eta\) have all been fixed independently.  It cannot calibrate
\(c_\eta\), prove branch membership, prove the averaged-Jacobian inverse,
choose the accepted nonlinear root, or turn fixed production tolerances into
an asymptotic solver policy.
\end{lemma}

\begin{proof}
Lemma~\ref{lem:inexact-newton} needs two separate inputs for each accepted
transition: membership in the local Newton branch ball and the proof-norm
residual inequality.  The reported quantity
\(\eta_h^{\rm reported}\), denoted \(\widehat\eta_h^{\rm log}\) in
diagnostic tables, records only a maximum over a finite set of already
computed trajectory transitions.  It is therefore a specialization of the
residual-size clause after the theorem instance has selected the branch,
norm, row convention, and constant.  It is not a quantifier over the compact
tube and it does not supply the local inverse that converts residual size to
stage error.

Consequently the logical direction is one-way.  A predeclared same-branch
stopping rule can instantiate the residual-size part of P6 for the
transitions it accepts, while the retained compact-tube branch hypotheses
still supply membership and local invertibility.  Reading the finite log in
the opposite direction would choose \(c_\eta\), the residual norm, or the
branch from the evidence that the theorem is meant to bound.  That reverse
reading is disallowed here.  The theorem-level condition remains
\(\eta_h^{\rm tube}\le c_\eta h^7\); the finite reported-grid log remains
diagnostic-grade solver information, not a P6 proof.
\end{proof}

\noindent\textbf{Proof-level solver-policy instantiation.}
The theorem can be invoked with a predeclared proof-level Newton acceptance
rule: before the run, fix the accepted branch, proof residual norm,
row scaling, \(c_\eta\), and the compact-tube constants, and accept a
transition for the theorem only when the branch-membership hypotheses of
Lemma~\ref{lem:inexact-newton} are retained and the proof-norm residual
satisfies \(\|F_{A,h}(\widetilde Z_A;y)\|\le c_\eta h^7\).  Under those
pre-fixed conditions the residual-size part of P6 is instantiated by the
stopping rule itself for the accepted transition set.  This does not prove
branch membership, averaged-Jacobian invertibility, or all-transition
termination; those remain the retained compact-branch solver hypotheses.  It
also does not promote fixed production tolerances, backend convergence flags,
post-hoc residual filtering, or source-policy runs into theorem-level
solver-policy proof.

\begin{lemma}[Residual-size subcondition under a predeclared P6 instance]
\label{lem:p6-residual-size-subcondition}
Fix a theorem instance before the run: the accepted Gauss-predictor branch,
proof residual norm, row scaling, compact-tube constants, and scalar
\(c_\eta\) are fixed, and the branch-ball and strong-local-inverse
hypotheses of Lemma~\ref{lem:inexact-newton} are retained.  Suppose that for
each accepted transition in the reported theorem domain the Newton iterate is
accepted only after the predeclared proof-norm residual test in
Eq.~\eqref{eq:p6-proof-norm-residual-test},
\begin{equation}
\label{eq:p6-proof-norm-residual-test}
  \|F_{A,h}(\widetilde Z_A;y)\|_{\rm proof}\le c_\eta h^7
\end{equation}
is satisfied.  Then the
residual-size part of the retained P6 clause is
instantiated for those accepted transitions, and
Lemma~\ref{lem:inexact-newton} contributes at most \(C_Nc_\eta h^7\) to the
local defect.  This verifies only the displayed proof-norm residual
inequality on the accepted transition set; it does not promote P6 from a
retained compact-tube solver hypothesis to a theorem-level solver-policy
certificate unless branch-ball membership, strong local inverse, and the
uniform all-transition envelope are proved separately.  This lemma does not
prove branch-ball membership,
strong-local-inverse persistence, all-transition solver termination, global
Newton convergence, remote-root exclusion, or any fixed-tolerance
asymptotic theorem.
\end{lemma}

\begin{proof}
The stopping rule is stated in the same proof residual norm and with the same
constant \(c_\eta\) used by the theorem instance.  Therefore, on every
accepted transition for which the retained branch and local-inverse
hypotheses hold, the inexact-Newton input satisfies
\(\eta_h\le c_\eta h^7\).  Substituting this inequality into
Lemma~\ref{lem:inexact-newton} gives Eq.~\eqref{eq:p6-inexact-newton-residual-size-bound},
\begin{equation}
\label{eq:p6-inexact-newton-residual-size-bound}
  \|\mathcal P_h(\widetilde Z_A)-\mathcal P_h(Z_A)\|
  \le C_N\eta_h\le C_Nc_\eta h^7 .
\end{equation}
The conclusion is only a residual-size instantiation on the preselected
branch.  It uses the branch-membership and local-inverse assumptions rather
than proving them, and it makes no statement about a backend fixed tolerance
or about solver behavior outside the accepted theorem domain.
\end{proof}

\begin{lemma}[Local-to-global transfer on the accepted tube]
\label{lem:local-global-transfer}
Let $\Psi_h$ be a one-step map on the smooth reduced chart from
Assumption~\ref{ass:regularity}, let \(K\) be the compact proof tube, and
suppose the exact trajectory on \(0\le t\le T\) has distance at least
\(d_K>0\) from \(\partial K\).
\noindent\textbf{Time-augmentation convention for the transfer.}
If the reduced FullVA flow is nonautonomous because drives, loads, or
constraint right-hand sides depend on time, this lemma is applied to the
augmented chart in Eq.~\eqref{eq:local-global-time-augmentation}:
\begin{equation}
\label{eq:local-global-time-augmentation}
  \widehat y=(y,\tau),\qquad \dot\tau=1 .
\end{equation}
Then \(\varphi_h\) denotes the autonomous augmented flow, \(\Psi_h\) advances
\(\tau\mapsto\tau+h\) exactly, and the reported theorem norm is obtained by
projecting back to the physical reduced variables; the time component
contributes zero grid error because both numerical and exact times are
\(t_n\) on the reported grid.  Equivalently, in unaugmented notation
\(\varphi_h(y)\) abbreviates the two-parameter flow
\(\varphi_{t+h,t}(y)\) at the fixed reported input time \(t=t_n\).  This
time convention is fixed before the local-defect and Gronwall estimates, so
no time-phase mismatch, hidden interpolation, or terminal-time correction is
inserted into the transfer.
Suppose that, for all \(y,\bar y\in K\), the local defect input is
Eq.~\eqref{eq:local-global-lemma-local-defect-input}:
\begin{equation}
\label{eq:local-global-lemma-local-defect-input}
  \|\Psi_h(y)-\varphi_h(y)\|\le C_\ell h^{p+1}
\end{equation}
and the stability input is Eq.~\eqref{eq:local-global-lemma-stability-input}:
\begin{equation}
\label{eq:local-global-lemma-stability-input}
  \|\Psi_h(y)-\Psi_h(\bar y)\|\le (1+C_s h)\|y-\bar y\|
\end{equation}
for $0<h\le h_0$, where \(C_s\ge0\).
For the accepted branch-selected map in Theorem~\ref{thm:g6fullva-order},
the same lemma is used in the restricted-domain form in
Eq.~\eqref{eq:local-global-restricted-domain}: if
\(\Psi_h\) is defined only on the accepted theorem domain
\(\mathcal D_h\subset K\), it is enough that, before the first domain exit,
the exact comparison state \(Y_n=\varphi_{t_n,0}(y_0)\) and the numerical
state \(y_n\) satisfy
\begin{equation}
\label{eq:local-global-restricted-domain}
\begin{gathered}
  y_n,Y_n\in \mathcal D_h\cap K,\\
  \|\Psi_h(Y_n)-\varphi_{t_{n+1},t_n}(Y_n)\|
  \le C_\ell h^{p+1},\\
  \|\Psi_h(y_n)-\Psi_h(Y_n)\|
  \le (1+C_s h)\|y_n-Y_n\|
\end{gathered}
\end{equation}
on the reported transition indices \(0\le n<N\).  This restricted form is
the version consumed by the theorem; it does not require an arbitrary
extension of the accepted map to all of \(K\).  The first-exit argument below
closes the compact-tube part of this condition, while branch-ball,
endpoint-ball, proof-norm, and P6 solver-envelope membership remain the
predeclared admissibility conditions defining \(\mathcal D_h\).
\noindent\textbf{Synchronized-time stability slice.}
In the augmented nonautonomous interpretation, this stability bound is
consumed only on the same reported input-time slice: the two arguments are
\((y,\tau_n)\) and \((\bar y,\tau_n)\), with \(\tau_n=t_n\). The accepted
map advances both time components exactly to \(\tau_{n+1}=t_{n+1}\), so the
stability estimate controls only the physical reduced-state difference and
does not hide a time-registration error. The proof never compares states with
different time phases, never fits \(C_s\) from time-shifted trajectory data,
and never absorbs interpolation or terminal-time mismatch into the stability
constant. Equivalently, in unaugmented notation the retained P2 stability
input is the one-step Lipschitz bound of the time-indexed map
\(\Psi_{h,t_n}\) on the physical state at fixed input time \(t_n\).
If a sharper negative contraction
coefficient is available, enlarge it to a nonnegative upper bound before
applying this lemma. Let \(t_n=n h\), and let \(N\) be any reported grid
length with \(0\le t_n\le T\) for \(0\le n\le N\); in particular,
\(Nh=T\) may be used when the final time is exactly divisible by \(h\).
Let the numerical trajectory be the accepted recurrence in
Eq.~\eqref{eq:local-global-accepted-recurrence}:
\begin{equation}
\label{eq:local-global-accepted-recurrence}
  y_{n+1}=\Psi_h(y_n),\qquad 0\le n<N,
\end{equation}
started from the same reduced-chart initial state \(y_0\) used by the exact
flow \(\varphi_{t_n}(y_0)\).
Local-to-global initial-state boundary: the lemma uses zero
initial error in the reduced chart. Any unreported initial projection,
warm-start, branch-selection, or chart-registration error would add a separate
initial-error term to the discrete Gronwall recursion and is not absorbed into
the \(C_\ell h^{p+1}\) one-step defect or the theorem's \(h^6\)
same-initial-state reported-grid estimate.
Define the Gronwall factor in
Eq.~\eqref{eq:local-global-lemma-gronwall-factor}:
\begin{equation}
\label{eq:local-global-lemma-gronwall-factor}
  \Gamma_s(T)=
  \begin{cases}
    \dfrac{e^{C_sT}-1}{C_s}, & C_s>0,\\[0.6em]
    T, & C_s=0 .
  \end{cases}
\end{equation}
Uniform-grid instance rule: the constants \(C_\ell\), \(C_s\), the tube
margin \(d_K\), the horizon \(T\), and the resulting \(\Gamma_s(T)\) are
chosen once for the accepted theorem instance before the reported grid is
considered.  The proof does not use transition-dependent constants
\(C_{\ell,n}\), \(C_{s,n}\), or a stepwise tube margin, and it does not
replace the recurrence after inspecting a local residual, endpoint diagnostic,
or observed order row.  If the accepted map is defined only on the theorem
domain \(\mathcal D_h\), this lemma is applied only to the preselected
accepted transition set described after the proof; Gronwall controls compact
tube retention but does not create new branch-ball, endpoint-ball, or P6
solver-envelope admissibility.
Then for \(h\) sufficiently small, the numerical trajectory remains in \(K\)
and satisfies the grid estimate in
Eq.~\eqref{eq:local-global-lemma-grid-bound}:
\begin{equation}
\label{eq:local-global-lemma-grid-bound}
  \max_{0\le n\le N}\|y_n-\varphi_{t_n}(y_0)\|
  \le C_T h^p,\qquad C_T=C_\ell\Gamma_s(T),
\end{equation}
with \(C_T\) independent of \(h\). This is an explicit grid-point estimate on
the accepted reduced-chart branch and a compact-tube transfer lemma on the
reported time grid \(t_n=n h\); exact final-time divisibility is not required.
The local one-step hypotheses are used only for transitions
\(0\le n<N\); the terminal index \(N\) is an output index, not an extra
solve or an additional local-map hypothesis.
If \(Nh<T\), the leftover interval \(T-Nh\) is outside this reported-grid
estimate; the lemma does not assert an error bound at an unreported terminal
time unless that time is added to the reported grid by a separate step.
It is not an invariant-region proof away from the accepted branch.
\end{lemma}

\begin{proof}
Let \(n_\ast\) be the first index at which the numerical trajectory would
leave \(K\), with \(n_\ast=N+1\) if no such index exists. For
\(n<n_\ast\), the estimates hold in \(K\). Let
\(\varepsilon_n=y_n-\varphi_{t_n}(y_0)\) and
\(e_n=\|\varepsilon_n\|\).  With the time-augmentation convention above,
the semigroup property used here is the autonomous augmented semigroup
identity; in unaugmented notation it is the identity
\(\varphi_{t_{n+1},0}(y_0)
=\varphi_{t_{n+1},t_n}(\varphi_{t_n,0}(y_0))\) at the reported grid time.
Adding and subtracting
$\Psi_h(\varphi_{t_n}(y_0))$ gives
Eq.~\eqref{eq:local-global-error-recursion-proof}:
\begin{equation}
\label{eq:local-global-error-recursion-proof}
  e_{n+1}
  \le (1+C_s h)e_n+C_\ell h^{p+1}.
\end{equation}
\noindent\textbf{Truncated first-exit recursion checkpoint.}
At this point the recursion is not yet iterated over the whole reported grid.
Set the truncated index in Eq.~\eqref{eq:local-global-truncated-index}:
\begin{equation}
\label{eq:local-global-truncated-index}
  m_\ast=\min\{N,n_\ast\}.
\end{equation}
For each output index \(0\le n\le m_\ast\), every transition used in the
Gronwall sum has index \(0\le j<n\le n_\ast\), hence \(j<n_\ast\).  The local
defect and stability hypotheses are therefore evaluated only before the
hypothetical first exit, and the proof does not use the local estimates after
a hypothetical exit.  Iterating the scalar recursion on this truncated index
set and using the discrete Gronwall argument in inequality form, with
\(n h\le T\), yields Eq.~\eqref{eq:local-global-gronwall-proof} for every
\(0\le n\le m_\ast\):
\begin{equation}
\label{eq:local-global-gronwall-proof}
  e_n
  \le C_\ell h^{p+1}\sum_{j=0}^{n-1}(1+C_s h)^j
  \le C_\ell h^p\,\Gamma_s(T),
\end{equation}
where the definition of \(\Gamma_s(T)\) supplies the limiting case
\(C_s=0\). Let \(C_T=C_\ell\Gamma_s(T)\). Choose \(h\) small enough that
\(C_T h^p<d_K/2\).
First-exit bootstrap domain: the local estimates are used only up to the
first possible exit index \(n_\ast\), and \(h_0\) is then shrunk after
the Gronwall constant \(C_T\) and the tube margin \(d_K\) are fixed.
Thus the proof does not assume global tube retention as the conclusion;
it proves retention by contradicting the first-exit alternative. The
first-exit index is a proof device, not a trajectory failure log. The
stability constant \(C_s\), tube margin \(d_K\), Gronwall factor
\(\Gamma_s(T)\), and \(C_T\) are fixed before any reported grid errors,
work-precision slopes, or exit diagnostics are inspected; shrinking \(h_0\)
after those choices cannot calibrate \(C_s\) or the tube margin from a
finite run.
Then every numerical point covered by the truncated induction stays within
distance \(d_K/2\) of the exact trajectory.  If \(n_\ast\le N\), the estimate
at \(n_\ast\) implies that \(y_{n_\ast}\) is still inside \(K\), contradicting
the definition of the first exit.  Hence \(n_\ast=N+1\), and only after this
contradiction is the estimate extended to every reported index
\(0\le n\le N\). This closes the bootstrap and proves the same-initial-state
reported-grid \(O(h^p)\) estimate on the accepted tube.
\end{proof}

\noindent\textbf{Accepted-domain specialization of the transfer.}
For Theorem~\ref{thm:g6fullva-order}, Lemma~\ref{lem:local-global-transfer}
is not applied to an arbitrary extension of the accepted map from
\(\mathcal D_h\) to all of \(K\).  The recurrence is read only on the
preselected accepted transition set of
Definition~\ref{def:accepted-branch-map}.  The first-exit proof closes the
compact-tube component of domain membership: once the local defect and
stability estimates have been fixed, the Gronwall bound keeps the numerical
and exact comparison states in the retained compact tube for \(0<h\le h_0\).
The other domain components--Gauss-predictor branch selection, branch-ball and
endpoint-ball membership, the fixed proof norm, and the P6 scaled
solver-envelope condition--are not consequences of Gronwall.  They remain the
same admissibility conditions on the same reported transition set.
Thus the transfer never propagates through an undefined accepted step and
never converts a solver-policy or branch-domain failure into a covered
lower-order case.

\begin{table}[htbp]
\centering
\caption{Local-to-global transfer table for
Lemma~\ref{lem:local-global-transfer} and
Theorem~\ref{thm:g6fullva-order}.}
\label{tab:local-global-transfer-ledger}
\scriptsize
\setlength{\tabcolsep}{2pt}
\renewcommand{\arraystretch}{0.9}
\begin{tabular}{P{0.18\linewidth}P{0.30\linewidth}P{0.27\linewidth}P{0.17\linewidth}}
\toprule
Transfer layer & Input consumed & Output produced & Explicit non-use \\
\midrule
Local defect input &
Uniform branch-selected local defect \(C_{\rm loc}h^7\) on the accepted
compact proof tube. &
The lemma input \(C_\ell h^{p+1}\) with \(p=6\). &
Residual/reaction tables, finite \(h\)-sweep fits, fixed production
tolerances. \\
Stability and tube bootstrap &
The compact-tube stability scale \(1+C_s h\) and the tube-retention margin
\(d_K\). &
The accepted numerical trajectory remains in \(K\) after shrinking \(h_0\). &
Invariant-region claims away from the accepted branch or remote chart
claims. \\
Discrete Gronwall recursion &
The accepted recurrence \(y_{n+1}=\Psi_h(y_n)\), started from the same
initial reduced-chart state, on reported transition steps \(0\le n<N_h\). &
\(C_{\rm red}=C_{\rm loc}\Gamma_s(T)\) and the reduced-chart grid-point
estimate. &
Resampled or postprocessed sequences, unreported initial projection errors,
exact final-time divisibility, an extra terminal solve, or an additional
local-map hypothesis at \(T\). \\
Reported-grid output &
The reported output indices \(0\le n\le N_h\), \(N_hh\le T\), followed by
the reporting map in the theorem. &
The reduced-chart \(h^6\) bound and the later position--velocity reporting-map
consequence. &
Unreported leftover interval claims, residual-to-error promotion, or
source-policy/full-\tfe{} readiness. \\
\bottomrule
\end{tabular}
\end{table}

Local-to-global transfer rule: the transfer consumes
only the uniform branch-selected local defect, the compact-tube stability
scale, and the tube-retention margin before the transfer step.
\mbox{discrete Gronwall on reported transition steps} is then applied to the accepted recurrence
\(y_{n+1}=\Psi_h(y_n)\) from the \mbox{same initial reduced-chart state}, not a
resampled or postprocessed sequence, and it consumes zero initial error rather
than hiding an initial projection or branch-registration defect. It produces a
reduced-chart grid-point estimate on reported output indices; the
position--velocity estimate is a later reporting-map consequence. The table
does not use residual/reaction tables, finite \(h\)-sweep fits, fixed
production tolerances, source-policy rows, unreported initial projection
errors, an unreported leftover interval, or an extra terminal solve at \(T\).
It does not prove a theorem-level solver-policy theorem, a P7 residual-to-error
transfer, source-policy readiness, or full-\tfe{} readiness. Current
certificates prove the direct PC2 residual bridge under the retained P6
hypothesis; they do not prove P6 or P7 themselves.

\noindent\textbf{Stability-map binding checkpoint.}
In the invocation of Lemma~\ref{lem:local-global-transfer} inside
Theorem~\ref{thm:g6fullva-order}, the abstract one-step map \(\Psi_h\) is
the endpoint-closed inexact accepted map
\(\Psi_h^{\mathrm{G6FVA}}=\mathcal P_h(\widetilde Z_A)\), restricted to the
same accepted compact branch and proof norm used in the local-defect
estimate.  The maps \(\Psi_h^G\), \(\Psi_h^A\), and \(\Psi_h^E\) are used
only to decompose the one-step local defect before the recurrence is written;
their stability is not substituted into the global recursion.  The stability
constant \(C_s\) is therefore the retained same-branch stability scale for the
accepted map under the P2 compact-tube and P6 solver-scale interfaces, not a
constant borrowed from the Gauss map, exact-stage root, or endpoint-closed
exact-stage map.  After the four local perturbation terms have produced
\(\|\Psi_h^{\mathrm{G6FVA}}(y)-\varphi_h(y)\|\le C_{\rm loc}h^7\), the
Gronwall step uses one recurrence only:
\(y_{n+1}=\Psi_h^{\mathrm{G6FVA}}(y_n)\).  This prevents a hidden
map-switch in the local-to-global proof and keeps endpoint closure and
inexact Newton inside the same branch-selected map whose grid error is
reported.

\begin{lemma}[Reporting-map norm consequence on the accepted chart]
\label{lem:reporting-map-norm-consequence}
Let \(K\) be the retained compact proof tube in the reduced chart, and let
\(\mathcal R:K\to\mathbb R^{m_{qv}}\), \(\mathcal R(y)=(q(y),v(y))\), be the
reported position--velocity map.  Assume \(\mathcal R\) is \(C^1\) on a
neighborhood of \(K\).  Choose once a fixed convex coordinate neighborhood
\(U_K\) whose closure lies in that chart neighborhood, contains \(K\) with
positive margin, and satisfies the derivative bound in
Eq.~\eqref{eq:reporting-map-derivative-bound},
\begin{equation}
\label{eq:reporting-map-derivative-bound}
  \sup_{y\in U_K}\|D\mathcal R(y)\|\le C_{\mathcal R}.
\end{equation}
\noindent\textbf{Convex chart-neighborhood checkpoint.}
The reporting step does not assume that \(K\) itself is convex.  The line
segments used by the integral mean-value formula are taken only in the fixed
Euclidean coordinates of \(U_K\), and the local-to-global tube is shrunk, if
needed, so that the reported numerical and exact reduced-chart points and all
segments between them stay inside \(U_K\).  Thus the reporting-map constant is
a compact-neighborhood derivative bound chosen before the reported grid,
output formatting, or fitted errors are inspected; it is not an atlas switch,
not a remote-chart norm equivalence, and not a fitted norm-equivalence
constant.
Reported-component selection rule: the codomain and component ordering of
\(\mathcal R(y)=(q(y),v(y))\), together with the norm weights used for the
reported position--velocity estimate, are fixed before the theorem instance
and before any numerical table is read.  The map does not include
accelerations, multipliers, reconstructed reactions, constraint residuals, or
stage residual monitors.  Adding any of those quantities to the reported
output would define a different map \(\mathcal R_{\rm ext}\) and would require
its own compact-tube derivative bound, admissible chart neighborhood, and
local-to-global input estimate; it cannot be obtained by appending columns to
the numerical result table after Theorem~\ref{thm:g6fullva-order} has been
invoked.
\noindent\textbf{SO(3) reporting convention.}
When a component of \(q\) or \(v\) contains orientation or angular velocity,
the displayed reporting norm is the norm induced by the same local
retraction/log chart and component weights fixed above.  The estimate is
therefore a local bound on the chart-reported rotational coordinates (or on the
smooth rotational output already included in \(\mathcal R\)), not an
atlas-wide \(SO(3)\) geodesic distance, not a statement through a chart cut or
antipodal parametrization, and not invariant under a post-hoc change of
orientation metric.  Replacing the reported orientation metric, switching
charts, or comparing a remote attitude representation would require a new
smooth output map on the same compact tube and a new derivative constant fixed
before the theorem instance.
If two reported-grid reduced-chart sequences \(y_n\) and \(\bar y_n\) lie in
\(K\) and satisfy the reduced-chart input estimate in
Eq.~\eqref{eq:reporting-map-reduced-input},
\begin{equation}
\label{eq:reporting-map-reduced-input}
  \max_{0\le n\le N_h}\|y_n-\bar y_n\|\le C_{\rm red}h^6,
\end{equation}
then, on the same reported output indices \(0\le n\le N_h\),
Eq.~\eqref{eq:reporting-map-output-bound} follows:
\begin{equation}
\label{eq:reporting-map-output-bound}
  \max_{0\le n\le N_h}
  \|\mathcal R(y_n)-\mathcal R(\bar y_n)\|
  \le C_{qv}h^6,\qquad C_{qv}=C_{\mathcal R}C_{\rm red}.
\end{equation}
This is a local chart-to-output norm consequence on the accepted compact
tube. It is not a residual-to-error transfer theorem, not a global atlas
equivalence, not a remote-branch statement, and not a terminal-time or
dense-output estimate outside the reported grid.
\end{lemma}

\begin{proof}
For each reported index \(n\), the segment
\(\bar y_n+\theta(y_n-\bar y_n)\), \(0\le\theta\le1\), lies in the fixed
convex chart neighborhood \(U_K\) by the checkpoint above. The integral
mean-value formula in this one fixed Euclidean chart gives
Eq.~\eqref{eq:reporting-map-mean-value},
\begin{equation}
\label{eq:reporting-map-mean-value}
  \mathcal R(y_n)-\mathcal R(\bar y_n)
  =
  \int_0^1
  D\mathcal R\bigl(\bar y_n+\theta(y_n-\bar y_n)\bigr)
  (y_n-\bar y_n)\,d\theta .
\end{equation}
Taking norms and the compact-tube derivative supremum yields
Eq.~\eqref{eq:reporting-map-point-bound},
\begin{equation}
\label{eq:reporting-map-point-bound}
  \|\mathcal R(y_n)-\mathcal R(\bar y_n)\|
  \le C_{\mathcal R}\|y_n-\bar y_n\|.
\end{equation}
Maximizing over exactly the same reported indices and substituting the
reduced-chart grid estimate gives the displayed \(C_{qv}h^6\) bound. The
argument uses only the fixed local chart, fixed convex coordinate
neighborhood \(U_K\), compact derivative bound, and reported-grid
reduced-chart estimate; it does not use residual/reaction tables,
source-policy rows, fitted constants, output interpolation, or a chart switch.
\end{proof}

Theorem~\ref{thm:g6fullva-order} is therefore a conditional
local-defect-to-global-error theorem: sixth order follows from the stated FullVA lift,
implemented residual-defect, endpoint-closure, compact-tube stability, and
inexact-Newton contracts. The active direct stage-residual proof route, and
only that route, is closed by the 96-row non-dynamic certificate plus the D5
direct-substitution certificate for the 36 Newton--Euler rows. The
AD-expanded certificate provides implementation-path and derivative-cell
closure only for those implemented direct-route identities. This AD-expanded
certificate is not the primitive symbolic defect certificate; the primitive
dynamic symbolic oracle remains open outside the active direct route. AD-expanded
implementation-path boundary: the closed certificate supplies AD-expanded
implementation-path/derivative-cell closure for the active direct residual
route only; it does not close the primitive dynamic symbolic oracle,
source-paper residual replacement, or source-policy package readiness. What
remains conditional is not the row-local substitution itself, but the
regularity tube, the theorem-level scaled nonlinear-solver hypothesis, and any
separate residual-to-error promotion beyond the accepted smooth branch. The
numerical h-sweeps below support the scale of the claim, and the solver
diagnostics are deliberately finite: a one-step scaled-tolerance probe, a
short-trajectory scaled-tolerance probe, and a finite tolerance-regime sweep
exercise reported-grid residual maxima on smooth windows. These finite
records are the probes referenced by the P6 evidence boundary. Those probes are
finite solver-policy diagnostics only: they are not proof inputs, they are not a
theorem-level tolerance sweep over every reported trajectory, and they are not
a substitute for the explicit \(\eta_h^{\rm tube}\le c_\eta h^7\) theorem
condition.
Current certificates prove only the direct PC2 residual bridge under the retained
P6 interface; they do not discharge P6 as a theorem-level solver-policy result,
and they leave P7 as a separate residual-to-error transfer problem.
Tables~\ref{tab:newton-euler-obligations}--\ref{tab:dynamic-proof-closure-matrix}
record the row-level implementation inputs and the remaining solver
boundary. For validation scope, no residual-to-trajectory-error transfer
theorem is accepted in this manuscript.
Residual tables and residual-to-error ratios are diagnostics only.
The residual-to-error route remains separate: residual rows
for four-link and slider-crank stay mechanism-coverage evidence, not accepted
dynamic order rows, until the seven residual-to-error obligations are
discharged.
Table~\ref{tab:p7-residual-to-error-obligation-ledger} lists those
obligations separately so that the transfer requirements are explicit.

\begin{table}[htbp]
\centering
\caption{Residual-to-error transfer requirements for mechanism residual rows.}
\label{tab:p7-residual-to-error-obligation-ledger}
\scriptsize
\setlength{\tabcolsep}{2pt}
\renewcommand{\arraystretch}{0.9}
\begin{tabular}{P{0.12\linewidth}P{0.30\linewidth}P{0.16\linewidth}P{0.34\linewidth}}
\toprule
Item & Required before residual promotion & Condition state & Manuscript consequence \\
\midrule
P7-1 & Trajectory-level residual identity relating the fixed residual operator
to closed-loop trajectory error on the same reported branch; this is not the
stage identity at \(Z_G\) & Open &
Residual tables do not by themselves define a same-branch trajectory-error
residual identity. \\
P7-2 & Residual-consistency rate \(O(h^7)\) on that branch, with frozen row
scaling & Open &
Observed residual magnitudes are not a uniform residual-consistency theorem in
a fixed proof norm. \\
P7-3 & Closed-loop DAE stability or inf-sup control for that fixed operator &
Open &
No uniform stability or inf-sup inverse is accepted for the residual-only
rows. \\
P7-4 & Calibrated non-floor-limited residual-to-error estimator in the reported
trajectory norm & Open &
No estimator is calibrated away from the reference floor in the promoted
trajectory norm. \\
P7-5 & Reference-floor exclusion & Open &
Floor diagnostics do not promote residual rows to dynamic-order rows. \\
P7-6 & Coarse-first accepted dynamic-order campaign & Open &
Four-link, slider-crank, and source-policy residual rows are not accepted
dynamic-order campaigns. \\
P7-7 & Manuscript transfer theorem/proof & Open &
No theorem in this manuscript transfers those residual bounds to trajectory
error. \\
\bottomrule
\end{tabular}
\end{table}

Residual-to-error boundary: all
seven P7 dependency obligations remain open. Residual tables and
residual-to-error ratios are diagnostics only; no mechanism-coverage residual
row, closed-loop reaction table, common-reference row, or source-policy row is
promoted to trajectory-error/order evidence until dynamic residual identity,
\(O(h^7)\) residual consistency, stability or inf-sup control, calibrated
estimator, reference-floor exclusion, \mbox{coarse-first campaign}, and manuscript
transfer theorem close together.
The closed-loop reaction table remains diagnostic in this P7 sense.
The seven items are conjunctive: P7-1--P7-3 are the mathematical transfer
operator requirements, P7-4--P7-6 are the estimator and campaign requirements,
and P7-7 is the manuscript theorem requirement. Closing only one group would
not promote a residual row.
Until all seven are closed together, residual magnitudes remain diagnostics
and cannot replace the accepted smooth-branch order proof above.
Equivalently, a residual-to-error \mbox{implication would have to prove a bound} of
the form in Eq.~\eqref{eq:p7-residual-transfer-template},
\begin{equation}
\label{eq:p7-residual-transfer-template}
  \|e_h\|_{\rm traj}
  \le C_{\rm P7}\Bigl(
     \|\mathcal R_h^{\rm mech}\|_{\rm res}
     +\rho_{\rm res}(h)
     +\varepsilon_{\rm floor}(h)
     +\varepsilon_{\rm est}(h)
  \Bigr)
\end{equation}
with a uniform stability or inf-sup constant on the reported branch, an
a priori residual-consistency rate \(\rho_{\rm res}(h)\), an estimator
calibrated away from the reference floor, and the same coarse-first
dynamic-order campaign. The present residual rows do not supply that
implication: small residual norms alone are not a trajectory-error theorem,
and residual rows remain diagnostics even when their measured norms are small.
The reason is logical rather than empirical:
an \(O(h^7)\) residual trend, even on the reported branch, is not accepted as
an \(O(h^6)\) trajectory-error statement unless an \(h\)-independent transfer
constant, a same-branch residual identity, a reference-floor exclusion, and
the manuscript transfer proof are all present. Such a transfer theorem would
be a separate P7 boundary theorem; it is not hidden inside the local stage-residual
defect, the Newton--Euler row certificate, or the common-reference tables.
\mbox{P7 fixed-operator requirement}: before any residual table could
be promoted, the proof would have to freeze the reported branch, chart,
residual row scaling, trajectory norm, and same-branch residual transfer
operator. A \mbox{fixed linearized inverse} must be declared before the
residual data are read. In a standard route this would be such an inverse or
inf-sup estimate such as Eq.~\eqref{eq:p7-linearized-infsup-template},
\begin{equation}
\label{eq:p7-linearized-infsup-template}
  \|\delta z\|_{\rm traj}
  \le C_{\rm P7}\,
     \|D\mathcal R_h^{\rm mech}(z_h^\star)\delta z\|_{\rm res}
\end{equation}
with \(C_{\rm P7}\) independent of \(h\), followed by a nonlinear remainder
estimate plus reference-floor and estimator closure on the same branch. The
constant is fixed before observed residual data are read; it is not obtained
by dividing measured trajectory errors by measured residuals. Without this
fixed inverse, residual-size convergence can be a consistency diagnostic but
not an order theorem.
This is the \mbox{fixed transfer operator} requirement for P7.
\noindent\textbf{No empirical ratio calibration.}
Residual-to-error ratios printed for four-link, slider-crank, or source-policy
diagnostics are posterior measurements, not transfer constants.  They use the
same finite reference calculation and reported output norm that a transfer
theorem would have to justify, so they cannot define \(C_{\rm P7}\), choose a
mechanism-specific residual weight, remove a reference floor, or select a
different trajectory norm after the run.  A valid P7 route would need one
predeclared operator/norm pair and one \(h\)-independent stability or inf-sup
constant on the accepted branch before any residual table is read; otherwise
the ratio is only a diagnostic consistency number.

P7 residual-to-error operator boundary: the separate P7 theorem is not a
scalar residual-size check. It would have to define, on the same reported
closed-loop branch, a mechanism residual operator
\(\mathcal R_h^{\rm mech}\), a trajectory norm \(\|e_h\|_{\rm traj}\)
matching the reported dynamic-order rows, and a uniform stability or inf-sup
inverse before reference-floor or estimator terms are added. The present
four-link and slider-crank rows check constraint closure, reaction
consistency, and local residual scale only. They do not identify a same branch
residual operator with a uniform inverse, and they do not prove that small
reaction or constraint residuals control the trajectory error.
P7 fixed-instance rule: any residual-to-error theorem used for promotion must predeclare,
before reading the residual table, the residual operator, residual norm,
reported trajectory components, trajectory norm, branch selector, reference
policy, floor-exclusion rule, estimator form, and stability or inf-sup
constant family. Changing any one of these after the four-link, slider-crank,
or source-policy residual rows are inspected defines a different diagnostic
row, not a promotion of the existing row. In particular, a constraint residual
norm, reaction residual norm, common-reference trajectory error, and
source-policy timing row cannot be combined after the fact into a single
P7 theorem unless a same-branch operator and one \(h\)-independent transfer
constant bind the entire tuple.

\begin{lemma}[Same-codomain obstruction for residual-to-error promotion]
\label{lem:p7-same-codomain-obstruction}
The obstruction is not merely that the present residual tables are finite.
They live in a different proof codomain from the trajectory-error conclusion.
Let \(\mathcal R_h^{\rm mech}:\mathcal X_h\to\mathcal Y_h^{\rm res}\)
denote a candidate mechanism residual/reaction diagnostic on a closed-loop
branch, and let \(E_h:\mathcal X_h\to\mathcal Y_h^{\rm traj}\) denote the
reported position/velocity error map used by the dynamic-order claim.  A proof
of P7 would have to fix, before the data are read, the branch chart, row
weights \(W_h\), residual norm, trajectory norm, and reference/floor term
\(\rho_h\), and then prove a uniform transfer estimate such as
Eq.~\eqref{eq:p7-same-codomain-transfer},
\begin{equation}
\label{eq:p7-same-codomain-transfer}
  \|E_h z-E_h z_h^\star\|_{\rm traj}
  \le
  C_{\rm P7}
  \|W_h(\mathcal R_h^{\rm mech}(z)-\mathcal R_h^{\rm mech}(z_h^\star))\|_{\rm res}
  + C_{\rm P7}\rho_h ,
\end{equation}
for all \(z\) in the same proof tube, with \(C_{\rm P7}\) independent of
\(h\) and with \(\rho_h\) bounded in the same asymptotic scale required by
the promoted trajectory claim.  Without this same-codomain transfer, a small
constraint residual, a small reaction residual, or a small posterior
residual-to-error ratio is an element of \(\mathcal Y_h^{\rm res}\), not a
bound in \(\mathcal Y_h^{\rm traj}\).  The four-link, slider-crank, and
source-policy residual rows therefore remain mechanism-coverage diagnostics
and cannot be substituted into Theorem~\ref{thm:g6fullva-order}.
\end{lemma}

\begin{proof}
Equation~\eqref{eq:p7-same-codomain-transfer} is the missing bridge from the
residual diagnostic codomain to the trajectory-error codomain.  Without a
predeclared same-branch operator, row weighting, trajectory norm, and
\(h\)-independent constant \(C_{\rm P7}\), the norm of an element in
\(\mathcal Y_h^{\rm res}\) supplies no inequality in
\(\mathcal Y_h^{\rm traj}\).  Posterior ratios computed from the same finite
reference data cannot define this operator or constant, because doing so
would choose the transfer after the residual table has already been inspected.
Thus residual-size diagnostics can corroborate mechanism coverage, but they
cannot enter the local-defect, Gronwall, or reporting-map chain of
Theorem~\ref{thm:g6fullva-order} as trajectory-error/order evidence.
\end{proof}

\noindent\textbf{Kernel and observability obstruction.}
The required P7 estimate is also an observability statement, not only a
residual-rate statement.  Even a preasymptotic \(O(h^7)\) trend in
\(\|\mathcal R_h^{\rm mech}\|\) would not control the reported
trajectory components if the linearized residual has an \(h\)-uniform
near-null direction whose image under \(E_h\) is nonzero.  A valid P7 proof
must therefore fix the branch chart, gauge or quotient for algebraic
variables, row weights, and trajectory projection, and then prove that every
trajectory-visible perturbation is controlled by the residual operator up to
the declared floor/estimator term.  Reaction residuals, constraint closure,
or posterior residual-to-error ratios do not by themselves exclude such
trajectory-visible kernel directions.  This observability or inf-sup estimate
is exactly the open stability component of P7-3, and it must be fixed
before the residual table is read.

The proof uses two deliberately separate implementation routes. The active
route is the direct 132-row residual-defect route: the 96 non-dynamic rows
have an \(O(h^7)\) defect by the kinematic/lower-pair certificate, and the
36 Newton--Euler rows vanish under the D5 direct-substitution certificate on
the lifted Gauss stage. This is the route used in
Theorem~\ref{thm:g6fullva-order}. The route-level use is
anti-circular in the following narrow sense: the D5 certificate is evaluated
at the lifted Gauss stage \(Z_G\) and at the implemented residual rows before
the local perturbation and local-to-global transfer steps, so it does not use
the computed trajectory error or the conclusion of
Theorem~\ref{thm:g6fullva-order}.

\begin{table}[htbp]
\centering
\caption{Direct-route anti-circularity table for
Theorem~\ref{thm:g6fullva-order}.}
\label{tab:direct-route-anticircularity-ledger}
\scriptsize
\setlength{\tabcolsep}{2pt}
\renewcommand{\arraystretch}{0.9}
\begin{tabular}{P{0.18\linewidth}P{0.28\linewidth}P{0.26\linewidth}P{0.20\linewidth}}
\toprule
Proof object & Fixed before & May depend on & Explicit non-dependence \\
\midrule
Lifted Gauss stage \(Z_G\) &
D5 evaluation and local perturbation. &
Retained smooth-lift and compact-tube hypotheses. &
Inexact Newton iterates, computed trajectory error, or the theorem conclusion. \\
Implemented residual rows &
D5 row substitution and PC2 consumption. &
The implemented 132-row formula map, row labels, and runtime AD binding. &
Source-paper residual replacement, residual/reaction diagnostics, or
source-policy rows. \\
D5 direct-substitution certificate &
Stage-root perturbation, local-defect summation, and local-to-global transfer. &
Row-local Newton--Euler identities and the direct \(Z_G\) substitution. &
Discrete Gronwall, reporting-map consequences, observed \(h\)-sweeps, or P7
residual-to-error promotion. \\
Local perturbation/local defect &
The theorem's grid-point conclusion. &
PC1--PC2 plus the retained compact-tube, endpoint-closure, and inexact-Newton
interfaces, with retained P6 and PC4 as the solver-scale interface and
residual-to-error exclusion. &
Separate source-policy/full-\tfe{} readiness or reproducibility package readiness evidence. \\
Theorem conclusion &
No antecedent input; it is the downstream output. &
The preceding local-defect, stability, local-to-global, and reporting-map
steps. &
Feedback into \(Z_G\), D5 identities, P5, or the primitive symbolic route. \\
\bottomrule
\end{tabular}
\end{table}

Direct-route anti-circularity rule: the direct route
is admissible because D5 direct-substitution identities are fixed before local
perturbation, local-to-global transfer, and the theorem conclusion; \(Z_G\)
and the implemented residual rows are fixed at the same pre-perturbation level.
The theorem conclusion
does not justify row identities; finite \(h\)-sweeps, residual/reaction
tables, source-policy rows, and the primitive symbolic route are not sources
for D5. This table does not prove P1/P2/P6/P7, the dynamic symbolic oracle,
the primitive/Taylor route, source-policy/full-\tfe{} readiness, or
reproducibility package readiness evidence.
The nonuse rule is one-way: row identities feed the theorem, while the
theorem's downstream order estimate cannot feed back into row identity,
implementation binding, or D5 closure.

The primitive symbolic
Taylor route is retained only as a conditional decomposition of the same
dynamic rows into state-, acceleration-, multiplier-, geometry-, and
gyroscopic-lift obligations. Because those primitive lift obligations remain
separate assumptions, the primitive route is not used to close the theorem.
This separation prevents the open primitive symbolic route from weakening the
accepted direct proof, and also prevents the direct proof from being
misreported as a complete source-paper \tfe{} residual replacement.

The primitive-route status table is one-closed/five-open:
\(P_{\mathrm{tube}}\) is closed, while \(P_{\mathrm{state}}\),
\(P_{\mathrm{acc}}\), \(P_{\lambda}\), \(P_{\mathrm{geom}}\), and
\(P_{\mathrm{gyro}}\) remain open.  The D5 term inventory therefore has no
certified induced primitive-route Taylor subterm bounds in the active theorem
route.
Consequently, the theorem does not import a primitive/Taylor bound as PC2
evidence; it consumes the direct 132-row route and carries the compact-tube
solver condition separately.

\begin{table}[htbp]
\centering
\caption{Implementation-route/certificate separation table for
Theorem~\ref{thm:g6fullva-order}.}
\label{tab:implementation-route-oracle-separation-ledger}
\scriptsize
\setlength{\tabcolsep}{2pt}
\renewcommand{\arraystretch}{0.9}
\begin{tabular}{P{0.18\linewidth}P{0.29\linewidth}P{0.27\linewidth}P{0.18\linewidth}}
\toprule
Route or certificate & Theorem role & Closed evidence & Explicit non-use \\
\midrule
Active direct residual-bridge route &
The stage-residual input used by Theorem~\ref{thm:g6fullva-order}. &
The 96-row \(O(h^7)\) non-dynamic certificate plus the 36
Newton--Euler direct-substitution zero rows give the implemented 132-row
route. &
Primitive/Taylor subterm closure or source-paper residual replacement. \\
Runtime formula/AD binding &
P3 identifies the implemented residual/Jacobian path being solved and
differentiated. &
Full 132-row formula map and finite AD checks bind the implementation path. &
A global source-code equivalence proof or theorem-level solver-policy
theorem. \\
AD-expanded certificate &
Closes the implementation-path/derivative-cell identities needed for the active
direct proof boundary. &
The AD-expanded closure covers the Newton--Euler derivative cells for the
implemented direct route. &
Primitive symbolic defect certificate or dynamic symbolic oracle completion. \\
Primitive symbolic route &
Diagnostic decomposition of the same 36 dynamic rows into primitive lift
obligations. &
One primitive obligation is closed and five remain open; no induced
primitive-route Taylor subterm bounds are invoked by the theorem. &
Not an input to the PC2 stage-residual condition and not a contradiction of the
active direct-substitution route. \\
Source-paper residual replacement &
No theorem input. &
Kept outside the accepted proof and outside current source-policy readiness. &
Full-\tfe{} replacement, source-policy reproduction, or reproducibility package readiness claims. \\
\bottomrule
\end{tabular}
\end{table}

Implementation-route/certificate separation: the
accepted proof consumes the implemented 132-row direct route, runtime formula/AD
binding, and implementation-path formula certificate only for the
stage-residual argument discharged by the direct route. The open primitive
symbolic route and the non-active primitive symbolic-oracle completion record
remain diagnostic records for separate proof extensions; they are
not inputs to the theorem's stage-residual condition, not contradictions of the direct D5 substitution closure,
not source-paper residual replacement, not source-policy/full-\tfe{}
readiness, and not reproducibility package readiness evidence.
The direct D5 substitution closure is consumed only through this implemented
132-row direct-route proof path.

The proof dependencies used by the theorem are:
the PC2 stage-residual condition consumes the 96-row non-dynamic certificate and the 36-row D5
direct-substitution certificate; P3 consumes only runtime formula/AD binding;
P6 consumes the compact-tube branch-solver hypothesis, while finite solver probes
attach only to P6 and do not prove the theorem-level solver condition; and P7
records the residual-to-error boundary, while residual tables
are recorded only as diagnostics for the open P7 residual-to-error boundary and do not prove a separate residual-to-error transfer
theorem. The primitive symbolic route is a separate primitive-route record and is
not an input to the PC2 stage-residual condition. The PC2 discharge therefore
does not prove P6 or P7; proving either statement would require independent
proof. Likewise, an open primitive symbolic route cannot
invalidate the accepted direct-substitution closure.
For the theorem dependency boundary, \mbox{finite solver probes attach only to P6}.
Equivalently, the primitive symbolic route \mbox{is not an input to the PC2
stage-residual condition}.

\noindent\textbf{Scope of theorem.}
The mathematical claim of
Theorem~\ref{thm:g6fullva-order} is a conditional order-six theorem under the
retained P1--P3 theorem interfaces, the retained P4 binding side, and the P6
solver-scale interface; P4's 96-row certificate and P5's direct
Newton--Euler route supply the proved residual rows, while P7 is only the
open residual-to-error boundary. Equivalently, it does not claim an unconditional
theorem without theorem-domain interfaces, a closed
\(\eta_h^{\rm tube}\le c_\eta h^7\) solver-policy condition,
fixed-tolerance asymptotic proof, an accepted residual-to-error transfer
theorem for mechanism rows, or source-policy/full-TFE package readiness.
Any paper-facing shorthand about direct residual-bridge discharge is read only
in this direct-PC2 sense: it refers to the same-branch 132-row residual-value
bridge, not to the primitive/Taylor lane, the retained P6 solver-policy
theorem, the open P7 residual-to-error boundary, source-policy readiness, or
a full-\tfe{} replacement.

\begin{table}[htbp]
\centering
\caption{Scope of auxiliary evidence for Theorem~\ref{thm:g6fullva-order}.}
\label{tab:theorem-use-rule}
\scriptsize
\setlength{\tabcolsep}{2pt}
\renewcommand{\arraystretch}{0.9}
\begin{tabular}{P{0.22\linewidth}P{0.34\linewidth}P{0.36\linewidth}}
\toprule
Evidence class & Permitted theorem use & Explicit non-use \\
\midrule
Accepted direct stage-row route &
The 96-row non-dynamic certificate and the 36-row D5 direct-substitution
certificate supply the stage-residual input required by PC2 on \(Z_G\). &
This is not a primitive symbolic certificate, source-paper residual
replacement, or full-\tfe{} stage replacement. \\
Finite solver probes &
They diagnose branch-scale consistency for the retained P6 solver interface. &
They do not prove a theorem-level \(\eta_h\) solver-policy theorem and do not
turn fixed tolerances into an asymptotic proof. \\
Residual and reaction tables &
They support mechanism coverage and the open P7 residual-to-error diagnostic
route. &
They do not promote four-link, slider-crank, or source-policy residual rows to
trajectory-error/order evidence. \\
Source-policy and local reproducibility records &
They record reproducibility provenance for the narrowed claim and current
reproducibility record state. &
They do not prove full source-policy reproduction, full source-policy runner
package readiness, or external-superiority claims. \\
\bottomrule
\end{tabular}
\end{table}

Auxiliary-evidence scope: theorem inputs are restricted to the retained
P1/P2/P3 conditions, the retained P4 binding interface, the P6 solver-scale
condition, and the P5 direct Newton--Euler route. Finite solver
probes, residual/reaction tables, source-policy rows, and local reproducibility record
evidence are diagnostic or reproducibility-boundary evidence only; they are not promoted
to theorem-input status, residual-to-error transfer, a theorem-level
\(\eta_h\) solver-policy result, or source-policy/full-TFE readiness.
The compact theorem-input tuple is exactly P1/P2/P3 conditions, the retained
P4 binding interface, the P6 solver-scale condition, and the P5 direct route;
all other records in this table remain auxiliary evidence.

\begin{table}[htbp]
\centering
\caption{Quantifier/domain table for Theorem~\ref{thm:g6fullva-order}.}
\label{tab:quantifier-domain-ledger}
\scriptsize
\setlength{\tabcolsep}{2pt}
\renewcommand{\arraystretch}{0.92}
\begin{tabular}{P{0.17\linewidth}P{0.31\linewidth}P{0.26\linewidth}P{0.18\linewidth}}
\toprule
Quantifier layer & What is fixed or quantified & Valid domain & Explicit exclusions \\
\midrule
Fixed proof data &
Fix \(T\), the accepted compact proof tube, the reduced chart, the
Gauss-predictor branch rule, the retained P1--P3 contracts, the P4 binding
contract, the P6 solver-scale contract, the P5 direct-route certificate, and
the P7 output boundary;
then choose \(h_0\) and the theorem constants. &
The smooth accepted branch specified by Assumption~\ref{ass:regularity}. &
Finite \(h\)-sweep fits, source-policy rows, residual-only mechanism rows,
and local reproducibility record state. \\
One-step estimate &
For every \(0<h\le h_0\) and every \(y\) in the retained compact tube
satisfying the branch-solver hypothesis. &
The accepted branch-selected \method{} one-step map initialized from the
Gauss predictor. &
Arbitrary Newton initializations, remote nonlinear roots, remote charts, and
nonsmooth contact events. \\
Reported-grid domain qualification &
For any reported \(N_h\) with \(N_hh\le T\), \(t_n=n h\), and transition
steps \(0\le n<N_h\) that satisfy the predeclared accepted theorem-domain
conditions in \(\mathcal D_h\) before the estimate is invoked. Manuscript order rows instantiate the exact-final-time case
\(N_hh=T\) when reported; the \(N_hh<T\) case is only a prefix-grid estimate. &
The grid maximum over output indices \(0\le n\le N_h\). &
Any error claim on the leftover interval \(T-N_hh\), dense-output terminal
value at \(T\), or extra terminal solve outside the reported transitions.
Exact-final-time grids are included as the \(N_hh=T\) case but are not
required by the theorem. \\
Solver residual scale &
The theorem uses \(\eta_h^{\rm tube}\le c_\eta h^7\) uniformly on the
compact tube; the reported \(\eta_h^{\rm reported}=\max_{0\le n<N_h}\eta_{h,n}\) is only the
trajectory specialization. &
Branch-selected Newton residuals on the same accepted compact-tube branch. &
Fixed production tolerances or finite solver probes as asymptotic proof. \\
\bottomrule
\end{tabular}
\end{table}

Quantifier/domain rule: the theorem quantifiers are
ordered as fixed proof data, then \(h_0\) and constants, then all
\(0<h\le h_0\), all retained-tube states \(y\), and any reported grid
\(N_hh\le T\) on the accepted branch. They do not quantify over arbitrary
Newton roots, remote charts, residual-only mechanism rows, source-policy rows,
unreported terminal times, or local reproducibility record states. The table records
the theorem domain; it does not prove a theorem-level solver-policy result,
a P7 residual-to-error boundary, source-policy readiness, or full-TFE readiness.
The theorem constants as suprema rule is fixed at this point: all theorem
constants are compact-tube suprema chosen before any reported branch, observed
slope, residual log, or work-precision table is read.

\noindent\textbf{Reader-facing theorem summary.}
Before the admissibility clauses below, the mathematical content is a standard
local-to-global implication.  On one fixed compact branch, the implemented
\(132\)-row FullVA residual receives a same-branch direct \(96+36\) residual
certificate at the lifted Gauss stage \(Z_G\); the endpoint closure and
inexact Newton solve contribute controlled \(O(h^7)\) perturbations; and
accepted repeated steps on the same reported grid therefore satisfy the
local-defect estimate in Eq.~\eqref{eq:g6fva-local-defect-theorem}, the
reduced-grid estimate in Eq.~\eqref{eq:g6fva-reduced-grid-bound}, and the
reported position--velocity estimate in Eq.~\eqref{eq:g6fva-reported-grid-bound}.
The remaining clauses in this subsection fix the branch, norm, row scaling,
solver envelope, and non-claims needed to make that implication unambiguous;
they are not additional evidence sources and they do not add a second theorem.

\begin{definition}[Accepted branch-selected one-step map]
\label{def:accepted-branch-map}
The accepted branch-selected one-step map
\(\Psi_h^{\mathrm{G6FVA}}\) is defined on the accepted subset of the retained
compact proof tube. Its theorem domain is the subset
\(\mathcal D_h\) of retained-tube states for which all branch, local-inverse,
endpoint-ball, tube-retention, and scaled solver-residual conditions below
hold with the fixed proof norm and constants.  More precisely,
\(\mathcal D_h\) is the state projection of a predeclared admissible graph:
each accepted input state is paired with the selected Gauss-predictor branch,
local stage root, endpoint-output data, proof norm, and solver-residual
certificate for that transition.  Thus \(y\in\mathcal D_h\) means that this
same input state and its selected transition data satisfy the branch,
endpoint, and P6 predicates; P6 is not inferred from state membership alone.
For \(y\in\mathcal D_h\), the
map is obtained by: (i) initializing the stage solve from the Gauss predictor
in the fixed local chart; (ii) selecting the local root of the implemented
132-row residual in the same chart, with fixed row ordering, finite row
scaling, and proof norm, and binding any computed iterate to that local
branch; (iii) applying the endpoint-output map
\(\mathcal P_h=\mathcal C_h\circ\mathcal E_h\), where \(\mathcal E_h\) is the
Lie endpoint reconstruction and \(\mathcal C_h\) is the accepted local
endpoint-closure map; and (iv) accepting the transition only when the computed
Newton iterate on that branch has residual satisfying the compact-tube solver
hypothesis. The rule is: \mbox{map is undefined for theorem purposes} whenever any
branch, tube, or solver-scale condition fails; the theorem makes no claim
about that step. A reported trajectory
is repeated application only over defined transitions \(0\le n<N_h\);
equivalently, a grid invocation requires \(y_n\in\mathcal D_h\) for every
transition index \(0\le n<N_h\). The exact comparison states
\(y(t_n)=\varphi_{t_n}(y_0)\) at which the local defect is evaluated must also
be admissible as exact-state inputs for that local-defect comparison. We
denote this same fixed input-domain condition by \(y(t_n)\in\mathcal D_h\)
below; it is a statement about the branch-selected map evaluated at the exact
state, not a post-run reacceptance test. The compact-tube part is supplied by
the same sub-tube/first-exit bootstrap, while
the branch and solver-scale parts remain retained theorem conditions for the
one-step map evaluated at those inputs; otherwise the term
\(\Psi_h^{\mathrm{G6FVA}}(y(t_n))-\varphi_h(y(t_n))\) is not a theorem
object. The compact-tube membership is supplied by the first-exit bootstrap
after \(h_0\) is fixed, not by an a posteriori deletion of failed steps; the
scaled residual part is not supplied by first exit and remains the retained
P6 hypothesis.
\(\eta_h^{\rm reported}\) is the maximum over those accepted transitions, not an extra
terminal solve or remote nonlinear root.
\end{definition}

First-exit/P6 separation: Lemma~\ref{lem:local-global-transfer} proves
compact-tube retention for the accepted numerical and exact comparison states.
It does not prove the scaled residual part of \(\mathcal D_h\).  The P6
condition \(\eta_h^{\rm tube}\le c_\eta h^7\) is a separate branch-solver
hypothesis for each accepted transition; if that condition fails, the
transition is outside the theorem instance rather than a lower-order case
covered by the same constants.

\paragraph{Common admissible-window checkpoint}
The theorem uses one common admissible small-step window, not a collection of
per-lemma or per-run windows.  After the compact tube, chart, row ordering,
row scaling, proof norm, Gauss-predictor branch rule, endpoint-anchor family and
right-inverse constants, stage inverse constants, strong residual-inverse
constants, stability constant, and \(c_\eta\) have been fixed, define the
common window in Eq.~\eqref{eq:common-admissible-window},
\begin{equation}
\label{eq:common-admissible-window}
  h_\star=
  \min\{h_{\rm G},h_{\rm res},h_{\rm K},h_{\rm E},h_{\rm N},
  h_{\rm stab},h_{\rm rep},\bar h\},
\end{equation}
where the terms are respectively the windows for the reduced Gauss local
defect, the 132-row residual bridge, the Kantorovich stage-root ball, the
endpoint-closure ball, the inexact-Newton branch ball, the stability/tube
bootstrap, the reporting-map chart neighborhood, and the preselected
compact-tube residual-map window.  The theorem then chooses
\(h_0\le h_\star\), with the additional explicit Kantorovich restrictions
in Eq.~\eqref{eq:common-admissible-kantorovich},
\begin{equation}
\label{eq:common-admissible-kantorovich}
  2MC_Rh_0^7\le r,\qquad
  2M^2LC_Rh_0^7\le \frac12 ,
\end{equation}
and the endpoint/inexact-Newton ball margins required in
Lemmas~\ref{lem:endpoint-closure} and \ref{lem:inexact-newton}.  This is a
single same-branch admissibility rule: if one component window fails, the
transition is outside Definition~\ref{def:accepted-branch-map}; it is not
covered by a different constant, a different Newton root, or a later
diagnostic row.  Shrinking \(h_0\) after the proof constants are fixed is
allowed, but selecting a new tube, inverse constant, endpoint-anchor family,
residual norm, row scaling, or solver-envelope constant after reading a run
would define a different theorem instance.

\noindent\textbf{Solver-envelope admissibility criterion.}
For the fixed theorem data, define the solver-admissible transition set
in Eq.~\eqref{eq:solver-admissible-set},
\begin{equation}
\label{eq:solver-admissible-set}
  \mathcal S_h(c_\eta)=
  \{\,n:\ y_n\in K,\ \widetilde Z_{A,n}\in B_{r_A}(Z_{A,n}),\
  \|F_{A,h}(\widetilde Z_{A,n};y_n)\|\le c_\eta h^7\,\}.
\end{equation}
The theorem is applied only on reported transition indices
\(\{0,\ldots,N_h-1\}\subset\mathcal S_h(c_\eta)\), with the norm, row
scaling, branch rule, local ball, and \(c_\eta\) fixed before the run.  This
criterion is a sufficient trajectory-level verification condition for the
residual-size part of P6 once branch-ball membership is also certified; it is not a
postprocessing filter.  Shrinking \(h_0\) can satisfy small-step requirements
after the constants are fixed, but it cannot turn a fixed absolute tolerance,
a residual ratio observed after the run, or a failed branch-ball membership
into the compact-tube solver envelope required by
Lemma~\ref{lem:inexact-newton}.

Accepted-map definition boundary: \(\Psi_h^{\mathrm{G6FVA}}\)
is a branch-selected accepted-domain map tied to the Gauss predictor, the local accepted
stage root of the implemented 132-row residual, the fixed proof norm, and the
endpoint-output map \(\mathcal P_h=\mathcal C_h\circ\mathcal E_h\). It is
\mbox{not a global nonlinear-solver selection rule}, not a remote-chart equivalence claim,
not a residual-only mechanism-row error theorem, not a source-policy
reproduction result, and not a terminal-time correction outside the accepted
transition set.
The phrase branch-selected accepted-domain map tied to the Gauss predictor
means that the map, the local stage root, the endpoint-output map, and the
scaled solver-residual condition are one same-branch theorem object.

Local-defect map-chain definition: on every defined accepted
transition, the proof compares exactly one chain of maps on the same compact
branch: \(\varphi_h\), the exact reduced flow; \(\Psi_h^G\), the reduced Gauss
collocation endpoint map; \(\Psi_h^A=\mathcal E_h(Z_A)\), the endpoint
reconstructed from the \mbox{exact local root of the implemented} 132-row residual
before endpoint closure; \(\Psi_h^E=\mathcal C_h(\Psi_h^A)\), the
endpoint-closed exact-stage reference for the inexact Newton perturbation; and
\(\Psi_h^{\mathrm{G6FVA}}=\mathcal P_h(\widetilde Z_A)\), the
endpoint-closed inexact Newton map. All four
local-defect terms are measured in the same reduced-chart endpoint norm after
the fixed row ordering, finite row scaling, and proof norm have already
defined the residual-to-stage perturbation input. The triangle inequality
does not change branch, chart, norm, or transition index, and none of its
terms is supplied by mechanism residual/reaction tables, source-policy rows,
or a terminal-time correction outside the reported transition set.

Algorithm-order consistency: Algorithm~1 computes the inexact stage iterate
\(\widetilde Z_A\) first, reconstructs its Lie endpoint by \(\mathcal E_h\),
and then applies the accepted local endpoint closure \(\mathcal C_h\). The
proof uses the same implemented output order in
\(\Psi_h^{\mathrm{G6FVA}}=\mathcal P_h(\widetilde Z_A)\), with
\(\mathcal P_h=\mathcal C_h\circ\mathcal E_h\). The auxiliary map
\(\Psi_h^E=\mathcal P_h(Z_A)\) is only the exact-root reference obtained by
applying that same endpoint-output map to \(Z_A\); it is not a reordered
algorithm and it is not an endpoint closure performed before Newton. Thus the
\(C_N\eta_h^{\rm tube}\) term compares two already endpoint-closed outputs,
whereas the \(C_Eh^7\) term compares the exact-stage unclosed and closed
outputs. The local-defect sum therefore never mixes an unclosed inexact
endpoint with a closed exact endpoint.

\noindent\textbf{Non-vacuity boundary.}
Theorem~\ref{thm:g6fullva-order} is an abstract conditional consistency
theorem for the nonempty compact branch specified in
Assumption~\ref{ass:regularity} and
Definition~\ref{def:accepted-branch-map}. Here nonempty means that the exact
comparison states lie in the state projection \(\mathcal D_h\) of the
predeclared admissible graph, and their lifted Gauss-predictor stages and
selected transition data lie in the associated local stage chart, compact
tube, endpoint ball, proof norm, and P6 solver-envelope predicates for
\(0<h\le h_0\); it is a
theorem-domain condition, not a set obtained by deleting failed reported
steps after a run. The finite smooth-branch runs and
endpoint/solver diagnostics are reported only as finite diagnostics for
branch-model relevance on the accepted benchmark, not as theorem evidence for
P6. They do not prove global branch selection, remote-root exclusion, or P6
solver-policy theorem. Thus the theorem is not made unconditional by the
experiments, and the experiments are not used to define away the theorem
domain after the fact.

\noindent\textbf{Direct residual-bridge scope convention.}
In this proof, direct residual-bridge closure denotes only the
residual-bridge/Kantorovich input used by
Theorem~\ref{thm:g6fullva-order}: the 96 non-dynamic implementation rows, the
direct D5 36-row Newton--Euler identity at the lifted Gauss stage, and the
local perturbation route that consumes them.  It does not prove a
solver-policy theorem or a P7 residual-to-error boundary, does not complete the separate
primitive/Taylor route, and does not establish multiplier/reaction output-order
claims, source-policy reproduction, or full-\tfe{} replacement. This is a mathematical
scope convention for the direct route; it is not an additional theorem
assumption and it does not enlarge the theorem conclusion. The exclusion is
part of this theorem instance, not a post-hoc classification that can be
changed after reading numerical tables or auxiliary records.
Only the direct PC2 residual-value bridge is discharged here.  This
reader-facing discharge does not close the primitive/Taylor route, P6, P7,
source-policy, multiplier/reaction, or full-\tfe{} interfaces.

\noindent\textbf{Theorem input-output reading rule.}
The theorem is read as one forward implication on a fixed compact branch.  Let
\(\mathcal I_h(K)\) collect the retained admissibility interfaces: P1, P2,
and P3 on the accepted compact tube; the separate P6 solver-scale interface;
and the branch, row-ordering, row-scaling, and proof-norm binding needed to
interpret P4.  Let
\(\mathcal B_h^{96+36}(Z_G)\) denote the proved P4 96-row non-dynamic
certificate plus the P5 direct 36-row identity at the lifted Gauss stage.  In
this notation P4 is not hidden inside the retained interface as an assumed
defect estimate: its row-local \(O(h^7)\) certificate is a discharged input,
but it is meaningful only after those admissibility interfaces have
fixed the same branch and residual convention used by the theorem. Then the
proof chain is Eq.~\eqref{eq:theorem-forward-chain}:
\begin{equation}
\label{eq:theorem-forward-chain}
\begin{aligned}
\mathcal I_h(K)+\mathcal B_h^{96+36}(Z_G)
&\Longrightarrow
  \|F_{A,h}(Z_G;y)\|\le C_Rh^7\\
&\Longrightarrow
  \|\Psi_h^{\mathrm{G6FVA}}(y)-\varphi_h(y)\|\le C_{\rm loc}h^7\\
&\Longrightarrow
  \max_{0\le n\le N_h}\|y_n-y(t_n)\|\le C_{\rm red}h^6,\\
&\hspace{2.3em}
  \max_{0\le n\le N_h}\|\mathcal R(y_n)-\mathcal R(y(t_n))\|
  \le C_{qv}h^6 .
\end{aligned}
\end{equation}
No arrow in this chain is used backwards.  In particular, observed grid
orders, residual tables, finite solver logs, or a discharged direct residual bridge do
not prove the retained compact-tube interfaces, turn the solver-scale
condition into a separate solver theorem, prove a residual-to-error theorem
for reported mechanism rows, or certify the separate primitive/Taylor route.
The same instance discipline governs numerical reporting: after
\(\mathcal I_h(K)\), \(\mathcal B_h^{96+36}(Z_G)\), the branch, proof norm,
row scaling, residual envelope, \(c_\eta\), and \(h_0\) have been fixed, a
reported run may only specialize that theorem instance. It cannot
retroactively change the theorem branch, norm, row scaling, or residual
envelope to make one of the displayed implications true.

\begin{lemma}[P6/P7 non-closing arrow contract]
\label{lem:p6-p7-nonclosing}
For the theorem instance fixed by the compact branch, residual convention,
proof norm, endpoint-output map, and solver envelope, P6 and P7 are
forward-only boundary clauses.  The P6 clause is an input-domain clause:
Eq.~\eqref{eq:p6-forward-clause},
\begin{equation}
\label{eq:p6-forward-clause}
  \eta_h^{\rm tube}\le c_\eta h^7
  \Longrightarrow C_N\eta_h^{\rm tube}\le C_Nc_\eta h^7 ,
\end{equation}
so it only bounds the inexact-Newton summand after the compact-tube branch,
proof norm, row scaling, and residual envelope have been fixed. The reverse
arrow is not used: a small reported residual, a tolerance sweep, or an
observed sixth-order slope does not prove the compact-tube P6 condition. The
P7 clause is an output-exclusion clause:
Eq.~\eqref{eq:p7-nonpromotion-clause},
\begin{equation}
\label{eq:p7-nonpromotion-clause}
\begin{aligned}
  \max_n\|y_n-y(t_n)\|=O(h^6)
  &\not\Longrightarrow
  \{\text{residual-only mechanism rows are}\\
  &\hspace{5.0em}\text{trajectory-error/order evidence}\}.
\end{aligned}
\end{equation}
Residual and reaction rows therefore remain mechanism-coverage diagnostics
until a separate branch-compatible residual-to-error theorem supplies the
fixed transfer operator and \(h\)-independent constant required in
Table~\ref{tab:p7-residual-to-error-obligation-ledger}. These two
non-closing arrows are part of the theorem boundary, not additional
conclusions.
\end{lemma}

\begin{proof}
The P6 statement is exactly Lemma~\ref{lem:inexact-newton} applied after the
accepted branch, residual norm, row scaling, endpoint-output map, branch ball,
and residual envelope have already been fixed.  That lemma maps the
proof-norm residual bound to the inexact-Newton endpoint perturbation; it
does not map a finite residual log, tolerance sweep, or observed grid slope
back to the compact-tube branch membership, strong local inverse, or
\(h^7\) residual envelope required for P6.  Reversing that arrow would change
the theorem domain after inspecting diagnostics and would therefore be a
different theorem instance.

The P7 statement is separate from the local-defect and Gronwall chain.  The
theorem output is a reduced-chart and reported position--velocity grid
estimate for the accepted branch-selected map.  A residual or reaction row is
an element of a different codomain: it is a mechanism-coverage diagnostic
unless a branch-compatible residual-to-error operator with an
\(h\)-independent constant is proved on the same reported transition set.
The seven entries in Table~\ref{tab:p7-residual-to-error-obligation-ledger}
are precisely the required data for such a transfer theorem.  Hence P6 cannot
be proved by finite solver diagnostics, and P7 cannot be proved by residual
or reaction tables without adding the separate P7 boundary-theorem ingredients
named above.
\end{proof}

\noindent\textbf{Pre-data theorem-instance discipline.}
Fix a compact tube \(K\), a reduced chart, the Gauss-predictor branch
selector, the row-scaled residual \(F_{A,h}\), the proof residual norm, the
endpoint-output map, the reporting map, constants \(c_\eta\) and \(h_0\), and
the local inverse, endpoint, stability, and norm-equivalence constants before
any solver log, residual table, fitted order, or source-policy row is read.
For this fixed theorem instance, Theorem~\ref{thm:g6fullva-order} may consume
only statements attached to the same instance:
the \(96+36\) row residual bridge at \(Z_G\), the compact-tube inverse and
endpoint hypotheses, the branch-selected inexact-Newton envelope
\(\eta_h^{\rm tube}\le c_\eta h^7\), and the local-to-global stability
estimate.  A reported run can specialize the instance only by showing that
its accepted transitions use the predeclared branch, norm, row scaling, and
solver envelope.  It cannot define these objects after the residual or error
tables are inspected.

\noindent\emph{Argument.}
The proof of Theorem~\ref{thm:g6fullva-order} is a composition of maps with a
fixed domain and codomain.  The stage-root perturbation uses
\(D_ZF_{A,h}\) and the residual norm fixed on the compact tube; the endpoint
closure uses the predeclared endpoint-output map; the inexact-Newton term
uses the same residual norm through Lemma~\ref{lem:inexact-newton}; and the
Gronwall step uses the fixed reduced-chart stability family.  The constants
entering the local defect are the pre-data tuple in
Eq.~\eqref{eq:local-defect-constants},
\begin{equation}
\label{eq:local-defect-constants}
  C_G,\quad C_A,\quad C_E,\quad C_Nc_\eta ,
\end{equation}
are therefore chosen before the reported trajectory is evaluated.  If a later
diagnostic changes the branch selector, residual row weights, residual norm,
trajectory norm, reference policy, floor-exclusion rule, or \(c_\eta\), the
constants above and hence the theorem instance have changed.  Such a changed
diagnostic can only motivate a separate theorem instance; it is not a proof of the
already stated instance.

This pre-data order is exactly what separates the theorem-grade hypotheses
from finite diagnostics.  Solver logs may verify the recorded residuals for a
predeclared P6 instance; they do not create the compact-tube inverse or the
\(h^7\) envelope.  Residual and reaction tables may verify mechanism
coverage for a predeclared residual vocabulary; they do not create the fixed
P7 boundary operator or an \(h\)-independent residual-to-error constant.
Thus no diagnostic row can be moved upstream into the theorem proof unless it
is first reformulated as a predeclared theorem assumption or as an independent
lemma on the same fixed object.

\begin{lemma}[Same-object theorem-composition criterion]
\label{lem:same-object-composition}
Every lemma used in Theorem~\ref{thm:g6fullva-order} is attached to the same
mathematical object: the accepted Gauss-predictor branch, the row-scaled
\(132\)-row residual \(F_{A,h}\), the fixed proof norm, the endpoint-output
map \(\mathcal P_h=\mathcal C_h\circ\mathcal E_h\), and the compact-tube
solver envelope \(\eta_h^{\rm tube}\).  The P4 non-dynamic certificate, the
P5 Newton--Euler identity, Lemma~\ref{lem:stage-residual-defect},
Lemma~\ref{lem:endpoint-closure}, and Lemma~\ref{lem:inexact-newton} are
therefore composable only as a jointly admissible tuple whose branch, row
convention, proof norm, endpoint map, and solver envelope agree before the
constants are chosen.  A certificate for a different residual convention, a
solver log in a different norm, or a residual/reaction row measured after a
different endpoint operation is not an input to the theorem instance.
Auxiliary proof records, including the separate D5 primitive/Taylor inventory
and the P7 residual-to-error table, may record scope boundaries, diagnostics,
or separate-route material only after this tuple has been fixed: the separate
D5 primitive/Taylor inventory describes a separate sufficient proof route, whereas
P7 and residual/reaction tables remain diagnostic material outside the theorem estimate
until their own transfer theorem closes.  None of these records can
choose the tuple, supply an additional constant, or change the codomain of a
lemma already used in the composition.
\end{lemma}

\begin{proof}
The proof of Theorem~\ref{thm:g6fullva-order} uses the full-residual bridge
to supply a defect value for the same map \(F_{A,h}\) that appears in
Lemma~\ref{lem:stage-residual-defect}; that perturbation lemma then selects a
local stage root on the same Gauss-predictor branch and passes that stage
root through the same endpoint-output map used by Lemma~\ref{lem:endpoint-closure}.
Lemma~\ref{lem:inexact-newton} perturbs this accepted root only in the fixed
proof norm and only under the same compact-tube solver envelope.  Hence the
branch, row convention, norm, and endpoint map agree before the constants are
chosen and before the estimates are composed.

Equivalently, the reference-style C1--C4 slots are admissible only when their
outputs live in this same theorem object.  The C1 local-error slot may feed
the C3 perturbation slot only as the bridged residual value for the declared
\(F_{A,h}\); the C2 constrained-algebraic slot may feed C3 only as endpoint,
right-inverse, and lower-pair data for the same \(\mathcal P_h\) and compact
tube; and C3 may feed C4 only as a one-step map in the same reduced-chart
norm used by the stability recurrence.  A proof fragment whose natural
object is a BLieDF multistep update, a primitive \(162\)-subterm inventory, a
residual/reaction diagnostic, or a source-policy output row is therefore a
different typed object.  It requires an explicit equivalence or transport
lemma before it can enter the theorem composition.  Without that lemma, the
BLieDF fragment is reference guidance, the primitive inventory is
separate-route material, and residual or source-policy rows are outside the
theorem inputs.

If a P4 row certificate, P5 dynamic identity, endpoint-closure estimate, or
solver-residual bound were proved on mutually different stage branches,
residual conventions, proof norms, or endpoint maps, the codomain and
constant of one estimate would not be the domain data of the next estimate.
Such a collection cannot be spliced into a single theorem instance and could
not be substituted for the tuple consumed by the theorem without adding a
separate equivalence or transfer lemma.  This is the present FullVA analogue
of the fixed-discrete-object discipline in the
Wieloch--Arnold constrained-BDF proof: the reference motivates the order in
which objects are fixed and estimates are consumed, while the constants and
estimates used here are exactly those of the implemented Gauss6/FullVA
residual.
\end{proof}

\noindent\textbf{Theorem reading guide.}
The theorem below has four layers.  First, P1, P2, P3, and P6 fix the
domain of the branch-selected statement: smooth lift, compact tube, endpoint
map, implementation path, proof norm, stability, and solver residual
envelope.  Second, P4 and P5 supply the proved residual data on that fixed
object: the 96 non-dynamic FullVA rows contribute \(O(h^7)\), and the 36
Newton--Euler rows vanish at the lifted Gauss stage \(Z_G\).  Third, the
stage-root, endpoint, inexact-Newton, and local-to-global lemmas turn that
single \(132\)-row residual bridge into an \(O(h^7)\) local defect and an
\(O(h^6)\) same-initial-state reported-grid bound.  Fourth, the reporting
map converts that bound only to the displayed reduced-chart and
position--velocity output; residual, reaction, multiplier, constraint-output,
source-policy, and full-\tfe{} replacement claims are not theorem outputs.  The separate
primitive/Taylor inventory, residual-to-error promotion records, and
external-comparison/package records are outside these four layers.  The
primitive/Taylor propositions in this section are therefore conditional
auxiliary implications: they describe what a primitive-route proof for that
route must prove, but they cannot be read backward as evidence for P4, P5, P6, the
residual bridge, or the output scope consumed by this theorem.
In particular, the theorem's Taylor object is the fixed full residual map
\(F_{A,h}\) and the increment \(Z-Z_G\); the primitive \(162\)-term inventory is
a stricter alternative factorization of the dynamic block, not a hidden
premise of the direct \(96+36\) bridge.
Equivalently, Theorem~\ref{thm:g6fullva-order} consumes T1 and T2 only; T3 is
neither a premise nor a consequence of the theorem.
The Wieloch--Arnold constrained-BDF proof is invoked as proof-order discipline only for
this reading guide: no BLieDF local-error, multiplier, or recursion estimate
from that paper is a premise of the theorem below.

\noindent\textbf{Load-bearing proof spine.}
To separate the theorem's mathematical core from its boundary guards, the
proof below uses only the following five forward steps.  S1 applies the
classical three-stage Gauss local defect in the reduced FullVA chart.  S2
inserts the lifted Gauss stage into the fixed \(132\)-row implemented residual
and uses the direct \(96+36\) bridge to obtain
\(\|F_{A,h}(Z_G;y)\|\le C_Rh^7\).  In that step, the P5 direct substitution
identity gives the zero dynamic block at the same lifted Gauss stage and in
the same row ordering.  S3 applies the same-object
Taylor/Kantorovich, endpoint, and inexact-Newton perturbation estimates to
obtain the four-term local defect with
\(C_{\rm loc}=C_G+C_A+C_E+C_Nc_\eta\).  S4 propagates this local defect on
the accepted branch by the first-exit/Gronwall recurrence.  S5 applies the
fixed reporting map to the same reported grid.  P6 contributes only the
inexact-Newton summand in S3 through the retained envelope
\(\eta_h^{\rm tube}\le c_\eta h^7\); P7, the primitive/Taylor T3 route,
fixed-tolerance diagnostics, and source-policy records are boundary objects
and never feed S1--S5.

\noindent\textbf{Theorem normal form.}
Equivalently, the result is a numerical-analysis implication: fixed compact
branch hypotheses, one same-branch \(132\)-row residual certificate, endpoint
and solver perturbation estimates, and accepted one-step stability imply an
\(O(h^7)\) local defect and an \(O(h^6)\) same-initial-state reported-grid
position--velocity error.  The surrounding interface language identifies
which hypotheses and certificates instantiate this implication for the
implemented \method{} residual; it is not a second theorem and it does not
add conclusions beyond Eqs.~\eqref{eq:g6fva-local-defect-theorem},
\eqref{eq:g6fva-reduced-grid-bound}, and
\eqref{eq:g6fva-reported-grid-bound}.
Thus a reader may treat the interface labels as bookkeeping once their
provenance has been checked: the
mathematical statement remains the standard local-defect plus stability plus
reporting-map convergence argument, with the direct FullVA residual bridge
supplying the only nonstandard residual-defect input.

\noindent\textbf{Theorem input map.}
Tables~\ref{tab:p-interface-satisfaction-ledger}
and~\ref{tab:theorem-assumption-anchor-ledger} give the role of the named
interfaces consumed by the theorem.  Retained inputs are the P1/P2/P3
compact-branch and implementation-domain interfaces, the binding side of P4,
and the P6 solver-scale envelope.  Discharged residual inputs are only the
P4 96-row non-dynamic certificate and the P5 36-row Newton--Euler identity,
assembled on the same lifted Gauss stage.  Excluded from the theorem inputs
are a P7 residual-to-error boundary, the separate primitive/Taylor T3 route,
fixed-tolerance diagnostics, source-policy rows, and full-\tfe{} replacement
records.  These tables specify the theorem inputs for the invocation; they
are not an additional proof argument.

\noindent\textbf{Reference-style proof invariant.}
The Wieloch--Arnold constrained-BDF proof is invoked here only as a proof-order
invariant.  Its local-error theorem first inserts the exact constrained
solution into already fixed discrete equations; its global theorem then
propagates those local defects through the appropriate stability and
constraint-recursion slots.  The corresponding invariant for
Theorem~\ref{thm:g6fullva-order} is
\begin{equation}
\label{eq:reference-style-proof-invariant}
  \begin{gathered}
  \hbox{fixed FullVA residual and branch}
  \rightarrow
  \hbox{same-branch residual defect at } Z_G,\\
  \hbox{same-branch residual defect}
  \rightarrow
  \hbox{Taylor/Kantorovich local map},\\
  \hbox{Taylor/Kantorovich local map}
  \rightarrow
  \hbox{endpoint/solver perturbations},\\
  \hbox{endpoint/solver perturbations}
  \rightarrow
  \hbox{same-initial-state reported-grid estimate}.
  \end{gathered}
\end{equation}
Every theorem-facing arrow in
Eq.~\eqref{eq:reference-style-proof-invariant} is either proved by the cited
lemma or certificate, or retained explicitly as a P-interface hypothesis, for
the present row-scaled Gauss6/FullVA object.  No BLieDF local-error constant, BDF
zero-stability theorem, hidden-constraint multiplier estimate, or coupled
multistep recursion is imported.  Conversely, each downstream reported-grid estimate,
finite residual diagnostic, source-policy row, and primitive/Taylor status
table is not allowed to define an upstream arrow in the chain.

For this theorem invocation the named clauses are fixed before the theorem is
used.  P1 fixes the smooth lift and chart tube.  P2 fixes the compact stage
inverse, endpoint raw-defect/right-inverse interface, endpoint ball, and
stability/tube-retention data.  P3 fixes the implementation path, row
ordering, row scaling, and proof norm.  P4 supplies only the 96 non-dynamic
row-local certificate, and only under the P1--P3 binding.  P5 supplies the 36
Newton--Euler rows only through the direct substitution identity on the same
lifted Gauss stage.  P6 supplies only the branch-selected solver residual
envelope.  None of these clauses is inferred from observed slopes, residual
logs, source-policy rows, or post-hoc filtered solver logs.

\noindent\textbf{Reference technical-slot substitution.}
The constrained BLieDF proof of Wieloch--Arnold~\cite{wieloch2021bdf}
uses three technical slots before its final convergence statement: exact
solution insertion into fixed discrete equations to define local truncation
errors, a neighbourhood/start-value regime that keeps the Lie-group
coordinates and multiplier variables in the admissible proof tube, and a
coupled error recursion that propagates the local estimates.  In
Theorem~\ref{thm:g6fullva-order} these slots are not imported; they are
replaced by the present FullVA objects.  Exact solution insertion becomes the
insertion of the smooth Gauss lift \(Z_G\) into the fixed row-scaled residual
\(F_{A,h}\).  The neighbourhood/start-value regime becomes the retained
accepted-domain conditions on \(\mathcal D_h\): branch, tube, endpoint ball,
proof norm, implementation binding, and P6 solver envelope are fixed before
any run is inspected.  The start-value component of the BLieDF proof has no
multistep prefix analogue here: it is replaced by the same-reduced-chart
initial-state clause and by the one-step start-up boundary in the proof of
Theorem~\ref{thm:g6fullva-order}.  Any projection, chart-registration, or
branch-registration mismatch is therefore an explicit initial-error term
outside the local defect, not an unspoken finite-prefix estimate.  The
multistep coupled recursion becomes the one-step same-initial-state stability
and first-exit/Gronwall transfer of Lemma~\ref{lem:local-global-transfer}.
This is a slot-by-slot replacement of proof roles, not a transfer of BLieDF
estimates, hidden-constraint multiplier bounds, start-up estimates, or BDF
zero-stability constants.

\begin{theorem}[Conditional sixth-order error theorem for the accepted branch-selected \method{} map]
\label{thm:g6fullva-order}
Let the theorem instance be fixed before any numerical evidence is read: the compact
proof tube \(K\), accepted Gauss-predictor branch, reduced chart, row-scaled
implemented residual \(F_{A,h}\), proof norm, endpoint-output map, and
reporting map are those of Assumption~\ref{ass:regularity} and
Definition~\ref{def:accepted-branch-map}.  The retained/discharged partition
is part of the theorem statement.  For nonautonomous or driven rows this
means the time-indexed residual \(F_{A,h,t_n}\), the Gauss stage times
\(t_{n,i}=t_n+c_i h\), the endpoint map, and the proof norm are fixed on the
same transition slice before any estimate or diagnostic row is read; all
compact-tube constants used below are uniform over those accepted slices.
Retain the P1--P3 theorem interfaces, the
binding side of P4, and P6 on this branch: the smooth FullVA lift, uniform
stage inverse and derivative bounds, endpoint raw-defect/right-inverse and
endpoint-ball conditions, strong local residual inverse for the inexact-Newton
perturbation, stability with tube-retention bootstrap, implementation-path
binding, and the branch-solver envelope
\(\eta_h^{\rm tube}\le c_\eta h^7\) are fixed as branch-domain hypotheses.
P4 is split in this statement: its branch/norm/implementation-binding side is
retained as a theorem interface, while its proved 96-row non-dynamic
row-local certificate is consumed only below in the single residual-certificate
slot together with P5's 36 Newton--Euler rows.
The P6 clause is a residual-norm envelope for the selected Newton iterate in
the fixed proof norm; it is not an assumption that the endpoint state,
reported trajectory, or grid error is already \(O(h^7)\) or \(O(h^6)\).
Thus P6 is a retained theorem-domain condition, not an unresolved step in the
direct residual proof: the theorem proves the order of the accepted map
conditional on this envelope and does not claim an automatic Newton-policy
theorem.
The P7 residual-to-error boundary is not part of these retained
branch-domain hypotheses; it is an output nonclaim boundary for residual,
reaction, and mechanism-coverage rows.
Under these retained branch-domain hypotheses and the active same-branch
residual certificate, the accepted map satisfies exactly the local and
reported-grid estimates in
Eqs.~\eqref{eq:g6fva-local-defect-theorem}--\eqref{eq:g6fva-reported-grid-bound}.
\emph{Residual-certificate slot.}  For this theorem instance the active
residual-value certificate is fixed by Eq.~\eqref{eq:residual-certificate-slot}:
\begin{equation}
\label{eq:residual-certificate-slot}
  \mathsf R_{\rm active}(h,y)\equiv \mathsf R_{\rm dir}(h,y),
  \qquad
  \mathsf R_{\rm dir}(h,y)\Longrightarrow
  \|F_{A,h}(Z_G;y)\|\le C_Rh^7 .
\end{equation}
Here \(\mathsf R_{\rm dir}\) is the same-branch direct certificate of
Corollary~\ref{cor:route-separation-invariant}.  The theorem has a single
residual-certificate slot.  The separate primitive predicate
\(\mathsf R_{\rm prim}\) is not an additional premise; until it independently
proves all five primitive lift/reduction inputs and all \(162\) primitive
subterm bounds on the same tuple \((h,y,Z_G,F_{A,h})\), it supplies no theorem
input.  Equivalently, a replacement primitive/Taylor route must prove all five
primitive lift/reduction inputs and all \(162\) primitive subterm bounds before
it supplies theorem input; the present route status is exactly
Eq.~\eqref{eq:t3-nonclosure-status-card}.  If completed, a future primitive/Taylor certificate
may replace \(\mathsf R_{\rm active}\) only after proving a new same-tuple \(132\)-row
residual-value bound for the fixed object \((h,y,Z_G,F_{A,h})\); it may not combine with the accepted direct certificate to lower constants, discharge P6,
or prove P7.
In particular, a primitive/Taylor proof of the \(36\) Newton--Euler rows would
not enter the theorem merely as a \(36\)-row statement.  It would first have
to be reassembled with the same \(96\)-row non-dynamic certificate, row
ordering, row scaling, block norm, branch, and endpoint convention that define
the implemented \(132\)-row residual certificate.
Thus the theorem-facing Taylor assertion is a full-map assertion: after the
single same-tuple residual value for \(F_{A,h}(Z_G;y)\) is available,
Lemma~\ref{lem:stage-residual-defect} applies Taylor's formula to the fixed
implemented residual map and the increment \(Z-Z_G\).  It is not a separate
claim that each primitive Newton--Euler subterm has its own \(O(h^7)\) bound.
This matches the reference proof discipline in which local defects are formed
by inserting the smooth constrained solution into fixed discrete equations
before any multiplier or global-recursion estimate is used, while the
estimates themselves remain those of the present FullVA residual.
In the present theorem, \(\mathsf R_{\rm dir}\) is instantiated, not
supplemented, by the two discharged residual inputs on the same lifted Gauss
stage: P4's 96-row non-dynamic certificate gives the \(O(h^7)\)
non-dynamic block, and P5's direct Newton--Euler substitution identity gives
the 36 dynamic rows at \(Z_G\).  Consequently,
Lemma~\ref{lem:full-132-row-residual-bridge} supplies the residual-value
hypothesis \(\|F_{A,h}(Z_G;y)\|\le C_Rh^7\) required by
Lemma~\ref{lem:stage-residual-defect}.  Then
there exist \(h_0>0\) and constants
\(C_{\rm loc},C_{\rm red},C_{qv}\), independent of \(h\) and the reported grid length
\(N_h\), such that for \(0<h\le h_0\) and any fixed uniform grid
\(t_n=n h\), \(0\le n\le N_h\), with \(N_hh\le T\), whose every transition
satisfies the predeclared accepted branch, endpoint-ball, proof-norm, and P6
solver-envelope hypotheses before the estimate is invoked, the accepted
branch-selected \method{} one-step map satisfies, uniformly for
\(y\in\mathcal D_h\subset K\), where \(\mathcal D_h\) is the accepted theorem
domain of Definition~\ref{def:accepted-branch-map}, under the compact-tube scaled
branch-solver hypothesis \(\eta_h^{\rm tube}\le c_\eta h^7\),
\begin{equation}
\label{eq:g6fva-local-defect-theorem}
  \|\Psi_h^{\mathrm{G6FVA}}(y)-\varphi_h(y)\|\le C_{\rm loc}h^7 .
\end{equation}
Here \(\eta_h^{\rm tube}\) is the a priori compact-tube envelope in the fixed
proof norm; the reported
\(\eta_h^{\rm reported}=\max_{0\le n<N_h}\eta_{h,n}\) may only specialize this envelope after
the accepted branch-selected transitions have been fixed.
The corresponding reduced-chart grid values are generated from the same
reduced initial state by
\(y_{n+1}=\Psi_h^{\mathrm{G6FVA}}(y_n)\) on the accepted branch for
\(0\le n<N_h\).  For each transition, \(y_n\in\mathcal D_h\), and the
exact-state input used in the local-defect comparison satisfies the
corresponding admissibility condition, denoted \(y(t_n)\in\mathcal D_h\) in
the same admissible-graph sense.  This means the selected transition data for
the exact input also satisfy the retained branch, endpoint, proof-norm, and
P6 predicates; it is not an assertion that P6 follows from the exact state
alone.
The same-initial-state clause is part of the theorem statement: an initial
projection, chart-registration, or branch-selection discrepancy would enter
the Gronwall recurrence as a separate initial-error term and is not absorbed
into the constants \(C_{\rm red}\) or \(C_{qv}\).
The compact-tube part of this domain is handled by the first-exit
tube-retention bootstrap of Lemma~\ref{lem:local-global-transfer} after the
compact tube and \(h_0\) are fixed.  Branch-ball, endpoint-ball, proof-norm,
and P6 solver-envelope membership remain retained transition hypotheses on
the accepted branch; they are not inferred from the reported-grid estimate.  Under
these domain conditions, the local estimate in
Eq.~\eqref{eq:g6fva-local-defect-theorem} gives the reduced-grid bound
\begin{equation}
\label{eq:g6fva-reduced-grid-bound}
  \max_{0\le n\le N_h}\|y_n-y(t_n)\|\le C_{\rm red}h^6 ,
\end{equation}
and the reported position--velocity map
\(\mathcal R(y)=(q(y),v(y))\) satisfies
\begin{equation}
\label{eq:g6fva-reported-grid-bound}
  \max_{0\le n\le N_h}\|\mathcal R(y_n)-\mathcal R(y(t_n))\|
  \le C_{qv}h^6 .
\end{equation}
The map \(\mathcal R\) in Eq.~\eqref{eq:g6fva-reported-grid-bound} is only
the chart-to-position--velocity reporting map; it is not the P7 mechanism
residual operator \(\mathcal R_h^{\rm mech}\) used in the residual-to-error
certificate template.
Equations~\eqref{eq:g6fva-reduced-grid-bound}
and~\eqref{eq:g6fva-reported-grid-bound} are the only same-initial-state
reported-grid conclusions asserted by the theorem.
The constants may depend on the compact proof tube, \(T\), the regularity,
inverse, stability, endpoint-closure, and branch-solver scale constants in
Assumption~\ref{ass:regularity}, including \(c_\eta\), but not on \(h\),
\(N_h\), the dense/sparse backend, or the fixed production tolerances used in
finite diagnostic runs.
The quantifier order is fixed in the same sense: the compact tube, branch,
proof norm, row scaling, inverse constants, endpoint/right-inverse constants,
stability constant, and \(c_\eta\) are chosen first; then \(h_0\) is chosen
small enough for all local lemmas and the tube-retention bootstrap to hold;
only after this choice may a particular grid \(0<h\le h_0\) and reported
accepted transition set be considered.
Grid-scope boundary: the displayed maxima are over the same reported indices
\(0\le n\le N_h\), with accepted transitions \(0\le n<N_h\). Exact-final-time
reported order rows instantiate the theorem with \(N_hh=T\). If \(N_hh<T\),
the estimate is only a prefix-grid estimate; the leftover interval \(T-N_hh\),
any dense-output value, and any separate terminal-time value at \(T\) require
an additional argument outside this theorem. Undefined steps, branch changes,
or residuals measured outside the fixed proof norm are outside the theorem
domain rather than cases covered by the constants above.
\end{theorem}

\paragraph{Scope and non-claims of Theorem~\ref{thm:g6fullva-order}}
Under these compact-branch, endpoint/stability, implementation-binding, and
solver-scale conditions above, the accepted branch-selected map is order six
only as the same-initial-state reduced-chart bound and the reported
position--velocity bound on the same reported grid.  The theorem does not prove residual, reaction,
multiplier, or constraint-output error estimates; does not prove a P7
residual-to-error transfer; does not supply source-policy reproduction; and
it \mbox{does not turn fixed-tolerance finite runs} into an asymptotic
solver-policy proof.  Reported-grid domain qualification is therefore a specialization
rule, not a post-hoc acceptance rule: a reported run may instantiate the
theorem only when its transitions already satisfy the predeclared
\(\mathcal D_h\) branch, tube, endpoint, proof-norm, and P6 solver-envelope
conditions with the residual norm, row scaling, branch, and \(c_\eta\) fixed
before the data are read.
Accordingly, the reader-facing shorthand ``any reported fixed time
grid'' means only a reported fixed time grid whose transitions already satisfy
those predeclared conditions; it is not an unfiltered or post-hoc reported
grid.
Only one residual-value certificate is consumed in this theorem invocation.
A future primitive/Taylor certificate may replace the accepted direct
certificate only by proving a new same-tuple \(132\)-row residual-value bound;
it cannot be appended to the direct certificate to lower constants, remove
the solver-scale condition, or prove a residual-to-error theorem for reported
mechanism rows.
Closed-loop residual and reaction-consistency rows therefore remain
mechanism-coverage or diagnostic evidence unless a separate P7
residual-to-error theorem is proved for the same residual, branch, norm, and
reporting map.

\begin{lemma}[No reverse Taylor inference from the reported-grid estimate]
\label{lem:no-reverse-taylor-inference}
For the theorem instance fixed in Theorem~\ref{thm:g6fullva-order}, the
conclusion in Eq.~\eqref{eq:no-reverse-grid-conclusion}
\begin{equation}
\label{eq:no-reverse-grid-conclusion}
  \max_{0\le n\le N_h}\|y_n-y(t_n)\|=O(h^6)
\end{equation}
and any observed finite-grid order, residual history, reaction table, or
source-policy row do not imply any of the primitive/Taylor inputs that are
kept outside the active route.  In particular, they do not construct a
primitive projection \(\Pi_{\rm prim}\), do not prove the five open primitive
lift or bilinear estimates
\(P_{\mathrm{state}},P_{\mathrm{acc}},P_{\lambda},
P_{\mathrm{geom}},P_{\mathrm{gyro}}\), do not certify any of the \(162\)
primitive D5 Taylor subterm bounds, and prove neither P6 nor P7.  Hence the
separate T3 primitive/Taylor route remains a separate sufficient route
unless those primitive inputs and the associated residual-to-root, endpoint,
solver, and global-transfer chain are proved on the same compact branch.
\end{lemma}

\begin{proof}
The proof of Theorem~\ref{thm:g6fullva-order} is a one-way composition.  The
predeclared full residual \(F_{A,h}\), branch selector, proof norm, row
scaling, endpoint-output map, and compact-tube constants determine the
local-defect estimate first.  The Gronwall step then maps that local defect to
a grid error by summing over at most \(T/h\) transitions, and the reporting map
passes from reduced-chart error to position--velocity error.  Neither the
Gronwall summation nor the reporting map is an inverse theorem for the stage
residual or for the primitive dynamic subterms.

To reverse the conclusion one would need additional data not present in the
theorem: an \(h\)-independent inverse from reported grid errors back to the
accepted stage perturbation, an inverse from the full residual value to a
primitive \(162\)-subterm decomposition, and uniform lift estimates for the
five primitive inputs listed above.  Those are exactly the open
primitive-route obligations recorded in the D5 tables.  A finite observed
order or a small residual table can be consistent with these objects, but it
does not prove their existence with the branch, norm, scaling, and constants
fixed before the run.  Therefore the global \(O(h^6)\) output estimate and
finite diagnostics cannot be used backward to close the primitive/Taylor
route, P6, or P7.
\end{proof}

\begin{proof}[Proof of Theorem~\ref{thm:g6fullva-order}]
Let $\varphi_h$ denote the exact reduced constrained flow in the smooth chart
from Assumption~\ref{ass:regularity}, and let $G_h$ be the three-stage Gauss
collocation map applied to that reduced flow. The one-step estimate is local
in a retained-tube state \(y\) and uses only the compact-tube solver envelope
\(\eta_h^{\rm tube}\).  The proof's load-bearing solver object remains this
compact-tube envelope.  This paragraph records the retained P6 notation before
the local-defect sum; it is consumed only as the fourth summand after the
direct residual bridge, the full-map Taylor/Kantorovich stage-root
perturbation, and the endpoint-closure estimate have been fixed.  For each
transition in the theorem domain, the
accepted iterate is already on the Gauss-predictor branch ball, the proof
residual norm and row scaling are fixed, and the retained P6 envelope gives
\begin{equation}
\label{eq:accepted-newton-residual-envelope}
  \|F_{A,h}(\widetilde Z_A;y)\|
  \le \eta_h^{\rm tube}\le c_\eta h^7 .
\end{equation}
When Eq.~\eqref{eq:accepted-newton-residual-envelope} is enforced by a
predeclared proof-norm stopping rule, Lemma~\ref{lem:p6-residual-size-subcondition}
records only this residual-size instantiation of P6; the theorem proof still
retains the branch-ball, local-inverse, and compact-tube hypotheses before
using the perturbation lemma.
Lemma~\ref{lem:inexact-newton} then consumes exactly this same-branch
residual-size input from Eq.~\eqref{eq:accepted-newton-residual-envelope} and
gives
\begin{equation}
\label{eq:inexact-newton-output-perturbation}
  \|\mathcal P_h(\widetilde Z_A)-\mathcal P_h(Z_A)\|
  \le C_N\eta_h^{\rm tube}\le C_Nc_\eta h^7 .
\end{equation}
Equation~\eqref{eq:inexact-newton-output-perturbation} is the sole
inexact-Newton contribution to the local-defect sum below.
Only after a reported accepted transition set has been fixed may the
per-transition accepted Newton residuals be recorded as
\(\eta_{h,n}\), \(0\le n<N_h\), with reported-grid specialization
\(\eta_h^{\rm reported}=\max_{0\le n<N_h}\eta_{h,n}\). This is the
reported-trajectory specialization of the predeclared envelope; it is not an additional local
hypothesis, not a fitted residual sequence, and not fitted from the residual
logs. It is the maximum over the accepted branch-selected Newton solves only
after the reported transition set has been fixed. Branch consistency and the
retained P6 condition then give, for those recorded transitions and only in
the fixed proof norm,
\begin{equation}
\label{eq:reported-residual-specialization}
  \eta_{h,n}\le \eta_h^{\rm tube}\le c_\eta h^7 .
\end{equation}
Equation~\eqref{eq:reported-residual-specialization} is only a reported-grid
specialization of the predeclared compact-tube envelope.
The residual norm in this statement is the fixed proof norm of the accepted
\(132\)-row residual after the chart, row ordering, and row-scaling
conventions have been fixed; implementation logs can specialize the
per-transition values \(\eta_{h,n}\), and hence
\(\eta_h^{\rm reported}\), only through that same residual map and an
\(h\)-independent compact-tube norm equivalence. The proof does not use
hidden \(h\)-dependent row weights, and
any logged-norm comparison is through an \(h\)-independent
finite-dimensional norm-equivalence constant after the explicit residual
\(h\)-factors have already been defined. Finite logs and scaled-tolerance
probes can record a reported-grid diagnostic specialization of the displayed
maximum after the branch and proof norm are fixed, but they do not discharge
the tube-envelope hypothesis or create an independent solver-policy proof. Thus
the proof does not fit the constants \(C_N\), \(c_\eta\), or \(C_{\rm loc}\)
from reported Newton logs, and a small residual in a different norm, branch,
row scaling, or terminal solve is not an input to the estimate. Throughout
the proof, \(\Psi_h^{\mathrm{G6FVA}}\) denotes this branch-selected accepted map:
Newton is initialized from the Gauss predictor on the same compact-tube
branch, and the estimate is not a statement about arbitrary nonlinear roots.
Solver-envelope quantifier discipline: the scalar solver-envelope constant and fixed
proof norm are chosen before \(h\), \(N_h\), or the finite logs; changing that
constant after inspecting residuals changes the theorem assumption and
\(C_{\rm loc}\), not the validity of an already accepted proof instance.
Classical collocation theory gives a
uniform Gauss truncation constant \(C_G\), after shrinking \(h_0\) if needed,
such that Eq.~\eqref{eq:gauss-local-truncation} holds uniformly for \(y\) in
the compact trajectory tube \cite{hairer1993solving}:
\begin{equation}
\label{eq:gauss-local-truncation}
  \|\Psi_h^G(y)-\varphi_h(y)\|\le C_Gh^7
\end{equation}
Gauss truncation constant boundary: \(C_G\) is the analytic collocation
constant obtained from \(C^7\) reduced-vector-field bounds on the fixed
compact tube. It is fixed before implementation residual diagnostics,
endpoint-closure logs, work-precision tables, or observed smooth \(h\)-sweep
slopes are inspected; those slopes may corroborate the scale but cannot
choose \(C_G\), shrink the tube, or select the Gauss-predictor branch.
The standard Gauss theorem is applied only after the constrained system has
been represented in the smooth reduced chart of
Assumption~\ref{ass:regularity}. If the reduced field \(f(y,t)\) is treated as
nonautonomous, the usual autonomous augmentation by \(t'=1\) is used locally.
The Lie-group coordinates, lower-pair constraints, and chart-to-configuration
norm equivalences enter only through the fixed compact-tube bounds and do not
change the \(h^7\) local Gauss defect. No remote chart, nonsmooth event, or
global Lie-group order statement is invoked here.
\noindent\textbf{Stage-time binding for nonautonomous rows.}
For a reported transition starting at \(t_n\), the proof uses the
time-indexed residual \(F_{A,h,t_n}\) defined by the method equations with
stage times \(t_{n,i}=t_n+c_i h\).  The lifted Gauss stage is therefore
\(Z_G(y,t_n)\), obtained by inserting the exact reduced solution and its
FullVA lift at those same stage times.  Driven constraint right-hand sides,
loads, and prescribed kinematic functions are evaluated at \(t_{n,i}\)
before the residual map, row scaling, proof norm, and Newton branch are fixed
for that transition.  They are not unknown columns, not endpoint
post-corrections, and not values reconstructed from a terminal output or a
finite diagnostic table.  Thus the local defect compares
\(\Psi_{h,t_n}^{\mathrm{G6FVA}}(y)\) with
\(\varphi_{t_n+h,t_n}(y)\) on the same time slice, while the autonomous
augmentation above is only a notation for this synchronized time-indexed
one-step statement.
\noindent\textbf{Nonautonomous same-time invariant.}
For this transition, every object entering the local defect is keyed by the
same input time \(t_n\) and the same Gauss stage abscissae
\(t_{n,i}=t_n+c_i h\).  The Gauss lift \(Z_G(y,t_n)\), the implemented
residual \(F_{A,h,t_n}\), the direct P4/P5 residual certificate, the
endpoint map \(\mathcal E_{h,t_n}\), the endpoint-closure map
\(\mathcal C_{h,t_n}\), and the inexact-Newton perturbation map
\(\mathcal P_{h,t_n}\) are all formed for that same transition.  The
four-term local-defect sum therefore compares
\(\Psi_{h,t_n}^{\mathrm{G6FVA}}(y)\) with
\(\varphi_{t_n+h,t_n}(y)\) without mixing a residual evaluated at another
input time, a force value reconstructed from the terminal output, an
interpolated drive sample, or a dense-output endpoint.  The later Gronwall
recursion only concatenates these synchronized one-step maps at
\(t_n=n h\); it does not insert an additional time-registration or terminal
correction term.
The SO(3) BCH and adjoint effects are represented only through the retained P2
reduced-chart stability and reporting-map constants chosen for the present
compact tube. No BCH constant or Lie-algebra recurrence estimate from
Wieloch--Arnold is imported into the bound. This is the local Lie-algebra
error-recursion slot corresponding to the Wieloch--Arnold constrained-BDF proof, not
an imported estimate or an additional theorem conclusion.
Lemma~\ref{lem:fullva-stage-lift} identifies the mathematical FullVA stage
rows with the lifted reduced Gauss equations and pointwise constrained
dynamics. The accepted proof then evaluates the defined implemented
\(132\)-row formula residual at the lifted Gauss stage, after the local chart,
row ordering, and row scaling have been fixed. The stage-residual bound
\(\|F_{A,h}(Z_G;y)\|\le C_Rh^7\) is supplied by
Lemma~\ref{lem:full-132-row-residual-bridge}: the 96 non-dynamic rows
contribute the \(O(h^7)\) block and the 36 Newton--Euler rows vanish by the D5
direct-substitution certificate. This is a discharged formula-residual input
to the perturbation criterion, not a compact-tube regularity assumption and
not an empirical residual fit from the reported trajectory.
This is also the precise point at which the Wieloch--Arnold constrained-BDF
proof architecture~\cite{wieloch2021bdf} is invoked as an
organizational analogue.  The transfer is a proof-order transfer, not an
estimate transfer: first the
discrete Lie-group object is fixed, then the smooth branch is inserted into
that fixed object, then a local defect is proved, and only after that is any
global error recursion invoked.  Lemma~\ref{lem:taylor-route-separation}
identifies the active Taylor object as the full scaled map \(F_{A,h}\) in the
stage increment \(Z-Z_G\).  The proof therefore does not import the reference
paper's BLieDF hidden-constraint estimates, BDF \(p=k\) theorem, or any
primitive FullVA subterm bound; those would be different premises from the
132-row residual bridge used here.  The only structural correspondence is the
obligation order: fixed discrete object, smooth-branch insertion,
local-defect estimate, and then local-to-global propagation.
Concretely, the cited constrained-BDF analogue is the Section~6.2
Theorems~2--4 dependency order: Theorem~2 establishes local truncation
estimates, Theorem~3 forms the configuration/velocity global-error recursion,
and Theorem~4 uses hidden-constraint difference approximations and a
multiplier-error estimate to close the coupled recursion and final order
statement after those local estimates are available.  The present
theorem follows the same dependency order at the level of proof obligations:
the same-tuple \(132\)-row residual value is certified before the
Taylor/Kantorovich stage-root perturbation is invoked, and only then is the
compact-tube stability/Gronwall transfer applied to the reported grid.
In the reference proof, the local truncation error is generated by inserting
the analytical solution into the already specified BLieDF equations.  The
FullVA analogue used here is the single insertion of the lifted Gauss branch
into the already specified residual \(F_{A,h}\).  Therefore all later Taylor
and Kantorovich constants are constants of this fixed residual object; they
are not constants of a reference BDF equation, a primitive D5 factorization,
or a table of finite residual diagnostics.
\noindent\textbf{Reference-style insertion checkpoint.}
The constrained-BDF proof order fixes the discrete equations first, forms
local truncation errors by substituting the exact constrained solution into
those equations, and only then derives hidden-constraint, multiplier, and
global recursion estimates.  The corresponding FullVA checkpoint is the
one-way chain in Eq.~\eqref{eq:reference-style-insertion-chain}:
\begin{equation}
\label{eq:reference-style-insertion-chain}
\begin{aligned}
  Z_G &\xrightarrow{\text{insert}} F_{A,h}(Z_G;y),\\
  F_{A,h}(Z_G;y) &\xrightarrow{\text{direct bridge}} O(h^7),\\
  O(h^7) &\xrightarrow{\text{Taylor/Kantorovich}}
    \|Z_A-Z_G\|=O(h^7)
    \xrightarrow{\mathcal P_h} O(h^7).
\end{aligned}
\end{equation}
Every arrow in this checkpoint is an assertion about the implemented
Gauss6/FullVA residual and endpoint map fixed before the proof begins.  No
BLieDF hidden-constraint difference estimate, multiplier recursion,
BDF zero-stability norm, or primitive D5 decomposition is a premise of any
arrow in this chain, and no later grid error, finite residual row, or
source-policy row can be read backward to create such a premise.
\noindent\textbf{Proof-order instantiation.}
For this theorem the reference ordering is instantiated in the following
forward-only sequence.
\begin{enumerate}
  \item Fix the discrete object: the accepted Gauss-predictor branch,
  \(F_{A,h}\), row ordering, row scaling, proof norm, and endpoint-output map
  are fixed before any estimate or diagnostic row is used.  The reference
  correspondence table is an exposition map for this order of
  obligations; it is not itself a mathematical input.
  \item Prove the local defect on that object: the Gauss \(C_Gh^7\) defect is
  combined with Lemmas~\ref{lem:full-132-row-residual-bridge} and
  \ref{lem:stage-residual-defect}; the separate primitive 162-term Taylor
  closure is not used.
  \item Apply the constrained and solver interfaces: endpoint closure is
  handled by Lemma~\ref{lem:endpoint-closure}, while P6 remains the retained
  compact-tube condition \(\eta_h^{\rm tube}\le c_\eta h^7\). Fixed-tolerance
  logs and finite solver probes do not discharge P6.
  \item Propagate and report only after the local defect is fixed:
  Lemma~\ref{lem:local-global-transfer} and
  Lemma~\ref{lem:reporting-map-norm-consequence} give the grid and
  position--velocity estimates. P7 residual rows remain nonpromoted
  diagnostics.
\end{enumerate}
Every constant used after this ordering is supplied by
Assumption~\ref{ass:regularity} or by the displayed lemmas on the fixed
FullVA object.  No constant, Taylor order, or residual estimate is imported
from the reference paper, from Table~\ref{tab:reference-proof-order-correspondence},
or from a finite diagnostic trajectory.

\noindent\textbf{Direct-route sufficiency.}
From this point to the Gronwall step the proof uses exactly four positive
inputs: the Gauss local defect \(C_Gh^7\), the direct \(132\)-row residual
defect \(C_Rh^7\) assembled from the 96 non-dynamic rows plus the 36 P5
Newton--Euler rows, the endpoint-closure perturbation \(C_Eh^7\), and the
inexact-Newton perturbation \(C_N\eta_h^{\rm tube}\) with the retained P6
bound.  It uses no primitive \(162\)-term Taylor closure, no P7
residual-to-error promotion, no source-policy row, and no fixed-tolerance log
as a theorem input.  Thus the PC2 discharge below is a direct-route
sufficiency statement for the implemented \(132\)-row residual,
not a closure of the separate primitive/Taylor route or a promotion of
diagnostic residual rows.

\noindent\textbf{Route-exclusivity refinement.}
These four inputs are consumed as one same-branch proof tuple after the
compact tube, branch rule, row scaling, proof norm, and endpoint-output map
have been fixed.  If a separate T3 primitive/Taylor proof replaces the
dynamic-row residual-value certificate, it must first produce its own
\(C_R^{\mathrm{T3}}h^7\) bound for the same \(132\)-row residual map on that
same tuple.  It cannot be combined with the accepted direct \(C_Rh^7\) term to
lower constants, remove the retained P6 solver-scale condition, use P7
residual-to-error rows, or turn primitive-route status tables into theorem
inputs.  Thus the active theorem consumes one residual-value certificate,
then one T2 full-map Taylor/Kantorovich perturbation; it never mixes partial
certificates from the direct and primitive routes.

\noindent\textbf{Same-object direct-substitution calculation.}
For this theorem invocation, write the row-scaled residual at the lifted
Gauss stage as the fixed block vector in
Eq.~\eqref{eq:direct-residual-block-vector}:
\begin{equation}
\label{eq:direct-residual-block-vector}
  R_h=F_{A,h}(Z_G;y)
  =\big(F^{\rm nd}_{A,h}(Z_G;y),F^{\rm dyn}_{A,h}(Z_G;y)\big),
\end{equation}
where the first block contains the 96 non-dynamic rows and the second block
contains the 36 Newton--Euler rows in the row ordering fixed by P3.
For a nonautonomous transition this is shorthand for
\begin{equation}
\label{eq:nonautonomous-residual-slice}
  R_{h,n}=F_{A,h,t_n}(Z_G(y,t_n,h);y),
\end{equation}
where both row blocks are evaluated on that same time-indexed residual
slice.  Equation~\eqref{eq:nonautonomous-residual-slice} fixes the
time-indexed residual slice used for the direct block calculation.
The constants used below are uniform over the reported accepted time slices;
no non-dynamic row, Newton--Euler row, Jacobian, or endpoint map is imported
from a different \(t_n\) or from a terminal-output reconstruction.  The P4
certificate gives
\(\|F^{\rm nd}_{A,h}(Z_G;y)\|\le C_{\rm nd}h^7\), and the P5 direct
substitution identity gives
\begin{equation}
\label{eq:direct-dynamic-block-zero}
  F^{\rm dyn}_{A,h}(Z_G;y)=0 .
\end{equation}
Hence, for the fixed finite-dimensional block norm,
\begin{equation}
\label{eq:direct-residual-block-bound}
  \|R_h\|
  \le C_{\rm blk}\|F^{\rm nd}_{A,h}(Z_G;y)\|
  \le C_R h^7 ,
  \qquad C_R=C_{\rm blk}C_{\rm nd}.
\end{equation}
Equations~\eqref{eq:direct-dynamic-block-zero}
and~\eqref{eq:direct-residual-block-bound} are the direct-route residual
certificate consumed by Lemma~\ref{lem:stage-residual-defect}.
The phrase fixed block vector denotes this residual vector with this row
order, scaling, norm, branch, and stage variable.
This is the exact theorem-level residual value inserted into the
Taylor/Kantorovich perturbation, not a primitive \(162\)-term estimate and
not a residual measured at \(\widetilde Z_A\).  With
\(J_h=D_ZF_{A,h}(Z_G;y)\), the stage-root equation in
Lemma~\ref{lem:stage-residual-defect} is
\begin{equation}
\label{eq:stage-root-equation}
  \Delta=-J_h^{-1}R_h-J_h^{-1}Q_h(\Delta),
  \qquad \Delta=Z-Z_G .
\end{equation}
Here \(Q_h\) is the same full-map integral Taylor remainder defined in
Eq.~\eqref{eq:stage-pert-remainder} and bounded there by
\(\|Q_h(\Delta)\|\le (L/2)\|\Delta\|^2\).  The constants
\(M,L,r\) are therefore the retained compact-tube constants of the frozen
\(132\)-row residual instance, not constants recomputed from the direct
dynamic certificate, the primitive inventory, or the reported trajectory.
On the retained branch, the same compact-tube inverse and quadratic-remainder
constants give the accepted fixed point \(\Delta_A=Z_A-Z_G\) with
\begin{equation}
\label{eq:stage-root-perturbation-bound}
  \|Z_A-Z_G\|=\|\Delta_A\|\le 2MC_Rh^7=C_Zh^7 .
\end{equation}
The endpoint reconstruction Lipschitz bound in the same chart then gives
\begin{equation}
\label{eq:endpoint-stage-perturbation-bound}
  \|\Psi_h^A(y)-\Psi_h^G(y)\|
  =\|\mathcal E_h(Z_A)-\mathcal E_h(Z_G)\|
  \le M_EC_Zh^7=C_Ah^7 .
\end{equation}
Equations~\eqref{eq:stage-root-equation},
\eqref{eq:stage-root-perturbation-bound},
and~\eqref{eq:endpoint-stage-perturbation-bound} are the same-object
Taylor/Kantorovich and endpoint-reconstruction handoff.
Thus the direct P4/P5 block substitution and the T2 full-map
Taylor/Kantorovich step are one same-object calculation: the block residual
value, the inverse \(J_h^{-1}\), the quadratic remainder \(Q_h\), and the
endpoint map \(\mathcal E_h\) all use the same branch, row scaling, proof
norm, and stage variable.  No constant in this calculation is imported from
the reference proof, fitted from a finite trajectory, or borrowed from the
separate primitive/Taylor route, summed with that route, or used to partially
close that route; the non-active route status remains
Eq.~\eqref{eq:t3-nonclosure-status-card}.
In the present proof the dynamic-block certificate slot is occupied by
\begin{equation}
\label{eq:direct-dynamic-certificate}
  \mathsf{D}_{\rm dir}:\qquad
  F^{\rm dyn}_{A,h}(Z_G;y)=0 .
\end{equation}
A separate primitive certificate would have the template
\begin{equation}
\label{eq:t3-dynamic-certificate-template}
  \mathsf{D}_{\rm T3}:\qquad
  (P_{\mathrm{state}}\wedge P_{\mathrm{acc}}\wedge P_{\lambda}
   \wedge P_{\mathrm{geom}}\wedge P_{\mathrm{gyro}})
  \Longrightarrow
  \|F^{\rm dyn}_{A,h}(Z_G;y)\|
  \le C_{\rm dyn}^{\rm T3}h^7 .
\end{equation}
Equations~\eqref{eq:direct-dynamic-certificate}
and~\eqref{eq:t3-dynamic-certificate-template} record the accepted direct
certificate and the non-active primitive/Taylor replacement template.
The latter would be an alternative input to the same full-residual bound, not
an additional summand.  No instance of \(\mathsf{D}_{\rm T3}\) is available
in this manuscript because those five antecedent primitives remain open and
the primitive Taylor-bound certificate is not established on the theorem
branch.  The theorem consumes
\(\mathsf{D}_{\rm dir}\) as the proved direct certificate, does not assume
\(\mathsf{D}_{\rm T3}\), and never combines the two certificates.

Lemma~\ref{lem:stage-residual-defect} converts that same-object direct residual defect
into the explicit stage-root and endpoint perturbation bound
\(\|\Psi_h^A(y)-\Psi_h^G(y)\|\le C_Ah^7\).
Lemma~\ref{lem:endpoint-closure} gives the endpoint-closure contribution under
its right-invertibility and raw-defect hypotheses, with the explicit constant
\(C_E=4M_{E,\mathrm{ri}}C_{E,\mathrm{raw}}\), and
Lemma~\ref{lem:inexact-newton} keeps the inexact Newton contribution at the
same scale under the compact-tube branch-solver hypothesis
\(\eta_h^{\rm tube}\le c_\eta h^7\).
More explicitly, let $\Psi_h^G$ be the endpoint map reconstructed from the
reduced Gauss stage, let
\(\Psi_h^A=\mathcal E_h(Z_A)\) be the exact branch-selected accepted
local stage-root map after Lie endpoint reconstruction and before endpoint
closure, let \(\Psi_h^E=\mathcal C_h(\Psi_h^A)\) be the endpoint-closed
exact-stage map, and let
\(\Psi_h^{\mathrm{G6FVA}}=\mathcal P_h(\widetilde Z_A)\) with
\(\mathcal P_h=\mathcal C_h\circ\mathcal E_h\) be the endpoint-closed inexact
Newton map. The
preceding paragraph and lemmas give constants \(C_G,C_A,C_E,C_N\),
independent of \(h\), \(y\), and the reported accepted step index
\(0\le n<N_h\), such that,
uniformly on the
compact proof tube,
\begin{equation}
\label{eq:four-local-bounds}
\begin{gathered}
  \|\Psi_h^G(y)-\varphi_h(y)\|\le C_G h^7,\quad
  \|\Psi_h^A(y)-\Psi_h^G(y)\|\le C_A h^7,
\\
  \|\Psi_h^E(y)-\Psi_h^A(y)\|\le C_E h^7,\quad
  \|\Psi_h^{\mathrm{G6FVA}}(y)-\Psi_h^E(y)\|\le C_N\eta_h^{\rm tube}.
\end{gathered}
\end{equation}
Equation~\eqref{eq:four-local-bounds} lists the four same-branch local
perturbation inputs before they are summed.
The constants in these four bounds are uniform in \(y\) on the compact proof
tube and uniform in the reported accepted step index \(0\le n<N_h\); the
nonlinear-solver term is controlled through the compact-tube branch-solver
policy, whose reported-grid specialization is the single accepted-branch
maximum \(\eta_h^{\rm reported}=\max_{0\le n<N_h}\eta_{h,n}\) over the reported
branch-selected Newton steps.
Their allowed dependencies are exactly the compact-tube regularity,
right-inverse, endpoint-closure, stability, and branch-solver scale bounds
listed in Assumption~\ref{ass:regularity}; no constant in the local defect
sum is chosen from the reported grid length, backend, or fixed production
tolerance.
The single small-step threshold used in the theorem is the minimum of the
windows required by the Gauss truncation estimate, the residual-bridge
certificate, the Kantorovich stage-root perturbation, the endpoint-closure
right-inverse estimate, the inexact-Newton branch-ball estimate, and the
local-to-global tube-retention bootstrap. These windows are all chosen after
the compact proof tube, proof norm, row scaling, branch rule, and constants in
Assumption~\ref{ass:regularity} have been fixed, and before any reported grid,
solver log, separate D5 primitive/Taylor inventory, P7 residual table, or
source-policy diagnostic is read. Hence the four local \(O(h^7)\) estimates hold
simultaneously on the same accepted branch and in the same norm; the proof
does not splice estimates that require incompatible step-size restrictions or
different nonlinear roots.

\noindent\textbf{Uniform-constant timeline.}
Equivalently, the theorem has a single quantifier order: first fix the
compact tube, chart, row-scaled residual map, proof norm, endpoint map,
branch rule, and solver-envelope convention; then choose the constants in the
four local estimates and the one-step stability estimate; then shrink one
common \(h_0\); only after that may a particular \(h\le h_0\), grid length
\(N_h\), and reported trajectory be considered.  The constants
\(C_G,C_R,M,L,r,C_A,C_E,C_N,c_\eta,C_s\) are therefore uniform in the
reported step index and independent of the finite grid selected afterward.
This is the one-step analogue of the reference constrained-BDF proof step
where constants in the error recursion are chosen independently of \(h\) and
\(n\) after the local truncation equations are fixed.  Rechoosing any of
these constants per reported step, per backend norm, or after seeing a
finite diagnostic would define a different theorem instance rather than
strengthen this one.

When \(\eta_h^{\rm tube}\le c_\eta h^7\), define the explicit local-defect constant
\(C_{\rm loc}=C_G+C_A+C_E+C_Nc_\eta\).
The four-term local defect decomposition is therefore
\begin{equation}
\label{eq:four-term-local-defect-sum}
\begin{aligned}
  \|\Psi_h^{\mathrm{G6FVA}}(y)-\varphi_h(y)\|
  &\le
  \|\Psi_h^G(y)-\varphi_h(y)\|
  +\|\Psi_h^A(y)-\Psi_h^G(y)\|  \\
  &\quad
  +\|\Psi_h^E(y)-\Psi_h^A(y)\|
  +\|\Psi_h^{\mathrm{G6FVA}}(y)-\Psi_h^E(y)\|  \\
  &\le C_Gh^7+C_Ah^7+C_Eh^7+C_N\eta_h^{\rm tube}  \\
  &\le C_{\rm loc}h^7 .
\end{aligned}
\end{equation}
Equation~\eqref{eq:four-term-local-defect-sum} is the triangle inequality over Gauss truncation, stage-root
perturbation, endpoint closure, and inexact Newton, with an explicit uniform
constant sum after the compact-tube \(\eta_h^{\rm tube}\) bound. It gives
the local estimate stated in Eq.~\eqref{eq:g6fva-local-defect-theorem}.
This displayed budget is also the only point where P6 enters the positive
order proof: P6 controls the fourth summand, while the direct residual bridge
and the full-map Taylor/Kantorovich step control the stage-root summand.
Equivalently, the proof has instantiated
Lemma~\ref{lem:route-local-defect-sufficiency} on the current transition:
the local-defect obligations supply the same-branch input consumed here, and
the remaining proof uses only the stability, tube-retention, and reporting-map
steps below. No later global estimate, finite run, or diagnostic row feeds back
into \(C_{\rm loc}\).
\noindent\textbf{Forward-only Taylor invariant.}
No \(h^6\) reported-grid estimate, observed slope, residual table, or finite Newton
log is used to define \(Z_G\), choose \(F_{A,h}\), select the row scaling or
proof norm, prove the D5 direct identity, or calibrate
\(C_R,M,L,r,C_A,C_E,C_N\).  The allowed proof order is fixed object and
smooth-branch insertion, then direct residual value at \(Z_G\), then the
Taylor/Kantorovich stage-root perturbation, then endpoint and inexact-Newton
perturbations, and only then Gronwall and reporting.  Thus every Taylor
estimate in the accepted route is a pre-global estimate on the fixed
implemented equations; a reverse reading from reported errors or finite
diagnostics back to the Taylor inputs would be circular and is excluded from
the theorem.

The forward-only exclusion used in the next step is
Lemma~\ref{lem:no-reverse-taylor-inference}: neither the reported-grid estimate nor any
finite diagnostic is used backward to certify primitive/Taylor inputs, P6, or
P7.

Proof structure: the load-bearing proof is the
four-term uniform local-defect sum followed by the stated local-to-global
transfer and reporting-map consequence. The proof-interface tables specify
domain and nonuse boundaries; they do not replace the local-defect estimate,
choose
constants from finite \(h\)-sweeps, or promote residual tables into
trajectory-error evidence.
The proof direction is therefore forward only.  The lifted Gauss stage,
implemented residual identity, endpoint correction, and solver-envelope constant
are fixed before the recurrence for \(e_n\) is written; the \(O(h^6)\)
reported-grid estimate is not used to choose them.  The only loss from the local \(h^7\)
defect to the same-initial-state reported-grid \(h^6\) estimate is the explicit Gronwall
accumulation over at most \(T/h\) accepted transitions.
The one-step stability scale in Assumption~\ref{ass:regularity} supplies the
\((1+C_s h)\) hypothesis of Lemma~\ref{lem:local-global-transfer}, and the
accepted compact tube supplies the tube-retention bootstrap used there.
To make the transfer explicit, define the reduced-chart grid error on the
reported indices by
\begin{equation}
\label{eq:reduced-grid-error-definition}
  \varepsilon_n=y_n-y(t_n),\qquad e_n=\|\varepsilon_n\| .
\end{equation}
Equation~\eqref{eq:reduced-grid-error-definition} defines the only propagated
state-error vector and scalar error norm in the theorem. For this theorem
invocation the first-exit index is the first reported index at which either
the numerical/exact states leave the retained compact tube or the next
accepted transition is no longer defined in \(\mathcal D_h\). The recurrence
below is used only before that combined compact-tube/domain exit. Here
\(\mathcal D_h\) denotes the state projection of the predeclared accepted-domain
graph carrying the retained branch-ball, endpoint-ball, proof-norm, and P6
solver-envelope requirements; membership in those non-compact-tube fiber
components is not inferred from the Gronwall recursion. The
small-step first-exit bootstrap closes the compact-tube component after the
constants are fixed; the branch-ball, endpoint-ball, proof-norm, and P6
solver-envelope components are retained admissibility conditions on the same
reported transition set and are not proved by the Gronwall recursion. Thus the
proof never extends
\(\Psi_h^{\mathrm{G6FVA}}\) outside its accepted theorem domain.
\noindent\textbf{Theorem-domain restriction.}
Domain-exit is not a lower-order conclusion: a branch-ball, endpoint-ball,
compact-tube, or P6 solver-envelope failure does not create a covered
lower-order transition. It means the stated theorem instance is undefined for
that transition until a different theorem domain is proved. The first-exit
argument is therefore a fixed-domain bootstrap for the compact-tube component:
it uses the retained tube margin and stability scale to keep the accepted
reported states inside the compact tube for \(0<h\le h_0\).  The other
components of \(\mathcal D_h\)--branch-ball membership, endpoint-ball
membership, proof-norm binding, and the P6 solver envelope--remain
predeclared transition-domain conditions on the same accepted transition set.
They are not proved by the first-exit bootstrap, and they are not recovered by
a post-run row deletion rule that would discard failed steps after inspecting a
trajectory.
On the accepted branch, after the branch-selection and tube-retention
hypotheses have been fixed, the local defect and stability estimates give the
error recurrence
\begin{equation}
\label{eq:error-recurrence}
  e_{n+1}\le (1+C_s h)e_n+C_{\rm loc}h^7,\qquad e_0=0,
  \qquad 0\le n<N_h .
\end{equation}
Equation~\eqref{eq:error-recurrence} is the scalar Gronwall recurrence
for the reduced-chart state-error norm.
Its forcing term is only the already proved same-branch \(C_{\rm loc}h^7\)
local defect from Eq.~\eqref{eq:four-term-local-defect-sum};
raw residual tables, P6 solver logs, and P7 mechanism residual diagnostics are
not additional right-hand-side terms in this recurrence.
Because \(e_n\) is the reduced-chart state-error norm and not a residual norm, this
recurrence is not the P7 residual-to-error boundary template in
Eq.~\eqref{eq:p7-residual-transfer-template}.
\noindent\textbf{Reference-style recursion analogue.}
This is the point in the present proof corresponding to the global recursion
step in the constrained Lie-group BDF reference.  The analogy is structural,
not an estimate transfer, and it matches only the proof order in which a
proved local defect is propagated.  It does not transfer the BLieDF propagated
variable set, companion matrices, multiplier recursion, or recursion constants.
In the reference proof, Lie-algebra configuration
errors, constraint differences, and multiplier estimates are propagated
together only after the local truncation defect has been proved for the fixed
BLieDF equations.  Here the lower-pair algebraic variables and stage
multipliers have already been consumed inside the fixed \(132\)-row
FullVA residual, the endpoint right-inverse, and the accepted branch map
\(\Psi_h^{\mathrm{G6FVA}}\).  After those same-object inputs have produced
the local defect \(C_{\rm loc}h^7\), the only quantity propagated by the
theorem is the reduced-chart state-error norm \(e_n\) through the displayed
one-step stability inequality.  There is no hidden multiplier recursion in
this step, no reaction-output norm inside \(e_n\), and no residual-to-error
operator being applied tacitly.  A multiplier or reaction error theorem would
therefore require a separate coupled output recursion with its own
branch-compatible constants; it is not a consequence of the scalar/vector
Gronwall inequality above.
Zero-initial-error rule: the displayed recurrence is the theorem instance
with the numerical and exact trajectories started from the same reduced-chart
initial state.  If an initial projection, chart-registration step,
warm-start correction, or branch-selection adjustment introduced
\(e_0\ne0\), the same Gronwall calculation would contain the additional term
\(\exp(C_sT)\|e_0\|\).  That term is not absorbed into
\(C_{\rm loc}h^7\) or \(C_{\rm red}h^6\).  Thus the theorem's order-six grid
estimate is a same-initial-state statement; any reported run with a separate
initialization error would need its own initial-error bound before invoking
the theorem.
One-step start-up boundary: unlike the constrained BLieDF reference proof,
there is no multistep start-up layer such as \(0\le n<2k-1\) in the theorem
instance above.  The accepted map is a one-step map, so the recurrence starts
at \(n=0\) from the exact reduced initial state and uses only the single
previous accepted state.  Any numerical start-up perturbation, projection
warm start, or branch-registration mismatch is therefore represented by the
initial-error term described above; it is not hidden in the local defect,
not treated as a separate finite prefix with its own constants, and not
estimated from the first few reported steps.
The induction hypothesis is that the numerical and exact states remain in the
retained compact tube for the reported transition under consideration; this is
the tube-retention induction supplied by Assumption~\ref{ass:regularity}, not
a condition read from trajectory plots after the run. In this theorem
invocation the induction also includes \(y_n\in\mathcal D_h\) for the
transition being estimated and \(y(t_n)\in\mathcal D_h\) for the exact-state
local-defect evaluation, so a branch, endpoint-ball, or solver-scale failure
is a theorem-domain exit rather than a lower-order case.
Equivalently, the theorem uses the first-exit bootstrap in
Lemma~\ref{lem:local-global-transfer}: the recurrence is applied only until a
hypothetical exit from the retained tube, and the small-step threshold is
chosen so that the resulting \(C_{\rm red}h^6\) bound stays inside the tube
margin while the retained branch/domain conditions remain in force. The
first-exit bootstrap closes the compact-tube component under the retained P2
tube-margin and stability hypothesis; branch-ball, endpoint-ball, proof-norm,
and P6 solver-envelope membership remain theorem-domain conditions for the
accepted transition set. None of these conditions is assumed after inspecting
the reported trajectory. Iterating the recurrence
before any reported residual or work-precision table is inspected gives
\begin{equation}
\label{eq:gronwall-sum-bound}
  e_n\le C_{\rm loc}h^7
  \sum_{j=0}^{n-1}(1+C_s h)^j
  \le C_{\rm loc}h^6\Gamma_s(T),
\end{equation}
where the Gronwall factor is
\begin{equation}
\label{eq:gronwall-factor}
  \Gamma_s(T)=
  \begin{cases}
    (\exp(C_sT)-1)/C_s, & C_s>0,\\
    T, & C_s=0 .
  \end{cases}
\end{equation}
Equations~\eqref{eq:gronwall-sum-bound}
and~\eqref{eq:gronwall-factor} are the only source of the one-power loss
from local \(h^7\) to grid \(h^6\).
Accumulation-factor boundary: the loss of one power of \(h\)
comes only from accumulating at most \(O(1/h)\) reported transition defects
over a fixed horizon. The bound uses \(N_hh\le T\) before the constant is
chosen, so the geometric sum is bounded by \(h^{-1}\Gamma_s(T)\) and the
constant \(C_{\rm red}\) is independent of the reported step count \(N_h\).
It is not fitted from the number of reported steps, an observed
work-precision slope, or a residual table. Likewise, \(C_s\), the tube
margin, \(\Gamma_s(T)\), and \(C_{\rm red}\) are inherited from
Assumption~\ref{ass:regularity} and
Lemma~\ref{lem:local-global-transfer} before reported grid errors are read.
The count \(N_h\) is used only through \(N_hh\le T\); it cannot fit a
stability constant or repair a \mbox{failed tube-retention row}.
Applying that lemma with \(p=6\), \(C_\ell=C_{\rm loc}\), and the reported
grid length denoted by \(N_h\) gives the explicit reduced-chart grid constant
\begin{equation}
\label{eq:reduced-grid-constant}
  C_{\rm red}=C_{\rm loc}\Gamma_s(T).
\end{equation}
Writing \(d_K\) for the retained distance of the exact reduced trajectory to
the compact-tube boundary in Lemma~\ref{lem:local-global-transfer}, the
first-exit closure uses one further pre-data restriction of the common
small-step threshold: after \(C_{\rm red}\) and \(d_K\) are fixed, but before
any reported grid is read, shrink \(h_0\) so that
\(C_{\rm red}h_0^6<d_K/2\).  Then for every \(0<h\le h_0\) the bound obtained
on the truncated first-exit interval keeps the numerical state strictly
inside the retained compact tube, contradicting any compact-tube first exit.
This is the only place where the tube margin enters the grid-error proof; it
does not repair branch-ball, endpoint-ball, proof-norm, or P6
solver-envelope failures.  Equation~\eqref{eq:reduced-grid-constant} is the
constant used in
the reduced-grid bound in Eq.~\eqref{eq:g6fva-reduced-grid-bound}
over the same reported time grid \(t_n=n h\), \(0\le t_n\le T\), with the
maximum taken over \(0\le n\le N_h\); exact final-time divisibility is not a
separate theorem condition. The terminal index \(N_h\) is only the last output
value; the branch-selected Newton solves and the solver maximum above are
indexed by the transition steps \(0\le n<N_h\). If \(N_hh<T\), the
unreported leftover interval \(T-N_hh\) is outside the stated grid maximum;
no terminal-time error at \(T\) is inferred without adding a separate
reported step. The final passage to
the reported position and velocity errors is the application of
Lemma~\ref{lem:reporting-map-norm-consequence}; it uses compact chart
norm-equivalence, not a new dynamic estimate. On the retained compact tube, the
chart-to-\((q,v)\) reporting map
\begin{equation}
\label{eq:reporting-map-definition}
  \mathcal R(y)=(q(y),v(y))
\end{equation}
is \(C^1\) with uniformly bounded derivative, say
\(\|D\mathcal R\|\le C_{\mathcal R}\). Define
\(C_{qv}=C_{\mathcal R}C_{\rm red}\). Hence
\begin{equation}
\label{eq:reporting-map-bound-proof}
  \max_{0\le n\le N_h}\|\mathcal R(y_n)-\mathcal R(y(t_n))\|
  \le C_{\mathcal R}\max_{0\le n\le N_h}\|y_n-y(t_n)\|
  \le C_{qv}h^6.
\end{equation}
Equations~\eqref{eq:reporting-map-definition}
and~\eqref{eq:reporting-map-bound-proof} define the reported output map and
its compact-chart norm consequence.
This step is local to the accepted chart and does not claim a global
equivalence across remote charts, nonsmooth contact events, or other branches
outside the compact proof tube; for rotational components it is the fixed
local retraction/log-chart reporting convention above, not a global
\(SO(3)\) geodesic-error theorem.
\end{proof}

\noindent\textbf{Scope of the theorem conclusion.}
The theorem conclusion is limited to the displayed local-defect and
same-initial-state reported-grid bounds for the accepted reduced-chart
branch-selected map.
In particular, residual tables are recorded only as diagnostics for the open
P7 residual-to-error boundary; they do not
prove P7, do not alter \(C_{\rm loc}\), and do not promote
\mbox{mechanism-coverage residual rows} to \mbox{trajectory-error/order evidence}. The
theorem \mbox{does not prove a P7 residual-to-error boundary}; the
corresponding exclusion rule is explicit: do not promote
mechanism-coverage residual rows to trajectory-error/order evidence unless a
separate residual-to-error theorem and branch-compatible dynamic-order
campaign have been supplied.
This P7 boundary is independent of the P6 solver-scale boundary: P6 controls
only the inexact Newton summand, while P7 would require a separate
residual-to-trajectory transfer operator.
The solver-scale clause is a retained hypothesis on the selected Gauss-predictor
branch; for a reported trajectory it applies only to transitions whose
selected Newton iterate remains on that branch and whose reported-grid maximum
\(\eta_h^{\rm reported}=\max_{0\le n<N_h}\eta_{h,n}\) satisfies
\(\eta_h^{\rm reported}\le c_\eta h^7\) as a trajectory specialization of
the retained compact-tube envelope \(\eta_h^{\rm tube}\). Equivalently,
\(\eta_{h,n}\le c_\eta h^7\) for every accepted reported transition. The
theorem does not choose \(c_\eta\), the residual norm, or the accepted branch
from reported residual ratios.
If the implementation stops by a fixed production tolerance or a backend
residual norm, that log specializes this clause only after the
same-branch proof-norm conversion in
Eq.~\eqref{eq:p6-backend-to-proof-norm-conversion} and the
predeclared \(c_\eta h^7\) envelope have been fixed; otherwise the log
remains diagnostic and is not a P6 theorem input.
In particular, the branch-ball membership certificate and the \(h^7\)
residual envelope are checked on the same transition; neither one is inferred
from the other.
Those objects are fixed before the grid, Newton logs, and tolerance sweeps;
if a computed transition fails the branch or scaled-residual condition, the
branch-selected map is undefined there rather than lower order.
Residual-only mechanism rows, fixed-tolerance logs, and comparison or
source-package records remain outside the local-defect estimate and outside
the theorem conclusion.

The stage-residual input is the implemented 132-row direct route. The
alternative primitive/Taylor decomposition of the Newton--Euler rows is a
separate dynamic-block residual-certificate route, not a premise of this
theorem. Its five uniform primitive inputs
\(P_{\mathrm{state}},P_{\mathrm{acc}},P_{\lambda},
P_{\mathrm{geom}},P_{\mathrm{gyro}}\) are not established on the accepted
theorem branch.  This is a scope boundary for the separate primitive route, as recorded in
Tables~\ref{tab:d5-primitive-blocker-ledger}
and~\ref{tab:d5-open-primitive-proof-interfaces} and the accompanying
conditional primitive/Taylor implication. The primitive/Taylor route is a
non-active replacement-certificate route rather than current theorem evidence;
more precisely, it is a separate dynamic-block residual-certificate route. It would become
a separate sufficient primitive-route implication for the dynamic-row PC2
input only after those five primitives are independently proved on the same
compact branch, proof norm, row scaling, and endpoint/output convention, and
the resulting dynamic-row value is reassembled with the same 96-row
non-dynamic certificate before being consumed by the unchanged T2 full-map
perturbation chain. It is not
a term in \(C_{\rm loc}\), not an additional
premise of the direct residual-bridge theorem, and the theorem conclusion
does not certify it. Equivalently, the theorem uses the implication from the certified
direct bridge to the full-residual defect; it does not use the converse
implication from theorem order to primitive subterm closure, and it does not
require the non-active separate primitive route to be refuted. Finite solver
probes, fixed-tolerance production logs, residual/reaction tables, and
common-reference or source-policy rows are therefore not post-theorem
discharge mechanisms. They may diagnose branch plausibility or record
diagnostic scope, but they do not prove a solver-policy theorem, do not prove
a P7 residual-to-error boundary, do not complete
the primitive/Taylor route, and do not enlarge the theorem beyond the accepted
branch-selected map stated above.  The same direct residual-bridge scope convention
applies throughout this section: it refers only to the
residual-bridge/Kantorovich input for this conditional theorem. It does not
prove a solver-policy theorem or a P7 residual-to-error boundary, does not
complete the primitive/Taylor route, and does not establish source-policy or
full-\tfe{} readiness.
The conditional primitive/Taylor implication therefore states two distinct mathematical
facts, summarized by the status card
Eq.~\eqref{eq:t3-nonclosure-status-card}.  Under the named primitive inputs,
all \(162/162\) reduction
templates have been assigned to row-level subterms. This is template
coverage only; the certified count of actual primitive Taylor bounds is \(0/162\) without
those inputs, so the certificate does not provide an actual primitive-route
Taylor bound for a T3 replacement certificate in the accepted theorem route.
The first fact shows that the primitive dependency map is complete as a
conditional implication. It is not theorem evidence unless the named
primitives are independently proved on the same compact branch. The second keeps the
separate primitive route out of the accepted theorem until its uniform lift
primitives are proved. These
two statements are not additive proof evidence: the conditional templates
cannot be spliced with the direct residual bridge to strengthen the theorem
conclusion or remove a retained interface.
A primitive-route theorem would have to prove the five open primitive
obligations on the same compact branch---the state and acceleration root
lifts, the multiplier lift propagated from them, and the downstream geometry
and gyroscopic reductions---and then redo the residual-to-root, endpoint,
solver, and global-transfer chain for that route.

\noindent\textbf{Domain of Validity.}
To invoke Theorem~\ref{thm:g6fullva-order}, the proof instance must use the
same reduced initial state, the same accepted compact-tube branch, the same
implemented 132-row residual and proof norm, the same endpoint-output map
\(\mathcal P_h=\mathcal C_h\circ\mathcal E_h\), the scaled branch-solver
condition \(\eta_h^{\rm tube}\le c_\eta h^7\), and the same reported grid
\(t_n=nh\), \(0\le n\le N_h\), with accepted transitions \(0\le n<N_h\)
and with \(y_n\in\mathcal D_h\) together with exact-state local-defect
admissibility, denoted \(y(t_n)\in\mathcal D_h\) in the same admissible-graph
sense, for each transition input.
This common accepted-domain membership is a theorem
hypothesis/bootstrap condition, not a filter chosen after a run.
No leftover terminal-time value, fixed production tolerance, residual-only
mechanism row, source-policy row, or reproducibility package state is part of this
invocation. In particular, no reported fixed-tolerance row instantiates P6
unless its branch membership and proof-norm residual maximum satisfy the
stated \(h^7\) envelope on the same accepted transition set. Failure of any
domain condition excludes that step or grid from the theorem domain; it does not
create an alternative lower-order theorem.

\begin{table}[htbp]
\centering
\caption{Accepted-branch consistency for
Theorem~\ref{thm:g6fullva-order}.}
\label{tab:accepted-branch-consistency-ledger}
\scriptsize
\setlength{\tabcolsep}{2pt}
\renewcommand{\arraystretch}{0.9}
\begin{tabular}{P{0.18\linewidth}P{0.28\linewidth}P{0.27\linewidth}P{0.19\linewidth}}
\toprule
Proof object & Branch identity used & Consistency source & Explicit exclusion \\
\midrule
Gauss predictor &
Same reduced-chart initial state \(y\) and same compact proof tube. &
Assumption~\ref{ass:regularity} and the reduced Gauss collocation map
\(G_h\). &
Remote charts and nonsmooth contact events. \\
Accepted stage root &
Local root selected from the Gauss-predictor ball on the compact tube. &
Compact-tube inverse condition and Lemma~\ref{lem:stage-residual-defect}. &
Arbitrary Newton roots or global nonlinear solver uniqueness. \\
Endpoint correction &
Endpoint closure applied to that accepted local root before the inexact
Newton term. &
Lemma~\ref{lem:endpoint-closure} and its local ball/right-inverse
hypotheses. &
Extra terminal solve or unreported leftover interval. \\
Inexact Newton iterate &
Newton iterate remains on the retained branch with
\(\eta_h^{\rm tube}\le c_\eta h^7\). &
Lemma~\ref{lem:inexact-newton} and retained P6 branch-solver condition. &
Fixed production tolerances as asymptotic proof. \\
Reported output grid &
Reported states \(y_n\) are generated by repeated accepted-branch
transitions on \(0\le n<N_h\). &
Lemma~\ref{lem:local-global-transfer} on the same reported grid. &
Residual/reaction tables, source-policy rows, and terminal-time claims at
\(T\). \\
\bottomrule
\end{tabular}
\end{table}

Accepted-branch consistency rule: every theorem
bound is tied to the same compact-tube branch selected by the Gauss predictor,
the local accepted stage root, endpoint closure, and retained inexact-Newton
solver scale. The branch-consistency statement does not claim remote-chart
equivalence, arbitrary Newton-root selection, global nonlinear solver
uniqueness, fixed-tolerance asymptotic proof, residual-to-error promotion,
source-policy reproduction, or terminal-time error at an unreported \(T\).

\begin{table}[htbp]
\centering
\caption{Nonlinear-solver scale conditions for
Theorem~\ref{thm:g6fullva-order}.}
\label{tab:nonlinear-solver-scale-ledger}
\scriptsize
\setlength{\tabcolsep}{2pt}
\renewcommand{\arraystretch}{0.9}
\begin{tabular}{P{0.18\linewidth}P{0.29\linewidth}P{0.27\linewidth}P{0.18\linewidth}}
\toprule
Solver-scale object & Theorem role & Mathematical condition and evidence boundary & Explicit non-use \\
\midrule
Compact-tube policy &
\(\eta_h^{\rm tube}\le c_\eta h^7\) uniformly on the retained
branch-selected compact tube. &
Retained P6 theorem condition used by Lemma~\ref{lem:inexact-newton}. &
Fixed production tolerances or finite probes as theorem-level proof. \\
Proof-norm binding &
The \(F_{A,h}\)-residual is measured in the fixed proof norm of the accepted
132-row residual after chart and row-scaling conventions are set. &
Implementation logs can be compared with \(\eta_h^{\rm tube}\) only under the
accepted implementation-path binding and compact-tube norm equivalence. &
A second residual budget, residual-log calibration of constants, or proof of
P6 from logged residuals. \\
Stopping proxy &
Any correction-norm, relative-residual, iteration-count, or backend success
criterion used by an implementation. &
It must be converted in advance to the same absolute
\(F_{A,h}\)-residual bound \(c_\eta h^7\) on the accepted branch. &
Proxy convergence flags, machine-precision floors, or adaptive stopping rules
as theorem inputs. \\
Reported-grid specialization &
\(\eta_h^{\rm reported}=\max_{0\le n<N_h}\eta_{h,n}\) over accepted branch-selected Newton
solves. &
Trajectory specialization only after the compact-tube policy, proof norm,
branch rule, and \(c_\eta\) have already been fixed. &
Using the reported maximum itself to prove the compact-tube policy, choose
\(c_\eta\), admit arbitrary Newton initializations, remote roots, or extra
terminal solves. \\
Finite solver diagnostics &
One-step scaled-tolerance probe, short-trajectory scaled-tolerance probe, and
tolerance-regime sweep. &
Diagnose branch-scale consistency and are attached only to P6. &
Closing the compact-tube policy for every retained state or the full reported
campaign. \\
Local-defect absorption &
\(C_N\eta_h^{\rm tube}\le C_Nc_\eta h^7\) enters
\(C_{\rm loc}=C_G+C_A+C_E+C_Nc_\eta\). &
Only after the retained compact-tube solver-scale condition is assumed. &
Fitting \(C_N\), \(c_\eta\), or \(C_{\rm loc}\) from residual tables or
h-sweeps. \\
Retained solver-policy condition &
The theorem is an order-six theorem under the stated compact-tube,
endpoint/stability, implementation-binding, and solver-scale interfaces. &
Recorded separately from direct stage-residual closure. &
Full source-policy archive status, source-policy closure, or residual-to-error
promotion. \\
\bottomrule
\end{tabular}
\end{table}

Nonlinear-solver scale rule: the proof consumes the
compact-tube \(\eta_h^{\rm tube}\le c_\eta h^7\) policy only as a retained P6
theorem condition in the fixed proof residual norm. Implementation logs bind
to that norm only through the accepted implementation path and compact-tube
norm equivalence. The reported-grid maximum \(\eta_h^{\rm reported}\) and the finite
scaled-tolerance probes are branch-scale diagnostics and trajectory
specializations only for a P6 instance that has already fixed the tube,
branch, norm, and \(c_\eta\). They are not a theorem-level solver-policy proof,
a fixed-tolerance asymptotic proof, global Newton convergence, arbitrary
Newton-root selection, a second residual budget, source-policy reproduction,
residual-to-error promotion, or reproducibility package readiness evidence.
Solver-envelope quantifier discipline: the scalar solver-envelope constant and fixed
proof norm are chosen before \(h\), \(N_h\), or finite logs. Changing that
constant after logs changes the theorem hypothesis and local-defect constant,
not a proof of P6.
The accepted branch and the strong-local-inverse ball are fixed at the same
time. Reported residual ratios can diagnose finite compatibility with the
retained P6 scale, but they cannot calibrate \(c_\eta\), select a different
Newton branch, or turn failed branch-scale rows into theorem inputs.

\begin{table}[htbp]
\centering
\caption{Local-defect decomposition for
Theorem~\ref{thm:g6fullva-order}.}
\label{tab:local-defect-decomposition-ledger}
\scriptsize
\setlength{\tabcolsep}{2pt}
\renewcommand{\arraystretch}{0.9}
\begin{tabular}{P{0.18\linewidth}P{0.27\linewidth}P{0.28\linewidth}P{0.19\linewidth}}
\toprule
Triangle term & Bound used & Proof source & Explicit non-use \\
\midrule
Gauss truncation &
\(\|\Psi_h^G(y)-\varphi_h(y)\|\le C_Gh^7\). &
Classical three-stage Gauss collocation on the smooth reduced flow. &
Implemented residual tables or finite-run h-sweeps. \\
Stage-root perturbation &
\(\|\Psi_h^A(y)-\Psi_h^G(y)\|\le C_Ah^7\). &
Lemma~\ref{lem:full-132-row-residual-bridge} supplies the \(C_Rh^7\)
residual from the 96-row non-dynamic certificate plus the 36-row D5
direct-substitution zero certificate; Lemma~\ref{lem:stage-residual-defect}
then uses the compact-tube inverse/Kantorovich step. &
Remote nonlinear roots, primitive dynamic symbolic oracle closure, or source-paper
residual reproduction. \\
Endpoint closure &
\(\|\Psi_h^E(y)-\Psi_h^A(y)\|\le C_Eh^7\). &
Lemma~\ref{lem:endpoint-closure} with local right inverse, endpoint-ball
margin, and raw endpoint defect. &
Global chart equivalence, residual-to-error promotion, or terminal-time
extra solve. \\
Inexact Newton termination &
\(\|\Psi_h^{\mathrm{G6FVA}}(y)-\Psi_h^E(y)\|\le
C_N\eta_h^{\rm tube}\le C_Nc_\eta h^7\). &
Lemma~\ref{lem:inexact-newton} plus the retained compact-tube branch-solver
scale, applied to the full endpoint-output map
\(\mathcal P_h=\mathcal C_h\circ\mathcal E_h\). &
Fixed production tolerances, global Newton convergence, or arbitrary
initial guesses. \\
Local-defect sum &
\(C_{\rm loc}=C_G+C_A+C_E+C_Nc_\eta\). &
Triangle inequality over the four accepted-branch terms above. &
Local-to-global grid transfer, source-policy rows, or residual/reaction
diagnostics. \\
\bottomrule
\end{tabular}
\end{table}

Local-defect decomposition rule: the local \(h^7\)
statement is obtained only by summing the four displayed accepted-branch
terms: Gauss truncation, stage-root perturbation, endpoint closure, and
inexact Newton termination under the compact-tube branch-solver scale. The
local-defect sum does not use finite-run h-sweeps, residual/reaction
tables, source-policy rows, remote nonlinear roots, global Newton
convergence, residual-to-error promotion, or full-\tfe{} readiness.

Norm/row-scaling bridge: the proof residual norm, stage-root
norm, reduced-chart norm, and reported \(q,v\) norm are not interchangeable.
The fixed 132-row residual norm is used only inside the
stage-residual perturbation and inexact-Newton estimates. The compact-tube
uniform inverse maps that residual norm to the local stage-root norm; the
endpoint-output map \(\mathcal P_h=\mathcal C_h\circ\mathcal E_h\) maps stage
perturbations to the reduced-chart one-step norm, including the local
endpoint-closure Lipschitz factor; the local-to-global transfer remains in
that reduced-chart norm; and the reporting map applies the separate
compact-tube factor
\(C_{\mathcal R}\) only after the reduced-chart grid estimate is proved. All
row scalings, chart norms, and norm-equivalence constants are fixed before
\(h_0\) and before \(C_{\rm loc},C_{\rm red},C_{qv}\) are chosen. In
particular, the row scaling is not reweighted with \(h\), \(N_h\), backend,
or residual logs; it cannot move powers of \(h\) between residual rows and
constants. No
residual-log norm, backend norm, mechanism table norm, or fitted
work-precision scale is substituted into the four-term local-defect sum.

\begin{table}[htbp]
\centering
\caption{Theorem output scope for
Theorem~\ref{thm:g6fullva-order}.}
\label{tab:theorem-output-scope-ledger}
\scriptsize
\setlength{\tabcolsep}{2pt}
\renewcommand{\arraystretch}{0.9}
\begin{tabular}{P{0.18\linewidth}P{0.29\linewidth}P{0.27\linewidth}P{0.18\linewidth}}
\toprule
Output layer & Proven statement & Proof source & Explicit non-output \\
\midrule
One-step accepted map &
\(\|\Psi_h^{\mathrm{G6FVA}}(y)-\varphi_h(y)\|\le C_{\rm loc}h^7\). &
Four-term local defect decomposition on the accepted compact-tube branch. &
Residual/reaction tables, source-policy rows, remote nonlinear roots. \\
Same-initial-state reduced grid &
\(\max_{0\le n\le N_h}\|y_n-y(t_n)\|\le C_{\rm red}h^6\). &
Local-to-global transfer on same-initial-state reported output indices with \(N_hh\le T\). &
Unreported terminal-time error, leftover interval, or extra terminal solve. \\
Reported position--velocity map &
\(\max_{0\le n\le N_h}\|\mathcal R(y_n)-\mathcal R(y(t_n))\|\le C_{qv}h^6\). &
The compact chart-to-\((q,v)\) reporting map after the reduced-chart estimate. &
Remote chart equivalence, nonsmooth contact events, or residual-to-error
promotion. \\
Method-order sentence &
Order six in the same-initial-state reported-grid sense under the retained branch,
solver-scale, and compact-tube conditions. &
The one-step, reduced-chart, and reporting-map bounds above. &
External superiority, full source-policy reproduction, full-\tfe{} readiness,
or fixed-tolerance asymptotic proof. \\
\bottomrule
\end{tabular}
\end{table}

Theorem output scope rule: the theorem outputs are
exactly the branch-selected one-step \(h^7\) local defect, the reduced-chart
reported-grid \(h^6\) estimate, the reported position--velocity \(h^6\)
estimate, and the resulting same-initial-state reported-grid order-six sentence under the retained
conditions. The theorem output scope does not include residual/reaction table
promotion, unreported terminal-time error, remote-chart equivalence, external
superiority, full source-policy reproduction, full-\tfe{} readiness, or
fixed-tolerance asymptotic proof.

\noindent\textbf{Constraint-multiplier and reaction-output fence.}
The multiplier step in the Wieloch--Arnold constrained-BDF proof is a proof slot for
its own hidden-constraint and multiplier-error estimate; it is not a theorem
output imported here.  In the present FullVA proof, the stage multipliers
\(\lambda\) are algebraic unknowns used to close the lower-pair and
Newton--Euler rows inside \(F_{A,h}\).  They enter the accepted theorem only
through the \(132\)-row residual bridge, the endpoint-closure lemma, and the
compact-tube right-inverse hypotheses.  The theorem does not assert an
\(h^6\) trajectory estimate for multiplier histories, reconstructed reaction
forces, reaction residuals, or constraint residual tables.  Those quantities
remain mechanism-coverage diagnostics unless a separate multiplier/reaction
reporting theorem supplies its own output map, uniform inverse or inf--sup
constant, branch/solver envelope, reference policy, and local-to-global
recursion.  Consequently, small reaction residuals can support consistency of
the implemented FullVA rows, but they cannot be read as a residual-to-error
promotion or as a multiplier-order claim of
Theorem~\ref{thm:g6fullva-order}.

\begin{table}[htbp]
\centering
\caption{Reporting-map/norm-equivalence table for
Theorem~\ref{thm:g6fullva-order}.}
\label{tab:reporting-map-norm-equivalence-ledger}
\scriptsize
\setlength{\tabcolsep}{2pt}
\renewcommand{\arraystretch}{0.9}
\begin{tabular}{P{0.18\linewidth}P{0.28\linewidth}P{0.27\linewidth}P{0.19\linewidth}}
\toprule
Reporting step & Input estimate & Uniform map property & Explicit non-use \\
\midrule
Reduced-chart grid &
\(\max_{0\le n\le N_h}\|y_n-y(t_n)\|\le C_{\rm red}h^6\). &
Local-to-global transfer on the accepted reduced-chart branch. &
Residual/reaction tables, source-policy rows, or sweep fits. \\
Reporting map definition &
\(\mathcal R(y)=(q(y),v(y))\) on the same compact proof tube. &
The accepted chart-to-\((q,v)\) map is \(C^1\) on that tube. &
Global atlas equivalence or remote chart switching. \\
SO(3) orientation convention &
Orientation or angular-velocity components inside \(q,v\), when reported. &
The same local retraction/log chart and component weights fixed before the theorem. &
Atlas-wide geodesic distance, chart cuts, or post-hoc metric changes. \\
Derivative bound &
\(\|D\mathcal R\|\le C_{\mathcal R}\). &
Compact-tube supremum of the reporting-map derivative. &
Fitted constants, dense/sparse backend behavior, or fixed tolerances. \\
Reported position--velocity estimate &
\(\max_{0\le n\le N_h}\|\mathcal R(y_n)-\mathcal R(y(t_n))\|
\le C_{qv}h^6\). &
\(C_{qv}=C_{\mathcal R}C_{\rm red}\) after the reduced-chart estimate. &
A new dynamic estimate or residual-to-error transfer theorem. \\
Reported grid scope &
The same output indices \(0\le n\le N_h\) and \(t_n=nh\). &
The reporting map is applied only to the proved grid-point estimate. &
Terminal-time claim at an unreported \(T\) or an extra terminal solve. \\
\bottomrule
\end{tabular}
\end{table}

Reporting-map/norm-equivalence rule: the final
reported \(q,v\) estimate consumes only the reduced-chart grid error and the
compact chart-to-\((q,v)\) \(C^1\) derivative bound on the accepted tube. It
is a local reporting-map consequence on the same reported grid, not a
residual-to-error transfer theorem, reaction diagnostic, global chart equivalence,
remote branch statement, source-policy reproduction, fitted-constant claim,
or terminal-time claim at an unreported \(T\).

The reported-grid conclusion is not an interpolation statement. The sequence
\(y_n\) is compared only with \(y(t_n)\) at the reported output indices
\(0\le n\le N_h\) chosen under \(N_hh\le T\). No bound is asserted for
\mbox{off-grid sampling times}, dense output, resampled work-precision curves, or a
\mbox{postprocessed terminal value} unless those times are added as reported grid
points and the same branch/tube hypotheses are verified there. The constants
\(C_{\rm red}\) and \(C_{qv}\) are selected before the reported trajectory is
read and do not depend on output formatting or interpolation policy.

\begin{table}[htbp]
\centering
\caption{Constant-dependency table for Theorem~\ref{thm:g6fullva-order}.}
\label{tab:constant-dependency-ledger}
\scriptsize
\setlength{\tabcolsep}{2pt}
\renewcommand{\arraystretch}{0.9}
\begin{tabular}{P{0.17\linewidth}P{0.29\linewidth}P{0.28\linewidth}P{0.18\linewidth}}
\toprule
Quantity & Proof source & Allowed dependencies & Explicit non-dependencies \\
\midrule
\(h_0\) &
Common small-step threshold for the compact tube, Gauss truncation,
stage-root inverse, endpoint closure, and branch-selected Newton estimates. &
Compact proof tube, regularity margins, local inverse radii, endpoint-ball
margin, and branch-selection radius. &
Reported grid length \(N_h\), dense/sparse backend, source-policy rows. \\
\(C_{\rm loc}\) &
\(C_G+C_A+C_E+C_Nc_\eta\), the four-term one-step defect sum. &
Gauss truncation, direct stage-root perturbation, endpoint closure,
strong local residual inverse, and the retained compact-tube solver scale. &
Finite h-sweep fits, fixed production tolerances, residual/reaction tables. \\
\(C_{\rm red}\) &
\(C_{\rm loc}\Gamma_s(T)\), the reduced-chart grid constant from the
local-to-global transfer lemma. &
The horizon \(T\), the nonnegative stability constant \(C_s\), and
the accepted tube-retention bootstrap. &
Exact final-time divisibility, unreported leftover intervals, external
comparison rows. \\
\(C_{qv}\) &
\(C_{\mathcal R}C_{\rm red}\), the reporting-map constant for
\((q,v)\). &
The \(C^1\) chart-to-\((q,v)\) map on the retained compact proof tube and
the reduced-chart estimate. &
Remote charts, nonsmooth contact events, local reproducibility record state. \\
\bottomrule
\end{tabular}
\end{table}

Constant-dependency rule: the theorem constants are
compact-tube constants only. They are not fitted from the finite h-sweep,
solver tolerance table, source-policy comparison, residual/reaction table, or
local reproducibility record. The table records analytic constant dependencies; it
does not prove a theorem-level \(\eta_h\) solver-policy theorem, a P7
residual-to-error transfer, source-policy readiness, or full-TFE readiness.
\mbox{Constant-dependency rule keyword: compact-tube constants only.}

Constant-selection hierarchy: first fix the mechanism data, the
reported smooth branch, the compact proof tube, the reduced chart, the
Gauss-predictor branch rule, the retained P1--P3 contracts, the P4 binding
contract, the P6 solver-scale contract, and the P5 direct-route certificate.
Then choose \(h_0\) small
enough for the Gauss truncation bound, the local
stage inverse, the endpoint right-inverse ball, the inexact-Newton
neighborhood, and the tube-retention bootstrap to hold on that same tube.
Then choose the theorem constants as suprema or algebraic combinations of
those fixed compact-tube bounds. Only after these choices is a reported grid
\(N_hh\le T\) considered. \mbox{No constant is re-fit} per step size, per backend,
per mechanism-coverage row, or from observed slopes, residual logs, or
work-precision tables.

\begin{lemma}[Order of the local paper-style \tfe{} target]
\label{lem:tfe-target-order}
The proof-level order target for the $m=3$ Gauss--Lobatto paper-style
\tfe{} formula used in this manuscript is five.
\end{lemma}

\begin{proof}
The source-paper family is compared through the temporal finite-element
formula-order convention
$p_{\mathrm{TFE}}(m)=2m-1$ in Eq.~\eqref{eq:tfe-order}. Substituting
$m=3$ gives $p_{\mathrm{TFE}}(3)=5$. This is the formula-order comparator
used here. It is distinct from the source paper's reported DAE pendulum
observation, where order reduction can make the implemented $m=3$ run appear
fourth order.
\end{proof}

\begin{proposition}[Conditional order comparison]
\label{prop:order-comparison}
Under Assumption~\ref{ass:regularity} and the current reproducibility
contract,
\method{} has higher formal order than the local paper-style
$m=3$ Gauss--Lobatto \tfe{} formula target in the narrow order-comparison
sense.
\end{proposition}

\begin{proof}
Theorem~\ref{thm:g6fullva-order} gives the conditional route to the
sixth-order method claim for \method{}. Lemma~\ref{lem:tfe-target-order}
gives the fifth-order local
paper-style $m=3$ comparator. This theorem--lemma pair is the formal-order
proof of the proposition. The generated h-sweep records observed smooth
position and velocity orders 7.161 and 7.066 only as consistency evidence
after that comparison; those finite-window slopes are not proof inputs. The
four-example validation criterion is a mechanism-coverage criterion, not a single order
theorem: the single- and
double-pendulum dynamic FullVA rows provide dynamic order evidence, while the
four-link and slider-crank rows are driven closed-loop kinematic/reaction
checks whose trajectory differences sit at the roundoff floor and are not
interpreted as method-order estimates. Therefore the conditional formal-order
comparison is supported only inside the smooth accepted dynamic path specified
in Assumption~\ref{ass:regularity} and Theorem~\ref{thm:g6fullva-order}. The
comparison inherits the compact-tube branch selection and solver-scale
restrictions of Theorem~\ref{thm:g6fullva-order}: Newton is initialized from
the Gauss predictor on the accepted branch, and the nonlinear residual budget
is the theorem-level \(\eta_h^{\rm tube}\le c_\eta h^7\) policy. It is not a
theorem about arbitrary Newton initializations, remote roots, fixed production
tolerances, or global solver convergence. The proposition does not compare
error constants, wall-clock performance against independently implemented
baselines, robustness in the nonsmooth sharp-friction regime,
four-link/slider-crank dynamic convergence order, or complete source-paper
residual reproduction. Those exclusions define the claim boundary; in
particular, the method is not a complete source-paper temporal finite-element
residual replacement.
\end{proof}

\begin{corollary}[Non-claim]
Proposition~\ref{prop:order-comparison} does not prove that the implemented
nonlinear residual is the complete source-paper temporal finite-element
residual.
\end{corollary}

\begin{proof}
The accepted method solves Eq.~\eqref{eq:residual}. A complete source-paper
residual reproduction would require replacing every stage row in that solve by
accepted paper-derived weak rows and passing the same validation criteria.
That stronger criterion remains open.
\end{proof}

\begin{definition}[P7 boundary-certificate template]
\label{def:p7-transfer-certificate-template}
For a mechanism-coverage row family on a fixed accepted branch, an accepted
P7 residual-to-error certificate would be a same-branch tuple
\begin{equation}
\label{eq:p7-certificate-tuple}
  \mathcal T_{\rm P7}
  =(\mathcal R_h^{\rm mech},\|\cdot\|_{\rm traj},C_{\rm P7},
    \rho_{\rm res},\varepsilon_{\rm floor},\varepsilon_{\rm est},
    \mathcal G_{\rm dyn}).
\end{equation}
and this tuple is fixed before the residual table is inspected, with the following properties.
Equation~\eqref{eq:p7-certificate-tuple} fixes the object whose transfer
bound is Eq.~\eqref{eq:p7-transfer-template-bound}.
First, \(\mathcal R_h^{\rm mech}\) is the residual operator actually reported
for the mechanism rows, with the same branch, time grid, row scaling, and
norm used in the table. Second, \(\|\cdot\|_{\rm traj}\) is the trajectory
error norm in which a dynamic-order row would be claimed. Third,
\(C_{\rm P7}\) is independent of \(h\) on the same compact branch. Fourth,
\(\rho_{\rm res}(h)\) is an a priori residual-consistency rate for inserting
the exact branch into the mechanism residual rows. Fifth,
\(\varepsilon_{\rm floor}(h)\) and \(\varepsilon_{\rm est}(h)\) bound,
respectively, reference-floor and estimator/calibration errors on the same
reported campaign. Sixth, \(\mathcal G_{\rm dyn}\) supplies the dynamic
identity, closed-loop DAE stability or inf-sup estimate, and output binding
needed to connect residual rows to trajectory variables. Seventh, the
dynamic-order campaign has a branch-compatible refinement set on which all
of these constants and rates are valid.
The stability or inf-sup clause is a nullspace exclusion.  It must rule out
same-branch perturbations that are invisible to the reported mechanism
residual but visible in the trajectory norm; otherwise a constraint or
reaction residual can be small while the reported position--velocity error is
not controlled.  Thus the residual operator, row weights, gauge conventions,
and trajectory output variables must be fixed together before any residual row
is promoted.

The certificate would have to imply Eq.~\eqref{eq:p7-transfer-template-bound},
uniformly for all sufficiently small
reported step sizes in that campaign,
\begin{equation}
\label{eq:p7-transfer-template-bound}
  \|e_h\|_{\rm traj}
  \le C_{\rm P7}\Bigl(
     \|\mathcal R_h^{\rm mech}\|_{\rm res}
     +\rho_{\rm res}(h)
     +\varepsilon_{\rm floor}(h)
     +\varepsilon_{\rm est}(h)
  \Bigr).
\end{equation}
Definition~\ref{def:p7-transfer-certificate-template} is a closure criterion,
not a theorem premise of Theorem~\ref{thm:g6fullva-order}: no such tuple is
proved or invoked in this manuscript.
\end{definition}

\begin{corollary}[Residual-to-error transfer boundary]
The four-link and slider-crank residual/reaction rows in the validation
section, and the common-reference or source-policy residual rows reported
elsewhere, do not become accepted dynamic-order or trajectory-error rows in
this manuscript.
\end{corollary}

\begin{proof}
This exclusion is not a residual-transfer theorem.
Those rows are mechanism-coverage evidence for the accepted FullVA residual
vocabulary, not residual-to-error evidence for the reported trajectory norm.
Promoting them to trajectory-error/order evidence would require
the seven residual-to-error obligations in
Table~\ref{tab:p7-residual-to-error-obligation-ledger} to close on the same
branch: dynamic residual identity, residual-consistency rate, closed-loop DAE
stability or inf-sup control, calibrated estimator, reference-floor exclusion,
coarse-first dynamic-order campaign, and a manuscript transfer theorem.  Equivalently,
the closure would have to instantiate the same-branch tuple in
Definition~\ref{def:p7-transfer-certificate-template}. Since that joint
criterion remains open, the rows remain diagnostics.
Equivalently, that separate promotion would require an additional same-branch
residual operator with a uniform inverse, fixed independently of the observed
residual table.
More formally, a proof of such a transfer would first have to define a same-branch residual
transfer operator on the reported closed-loop branch, a trajectory norm
\(\|e_h\|_{\rm traj}\) matching the claimed order rows, and an
\(h\)-independent constant \(C_{\rm P7}\) such that, for all sufficiently
small \(h\) in the claimed dynamic-order campaign, the transfer bound in
Eq.~\eqref{eq:corollary-p7-transfer-bound} would have to hold:
\begin{equation}
\label{eq:corollary-p7-transfer-bound}
  \|e_h\|_{\rm traj}
  \le C_{\rm P7}\Bigl(
     \|\mathcal R_h^{\rm mech}\|_{\rm res}
     +\rho_{\rm res}(h)
     +\varepsilon_{\rm floor}(h)
     +\varepsilon_{\rm est}(h)
  \Bigr),
\end{equation}
with the residual identity, residual-consistency rate, floor exclusion, and
estimator terms all attached to the same branch. The present residual and
reaction tables provide none of this transfer operator: they report small
mechanism residuals and consistent reactions, but they do not provide the
uniform inverse/stability constant or the reference-floor and estimator
closure needed to convert those residuals into trajectory-error order.
In particular, they do not exclude residual-null directions on the reported
closed-loop branch.  A perturbation that changes the reported trajectory
while leaving the chosen residual/reaction rows unchanged would violate any
uniform P7 transfer estimate, so finite residual smallness is not enough to
define the missing inverse.
Equivalently, PC4 is not a residual-to-error closure statement.  The
manuscript retains the scope rule that these rows are
not promoted; this retained boundary does not supply any stability,
estimator, floor-exclusion, campaign, or manuscript-theorem component of the
open P7 route.
\end{proof}

\begin{table}[htbp]
\centering
\caption{Proof-dependency map for Proposition~\ref{prop:order-comparison}.}
\label{tab:proof-traceability}
\tiny
\setlength{\tabcolsep}{2pt}
\renewcommand{\arraystretch}{0.68}
\begin{tabular}{P{0.28\linewidth}P{0.36\linewidth}P{0.20\linewidth}}
\toprule
Theorem component & Mathematical input & Role in the claim \\
\midrule
Accepted one-step map &
132-row residual, one-step algorithm, and endpoint closure in
Sections~\ref{sec:method}--\ref{sec:validation}. &
Definition fixed for the theorem and numerical protocol. \\
Endpoint closure perturbation &
Lemma~\ref{lem:endpoint-closure}: closed endpoint-anchor family, uniform
endpoint-ball margin, local right-inverse perturbation, and raw endpoint
defect. &
Retained; local \(O(h^7)\) map perturbation, not a substitute for the
stage-row defect certificate. \\
Sixth-order method path &
The mathematical input is the conditional criterion in
Theorem~\ref{thm:g6fullva-order}. &
Analytic criterion under the theorem hypotheses; the reported 7.161/7.066 slopes and
finite-run scale check are numerical consistency evidence only, not proof
closure or solver proof. \\
Local-to-global transfer &
Lemma~\ref{lem:local-global-transfer}: one-step \(O(h^7)\) defect,
\(1+C_s h\) stability, and tube retention for \(N_hh\le T\). &
Grid-point transfer; constants independent of \(h,N_h\); no unreported
terminal-time error when \(N_hh<T\). \\
Comparator order &
Lemma~\ref{lem:tfe-target-order} states the local $m=3$ expected order
$2m-1=5$ and separates it from source-paper residual reproduction. &
Formal-order comparator only. \\
Four examples &
The four ASME-style examples in Section~\ref{sec:four-examples} check the
accepted mechanism-coverage policy. &
Mechanism coverage; four-link/slider-crank rows are residual/reaction
checks, not order-proof rows. \\
Residual-to-error promotion &
No residual-to-error transfer theorem is accepted for residual-only mechanism
or source-policy rows. &
Excluded output boundary; residual tables are diagnostics, not smooth-branch theorem
inputs. \\
Non-claims &
Complete source-paper residual reproduction and full-\tfe{} stage
replacement remain excluded from the accepted theorem. &
Excluded stronger criterion. \\
Stage-row perturbation &
Lemmas~\ref{lem:fullva-stage-lift}--\ref{lem:stage-residual-defect} plus the
96-row certificate and 36-row direct Newton--Euler identity on \(Z_G\). &
Direct 132-row stage-residual input; symbolic defect certificate remains a
provenance record, not the route that discharges PC2. \\
Implementation fidelity &
One residual/Jacobian path, 132-row formula-row checks, finite-probe AD checks,
and AD-expanded Newton--Euler derivative-cell closure. &
Accepted implementation path for the stage-residual proof; not a primitive
symbolic defect certificate and not a source-paper residual replacement. \\
Branch selection &
Assumption~\ref{ass:regularity} and
Lemmas~\ref{lem:stage-residual-defect}--\ref{lem:inexact-newton} restrict the
accepted map to the compact-tube Gauss-predictor Newton branch. &
Local branch only; arbitrary Newton initializations, remote roots, and global
solver convergence are not claimed. \\
Inexact Newton error &
Lemma~\ref{lem:inexact-newton} states the local tolerance scale; finite probes
remain diagnostics. &
Theorem condition retained; fixed production runs do not close solver policy. \\
\bottomrule
\end{tabular}
\end{table}

\subsection{External comparison boundary}
\label{sec:external-comparison-boundary-pointer}

External source-paper comparison is kept outside the main method and proof
claims.  The accepted evidence below supports the branch-selected
Gauss6/FullVA method, the conditional sixth-order theorem, and method-side
mechanism coverage.  The \tfe{} and public-code source-policy rows are
retained only as diagnostic boundary records; they close no source-paper
superiority row and no theorem input.  The detailed same-test promotion
criterion, common-reference tables, and source-policy runner gaps are recorded
in Section~\ref{sec:cross-paper-boundary} so that the main numerical section
is read in the intended evidence order: dynamic-order evidence first,
mechanism-coverage evidence second, external-comparison diagnostics last.

\section{Numerical evidence}
\label{sec:numerical}

The numerical section should be read as a sequence of evidence blocks rather
than as a single undifferentiated benchmark table. First, the smooth
cylindrical-chain sweep tests the accepted 132-row \method{} stage solve and
is finite-run consistency evidence after the conditional order-six theorem;
it is not a premise that discharges the compact-tube, endpoint, or
solver-scale hypotheses.
Second, the four ASME-style mechanisms test the lower-pair FullVA vocabulary
across driven, open-chain, and closed-loop settings; only the single and
double pendula are used as accepted dynamic-order examples, while the
closed-loop examples supply constraint, reaction, and coarse-window trajectory
diagnostics only. Third, the source-paper and public-code rows are comparison
diagnostics whose role is to expose what remains open in a fair same-test
campaign. Finally, the work and solver tables make the evidence reproducible
by recording Newton counts, tolerances, and runtime, but they do not by
themselves turn the diagnostic baseline rows into superiority claims.

The recorded h-sweep uses $h=\{0.04,0.02,0.01\}$ and an $h=0.005$ reference
for the smooth cylindrical-chain benchmark. Table~\ref{tab:orders} gives the
headline order evidence.
Reference-solution boundary: every observed order reported from a finite
\(h_{\rm ref}\) row is a finite-reference consistency diagnostic, not a
replacement for the exact-flow comparison used in
Theorem~\ref{thm:g6fullva-order}.  The reference step, output norm, and
reported grid are fixed before the fit is computed.  A finite reference row
cannot discharge the compact-tube regularity, endpoint-closure, P6
solver-scale, or P7 residual-to-error boundaries, cannot be used to tune the
reported norm after the fact, and cannot exclude a reference floor without a
separate refinement or analytic-reference argument.

\begin{table}[htbp]
\centering
\caption{Observed orders for the accepted \method{} path.}
\label{tab:orders}
\begin{tabular}{P{0.30\linewidth}P{0.20\linewidth}P{0.18\linewidth}P{0.18\linewidth}}
\toprule
Case & Step sizes & Position order & Velocity order \\
\midrule
Smooth cylindrical chain & $0.04,0.02,0.01$ & 7.161 & 7.066 \\
Sharp friction, coarse sweep & $0.04,0.02,0.01$ & 1.894 & 2.683 \\
\bottomrule
\end{tabular}
\end{table}

\begin{table}[htbp]
\centering
\caption{Smooth cylindrical-chain projected-endpoint errors.}
\label{tab:smooth-errors}
\begin{tabular}{P{0.12\linewidth}P{0.16\linewidth}P{0.18\linewidth}P{0.18\linewidth}P{0.16\linewidth}}
\toprule
$h$ & Steps & Position error & Velocity error & Newton iterations \\
\midrule
0.04 & 2 & $1.761\times10^{-7}$ & $1.007\times10^{-5}$ & 12 \\
0.02 & 4 & $1.838\times10^{-9}$ & $2.954\times10^{-8}$ & 24 \\
0.01 & 8 & $8.602\times10^{-12}$ & $5.609\times10^{-10}$ & 48 \\
\bottomrule
\end{tabular}
\end{table}

The slopes 7.161 and 7.066 in Tables~\ref{tab:orders}--\ref{tab:smooth-errors}
should not be read as a seventh-order method claim. The theorem above claims
sixth order under an $O(h^7)$ local-defect contract. The larger fitted slopes
come from a short, smooth, projected-endpoint window in which the leading
sixth-order error coefficient is small relative to higher-order terms, while
the finest row is already close to reference, nonlinear-tolerance, and
floating-point floors. Consequently the paper treats these values as
post-theorem consistency evidence for the order-six behavior on the accepted
smooth path, not as a proof premise or a new asymptotic order. A wider
asymptotic sweep with a compact-tube
\(\eta_h^{\rm tube}\le c_\eta h^7\) nonlinear tolerance would be needed before
claiming more than order six, and the closed-loop four-link and slider-crank
rows are still not promoted to dynamic order estimates.

\begin{table}[htbp]
\centering
\caption{Order-estimate acceptance policy used in the numerical section.}
\label{tab:order-acceptance-policy}
\scriptsize
\setlength{\tabcolsep}{3pt}
\begin{tabular}{P{0.22\linewidth}P{0.22\linewidth}P{0.19\linewidth}P{0.23\linewidth}}
\toprule
Evidence row & Step/reference policy & Order information & Manuscript role \\
\midrule
Smooth cylindrical chain &
$h=0.04,0.02,0.01$ against $h=0.005$. &
Position/velocity slopes 7.161/7.066. &
Post-theorem finite-run consistency check for the conditional order-six
theorem; not a proof premise or seventh-order claim. \\
Single pendulum &
Accepted dynamic FullVA rows plus coarse public-window rows. &
Minimum accepted dynamic order above six; coarse local position order about
6.05. &
Accepted dynamic-order example and coarse-first comparison row. \\
Double pendulum &
Method-side FullVA dynamic rows plus coarse public-window rows. &
Minimum accepted method-side order above six; coarse public-window
position/velocity slopes 7.951/7.042. &
Accepted dynamic-order example and coarse-first comparison row. \\
Four-link and slider-crank &
Closed-loop kinematic/reaction rows and strict common-reference coarse rows. &
Local coarse trajectory-diagnostic candidates are near six, but floor-limited
endpoint rows are not used as asymptotic order estimates. &
Mechanism-coverage and work/precision evidence; not accepted asymptotic
dynamic-order proof or an external dynamic-order claim. \\
Paper-style \tfe{} all-row attempt &
Same smooth cylindrical-chain steps as the accepted row. &
Position/velocity slopes 4.046/4.420. &
Internal diagnostic only; not a source-paper same-test row and not an
accepted full-\tfe{} replacement. \\
Public baseline families &
Public-code and same-window rows where available. &
Useful order/time rows, including nonconvergence rows for some coarse
settings. &
External-baseline evidence in progress; no external-superiority claim. \\
\bottomrule
\end{tabular}
\end{table}

\begin{table}[htbp]
\centering
\caption{Non-promoted same-step diagnostic with a \tfe{}-style all-row attempt.}
\label{tab:tfe-diagnostic-comparison}
\scriptsize
\setlength{\tabcolsep}{3pt}
\begin{tabular}{P{0.10\linewidth}P{0.29\linewidth}P{0.33\linewidth}P{0.16\linewidth}}
\toprule
$h$ & Accepted \method{} & Paper-style all-row \tfe{} diagnostic & Diagnostic status \\
\midrule
0.04 &
Position $1.761\times10^{-7}$; velocity $1.007\times10^{-5}$; 12 Newton. &
Position $9.064\times10^{-8}$; velocity $1.238\times10^{-5}$; 13 Newton. &
Not accepted. \\
0.02 &
Position $1.838\times10^{-9}$; velocity $2.954\times10^{-8}$; 24 Newton. &
Position $3.378\times10^{-9}$; velocity $1.931\times10^{-7}$; 21 Newton. &
Not accepted. \\
0.01 &
Position $8.602\times10^{-12}$; velocity $5.609\times10^{-10}$; 48 Newton. &
Position $3.322\times10^{-10}$; velocity $2.701\times10^{-8}$; 37 Newton. &
Not accepted. \\
\bottomrule
\end{tabular}
\end{table}

Table~\ref{tab:tfe-diagnostic-comparison} is retained only as a
non-promotion record for the current source-paper comparator status. The
paper-style column is an implemented all-row \tfe{} diagnostic using the
terminal Lobatto-node output, a consistent initial $z_0$ policy, and an
equilibrated Newton linear solve. Its recorded smooth orders are 4.046 in
position and 4.420 in velocity. Because the start policy is not a
paper-derived recurrent policy and the full-\tfe{} replacement criterion is
not met by the current evidence package, this row is not an accepted fair
baseline and is not part of the main order evidence.

\begin{table}[htbp]
\centering
\caption{Work and solver diagnostics for accepted evidence rows.}
\label{tab:work-diagnostics}
\scriptsize
\setlength{\tabcolsep}{3pt}
\begin{tabular}{P{0.20\linewidth}P{0.19\linewidth}P{0.12\linewidth}P{0.34\linewidth}}
\toprule
Evidence row & Finest work & Runtime (s) & Finest accepted value \\
\midrule
Smooth cylindrical chain &
$h=0.01$; 8 steps / 48 Newton & 0.199 &
Position $8.602\times10^{-12}$; velocity $5.609\times10^{-10}$. \\
Single pendulum &
$h=0.05$; 4 steps / 8 Newton & 0.013 &
Position $7.374\times10^{-12}$; max stage residual
$3.664\times10^{-12}$. \\
Double pendulum &
$h=0.01$; 20 steps / 60 Newton & 3.374 &
Orientation $1.597\times10^{-15}$; angular velocity
$4.867\times10^{-15}$. \\
Four-link closed loop &
$h=0.005$; 40 steps / 164 Newton & 0.118 &
Max dynamics residual $1.338\times10^{-13}$. \\
Slider-crank closed loop &
$h=0.005$; 40 steps / 163 Newton & 0.122 &
Max dynamics residual $6.492\times10^{-15}$. \\
\bottomrule
\end{tabular}
\end{table}

\begin{table}[htbp]
\centering
\caption{Nonlinear solver tolerances used for the reported evidence rows.}
\label{tab:solver-tolerances}
\scriptsize
\setlength{\tabcolsep}{3pt}
\begin{tabular}{P{0.20\linewidth}P{0.31\linewidth}P{0.27\linewidth}P{0.12\linewidth}}
\toprule
Evidence row & Nonlinear solve & Stopping rule & Finest count \\
\midrule
Cylindrical chain &
132-row \method{} residual with dense AD Jacobian and dense linear solve;
rank diagnostics use tolerance $10^{-10}$. &
Stop when $\|R\|_2<10^{-11}$ or $\|\Delta Z\|_2<10^{-11}$; cap 80 Newton
iterations. &
48 Newton. \\
Single pendulum &
Absolute-coordinate driven FullVA residual with AD Jacobian and dense linear
solve. &
Stop when $\|R\|_2<10^{-11}$ or $\|\Delta Z\|_2<10^{-12}$; cap 32 Newton
iterations. &
8 Newton. \\
Double pendulum &
Imported double-revolute FullVA harness with three Gauss stages, AD Jacobian,
and dense linear solve. &
Stop when $\|R\|_2<10^{-11}$ or $\|\Delta Z\|_2<10^{-11}$; cap 64 Newton
iterations. &
60 Newton. \\
Four-link closed loop &
Local fully constrained kinematic FullVA solve for $\Phi(q,t)=0$ followed by
linear solves for $\Phi_q\dot q=\nu$ and $\Phi_q\ddot q=\gamma$. &
Position Newton stops when $\|\Delta q\|_2<10^{-12}$; cap 30 Newton
iterations. &
164 Newton; max 4/step. \\
Slider-crank closed loop &
Same local closed-loop kinematic FullVA policy as the four-link example. &
Position Newton stops when $\|\Delta q\|_2<10^{-12}$; cap 30 Newton
iterations. &
163 Newton; max 4/step. \\
\bottomrule
\end{tabular}
\end{table}

For the asymptotic theorem, Table~\ref{tab:solver-tolerances} should be read
with the inexact-Newton policy of Lemma~\ref{lem:inexact-newton}: a proof-level
h-sweep must use the nonlinear residual threshold in
Eq.~\eqref{eq:solver-table-p6-scale},
\begin{equation}
\label{eq:solver-table-p6-scale}
  \eta_h^{\rm tube} \le c_\eta h^7
\end{equation}
on the compact-tube branch, or an equivalent update criterion that implies
the same stage-solution error. The fixed absolute tolerances in the recorded
runs are retained for reproducibility of the reported runs and explain the
reported Newton counts; they are not by themselves a derivation of the
asymptotic solver-error bound. A solver-scale check of the existing summary
residuals records roundoff-scale linear and stage residuals, but the ratios to
$h^7$ grow at the finest reference scale. A separate short-trajectory
scaled-tolerance probe records the finite-log reported-grid maximum
\(\eta_h^{\rm reported}:=\max_{0\le n<N_h}\eta_{h,n}
=\widehat\eta_h^{\rm log}\le c_\eta h^7\)
with \(c_\eta=10^4\) at
$h=0.04,0.02,0.01,0.005$ through $T=0.08$; all four diagnostic rows and all
30 trajectory steps pass the finite diagnostic filter, and the maximum final
residual divided by $h^7$ is 210.891. This diagnostic pass is not a
theorem-level P6 proof. It is a useful finite P6 diagnostic record, broader
than a one-step smoke test, and it records reported-window consistency only
rather than solver-policy closure. A companion finite tolerance-regime sweep compares fixed
$10^{-10}$, $c h^7$, and $c h^8$ stopping rules against a local
$h=0.0025$ reference over the same window. The fixed and $h^7$ policies give
position/velocity fits 6.946/6.608, while the $h^8$ policy gives
6.946/6.608. The two shrinking-target policies are separately recorded as
eight finite h-scaled rows and 60 h-scaled trajectory steps, with
position/velocity order floor 6.946/6.608. Thus, on this finite diagnostic
window, the fixed-tolerance row and the \(h\)-scaled rows have the same
recorded order floor; this is a finite-window tolerance comparison only, not
P6 instantiation or a theorem-level solver-policy proof. The scaled policies
exercise the theorem's shrinking tolerance form on this finite window, but the
sweep is still not a theorem-level scaled-tolerance proof over every reported
trajectory and not a scaled-tolerance proof under a uniform all-transition
policy for the reported trajectory.
The proof therefore remains explicitly conditional on the
\(\eta_h^{\rm tube}\le c_\eta h^7\) compact-tube policy.  The reported-log
quantity is only
\(\eta_h^{\rm reported}=\max_{0\le n<N_h}\eta_{h,n}\), also denoted
\(\widehat\eta_h^{\rm log}\) in diagnostic tables, after the accepted branch,
proof norm, row scaling, and \(c_\eta\) have already been fixed; it cannot
select the branch, norm, or theorem constant.

Tables~\ref{tab:work-diagnostics}--\ref{tab:solver-tolerances} and
Table~\ref{tab:tfe-diagnostic-comparison} are solver-scale and diagnostic
comparison tables, not accepted fair baseline comparisons. Runtime values are
single-machine recorded timings, and the mechanism rows report the accepted
finest-step evidence for each validation path. A publication-level comparison still
requires accepted implemented baselines and cross-method work/precision curves
on the same problems; these tables and Figure~\ref{fig:convergence}
only prevent the accepted evidence from being reported without any work,
tolerance, Newton-iteration, or source-paper diagnostic context.
For the four-link and slider-crank rows, the finest-step residual and reaction
entries are read through the P7 boundary stated after
Theorem~\ref{thm:g6fullva-order}: they demonstrate closed-loop mechanism
coverage and consistency of the constructed FullVA rows, but they are not
trajectory-order rows unless an accepted residual-to-error transfer is proved.

\paperfig{Figure_1_convergence.png}
{Convergence and work/precision evidence for the accepted \method{} path and
the implemented source-paper all-row \tfe{} diagnostic. The top row compares
the smooth cylindrical-chain errors and Newton-work/error curves against the
diagnostic paper-style all-row attempt; that curve is not an accepted fair
baseline. The bottom row contrasts the sharp-friction coarse sweep with the
refined cost envelope that recovers high order only at additional runtime.}
{fig:convergence}

The sharp-friction coarse sweep is not used as the main order claim. It is
reported as a limitation: deeper fixed refinement recovers high order at
increased cost, while the practical coarse regime remains order-reduced.
The recorded sharp-friction diagnostic shows the practical tradeoff: the coarse
$h=\{0.04,0.02,0.01\}$ sweep gives only 1.894/2.683 position/velocity order,
while ultra refinement over $h=\{0.005,0.0025,0.00125\}$ against
$h=0.000625$ recovers 7.521/5.725 all-point position/velocity order. The
high-order recovery requires an approximately $7.25\times$ runtime increase
relative to the $h=0.01$ baseline, so the paper keeps this as a practical
cost caveat rather than a solved production setting.

\section{Four-example mechanism validation}
\label{sec:four-examples}

The method-side evidence covers four ASME-style examples, with different
accepted roles. This section states the mechanism content used by the
validation rows and keeps their evidentiary roles distinct. The examples are
not treated as interchangeable benchmark names: each one exercises a different
part of the FullVA residual.
These validation rows are downstream evidence for method behavior; they do
not feed back into the constants, hypotheses, or residual-certificate slot of
Theorem~\ref{thm:g6fullva-order}.  In particular, closed-loop residual and
reaction rows remain coverage/diagnostic rows unless the separate P7
residual-to-error boundary is discharged.

\subsection{Driven single pendulum}

The single pendulum is a rigid bar of length $4.0$ m, square side length
$0.05$ m, density $7800$ kg/m$^3$, and gravity $g=9.81$ m/s$^2$. The center
of mass is $2.0$ m from the fixed pivot. The accepted driven trajectory is
\begin{equation}
  \theta(t)=\frac{\pi}{2}+\frac{\pi}{4}\cos(2t),
  \label{eq:single-drive}
\end{equation}
Equation~\eqref{eq:single-drive} defines the accepted driven single-pendulum
trajectory, with the corresponding analytic $\dot\theta$ and $\ddot\theta$ used for
reference values. The constraint set consists of a fixed pivot, two DP1
revolute-axis constraints, and a DP1 drive constraint. The strongest accepted
row is the absolute-coordinate driven FullVA residual:
each Gauss stage solves for orientation, center-of-mass position, velocity,
acceleration, angular velocity, angular acceleration, pivot reaction force,
axis multipliers, and drive multiplier. Its square residual includes pivot
velocity and pivot acceleration rows, two revolute-axis rows, and the
acceleration-level drive row.

\subsection{Double pendulum}

The double pendulum uses two rigid bars with lengths $4.0$ m and $2.0$ m,
the same square side length, density, and gravity as the single-pendulum
case, and two revolute joints. The first body is pinned to ground and the
second body is pinned to the first body. Friction and external torques are
disabled in the accepted method-side comparison. Because the available
legacy rA reference shows an h-independent discrepancy under the coordinate
convention change, the accepted reference policy is a nested local
\method{} self-reference at $h=0.005$; the rA comparison is retained only as
a diagnostic reference-floor test.

\subsection{Four-link closed loop}

The four-link example is a driven closed-loop lower-pair graph with three
moving rigid bodies and $18$ scalar constraints. The constraint inventory is
$12$ coordinate-difference rows and $6$ DP1 rows. Since the graph is fully
constrained, the accepted evidence is not a free dynamic convergence claim.
At each output time the local FullVA solve enforces
Eq.~\eqref{eq:four-link-fullva-constraint-levels},
\begin{equation}
\label{eq:four-link-fullva-constraint-levels}
  \Phi(q,t)=0,\qquad \Phi_q(q,t)\dot q=\nu(q,t),\qquad
  \Phi_q(q,t)\ddot q=\gamma(q,\dot q,t),
\end{equation}
then reconstructs multipliers and checks Newton--Euler translational and
rotational residuals. The largest recorded dynamics residual is
$1.338\times 10^{-13}$.

\subsection{Slider-crank closed loop}

The slider-crank example is also a three-body fully constrained graph with
$18$ scalar constraints, but it includes prismatic geometry through DP2 and
distance rows. Its inventory is $6$ coordinate-difference rows, $7$ DP1 rows,
$4$ DP2 rows, and one distance row. The accepted solve follows the same
closed-loop FullVA policy as the four-link mechanism. The largest recorded
constraint norm is $1.204\times10^{-15}$ and the largest dynamics residual is
$6.492\times10^{-15}$.

\begin{table}[htbp]
\centering
\small
\setlength{\tabcolsep}{2pt}
\caption{Four-example validation summary.}
\label{tab:asme}
\begin{tabular}{P{0.16\linewidth}P{0.25\linewidth}P{0.19\linewidth}P{0.26\linewidth}}
\toprule
Example & Accepted evidence & Role & Headline metric \\
\midrule
Single pendulum &
Driven absolute-coordinate FullVA residual &
Dynamic-order evidence &
Minimum observed order 6.024; max stage residual
$3.664\times10^{-12}$. \\
Double pendulum &
Method-side FullVA reference policy &
Dynamic-order evidence &
Minimum orientation/velocity order 6.089; endpoint velocity
constraints below $8.323\times10^{-12}$. \\
Four link &
Closed-loop kinematic FullVA and reaction dynamics &
Mechanism coverage only &
Max constraint norm $1.052\times10^{-14}$; max dynamics residual
$1.338\times10^{-13}$. \\
Slider crank &
Closed-loop kinematic FullVA and reaction dynamics &
Mechanism coverage only &
Max constraint norm $1.204\times10^{-15}$; max dynamics residual
$6.492\times10^{-15}$. \\
\bottomrule
\end{tabular}
\end{table}

\begin{table}[htbp]
\centering
\scriptsize
\setlength{\tabcolsep}{3pt}
\caption{Problem setup values used by the four ASME-style validation examples.}
\label{tab:asme-setup}
\begin{tabular}{P{0.14\linewidth}P{0.30\linewidth}P{0.24\linewidth}P{0.16\linewidth}}
\toprule
Example & Physical and constraint setup & Time grid and reference policy & Accepted check \\
\midrule
Single pendulum &
One steel bar, length $4$ m, square side $0.05$ m, density
$7800$ kg/m$^3$, mass $78$ kg, and
$J=\operatorname{diag}(0.0325,104.016,104.016)$ kg m$^2$.
The graph has one ground revolute joint: CD(3) and DP1(3). &
Driven angle
$\theta(t)=\pi/2+(\pi/4)\cos(2t)$, analytic reference,
$t_f=0.2$, and $h=\{0.20,0.10,0.05\}$. &
Minimum observed order $6.024$; max stage residual
$3.664\times10^{-12}$. \\
Double pendulum &
Two steel bars, lengths $4$ m and $2$ m, same side and density as the
single pendulum, masses $78$ kg and $39$ kg, with two revolute joints and no
friction or external torque. &
$q(0)=(-\pi/2,0)$ and regularized angular speeds
$(10^{-12},-10^{-12})$ rad/s. The accepted reference is nested local
FullVA at $h=0.005$; tested $h=\{0.04,0.02,0.01\}$. &
Orientation and angular-velocity orders $6.101/6.089$; endpoint velocity
constraint below $8.323\times10^{-12}$. \\
Four-link &
Three moving bodies with masses $2,1,1$ kg and inertia diagonals
$(4,2,0)$, $(12.4,0.01,0)$, and $(4.54,0.01,0)$ kg m$^2$.
The closed loop has CD(12) and DP1(6). &
Driven constraint $\cos(\pi t+\pi/2)$, $t_f=0.2$,
$h=\{0.02,0.01,0.005\}$, and local kinematic FullVA reference
$h=0.001$. &
Max constraint norm $1.052\times10^{-14}$; max dynamics residual
$1.338\times10^{-13}$. \\
Slider-crank &
Crank, rod, and slider masses $0.12,0.5,2$ kg with inertia diagonals
$(10^{-4},10^{-5},10^{-4})$, $(4\times10^{-3},4\times10^{-4},
4\times10^{-3})$, and $(10^{-4},10^{-4},10^{-4})$ kg m$^2$.
The graph has CD(6), DP1(7), DP2(4), and D(1). &
Driven constraint $\cos(-2\pi t+\pi/2)$, $t_f=0.2$,
$h=\{0.02,0.01,0.005\}$, and local kinematic FullVA reference
$h=0.001$. &
Max constraint norm $1.204\times10^{-15}$; max dynamics residual
$6.492\times10^{-15}$. \\
\bottomrule
\end{tabular}
\end{table}

\subsection{Coordinate-level mechanism specification}

The setup table gives the physical scales; the following coordinate data make
the lower-pair graphs reproducible directly from the manuscript. Let
$e_x=(1,0,0)$, $e_y=(0,1,0)$, and $e_z=(0,0,1)$. All point vectors $s$ and
direction vectors $a$ are body-fixed unless the body is ground; each CD group
uses the three coordinate directions $e_x,e_y,e_z$. Driver functions listed
here are the functions actually used by the setup code after model-file
overrides.

\paragraph{Single pendulum}
The moving body is body 1 and ground is body 0. The ground revolute joint uses
CD rows with $s_1=(-2,0,0)$ and $s_0=(0,0,0)$, plus DP1 rows
$(a_1,a_0)=(e_x,e_x)$ and $(e_y,e_x)$. The driven DP1 row uses
$(a_1,a_0)=(e_y,-e_z)$ and
$f(t)=\cos(\pi/2+(\pi/4)\cos(2t))$.

\paragraph{Double pendulum}
The two links are bodies 1 and 2, with ground body 0. The ground revolute
joint for body 1 uses CD rows with $s_1=(-2,0,0)$ and $s_0=(0,0,0)$, plus DP1
rows $(e_x,e_x)$ and $(e_y,e_x)$. The inter-link revolute joint uses CD rows
between body 1 and body 2 with $s_1=(2,0,0)$ and $s_2=(-1,0,0)$. The second
axis pair is represented by DP1 rows on body 2 against ground:
$(a_2,a_0)=(e_x,e_x)$ and $(e_y,e_x)$.

\paragraph{Four-link closed loop}
The moving bodies are the rotor, link 2, and link 3, numbered 1, 2, and 3;
ground is body 0. Revolute joint A between body 1 and ground uses CD rows with
$s_1=(0,0,0)$ and $s_0=(0,0,0)$, plus DP1 rows $(e_x,e_x)$ and $(e_y,e_x)$.
Revolute joint D between body 3 and ground uses CD rows with
$s_3=(0,3.7,0)$ and $s_0=(-4,-8.5,0)$, plus DP1 rows
$((0,0,1),e_y)$ and $(e_y,e_y)$. The spherical joint between bodies 3 and 2
uses CD rows with $s_3=(0,-3.7,0)$ and $s_2=(0,-6.1,0)$. The universal joint
between bodies 1 and 2 uses CD rows with $s_1=(-2,0,0)$ and
$s_2=(0,6.1,0)$, plus the DP1 row
$((0,-0.662,0.75),e_z)$. The driving DP1 row uses
$((-1,0,0),e_y)$ and $f(t)=\cos(\pi t+\pi/2)$.

\paragraph{Slider-crank closed loop}
The crank, connecting rod, and slider are bodies 1, 2, and 3; ground is body
4. Revolute joint A between ground and crank uses CD rows with
$s_4=(0,0.1,0.12)$ and $s_1=(0,0,0)$, plus DP1 rows $(e_y,e_x)$ and
$(e_z,e_x)$. The spherical joint B between rod and crank uses CD rows with
$s_2=(-0.15,0,0)$ and $s_1=(0,0.08,0)$. The revolute/cylindrical joint C
between rod and slider uses one DP1 row $(e_y,e_z)$ and two DP2 rows with
rod point $s_2=(0.15,0,0)$, slider point $s_3=(0,0,0)$, and rod directions
$e_y$ and $e_z$. The translational joint D between slider and ground uses DP1
rows $(-e_x,e_x)$, $(e_y,e_x)$, and $(-e_x,e_y)$, plus DP2 rows from slider
point $s_3=(0,0,0)$ to ground point $s_4=(0.2,0,0)$ with slider directions
$-e_x$ and $e_y$. The distance row constrains the rod point $(0.15,0,0)$ and
slider point $(0,0,1)$ with $f(t)=1$. The driving DP1 row uses
$(e_y,e_y)$ and $f(t)=\cos(-2\pi t+\pi/2)$.

\paperfig{Figure_2_asme_lower_pair_graph_bridge.png}
{Coordinate-oriented mechanism schematics and lower-pair constraint inventory
for the four ASME-style examples. Each panel shows the global coordinate frame,
named joints, the principal moving bodies, the driver where applicable, and
the CD/DP1/DP2/D constraint vocabulary used by the FullVA rows. Exact
body-fixed point and direction vectors are listed in the text immediately
preceding the figure.}
{fig:asme-graph}

\paperfig{Figure_3_asme_closed_loop_kinematic_fullva.png}
{Closed-loop kinematic FullVA mechanism-coverage evidence for the four-link
and slider-crank examples.}
{fig:closed-loop}

\section{Lower-pair constraint formulas}
\label{sec:lower-pair-formulas}

This section makes explicit the scalar constraint rows used by the four
example mechanisms. Let body $i$ have center position $r_i$, rotation matrix
$A_i\in SO(3)$, translational velocity $v_i$, angular velocity $\omega_i$,
translational acceleration $a_i$, and angular acceleration $\alpha_i$. A
marked point and a marked direction fixed in body $i$ are
$x_i=r_i+A_i s_i$ and $u_i=A_i a_i$, respectively. The row families used in
the examples are listed in Eq.~\eqref{eq:lower-pair-row-families}:
\begin{align}
  C_{\mathrm{CD},k}(q,t)
    &= e_k^T\left(x_j-x_i\right)-d_k(t),\\
  C_{\mathrm{DP1}}(q,t)
    &= u_i^T u_j-c(t),\\
  C_{\mathrm{DP2}}(q,t)
    &= u_i^T\left(x_j-x_i\right)-d(t),\\
  C_{\mathrm{D}}(q,t)
    &= \left\|x_j-x_i\right\|^2-\ell^2(t).
  \label{eq:lower-pair-row-families}
\end{align}
Here CD denotes coordinate difference, DP1 constrains two body-fixed
directions, DP2 constrains a body-fixed direction against a point difference,
and D is a distance row. Revolute joints are assembled from CD and DP1 rows;
the slider-crank prismatic joint additionally uses DP2 and distance rows.

The velocity and acceleration formulas use the rigid-body identities in
Eq.~\eqref{eq:lower-pair-rigid-kinematic-identities}:
\begin{align}
  \dot x_i &= v_i+\omega_i\times A_i s_i,&
  \dot u_i &= \omega_i\times u_i,\\
  \ddot x_i &= a_i+\alpha_i\times A_i s_i+
    \omega_i\times\left(\omega_i\times A_i s_i\right),&
  \ddot u_i &= \alpha_i\times u_i+
    \omega_i\times\left(\omega_i\times u_i\right).
  \label{eq:lower-pair-rigid-kinematic-identities}
\end{align}
Writing $\Delta x=x_j-x_i$, $\Delta\dot x=\dot x_j-\dot x_i$, and
$\Delta\ddot x=\ddot x_j-\ddot x_i$, the row derivatives assembled by the
FullVA block are Eq.~\eqref{eq:lower-pair-velocity-rows} at velocity level
\begin{align}
  \dot C_{\mathrm{CD},k}
    &= e_k^T\Delta\dot x-\dot d_k,\\
  \dot C_{\mathrm{DP1}}
    &= \dot u_i^T u_j+u_i^T\dot u_j-\dot c,\\
  \dot C_{\mathrm{DP2}}
    &= \dot u_i^T\Delta x+u_i^T\Delta\dot x-\dot d,\\
  \dot C_{\mathrm{D}}
    &= 2\Delta x^T\Delta\dot x-2\ell\dot\ell,
  \label{eq:lower-pair-velocity-rows}
\end{align}
and Eq.~\eqref{eq:lower-pair-acceleration-rows} at acceleration level:
\begin{align}
  \ddot C_{\mathrm{CD},k}
    &= e_k^T\Delta\ddot x-\ddot d_k,\\
  \ddot C_{\mathrm{DP1}}
    &= \ddot u_i^T u_j+2\dot u_i^T\dot u_j+u_i^T\ddot u_j-\ddot c,\\
  \ddot C_{\mathrm{DP2}}
    &= \ddot u_i^T\Delta x+2\dot u_i^T\Delta\dot x+
       u_i^T\Delta\ddot x-\ddot d,\\
  \ddot C_{\mathrm{D}}
    &= 2\Delta\dot x^T\Delta\dot x+
       2\Delta x^T\Delta\ddot x-
       2\left(\dot\ell^2+\ell\ddot\ell\right).
  \label{eq:lower-pair-acceleration-rows}
\end{align}

\subsection{Scalar lower-pair Jacobian entries}

The scalar Jacobian rows used by Eq.~\eqref{eq:fullva-jacobian} are obtained
by varying the marked points and directions in spatial form. Define the
geometry variables by Eq.~\eqref{eq:scalar-lower-pair-geometry-definitions}:
\begin{equation}
\label{eq:scalar-lower-pair-geometry-definitions}
  \rho_i=A_i s_i,\qquad u_i=A_i a_i,\qquad
  \Delta x=x_j-x_i .
\end{equation}
Let $\vartheta_i$ be the spatial infinitesimal rotation associated with the
stage perturbation $\delta\eta_i$ through Eq.~\eqref{eq:pose-variation}. Then
\begin{align}
  \delta \rho_i &= \vartheta_i\times \rho_i,&
  \delta u_i &= \vartheta_i\times u_i,\\
  \delta x_i &= \delta r_i+\vartheta_i\times\rho_i,&
  \delta\Delta x &= \delta x_j-\delta x_i .
  \label{eq:point-direction-variation}
\end{align}
Equation~\eqref{eq:point-direction-variation} supplies the point and
direction variations used by the scalar lower-pair Jacobian rows.
The position-level scalar Jacobian entries are therefore
\begin{align}
  \delta C_{\mathrm{CD},k}
    &= e_k^T\delta\Delta x,\\
  \delta C_{\mathrm{DP1}}
    &= \vartheta_i^T(u_i\times u_j)
       +\vartheta_j^T(u_j\times u_i),\\
  \delta C_{\mathrm{DP2}}
    &= \vartheta_i^T(u_i\times\Delta x)
       +u_i^T\delta\Delta x,\\
  \delta C_{\mathrm{D}}
    &= 2\Delta x^T\delta\Delta x .
  \label{eq:position-scalar-jacobian}
\end{align}
Equation~\eqref{eq:position-scalar-jacobian} gives the translational columns
directly through $\delta r_i,\delta r_j$ and the rotational columns through
the map $\delta\eta\mapsto\vartheta$ in Eq.~\eqref{eq:pose-variation}. For
driven rows, the prescribed functions $d_k(t)$, $c(t)$, $d(t)$, and
$\ell(t)$ are fixed during the Newton solve at a stage time and therefore do
not add unknown columns.

The velocity-level Jacobian follows by differentiating the velocity rows with
the perturbations
\begin{align}
  \delta\dot x_i
    &= \delta v_i+\delta\omega_i\times\rho_i
       +\omega_i\times\delta\rho_i,\\
  \delta\dot u_i
    &= \delta\omega_i\times u_i+\omega_i\times\delta u_i,\\
  \delta\Delta\dot x
    &= \delta\dot x_j-\delta\dot x_i .
  \label{eq:velocity-scalar-perturbation}
\end{align}
Substitution of Eq.~\eqref{eq:velocity-scalar-perturbation} into
$\dot C_{\mathrm{CD}}$, $\dot C_{\mathrm{DP1}}$,
$\dot C_{\mathrm{DP2}}$, and $\dot C_{\mathrm{D}}$ gives the exact
$K^{\dot\Phi,q}$ and $K^{\dot\Phi,\dot q}$ blocks in
Table~\ref{tab:stage-jacobian-block}. The acceleration-level Jacobian uses
\begin{align}
  \delta\ddot x_i
    &= \delta a_i+\delta\alpha_i\times\rho_i
       +\alpha_i\times\delta\rho_i \notag\\
    &\quad+\delta\omega_i\times(\omega_i\times\rho_i)
       +\omega_i\times(\delta\omega_i\times\rho_i
       +\omega_i\times\delta\rho_i),\\
  \delta\ddot u_i
    &= \delta\alpha_i\times u_i+\alpha_i\times\delta u_i \notag\\
    &\quad+\delta\omega_i\times(\omega_i\times u_i)
       +\omega_i\times(\delta\omega_i\times u_i
       +\omega_i\times\delta u_i),\\
  \delta\Delta\ddot x
    &= \delta\ddot x_j-\delta\ddot x_i .
  \label{eq:acceleration-scalar-perturbation}
\end{align}
Substituting Eq.~\eqref{eq:acceleration-scalar-perturbation} into the four
$\ddot C$ formulas above supplies the $K^{\ddot\Phi,q}$,
$K^{\ddot\Phi,\dot q}$, and $K^{\ddot\Phi,\ddot q}$ blocks. These formulas
are the scalar CD/DP1/DP2/D Jacobian entries checked against the dense AD
Jacobian on the accepted path.

For any assembled scalar row $\Phi(q,t)=0$, these formulas are represented
in the compact FullVA residual as position, velocity, and acceleration
levels in Eq.~\eqref{eq:assembled-fullva-levels},
\begin{equation}
\label{eq:assembled-fullva-levels}
\begin{aligned}
  \Phi(q,t)&=0,\\
  \Phi_q(q,t)\dot q+\Phi_t(q,t)&=0,\\
  \Phi_q(q,t)\ddot q+
  \frac{d}{dt}\left(\Phi_q(q,t)\right)\dot q+
  \frac{d}{dt}\Phi_t(q,t)&=0.
\end{aligned}
\end{equation}
Equivalently, the velocity and acceleration rows are written as
$\Phi_q(q,t)\dot q=\nu(q,t)$ and
$\Phi_q(q,t)\ddot q=\gamma(q,\dot q,t)$, with nonzero right-hand sides for
driven constraints. This is the constraint vocabulary behind the single
pendulum, double pendulum, four-link, and slider-crank evidence in
Table~\ref{tab:asme}.

The Newton unknown at each collocation stage is ordered as
Eq.~\eqref{eq:stage-unknown-order},
\begin{equation}
\label{eq:stage-unknown-order}
  Z_i=(r_i,\eta_i,v_i,\omega_i,a_i,\alpha_i,\lambda_i),
\end{equation}
where $\eta_i$ is the local Lie-algebra rotation increment. The lower-pair
Jacobian rows therefore have the block form in
Eq.~\eqref{eq:lower-pair-residual-linearization},
\begin{equation}
\label{eq:lower-pair-residual-linearization}
  \delta R_{\mathrm{lp}} =
  R_q\,\delta q_i + R_{\dot q}\,\delta \dot q_i+
  R_{\ddot q}\,\delta \ddot q_i+R_\lambda\,\delta\lambda_i,
\end{equation}
with $R_q$ containing CD/DP1/DP2/D position derivatives, $R_{\dot q}$
containing the velocity-level blocks, $R_{\ddot q}$ containing the
acceleration-level blocks, and $R_\lambda=0$ for pure kinematic rows. The
Newton--Euler rows use the same $\Phi_q(q_i,t_i)^T\lambda_i$ columns in their
force and torque balances. Thus the paper's compact notation corresponds to
the actual implementation-level row groups: collocation rows, velocity
collocation rows, Newton--Euler balance rows, and FullVA lower-pair rows.

\section{Diagnostic boundaries and limitations}

This section states diagnostic boundaries; it is not a second contribution line
and not a prerequisite for the theorem or the accepted validation rows.  The
method claim can be evaluated from the residual definition,
Theorem~\ref{thm:g6fullva-order}, and Section~\ref{sec:numerical}; the
diagnostics below are retained only to prevent stronger source-paper,
full-\tfe{}, sparse-performance, and external-superiority claims from being
read into the main method evidence.
The supplementary evidence contains additional diagnostic work toward
source-paper temporal finite-element residual replacement. This evidence is
useful because it identifies why a stronger reproduction claim is outside the
present scope, but it is not needed for
Proposition~\ref{prop:order-comparison}. The central diagnostic result is
that several endpoint or weak-row variants can close one property at a time,
but the reported evidence does not provide a source-free, projection-free,
non-terminal-row-replacement 132-row \tfe{} residual that both closes terminal
endpoint velocity and preserves the smooth order target.

\paperfig{Figure_4_order_closure_blend.png}
{Order/closure blend diagnostic. The tested path separates order preservation
from terminal endpoint-velocity closure, which is why independent full-\tfe{}
replacement remains open.}
{fig:order-closure-blend}

\paperfig{Figure_5_velocity_compression.png}
{Velocity-compression diagnostic. Local row-space probes quantify the gap
between available weak-row features and the terminal bridge needed for
accepted full-\tfe{} stage replacement.}
{fig:velocity-compression}

Sparse automatic differentiation is also treated as diagnostic engineering
evidence rather than a primary paper claim. The accepted sparsity pattern is
correct and quantified, but the recorded wall-clock measurements do not show a
stable speed win over dense \texttt{jacfwd}.
The current code should therefore be read as a reference implementation. The
recorded small-system pattern has a $132\times132$ dense dimension, 2637
nonzeros, and 90/60 column/row colors; row-VJP is still 1.20--1.22 times the
dense runtime, requiring a 16.9--18.3\% row-runtime reduction merely to match
dense. No body-count sweep for $N=2,4,8,16,32$, memory scaling,
conditioning-versus-body-count, or Jacobian/linear-solve breakdown is claimed.

\paperfig{Figure_6_sparse_speed_gap.png}
{Sparse backend speed-gap diagnostic. The sparse structure is correct and
quantified, but dense \texttt{jacfwd} remains faster in the recorded
wall-clock evidence.}
{fig:sparse-speed}

\paperfig{Figure_7_strict_common_reference_work_precision.png}
{Strict common-reference work/precision rows for bounded diagnostics on the
closed-loop four-link and slider-crank mechanisms. Bounded common-reference
diagnostic work/precision rows use the same finite accepted-method endpoint reference at
$h=0.0125$ over the coarse
$h=0.1,0.05,0.025$ window. The plot removes one fixed-grid diagnostic caveat
for these coarse rows, but it is still not an external superiority claim because
the broader source-paper suites and manuscript-level review remain open.}
{fig:strict-common-reference}

\paperfig{Figure_8_claim_boundary_limitations.png}
{Claim-boundary and limitation map for the present paper. Green
items are accepted claims or accepted evidence rows; orange items are bounded
or partial evidence; red items remain open boundaries. The figure summarizes why
the manuscript supports a conditional sixth-order \method{} claim under the
retained theorem interfaces but still does not claim complete source-paper residual replacement, external
same-test superiority, or a closed residual-to-error transfer theorem for the
closed-loop mechanisms.}
{fig:claim-boundary-limitations}

\paperfig{Figure_9_coarse_baseline_work_precision.png}
{Coarse-first baseline and work/precision diagnostics. The top panels
compare observed order on the same coarse window for the single and double
pendula. The lower panels use the strict common-reference closed-loop
summaries for four-link and slider-crank, with all plotted rows evaluated
against the same finite accepted-method endpoint reference at $h_{\rm ref}=0.0125$.
These panels document coarse-first diagnostics and work/precision
availability, but they are not an external-superiority claim and do not
replace the unresolved original source-policy campaigns.}
{fig:coarse-baseline-work-precision}

\paperfig{Figure_10_closed_loop_true_dynamic_order.png}
{Closed-loop coarse-dynamics diagnostic evidence for the four-link and
slider-crank mechanisms. Closed-loop coarse-window trajectory diagnostics plot
endpoint position, orientation,
linear-velocity, and angular-velocity errors against the $O(h^6)$ guide over
$h=0.1,0.05,0.025$ with $h_{\rm ref}=0.0125$. The bottom panels summarize the
finite-window fitted slopes and the non-oracle Newton stage-solve work. These
rows are mechanism-coverage and coarse-window trajectory diagnostics only, not
accepted dynamic-order proof; they run without
using stage-time kinematic oracles, without making $10^{-4}$ a default policy,
without promoting four-link or slider-crank to accepted dynamic-order rows, and
without asserting external same-test superiority. They are finite-window
diagnostics: they do not instantiate the residual-to-error implication, do not
close P7, and do not promote four-link or slider-crank to accepted dynamic-order
examples.}
{fig:closed-loop-true-dynamic-order}

\paperfig{Figure_12_all_method_result_matrix.png}
{All-method common-reference result matrix for the four validation examples.
Each cell reports observed velocity order and finest-step velocity error from
the paper numerical result matrix over
$h=0.1,0.05,0.025$ with $h_{\rm ref}=0.0125$. The local \method{} row is
method-side order evidence. All nonlocal rows remain bounded diagnostic rows with
open source-policy status, so the figure does not authorize an external
source-paper superiority claim.}
{fig:all-method-result-matrix}

\paperfig{Figure_13_work_precision_compendium.png}
{Work/precision compendium for fair candidate diagnostics. The first two panels
use the frictionless TFE source-pendulum candidate route on a shared $T=1$
grid and RK4 reference, plotting velocity error against wall time and Newton
work. The Algorithm-1-literal $T=10$ panels plot velocity and coordinate error
against Newton work from the Algorithm-1-literal work/precision diagnostic rows; these
fixed-$h$ endpoint rows remain source-policy open because exact-$T$ error
sampling and a source-equivalent DAE/method runner are not established. The
RA2021 panel summarizes same-window pendulum finest-row runtime and
velocity-error ratios relative to public $rA$ rows, and the closed-loop panel
collects strict common-reference four-link and slider-crank work/precision
rows. These panels consolidate existing evidence for review; they remain
candidate/common-reference diagnostics, not source-policy superiority claims.}
{fig:work-precision-compendium}

\begin{table}[htbp]
\centering
\caption{Diagnostic boundaries separated from the accepted theorem.}
\label{tab:diagnostic-boundaries}
\begin{tabular}{P{0.30\linewidth}P{0.40\linewidth}P{0.14\linewidth}}
\toprule
Diagnostic boundary & Current evidence & Claim boundary \\
\midrule
Independent full-\tfe{} replacement &
No accepted 132-row paper-derived weak residual with source-free terminal
closure and smooth order recovery. &
Boundary retained. \\
Sparse backend speed &
Sparse pattern and colors are quantified, but dense \texttt{jacfwd} is still
faster in the recorded wall-clock run. &
Backend caveat. \\
Sharp-friction coarse regime &
Coarse sharp-friction sweep is order-reduced; deeper fixed refinement recovers
high order at increased cost. &
Regime caveat. \\
\bottomrule
\end{tabular}
\end{table}

\section{Limitations}

Six limitations define the boundary of the claim, ordered from theorem and
method conditions to external reproduction boundaries.
\begin{enumerate}
  \item \textbf{Conditional proof boundary.} The accepted stage-residual input
  is closed for the implemented FullVA direct residual route: the 96
  non-dynamic implementation rows are supplied by the non-dynamic certificate
  and the 36 Newton--Euler dynamic rows are supplied by the D5
  direct-substitution certificate. The implementation-path binding is used
  only to identify the formulas differentiated by the active residual map;
  a term-by-term primitive verification of the dynamic block is kept as a
  separate sufficient route and is not an input to the accepted direct proof.
  Thus the active theorem input is the direct residual bridge: the
  same-branch residual-value input at the lifted Gauss stage \(Z_G\) is proved
  by that route, giving the stage-residual \(O(h^7)\) input used below, while separate
  primitive or symbolic dynamic-block certificates are not claimed. The theorem remains
  conditional on the smooth regularity tube and on an inexact nonlinear solve
  satisfying the stated compact-tube \(\eta_h^{\rm tube}\le c_\eta h^7\)
  policy; residual-to-error promotion for additional
  four-link/slider-crank dynamic-order claims remains outside the accepted
  theorem.
  This separation keeps the proof route distinct from broader promotion
  claims. The direct residual bridge supplies the local-defect input used by
  the theorem, and the proof certificates supply the accepted direct-route
  residual/Jacobian binding used by the bridge. Compact-branch regularity,
  implementation binding,
  endpoint perturbation, and solver scale remain explicit theorem-domain
  interfaces. Promoting closed-loop residual or reaction rows to trajectory
  order would require a separate branch-compatible residual-to-error theorem
  with a fixed residual operator, trajectory norm, stability or inf-sup
  inverse, nonlinear remainder control, reference-floor exclusion, and its own
  dynamic-order campaign. Small constraint, reaction, or mechanism residuals
alone do not provide that theorem. In the proof-interface terminology
retained in the theorem tables, P7 remains a separate output nonclaim/residual-to-error boundary: residual-only mechanism rows require a separate residual-to-error transfer theorem and are not a theorem premise.
  \item \textbf{Nonlinear tolerance and endpoint closure.} The asymptotic
  theorem requires a compact-tube nonlinear residual policy
  \(\eta_h^{\rm tube}\le c_\eta h^7\), while the
  reported production runs use a fixed finite-run tolerance. The current
  scaled-tolerance and tolerance-regime probes are finite reported-window
  diagnostics around this policy; they diagnose this scale but do not close
  it, prove an asymptotic solver contract, or convert fixed production
  tolerances into theorem hypotheses. Likewise, endpoint closure is a conditional
  local \(O(h^7)\) perturbation under Lemma~\ref{lem:endpoint-closure}
  and the retained P2 endpoint-ball/right-inverse hypotheses; the raw
  residual/correction rows do not prove those retained hypotheses and are
  not hidden source-paper or reference-solution projections.
  Closing P6 rather than retaining it would require a theorem-level solver
  result on the same compact tube: the Gauss-predictor branch would have to
  be selected for every accepted transition, the computed Newton iterate would
  have to remain in the strong local residual-inverse ball, and the
  \(h^7\)-scaled residual stopping rule would have to hold uniformly with
  constants independent of the reported grid and backend. Such a result would
  discharge the solver-scale interface only; it would still not promote
  residual-only mechanism rows unless the separate P7 residual-to-error
  transfer theorem were also proved.
  \item \textbf{Full source-paper residual replacement.} Independent
  full-\tfe{} stage replacement is not accepted. The accepted theorem is tied
  to the Gauss/FullVA residual in Eq.~\eqref{eq:residual}; replacing those rows
  by paper-derived weak rows is a separate method-construction problem.
  \item \textbf{Performance and scalability.} The sparse automatic
  differentiation structure is correct and quantified, but dense
  \texttt{jacfwd} remains faster in the recorded wall-clock evidence. The
  sparse backend is therefore an implementation diagnostic, not a claimed
  performance improvement. The current implementation is a small-system
  reference implementation; a production scalability claim would require an
  $N=2,4,8,16,32$ chain benchmark with wall-clock time, Newton iterations,
  Jacobian-assembly time, linear-solve time, conditioning or rank diagnostics,
  and memory scaling.
  \item \textbf{Sharp-friction coarse regime.} The practical coarse
  sharp-friction regime is order-reduced, although deeper fixed refinement
  recovers high order at additional cost. The method claim therefore rests on
  the smooth contract and treats sharp friction as a quantified caveat.
  \item \textbf{External same-test comparison.} The exact same-test external
  comparison criterion in Section~\ref{sec:cross-paper-boundary} remains open
  after the public-code baselines and selected/coarse \method{} pilots.
  The Chaturvedi--Sandu--Sandu pendulum tests are not source-policy
  reproduced: the encoded source-shaped candidate smokes and bounded probes
  remain diagnostic; the Kissel/Negrut suites still lack complete \method{}
  dynamic order/work rows under the published policies; and the closed-loop
  surrogate rows do not change that accepted-order boundary. The floor
  analysis narrows this boundary to velocity/acceleration floor evidence plus
  a position/reference floor, but it does not close the source-policy
  dynamic-order comparison. The velocity-partitioning public-code path is also
  unresolved.
\end{enumerate}
These limitations do not invalidate the narrower conditional formal-order
comparison; they define the proof, implementation, and reproduction work still
needed before a broader paper-level claim would be justified.

\section{Conclusions}

This paper has presented a monolithic Lie-group Gauss6/FullVA integrator for
lower-pair mechanisms. The method combines SO(3) retractions, lower-pair
position/velocity/acceleration rows, and Newton--Euler dynamics in one
square stage residual, then reconstructs the endpoint on the accepted
branch. Under the stated compact-branch, endpoint/stability,
implementation-binding, and solver-scale contracts, this construction gives the
conditional sixth-order Lie-group FullVA theorem of
Theorem~\ref{thm:g6fullva-order}. The conditional qualifier fixes the compact
branch, endpoint, implementation-binding, and solver-scale domain of the
theorem; it is not an empirical replacement for the local-defect proof. The
numerical evidence records finite-window smooth position and velocity orders
7.161 and 7.066 as consistency evidence after the order-six theorem has been
stated. The single- and double-pendulum rows provide
method-side dynamic-order evidence, while the four-link and slider-crank rows
provide diagnostic mechanism-coverage evidence, not accepted dynamic-order
proof and not residual-to-error evidence, for the same FullVA residual.
Against the
local order-five $m=3$ Gauss--Lobatto temporal finite-element target, this is
a conditional formal-order comparison, not an external same-test superiority
claim.

The proof route is anchored by the implemented 132-row residual bridge: the
96 non-dynamic rows have an \(O(h^7)\) certificate on the lifted Gauss stage,
and the 36 Newton--Euler rows close by direct substitution with zero
residual on the lifted Gauss stage \(Z_G\). That bridge is then combined with
the analytic local-defect decomposition, endpoint perturbation, the retained
compact-tube inexact-Newton scale condition, and local-to-global transfer. In
that reading, the theorem is a conditional implication for a fixed smooth
branch: once the compact branch, endpoint map, implementation binding, and
\(\eta_h^{\rm tube}\le c_\eta h^7\) solver-scale condition are in force, the
direct 132-row bridge supplies the \(O(h^7)\) residual-defect input; the
stage-residual perturbation, endpoint, and retained solver-scale lemmas then
supply the \(O(h^7)\) local defect and the Gauss
stability transfer gives the \(O(h^6)\) same-initial-state reported-grid
estimate. The finite-window order
fits, finite solver probes, and residual/reaction tables are consistency or
diagnostic evidence and do not replace those mathematical inputs. The separate
primitive/Taylor decomposition of the Newton--Euler rows remains a stricter
separate sufficient-route specification and is not used to supply \(C_{\rm loc}\), close the
solver-scale interface, close the residual-to-error boundary, or act as a
theorem input. The unproved primitive-route inputs are therefore a non-active
replacement-certificate boundary for the dynamic-block provenance, not an
open algebraic term in the accepted direct residual-bridge proof. Remaining
work concerns a separate
uniform same-branch solver-policy theorem, residual-to-error promotion for
closed-loop residual rows, and external same-test reproduction; these are
outside the theorem proved here.

Broader comparison work remains separate from this method result. It should
rerun the original pendulum tests and the Kissel/Negrut public-code mechanism
suites under identical benchmark definitions, reporting order, error,
constraint drift, Newton work, and runtime side by side. The
velocity-partitioning code path is part of that separate published-test baseline
scope. Complete source-paper temporal finite-element residual replacement and
a runner-centered full reproduction package remain outside the accepted claim.

\section*{CRediT authorship contribution statement}

Jingquan Wang: Conceptualization, Methodology, Software, Validation, Formal
analysis, Writing - original draft, Writing - review and editing.

\section*{Declaration of competing interest}

The author is affiliated with the University of Wisconsin--Madison. The author
declares no competing interests relevant to this work.

\section*{Funding}

No external funding was received for this research.

\section*{Data availability}

The journal source archive is the LaTeX archive: it contains the manuscript
sources, figures, and declarations. This is an editorial-source archive rather
than a full source-policy runner archive.
The companion reproducibility package contains the reported-result records and
read-only check scripts used for this paper. It is scoped to the accepted
method-side evidence and proof-dependency records; external same-test runner
reproduction for source papers and full \tfe{} runner comparisons are reserved
for a separate package. An archival public repository or DOI can be attached
at journal production or revision if the editorial office requires a public
deposit.

\section*{Declaration of generative AI and AI-assisted technologies in the writing process}

During preparation of this work, the author used an AI coding assistant for
drafting support, consistency checking, and validation-script editing. The
author reviewed and edited the content and takes full responsibility for the
published article.

\appendix

\section{Cross-paper same-test comparison criterion: diagnostic appendix}
\label{sec:cross-paper-boundary}

This subsection is a diagnostic guardrail, not a second contribution claim
and not part of the proof chain. It can be read as a promotion rule for separate
external source-paper claims after the method definition, proof, and accepted
method-side validation have already been fixed. A publication-level
source-paper comparison must be made on the same numerical tests used by the
source papers, not only on mechanisms with similar names. The criterion is
strict:
the body parameters, initial conditions, prescribed drivers, friction model,
step-size sweep, nonlinear tolerance policy, reference solution, error norm,
and work metric must be matched or explicitly justified. Table
\ref{tab:cross-paper-boundary} records the required external comparisons, the
status of this evidence package, and the claim that is allowed before
the remaining source-suite decisions close. The current paper uses these rows
only to prevent overclaiming: the method claim, the conditional theorem, and
the accepted numerical evidence do not depend on source-policy promotion.

\paragraph{Why the source-policy criterion is separate}
The source-policy criterion is required only for a stronger external
apples-to-apples statement: that \method{} has been evaluated under the
source papers' own model, output, tolerance, reference, and work policies and
would then be eligible for a separate source-paper replacement or superiority
comparison.
It is not a premise of the Gauss6/FullVA residual definition, the retained
conditional sixth-order theorem, or the accepted local dynamic-order and
mechanism-coverage evidence.  Therefore an open source-policy row does not
weaken the bounded method/proof claim; it only blocks the broader external
same-test and full reproducibility package claims.

\paragraph{Source-policy progress boundary}
The detailed source-policy record is boundary evidence only. It records local
candidate availability, public-code setup progress, \tfe{} formula diagnostics,
and open runner gaps, but no external row is promoted to a source-policy
dynamic-order or source-paper superiority claim.
The current progress entries are read only as diagnostic inventory: public-code
baseline order/timing rows, source-identity checks, coarse double-pendulum
self-reference runs, bounded four-link/slider-crank mechanism-diagnostic candidates, and
published-formula candidates record implementation progress but do not
establish external same-test reproduction.  No external source-policy
dynamic-order row is accepted, and the local method rows under external
source-policy dynamic-order rules remain unpromoted; therefore no
source-paper superiority, full source-policy work-precision, or full-\tfe{}
replacement claim is made.

\begin{table}[htbp]
\centering
\caption{Required same-test comparisons and current external-suite claim boundary.}
\label{tab:cross-paper-boundary}
\scriptsize
\renewcommand{\arraystretch}{0.88}
\setlength{\tabcolsep}{3pt}
\begin{tabular}{P{0.21\linewidth}P{0.24\linewidth}P{0.22\linewidth}P{0.19\linewidth}}
\toprule
External source & Published or code-backed tests & Current evidence & Claim allowed now \\
\midrule
Chaturvedi--Sandu--Sandu Lie-group \tfe{} paper
\cite{chaturvedi2026tfe} &
Rigid pendulum without and with joint friction; $m=1,2,3$ \tfe{};
Newmark--$\beta$ and trapezoidal comparators. &
Specification and source-shaped candidate smokes are encoded, including source
parameter/output policies and bounded \tfe{} probes; source-policy
absolute-coordinate DAE/method-runner rows are outside the current evidence
scope. &
Formula-order comparator only; no implemented source-paper superiority claim. \\
Kissel--Taves--Negrut absolute-coordinate $rA/rp/r\epsilon$ suite
\cite{taves2020exponential,kissel2021dwelling,kissel2022absolute} &
Single pendulum, double pendulum, slider-crank, and four-link mechanisms from
the public C1/C2 Python code and model files. &
Public baseline rows and timing rows are complete. \method{} has
single-pendulum public rows, double-pendulum coarse rows, and closed-loop
strict-common-reference work/precision rows, but accepted local
source-policy dynamic-order evidence is not promoted for any of the four
tracked examples. &
Coarse-first and common-reference evidence only; no external superiority
claim. \\
Fang--Kissel--Zhang--Negrut half-implicit suite
\cite{fang2023halfimplicit} &
Double pendulum, slider-crank, four-link, $N$-pendulum scaling, and
slider-crank joint-friction examples. &
Bounded $T=0.1$ $rA/rA_{\mathrm{half}}$ pilot complete; full $T=8$ policy and
\method{} rows open. &
Bounded pilot only; not a full half-implicit policy comparison. \\
Kissel--Bakke--Negrut velocity-partitioning Lie-group ODE suite
\cite{kissel2024velocitypartitioning} &
Single pendulum, double pendulum, four-link, and slider-crank tests reported
for velocity partitioning. &
The coordinate-partitioning proxy is identified through the implemented $rA$
wrapper; a distinct public VP code path remains unresolved and demoted. The
source-code path is not separately reproduced. &
Common-reference diagnostic rows only; no source-policy VP superiority claim. \\
\bottomrule
\end{tabular}
\end{table}

\noindent\begin{minipage}{0.98\linewidth}
\paragraph{Claim-policy implication}
The source-policy promotion criterion is interpreted as a logical promotion rule, not as a
visual preference over plots.  A paper-level external-superiority statement
requires the source benchmark, method-set, numerical-control, and work-metric
rows in Table~\ref{tab:tfe-same-test-protocol} to close for the
source-policy campaign being cited.  The current evidence record reports no
strict external error-claim rows, does not establish any source-policy-closed
external same-test error row, and leaves the same-test promotion conditions
open.
Therefore the common-reference matrix and bounded pilots may be used only as
diagnostics and cannot be used as premises in any implemented source-paper
superiority claim.
\end{minipage}\par

\begin{table}[htbp]
\centering
\caption{Minimum same-test \tfe{} acceptance protocol before any source-paper superiority claim.}
\label{tab:tfe-same-test-protocol}
\scriptsize
\setlength{\tabcolsep}{3pt}
\renewcommand{\arraystretch}{0.88}
\begin{tabular}{P{0.20\linewidth}P{0.47\linewidth}P{0.21\linewidth}}
\toprule
Promotion item & Required evidence & Evidence boundary \\
\midrule
Source benchmarks &
Reproduce the Chaturvedi--Sandu--Sandu frictionless pendulum and revolute
joint with Brown--McPhee friction using the source parameters, horizon,
initial data, drivers, and output variables. &
Specification and reference-feasibility diagnostics only; source rows not
accepted. \\
Method set &
Run source-equivalent \tfe{} $m=1,2,3$, Newmark--$\beta$, trapezoidal, and
\method{} on the same trajectories. &
Local formula/common-reference diagnostics only; no accepted source-policy
method matrix. \\
Numerical controls &
Use the same step sizes or a justified common refinement, matched nonlinear
tolerance policy, matched reference solution, and the same error norm and
constraint-drift reporting. &
Open. \\
Work metrics &
Report error versus wall-clock time, Newton iterations, Jacobian assembly
time, linear-solve time, and failed-step counts. &
Open. \\
Allowed claim &
Only after the rows above pass may the manuscript state a source-paper
accuracy or work-precision advantage. Until then, the \tfe{} comparison is a
formal-order-level comparison only; diagnostic rows remain non-claim evidence. &
No external-superiority claim. \\
\bottomrule
\end{tabular}
\end{table}

Consequently, the present numerical evidence supports only the method-side
Gauss/FullVA order claim and the local order-five \tfe{} formula-target
comparison. The method-side evidence consists of the smooth cylindrical-chain
benchmark plus the four ASME-style mechanisms in this package; the coverage
matrix records 48 rows and 32 completed rows, but no external superiority
claim. It does not support a stronger statement that \method{} is
better than the published Chaturvedi--Sandu--Sandu implementation or the
Kissel/Negrut public-code baselines on their own numerical tests. The extracted
public-code benchmark specification and row-level inventory are maintained in
the reproducibility appendix. The current evidence state is: specification
extracted; 2021 public baselines complete for the already extracted
order/timing evidence; 2021 public baseline order/timing evidence closed at
$12/12$ public order groups and $12/12$ timing rows; selected/coarse
\method{} pilot rows; a coarse public double-pendulum same-window baseline
run; and a single-pendulum coarse same-window tranche. The accepted local
\method{} source-policy dynamic-order set is nevertheless empty across the
four tracked examples: the double-pendulum row is coarse-first rather than the
public $10^{-4}$ policy,
and the four-link/slider-crank public-horizon rows are kinematic/reaction
residual rows rather than accepted dynamic-order rows. Because a separate
residual-to-error transfer theorem remains open, residual-only rows are not promoted to
order claims. The external same-test campaign remains incomplete:
full source-policy dynamic-comparison rows remain unaccepted, and the paper
makes no external-superiority claim.
Equivalently, the source-policy dynamic-order count is nevertheless 0/4 for
the external source-paper rows; this count is a progress boundary, not a
demotion of the accepted local dynamic-order evidence.

After the coarse common-reference replay was repaired to use fixed public-code
times $t_i=i h$, the result package also records a bounded
common-reference matrix. All 44 common-reference summary rows pass the
row-level policy checks, and the global comparison policy passes 14/14
checks, but a stricter all-example consistency check allows zero paper-level
external direct-error claims. Thus the replay is a shared coarse
final-state diagnostic, not a reproduction of the source papers' default
time horizons, output policies, error norms, or $10^{-4}$ policies.

\paragraph{Benchmark-policy dictionary}
The bounded common-reference policy used in
Tables~\ref{tab:common-reference-pack}--\ref{tab:common-reference-all-methods}
fixes $T=0.1$, $h\in\{0.1,0.05,0.025\}$, and
$h_{\rm ref}=0.0125$ for every runnable method.  The single-pendulum row uses
the analytic exact final state and a final $L^2$ error.  The double-pendulum
row uses the local \method{} $h_{\rm ref}=0.0125$ final state and a final
$L^\infty$ error.  The four-link and slider-crank rows use closed-loop
true-dynamic common-reference trajectories at $h_{\rm ref}=0.0125$ and then
apply the same order/error arithmetic to all runnable methods.  Source-policy
rows, by contrast, must follow each source paper or public code on its own
horizon, time grid, reference solution, error norm, tolerance, and output
variables.  This dictionary is the reason the common-reference table can be
arithmetically reconciled while source-policy external-superiority rows remain
open.
Table~\ref{tab:common-reference-pack} summarizes the nearest nonlocal
bounded common-reference error evidence, and
Table~\ref{tab:common-reference-all-methods} reports the full runnable
method-by-example order/error matrix.

We therefore separate two comparison statements. The finite-grid
common-reference diagnostic is assembled for arithmetic coverage of the
required matrix: all required method labels are accepted, alias-resolved, or
source-backed scope-excluded, and the 44-cell table is recomputed from 132 raw
rows with zero summary mismatches. The nonlocal cells remain bounded
common-reference diagnostics because they do not reproduce every source
paper's default horizon, tolerance, output policy, or error norm. These
diagnostic comparisons are not promoted to a CMAME paper-level
external-superiority claim. The stronger source-policy statement remains a
non-claim, and no common-reference count is promoted to a source-paper
superiority claim.

This paragraph separates local-candidate availability from source-policy
closure.  For
the 2021 absolute-coordinate suite, the public-code path, the
\texttt{body.r}/\texttt{body.dr}/\texttt{body.ddr} output map, and the endpoint
time-grid convention have been extracted and checked; in that bounded sense,
source identity is now closed only for the diagnostic row map, not for
source-policy closure.  The remaining gap is
promotion of local \method{} rows under the
same source-policy dynamic-order contract. Table~\ref{tab:ra2021-local-promotion-boundary}
records why the existing local candidates are not promoted: the single-pendulum
public-grid tranche is floor limited, the double-pendulum row is a coarse
self-reference run, and the four-link/slider-crank true-dynamic rows are
bounded $T=0.1$ diagnostics rather than the public $T=3$ source-policy rows.
For the 2022 half-implicit suite, the bounded $T=0.1$ replay has 24/24
successful rows, while the $T=8$ coarse replay has 18/24 successful rows and
only four of eight complete form--model groups.  For the original TFE pendulum
suite, bounded runner and endpoint-sensitivity diagnostics exist, but a
source-equivalent DAE/friction/output runner has not completed the original
rows.  Consequently no local method row is accepted under the external
source-policy dynamic-order contract, and no source-policy-closed external
same-test error row is established.

\begin{table}[htbp]
\centering
\caption{RA2021 local-candidate promotion boundary.}
\label{tab:ra2021-local-promotion-boundary}
\scriptsize
\setlength{\tabcolsep}{3pt}
\renewcommand{\arraystretch}{0.88}
\begin{tabular}{P{0.18\linewidth}P{0.38\linewidth}P{0.34\linewidth}}
\toprule
Example & Local candidate evidence & Why not promoted to RA2021 source-policy row \\
\midrule
Single pendulum &
Public $T=3$ row trio at $h=\{10^{-2},10^{-3},10^{-4}\}$ with
$h_{\rm ref}=10^{-3}$; finest velocity error is $7.01\times10^{-15}$. &
The row is roundoff/floor limited; observed order is not a stable
source-policy dynamic-superiority certificate. \\
Double pendulum &
Coarse $T=3$ self-reference row with $h=\{0.1,0.05,0.025\}$,
$h_{\rm ref}=0.0125$, and position/velocity order $7.951/7.042$. &
The step grid and reference do not match the RA2021 double-pendulum policy
$h=\{0.01,0.002,0.001\}$ with $h_{\rm ref}=10^{-4}$. \\
Four-link &
True-dynamic local row at $T=0.1$, $h=\{0.1,0.05,0.025\}$, with
position/velocity order $5.955/6.085$. &
The public $T=3$ rows are kinematic/reaction checks; the true-dynamic row is
not the RA2021 source-policy horizon. \\
Slider-crank &
True-dynamic local row at $T=0.1$, $h=\{0.1,0.05,0.025\}$, with
position/velocity order $6.164/7.341$. &
The public $T=3$ rows are kinematic/reaction checks; the true-dynamic row is
not the RA2021 source-policy horizon. \\
\bottomrule
\end{tabular}
\end{table}

\begin{table}[htbp]
\centering
\caption{Bounded common-reference velocity evidence from the current result pack; diagnostic only, not external-superiority evidence.}
\label{tab:common-reference-pack}
\scriptsize
\setlength{\tabcolsep}{3pt}
\renewcommand{\arraystretch}{0.88}
\begin{tabular}{P{0.18\linewidth}P{0.17\linewidth}P{0.17\linewidth}P{0.23\linewidth}P{0.17\linewidth}}
\toprule
Example & Local velocity order & Local finest velocity error & Nearest nonlocal method & Nearest nonlocal finest error \\
\midrule
Single pendulum & 5.801 & $1.155\times10^{-12}$ &
velocity partitioning $rA$ & $3.163\times10^{-11}$ \\
Double pendulum & 6.010 & $2.346\times10^{-13}$ &
2021 public $rA$ & $4.924\times10^{-5}$ \\
Four-link & 6.085 & $1.093\times10^{-11}$ &
velocity partitioning $rA$ & $1.087\times10^{-8}$ \\
Slider-crank & 7.341 & $9.189\times10^{-12}$ &
velocity partitioning $rA$ & $1.332\times10^{-7}$ \\
\bottomrule
\end{tabular}
\end{table}

\begin{table}[htbp]
\centering
\caption{All runnable common-reference velocity rows.  Each entry is observed velocity order / finest-step velocity error on the fixed-grid $h=\{0.1,0.05,0.025\}$ sweep with reference $h=0.0125$; the table is a diagnostic common-reference matrix, not a source-policy superiority table.}
\label{tab:common-reference-all-methods}
\scriptsize
\setlength{\tabcolsep}{3pt}
\renewcommand{\arraystretch}{0.88}
\resizebox{\linewidth}{!}{%
\begin{tabular}{lllll}
\toprule
Method & Single pendulum & Double pendulum & Four-link & Slider-crank \\
\midrule
\method{} &
5.801 / $1.155\times10^{-12}$ &
6.010 / $2.346\times10^{-13}$ &
6.085 / $1.093\times10^{-11}$ &
7.341 / $9.189\times10^{-12}$ \\
HI 2022 $rA$ &
1.005 / $7.725\times10^{-2}$ &
1.049 / $4.944\times10^{-5}$ &
1.015 / $1.970\times10^{-1}$ &
0.941 / $3.281\times10^{-2}$ \\
HI 2022 $rA_{\mathrm{half}}$ &
1.841 / $7.725\times10^{-2}$ &
-0.799 / $2.970$ &
1.590 / $1.970\times10^{-1}$ &
2.459 / $3.281\times10^{-2}$ \\
2021 public $rA$ &
0.617 / $1.305\times10^{1}$ &
1.052 / $4.924\times10^{-5}$ &
1.015 / $1.970\times10^{-1}$ &
0.941 / $3.281\times10^{-2}$ \\
2021 public $r\epsilon$ &
0.617 / $1.305\times10^{1}$ &
1.052 / $4.924\times10^{-5}$ &
2.894 / $1.458\times10^{-2}$ &
0.940 / $3.281\times10^{-2}$ \\
2021 public $rp$ &
0.517 / $1.499\times10^{1}$ &
0.361 / $5.241\times10^{-5}$ &
2.894 / $1.458\times10^{-2}$ &
0.940 / $3.281\times10^{-2}$ \\
TFE Newmark--$\beta$ &
0.901 / $1.479\times10^{1}$ &
1.787 / $4.246\times10^{-4}$ &
1.872 / $1.672\times10^{-2}$ &
2.051 / $1.531\times10^{-3}$ \\
TFE $m=1$ &
1.021 / $1.498\times10^{1}$ &
1.996 / $2.708\times10^{-4}$ &
3.866 / $5.253\times10^{-4}$ &
2.481 / $6.264\times10^{-4}$ \\
TFE $m=2$ &
0.054 / $1.102\times10^{2}$ &
1.870 / $8.711\times10^{-5}$ &
1.021 / $6.711\times10^{-4}$ &
2.290 / $9.881\times10^{-5}$ \\
TFE trapezoidal &
1.024 / $1.498\times10^{1}$ &
2.000 / $2.708\times10^{-4}$ &
3.888 / $5.306\times10^{-4}$ &
2.329 / $7.328\times10^{-4}$ \\
VP 2024 coordinate partitioning $rA$ &
0.000 / $3.163\times10^{-11}$ &
1.000 / $2.165\times10^{-3}$ &
-0.001 / $1.087\times10^{-8}$ &
-0.000 / $1.332\times10^{-7}$ \\
\bottomrule
\end{tabular}}
\end{table}

The full common-reference table is arithmetic-consistent but not, by itself, a
source-policy reproduction. An independent recomputation from the 132 raw
common-reference rows reproduces all 44 summary order/error entries with zero
summary mismatches. The remaining concern is interpretive rather than
arithmetical: 15 nonlocal method/example rows are flagged for source-policy
review, and 10 of those rows have small position error but large velocity error.
This pattern is consistent with a velocity-output, state-map, time-grid, or
reference-policy mismatch and therefore cannot support an external-superiority
claim. Table~\ref{tab:source-policy-diagnosis} records the resulting
claim boundary.

The flagged-row count should not be read as the complete set of open external
claim rows. All 40 nonlocal method/example rows remain bounded
common-reference diagnostics rather than source-policy-closed reproductions;
the 15 flagged rows are the anomaly or high-risk subset that requires special
attention before any separate external-superiority statement is considered. This
distinction keeps the order/error comparison and the source-policy claim
separate throughout the paper.

\begin{table}[htbp]
\centering
\caption{Source-policy diagnosis for flagged common-reference baseline rows.}
\label{tab:source-policy-diagnosis}
\scriptsize
\setlength{\tabcolsep}{3pt}
\renewcommand{\arraystretch}{0.88}
\begin{tabular}{P{0.31\linewidth}P{0.13\linewidth}P{0.21\linewidth}P{0.26\linewidth}}
\toprule
Diagnostic category & Rows & Affected examples & Claim-boundary interpretation \\
\midrule
Position aligned, velocity mismatched & 10 &
single pendulum, four-link &
Likely velocity-output, state-map, or reference-policy mismatch; not a
method-failure claim. \\
Single-pendulum velocity-output policy suspect & 7 &
single pendulum &
Original/Taves--Kissel--Negrut/TFE pendulum rows require source setup and
velocity-output policy closure. \\
Near reference floor; order not identifiable & 3 &
single pendulum, four-link, slider-crank &
Errors are too close to the shared reference floor for a meaningful observed
order statement. \\
Low order but small error; wider step window needed & 4 &
single pendulum, double pendulum, four-link, slider-crank &
The row is diagnostic only until a wider or source-policy step window is
reported. \\
Negative order with large error & 1 &
double pendulum &
The fixed-grid replay or selected form requires source-policy recheck before
comparison. \\
\bottomrule
\end{tabular}
\end{table}

The all-example source-policy check expands this diagnosis from category
counts to the actual method/example coverage. It checks all 15 flagged summary
rows and their 45 raw coarse rows at $h=\{0.1,0.05,0.025\}$. The flagged rows
cover all four examples: 8 single-pendulum rows, 2 double-pendulum rows,
4 four-link rows, and 1 slider-crank row. Table~\ref{tab:all-example-source-policy}
records the suite-level closure boundary. A stricter all-method/all-example
consistency check then checks the complete 44-cell common-reference table and its
132 raw rows; it keeps all 44 cells as bounded diagnostics but allows zero
strict external error-claim rows. In particular, the common-reference replay is
valid for the finite-grid diagnostic, but no suite has yet closed the stronger
source-policy reproduction required for an external-superiority claim.

\begin{table}[htbp]
\centering
\caption{All-example source-policy status for the flagged baseline rows.}
\label{tab:all-example-source-policy}
\scriptsize
\setlength{\tabcolsep}{3pt}
\renewcommand{\arraystretch}{0.88}
\begin{tabular}{P{0.22\linewidth}P{0.10\linewidth}P{0.33\linewidth}P{0.25\linewidth}}
\toprule
Suite & Rows & Common-reference row source & Closure boundary \\
\midrule
2021 absolute-coordinate public suite & 5 &
Public setup and $rA/rp/r\epsilon$ steppers, fixed-grid replay, and
\texttt{body.r/body.dr/body.ddr} output fields. &
Source identity is checked; promotion still requires local \method{}
source-policy dynamic rows, norm/output binding, runtime/Newton binding, and
an independent verification record. \\
Original \tfe{} pendulum formulas & 4 &
Surrogate common-reference row only: RA2021 public setup with locally installed
Newmark and \tfe{} formula proxies, not the source pendulum. &
Encode the original pendulum setup, source error norm, Brown--McPhee friction
law, output variables, and horizon, or keep these rows demoted. \\
HI 2022 half-implicit suite & 3 &
Public setup and half-implicit steppers under bounded fixed-grid replay. &
Run the full $T=8$ public policy or explicitly demote the bounded pilot rows. \\
VP 2024 coordinate partitioning & 3 &
Coordinate-partitioning proxy on the RA2021 public setup. &
Resolve the independent velocity-partitioning code path or keep the row as
proxy/common-reference evidence only. \\
\bottomrule
\end{tabular}
\end{table}

For the \tfe{} source evidence, we separate the friction evidence into two
levels.  The source text identifies the Brown--McPhee velocity-based
continuous friction model, reports $\mu_s=0.3$ and $\mu_d=0.2$, and uses
$h_{\rm ref}=10^{-4}$ for source reference solutions; its detailed
continuous-friction law is delegated to the source-paper friction references
\cite{chaturvedi2026tfe,browmcphee2016friction}.  The model check therefore
treats the encoded Brown--McPhee torque law as a candidate published formula
structure, not as a source-code-equivalent friction implementation.  The
transition-velocity, normal-load coupling, friction Jacobian, and DAE
multiplier policies remain source-policy boundaries.

\noindent Source-policy boundary summary:\par
\noindent Brown--McPhee velocity-based continuous friction model;\par
\noindent candidate published formula, not source-code-equivalent friction implementation;\par
\noindent source-policy boundaries;\par
\noindent Brown--McPhee formula is not fixed;\par
\noindent source-code-equivalent Brown--McPhee transition law;\par
\noindent transition velocity vt;\par
\noindent velocity scale in the v/vt Stribeck factor;\par
\noindent not the normal-load/multiplier coupling;\par
\noindent implementation choice, not source-policy evidence;\par
\noindent candidate-friction DAE trajectory contract records 12 method-grid rows and 56 step-residual rows with nonpositive friction power;\par
\noindent do not implement a monolithic source-policy DAE runner;\par
\noindent source-policy runner acceptance contract has six simultaneous parts;\par
\noindent transition law, including transition velocity;\par
\noindent velocity-based continuous friction model;\par
\noindent row-promotion rules binding error, order, runtime, and work
metrics;\par
\noindent monolithic absolute-coordinate DAE time integrator;\par
\noindent not an accepted source-policy reproduction.\par

\noindent
More concretely, the Brown--McPhee formula is not fixed by
\(\mu_s\), \(\mu_d\), and \(R\) alone.  It also depends on a
transition velocity \(v_t\), the velocity scale in the \(v/v_t\)
Stribeck factor, and on how the tangential velocity and normal reaction
are supplied by the constrained DAE.  The source-pendulum text records
\(\mu_s\), \(\mu_d\), \(R\), and \(h_{\rm ref}\), but not \(v_t\), not
the normal-load/multiplier coupling, and not the friction-Jacobian
convention.  Choosing the local surrogate's Stribeck velocities
\(0.05\), \(0.5\), or \(1.0\) would therefore be an implementation
choice, not source-policy evidence.
A bounded transition-velocity sensitivity check makes this boundary numerical
rather than merely documentary: over Stribeck velocities \(0.05\), \(0.5\),
and \(1.0\), with \(0.5\) as the baseline, the \(T=0.024\),
\(h_{\rm ref}=10^{-4}\) candidate DAE check changes the endpoint by at most
\(2.90\times 10^{-6}\) in coordinate and \(2.36\times 10^{-4}\) in velocity.
The check records three sensitivity rows, 36 contract rows, and zero
source-policy rows, so the missing transition velocity is treated as a
material unresolved source-law parameter rather than an innocuous
transcription detail.
The positive evidence from the surrogate-friction DAE layer is narrower:
the candidate-friction DAE trajectory contract records 12 method-grid rows
and 56 step-residual rows, all finite, with candidate DAE residuals below
\(10^{-9}\) and nonpositive friction power.  These checks certify local
dissipativity and absolute-coordinate residual consistency for the encoded
candidate law only.  They do not certify a source-code-equivalent
Brown--McPhee law, do not implement a monolithic source-policy DAE runner,
and do not close any source-policy row.

Consequently, the source-policy runner acceptance contract has six
simultaneous parts.  A promoted \tfe{} row would need a
source-code-equivalent Brown--McPhee transition law, including transition
velocity, normal-load coupling, friction Jacobian, and multiplier policy; a
monolithic absolute-coordinate DAE time integrator rather than a stepwise
residual replay; source grid, output norm, and endpoint/output-sampling
binding for \(T=10\) and \(h_{\rm ref}=10^{-4}\); nonlinear solve, Newton,
and tolerance-policy binding; method-runner equivalence for the \tfe{}
\(m=1,2,3\), Newmark--\(\beta\), trapezoidal, and local \method{}
comparator rows; and row-promotion rules binding error, order, runtime, and
work metrics.  Equivalently, the present package leaves two documentary
blocks open without requiring a new long numerical campaign---the
source-code-equivalent Brown--McPhee law and the full-\(T=10\)
endpoint/output-sampling policy for endpoint-incompatible step sizes---and
four execution-level blocks open: a monolithic source-policy DAE integrator,
source-policy method runners for the \tfe{} and one-step comparator rows, a
Gauss6/FullVA source-pendulum DAE comparator, and accepted work/precision row
binding.  Because these blocks are not jointly present, the candidate
DAE trajectory evidence remains a diagnostic record of what is missing, not an
accepted source-policy reproduction row.

The source check also includes the Appendix-B coefficient formulas for
\tfe{} $m=1$, $m=2$, and $m=3$ Gauss--Lobatto against the source equations
(40), (41), and (43).  The recorded coefficient certificate has three checked
rows and zero maximum absolute difference for the extracted source parameter
choices.  This closes coefficient-transcription evidence only; it does not
close source-policy method-runner equivalence or any source-policy reproduction
row.
Here, Appendix-B coefficient formulas denotes this coefficient-transcription
certificate only, not a closed source-policy method runner.
The diagnostic chain also includes a full-$T=10$ coarse formula probe for
the Appendix-B $m=3$ Gauss--Lobatto target.  That probe is finite and
residual-converged for one checked row and records the formal fifth-order
target, but it remains a non-source-policy formula diagnostic because it does
not invoke the source $h=10^{-4}$ reference campaign.
Restated as a source-policy boundary, the full-$T=10$ coarse formula probe
targets the Appendix-B $m=3$ Gauss--Lobatto target, records a formal
fifth-order target, and does not invoke the source $h=10^{-4}$ reference
campaign.
The source-reference side is therefore checked separately.  A full-$T=10$
frictionless source-reference feasibility probe integrates the extracted
source pendulum with $h_{\rm ref}=10^{-4}$ and compares it against a
$h=5\times10^{-5}$ check run, using 100000 and 200000 RK4 steps,
respectively.  The $h_{\rm ref}=10^{-4}$ final state is
$(\theta,\omega)=(-2.8790156956337563,-1.2836123173308356)$, and the
check-run mismatches are $3.82\times10^{-14}$ in coordinate and
$2.49\times10^{-13}$ in velocity.  This closes only source-reference
feasibility for the original pendulum time horizon.  It is not
source-policy method-runner equivalence for Newmark, trapezoidal, or
\tfe{}, it does not cover the Brown--McPhee friction row, and it does not
authorize external-superiority claims; no source-policy row is promoted.

The same source-pendulum candidate route now includes a \method{} row on
the extracted source output policy.  On the frictionless $T=1$ run with
$h_{\rm ref}=2.5\times 10^{-4}$ and $h=\{0.1,0.05,0.025\}$, the coordinate,
velocity, and Frobenius pairwise orders are
$(5.98,6.10)$, $(6.02,6.03)$, and $(5.98,6.10)$, respectively, with maximum
Newton residual below $10^{-12}$.  This is same-test candidate evidence for
the Gauss collocation order target on the source pendulum output policy.  The
candidate frictional row on the same grid shows reduced pairwise orders
(about $1.6$--$2.3$), so it is recorded only as a friction-feature diagnostic.
Neither row is an absolute-coordinate FullVA DAE source-policy reproduction,
and the \tfe{} source-policy closure count remains zero.
A separate same-test work/precision diagnostic runs Newmark--$\beta$,
trapezoidal, \tfe{} $m=1,2,3$, and \method{} on this same frictionless
source-pendulum $T=1$ grid.  It records 18 raw rows and six method summaries
with runtime and Newton-iteration work proxies; the diagnostic velocity-order
floors are about $6.02$ for the \method{} candidate and $5.19$ for the
\tfe{} $m=3$ Gauss--Lobatto formula target.  This diagnostic improves the
fair same-grid comparison, but it remains candidate-level evidence: the
source-policy method-runner equivalence flag and the external-superiority
claim flag are both false.

For the original \tfe{} pendulum suite, the source-model evidence also
includes a bounded active-row candidate smoke test. It uses the extracted
pendulum parameters and source output/error policy with
$h_{\rm ref}=10^{-4}$, but only on the short frictionless interval
$T=0.024$ and the source-shaped grid
$h=\{0.012,0.006,0.003\}$. The row is therefore useful for checking whether
the installed Newmark, trapezoidal, and \tfe{} formula runners have coherent
local order/error behavior, but it is not the full $T=10$ original-source
campaign and it closes zero source-policy rows. Table~\ref{tab:tfe-active-b2-smoke}
records this bounded diagnostic.

\begin{table}[htbp]
\centering
\caption{Bounded active-row candidate smoke for the original \tfe{} pendulum rows.  The run uses $T=0.024$, $h_{\rm ref}=10^{-4}$, and $h=\{0.012,0.006,0.003\}$; it is not a full $T=10$ source-policy reproduction.}
\label{tab:tfe-active-b2-smoke}
\scriptsize
\setlength{\tabcolsep}{3pt}
\renewcommand{\arraystretch}{0.88}
\resizebox{\linewidth}{!}{%
\begin{tabular}{P{0.20\linewidth}P{0.10\linewidth}P{0.20\linewidth}P{0.20\linewidth}P{0.13\linewidth}P{0.13\linewidth}P{0.18\linewidth}}
\toprule
Original-row method & Formula target & Velocity pair orders & Coordinate pair orders & Finest velocity error & Finest coordinate error & Claim boundary \\
\midrule
Newmark--$\beta$ & 2 & 1.973, 1.997 & 2.060, 2.022 & $1.650\times10^{-10}$ & $2.584\times10^{-12}$ & Bounded candidate smoke only \\
\tfe{} $m=1$ & 1 & 1.934, 1.915 & 1.000, 1.000 & $1.856\times10^{-10}$ & $5.742\times10^{-7}$ & Bounded candidate smoke only \\
\tfe{} $m=2$ & 3 & 4.000, 3.998 & 3.749, 3.609 & $6.482\times10^{-14}$ & $4.441\times10^{-15}$ & Bounded candidate smoke only \\
Trapezoidal & 2 & 1.973, 1.997 & 2.078, 2.028 & $1.650\times10^{-10}$ & $1.987\times10^{-12}$ & Bounded candidate smoke only \\
\bottomrule
\end{tabular}}
\end{table}

The same source-model evidence record also includes a non-heavy full-horizon candidate
probe at the original \tfe{} pendulum horizon $T=10$.  This diagnostic uses
$h=\{0.1,0.05,0.025\}$ and a coarse reference $h_{\rm ref}=0.0125$, so it
tests full-horizon large-step behavior without invoking the extracted
$h_{\rm ref}=10^{-4}$ source-policy reference.  All four active rows are
finite and Newton-residual stable, but the table remains a candidate-runner
diagnostic rather than a source-policy reproduction; it closes zero
source-policy rows.

\begin{table}[htbp]
\centering
\caption{Full-horizon coarse candidate probe for the original \tfe{} pendulum rows.  The run uses $T=10$, $h_{\rm ref}=0.0125$, and $h=\{0.1,0.05,0.025\}$; it does not invoke the source-policy $h_{\rm ref}=10^{-4}$ reference.}
\label{tab:tfe-full-t10-coarse-candidate}
\scriptsize
\setlength{\tabcolsep}{3pt}
\renewcommand{\arraystretch}{0.88}
\resizebox{\linewidth}{!}{%
\begin{tabular}{P{0.20\linewidth}P{0.10\linewidth}P{0.20\linewidth}P{0.20\linewidth}P{0.13\linewidth}P{0.13\linewidth}P{0.18\linewidth}}
\toprule
Original-row method & Formula target & Velocity pair orders & Coordinate pair orders & Finest velocity error & Finest coordinate error & Claim boundary \\
\midrule
Newmark--$\beta$ & 2 & 1.991, 1.998 & 2.028, 2.007 & $4.642\times10^{-3}$ & $1.956\times10^{-3}$ & Full-$T$ coarse candidate only \\
\tfe{} $m=1$ & 1 & 1.276, -0.077 & 2.623, 5.795 & $6.801\times10^{-3}$ & $5.226\times10^{-5}$ & Full-$T$ coarse candidate only \\
\tfe{} $m=2$ & 3 & 2.927, 2.952 & 2.694, 2.824 & $2.375\times10^{-6}$ & $3.346\times10^{-7}$ & Full-$T$ coarse candidate only \\
Trapezoidal & 2 & 1.993, 1.998 & 2.021, 2.005 & $3.558\times10^{-3}$ & $1.504\times10^{-3}$ & Full-$T$ coarse candidate only \\
\bottomrule
\end{tabular}}
\end{table}

\begin{table}[htbp]
\centering
\caption{Cross-paper evidence accumulated without treating it as an external superiority result.}
\label{tab:cross-paper-evidence}
\tiny
\setlength{\tabcolsep}{3pt}
\renewcommand{\arraystretch}{0.80}
\begin{tabular}{P{0.27\linewidth}P{0.16\linewidth}P{0.30\linewidth}P{0.15\linewidth}}
\toprule
Tranche & Rows & Main metric & Interpretation \\
\midrule
2021 public $rA/rp/r\epsilon$ baselines &
$27+9$ order rows; 12 timing rows &
Public $T=3$ policies; $h=10^{-2},10^{-3},10^{-4}$ except the
double-pendulum dynamic self-reference sweep
$h=10^{-2},2\times10^{-3},10^{-3}$. &
Baseline execution evidence only. \\
\method{} single public horizon &
3 rows &
$T=3$, $h=10^{-2},10^{-3},10^{-4}$; finest position/velocity errors
$2.584\times10^{-14}/7.012\times10^{-15}$. &
Public-policy rows closed, but order is reference-floor limited. \\
\method{} double public horizon, coarse &
3 rows &
$T=3$, $h=0.1,0.05,0.025$; reference $h=0.0125$; position/velocity orders
$7.951/7.042$. &
Coarse-first order/work evidence, not the public $10^{-4}$ policy. \\
Public double coarse same-window baselines &
9 rows, 6 ok &
$rA/r\epsilon$ at the same $T=3$, $h=0.1,0.05,0.025$, reference
$h=0.0125$ give position/velocity orders $0.703/0.754$; $rp$ fails Newton at
all three coarse steps. &
Large-step baseline evidence; not the source $10^{-4}$ policy. \\
\method{} and public single coarse same-window &
12 raw rows; 4 summary rows &
$T=3$, $h=0.1,0.05,0.025$, reference $h=0.0125$; \method{} position order
$6.054$; public velocity orders about $1.21$. &
Coarse-first evidence; no default $10^{-4}$ run or superiority claim. \\
\method{} closed-loop public mechanisms &
12 rows &
Selected and public-horizon residual trios for four-link/slider-crank; finest
public residuals $5.136\times10^{-13}/1.113\times10^{-14}$. &
Residual evidence, not dynamic order/work superiority. \\
Closed-loop surrogate and floor analysis &
2+2 rows &
Selected-window residual-to-error ratios for four-link/slider-crank; no
source-policy dynamic-order row is accepted; the floor analysis has two
velocity/acceleration evidence rows and two position/reference-floor limitations. &
Narrows the boundary; not external superiority. \\
2022 half-implicit pilot &
24 rows &
Bounded $T=0.1$, $h=0.02,0.01,0.005$ for $rA/rA_{\mathrm{half}}$ on four
examples. &
Source-family candidate; full $T=8$ policy remains open. \\
Closed-loop coarse trajectory diagnostics &
6 rows &
Four-link primary-state orders $5.955/5.955/6.085/5.971$; slider-crank
orders $6.164/6.159/7.341/6.426$ at
$T=0.1$, $h=0.1,0.05,0.025$, reference $h=0.0125$. &
Local coarse trajectory evidence only; not accepted dynamic-order proof or
public superiority. \\
Bounded common-reference work/precision diagnostic &
24 rows &
Local \method{} and public $rA/rp/r\epsilon$ rows use the same
finite accepted-method endpoint reference; the fixed-reference diagnostic gap count is zero. &
Figure~\ref{fig:strict-common-reference}; no external superiority claim. \\
Coarse-first readiness diagnostic &
4 rows &
Single and double pendulum are coarse same-window order/time ready; four-link
and slider-crank have local coarse trajectory diagnostics and strict
common-reference work/precision rows; the $10^{-4}$ policy is outside the
default evidence package. &
Readiness diagnostic, not a superiority claim. \\
Coverage matrix &
48 rows &
32 complete, 0 partial; no external-superiority claim. &
Diagnostic coverage table, not a paper claim. \\
\bottomrule
\end{tabular}
\end{table}


\section{Reproducibility package}
\label{sec:repro-appendix}

The flat LaTeX archive contains the editable Elsevier source, the compiled
manuscript, the thirteen figures, and declaration files. The companion
repository-facing reproducibility package contains the scripts needed to check
the reported tables. The paper-level checks read existing result records unless
a full numerical rebuild is explicitly requested; this keeps editorial checks
separate from computational campaigns. Companion records encode the accepted
claim scope and the non-claims, but the mathematical and numerical arguments
in the paper are the authority.

The current reproducibility package is therefore two-tiered.  The
reviewer-facing minimal candidate is a 10-file package with one 158-line
Python replay core; it is a shrink target for the reported-result records and
proof-dependency records, not a complete runner-centered source-policy
archive.  The local accepted-row companion is human-runnable and
self-contained for the four local examples and the 12 accepted local rows.
A broader external source-policy runner archive would require the 40
source-policy rows, the \tfe{} source-policy runner, and regeneration of both
the accepted local rows and any promoted source-policy rows from a standalone
runner.  That archive is outside the current paper package.

\begin{table}[htbp]
\centering
\caption{External reproduction and package boundaries retained by the paper package.}
\label{tab:objective-completion-blockers}
\scriptsize
\setlength{\tabcolsep}{3pt}
\renewcommand{\arraystretch}{0.9}
\begin{tabular}{P{0.16\linewidth}P{0.42\linewidth}P{0.30\linewidth}}
\toprule
Boundary & Current evidence state & Manuscript consequence \\
\midrule
External same-test reproduction &
External source-policy reproduction has not been established on the source
paper tests, and no external rows are treated as source-policy closed. &
No source-paper superiority or source-policy direct-error claim is made. \\
Original \tfe{} runner package &
A complete source-policy \tfe{} runner package would need the source pendulum
DAE runner, \tfe{}/comparator source-policy method runners, Gauss6
source-policy runner, full-T10 endpoint policy, and accepted work-precision
source rows. &
TFE rows are formula/common-reference diagnostics, not full-TFE or
source-policy runner evidence. \\
Runner archive boundary &
The local accepted-row runner package covers the bounded replay/provenance
role for the accepted rows; a full source-policy runner archive is outside
this package. &
The package supports accepted-row provenance while keeping full
source-policy runner evidence outside the accepted theorem and comparison
claims. \\
\bottomrule
\end{tabular}
\end{table}

External reproduction/package boundary: the three rows in
Table~\ref{tab:objective-completion-blockers} are scope boundaries, not
theorem hypotheses. They do not change the conditional theorem proof; they
keep source-policy reproduction, full-\tfe{} runner evidence, and a full
source-policy runner archive outside the accepted theorem and comparison
claims.
Notation boundary: theorem P6 denotes only the compact-tube solver-scale
interface \(\eta_h^{\rm tube}\le c_\eta h^7\), whereas the separate
original-\tfe{} source-policy DAE-runner/package gap is an external
reproduction boundary, not a theorem condition.

The implementation record identifies \nolinkurl{residual_cylindrical_chain},
\nolinkurl{R_VALUE}, and \nolinkurl{R_JAC} as the implemented FullVA residual
and Jacobian path. This supports independent reruns of the reported evidence;
it is not a substitute for the proof.

\begin{table}[htbp]
\centering
\caption{Algorithm-to-source map for the accepted \method{} execution path.}
\label{tab:algorithm-source-map}
\footnotesize
\setlength{\tabcolsep}{3pt}
\begin{tabular}{P{0.22\linewidth}P{0.34\linewidth}P{0.30\linewidth}}
\toprule
Algorithm block & Source location & Reproduction role \\
\midrule
Mechanism setup and examples &
\nolinkurl{v047_cylindrical_chain_pipeline/run_v047.py} &
Builds the cylindrical-chain data and the four ASME-style example rows used
by the accepted method-side evidence. \\
Residual and Jacobian path &
\nolinkurl{residual_cylindrical_chain}, \nolinkurl{R_VALUE},
\nolinkurl{R_JAC} &
Defines the 132-row FullVA residual and the single accepted
residual/Jacobian path used by Newton. \\
One-step solve and trajectory loop &
\nolinkurl{gauss_step}, \nolinkurl{integrate}, \nolinkurl{run_case} &
Implements the Algorithm~1 stage solve, endpoint reconstruction, endpoint
closure, and step diagnostics. \\
Method result records &
\nolinkurl{v047_cylindrical_chain_pipeline/results} &
Record the h-sweeps, endpoint closure diagnostics, four-example rows, and
accepted method-side order/coverage evidence. \\
Cross-paper comparison records &
\nolinkurl{v048_cross_paper_same_test_benchmarks/results} &
Store common-reference diagnostics and external-row status; these records do
not authorize source-paper superiority claims. \\
Package-level checks &
\nolinkurl{validate_paper_package.py --latex} and
\nolinkurl{validate_pipeline_outputs.py} &
Rebuild and check the manuscript package and repository inventory without
invoking a new full numerical campaign. \\
\bottomrule
\end{tabular}
\end{table}

\begin{table}[htbp]
\centering
\caption{Compact reproducibility contract used by the bounded paper package.}
\label{tab:artifact-contract}
\begin{tabular}{P{0.31\linewidth}P{0.53\linewidth}}
\toprule
Field & Paper-level value \\
\midrule
Accepted method & Gauss6/FullVA. \\
Conditional method order & Sixth-order path of
Theorem~\ref{thm:g6fullva-order}. \\
Smooth observed orders & Position/velocity 7.161/7.066. \\
Local method-side coverage examples & single pendulum, double pendulum, four-link, and
slider-crank. \\
Accepted dynamic order rows & Single and double pendulum. \\
Coverage-only for dynamic order & Four-link and slider-crank; their
closed-loop kinematic/reaction rows are not
accepted external dynamic-order rows. \\
Comparator & Local paper-style $m=3$ Gauss--Lobatto \tfe{} formula target. \\
Comparator expected order & 5. \\
Numerical comparison scope & Forty-four method/example common-reference
order-error cells, with external rows kept outside source-paper superiority. \\
Runner-package status & The local accepted-row companion is executable for
the 12 accepted local rows; a full source-policy runner archive is outside
the current package and would require separate source-policy closure and the
original-\tfe{} source-policy runner. \\
Non-claims & No full \tfe{} replacement, no source-paper superiority, and no
paper-level direct-error superiority against external methods. \\
\bottomrule
\end{tabular}
\end{table}

\noindent\textbf{Comparison reconciliation statement.}\par\noindent
The finite-grid common-reference diagnostic is arithmetically complete, but
the nonlocal cells remain diagnostic comparison records. \mbox{These rows remain}
diagnostic scope records and are not promoted to a CMAME paper-level
external-superiority claim.

\begin{table}[htbp]
\centering
\caption{Cross-example and review records retained with the source package.}
\label{tab:all-example-review-artifacts}
\begin{tabular}{P{0.31\linewidth}P{0.53\linewidth}}
\toprule
Record & Role \\
\midrule
All-example result matrix & Checks every one of the 44 method/example cells
and separates the 15 nonlocal rows that are sensitive to external-row
policy. \\
External-row table & Lists the row-level evidence still needed before any
external row can support source-paper superiority. \\
Case reconciliation record & Separates full-campaign not-run markers from the
bounded/coarse evidence already present for the RA2021 and HI2022 families. \\
Proof and package-boundary record & Keeps proof assumptions, dynamic-row closure
requirements, open obligations, all-example numerical sanity, and current
package boundary traceable while external reproduction items remain open. \\
\bottomrule
\end{tabular}
\end{table}

\begin{table}[htbp]
\centering
\caption{Source-package reproducibility and consistency checks.}
\label{tab:validators}
\footnotesize
\setlength{\tabcolsep}{4pt}
\renewcommand{\arraystretch}{0.92}
\begin{tabular}{P{0.34\linewidth}P{0.50\linewidth}}
\toprule
Check group & Role \\
\midrule
Source-format check &
Checks the Elsevier source, PDF, highlights, declarations, clean log, and
non-claim scope. \\
Archive-completeness check &
Checks the source archive, required figures, clean logs, and the rule that
the full numerical generator is not started by editorial checks. \\
Order-scope check &
Checks that single/double pendulum carry the accepted dynamic-order rows,
four-link/slider-crank remain coverage-only for dynamic order, and
$10^{-4}$ is not a default execution policy. \\
Numerical-matrix check &
Checks the 44-cell numerical matrix against the common-reference and reconciliation
records. \\
External-row check &
Checks the separated nonlocal rows, zero claim-ready external rows, and the
non-default-\(10^{-4}\) execution rule. \\
Proof-dependency check &
Checks proof-label dependency links, theorem-scope assumptions, and open
symbolic/primitive-route obligations that are not active direct-PC2
obligations. \\
Scope-consistency review &
Checks that the all-example matrix, proof-dependency record, and stated
non-claims use the same claim boundary. \\
Package check &
Runs the paper, source-archive, source-comparison, proof-matrix, four-example, and
accepted-method package checks without starting a new full numerical campaign. \\
Repository-provenance check &
Checks the broader repository inventory and records that the full numerical
generator was not started by the source-package checks. \\
\bottomrule
\end{tabular}
\end{table}

The source package includes the command file for these read-only checks.

\section{Supplementary source package}

The Elsevier source package contains the editable \nolinkurl{elsarticle}
manuscript source, the compiled manuscript PDF, the separate highlights file,
the declaration statements, and the thirteen figure files used in this paper.
A supporting order-scope record separates accepted dynamic order from
four-example mechanism coverage. A flat Editorial-Manager-oriented source
folder, \nolinkurl{cmame_submission_flat}, is also supplied. In that folder,
the submission TeX source and all referenced figure files live at the same
directory level, avoiding a dependency on the repository's figure
subdirectory.
The reproducibility package is therefore two-tiered: a compact 10-file
package with one 158-line Python replay core records the bounded local replay,
while the full development repository remains provenance for the broader
development record. The local accepted-row companion is human-runnable and
records 12 accepted local rows, but it is not a complete runner-centered
source-policy archive. A full source-policy runner archive is not included;
it would require closing the 40 source-policy rows under their source
policies.
The reviewer-facing runnable evidence for the narrowed current claim is the
\nolinkurl{cmame_narrowed_repro_bundle} directory. Its
\nolinkurl{run_narrowed_repro_bundle.py} entry point replays the paper result
matrix and launches the four self-contained local accepted-row runners, while
recording source-policy external-row closure and the full source-policy
runner archive as open boundaries. Thus the local reproducibility bundle
supports the present bounded claim boundary, but it is not a source-policy
runner archive and does not promote external-superiority rows.

\begin{thebibliography}{00}

\bibitem{chaturvedi2026tfe}
E. Chaturvedi, C. Sandu, A. Sandu,
Higher-order integration of index-3 DAE with friction using time finite
elements on Lie groups,
Multibody System Dynamics (2026).
doi:10.1007/s11044-026-10153-w.

\bibitem{kim2017higher}
W. Kim, J.N. Reddy,
A new family of higher-order time integration algorithms for the analysis of
structural dynamics,
Journal of Applied Mechanics 84 (7) (2017) 071008.
doi:10.1115/1.4036821.

\bibitem{hairer1993solving}
E. Hairer, G. Wanner,
Solving Ordinary Differential Equations II: Stiff and Differential-Algebraic
Problems,
Springer, Berlin, 1996.

\bibitem{gear1985automatic}
C.W. Gear, B. Leimkuhler, G.K. Gupta,
Automatic integration of Euler--Lagrange equations with constraints,
Journal of Computational and Applied Mathematics 12--13 (1985) 77--90.
doi:10.1016/0377-0427(85)90008-1.

\bibitem{munthe1998lie}
H. Munthe-Kaas,
Runge-Kutta methods on Lie groups,
BIT Numerical Mathematics 38 (1998) 92--111.

\bibitem{munthe1999high}
H. Munthe-Kaas,
High order Runge-Kutta methods on manifolds,
Applied Numerical Mathematics 29 (1999) 115--127.

\bibitem{bruls2012lie}
O. Bruls, A. Cardona, M. Arnold,
Lie group generalized-alpha time integration of constrained flexible
multibody systems,
Mechanism and Machine Theory 48 (2012) 121--137.
doi:10.1016/j.mechmachtheory.2011.07.017.

\bibitem{wieloch2021bdf}
V. Wieloch, M. Arnold,
BDF integrators for constrained mechanical systems on Lie groups,
Journal of Computational and Applied Mathematics 387 (2021) 112517.
doi:10.1016/j.cam.2019.112517.

\bibitem{terze2016singularity}
Z. Terze, A. Mueller, D. Zlatar,
Singularity-free time integration of rotational quaternions using
non-redundant ordinary differential equations,
Multibody System Dynamics 38 (2016) 201--225.
doi:10.1007/s11044-016-9518-7.

\bibitem{terze2020aircraft}
Z. Terze, D. Zlatar, V. Pandza,
Aircraft attitude reconstruction via novel quaternion-integration procedure,
Aerospace Science and Technology 97 (2020) 105617.

\bibitem{taves2020exponential}
J. Taves, A. Kissel, D. Negrut,
On an exponential map approach for rigid body kinematics and dynamics
analysis,
Technical Report TR-2020-08, Simulation-Based Engineering Laboratory,
University of Wisconsin--Madison (2020).

\bibitem{kissel2021dwelling}
A. Kissel, D. Negrut, J. Taves,
Dwelling on the connection between SO(3) and rotation matrices in rigid
multibody dynamics. Part 1: Description of an index-3 DAE solution approach,
in: Proceedings of the ASME 2021 International Design Engineering Technical
Conferences and Computers and Information in Engineering Conference,
IDETC-CIE2021, 2021.

\bibitem{kissel2022absolute}
A. Kissel, D. Negrut, J. Taves,
Constrained multibody kinematics and dynamics in absolute coordinates: a
discussion of three approaches to representing rigid body rotation,
Journal of Computational and Nonlinear Dynamics 17 (10) (2022) 101008.

\bibitem{fang2023halfimplicit}
L. Fang, A. Kissel, R. Zhang, D. Negrut,
On the use of half-implicit numerical integration in multibody dynamics,
Journal of Computational and Nonlinear Dynamics 18 (1) (2023) 014501.
doi:10.1115/1.4056183.

\bibitem{kissel2024velocitypartitioning}
A. Kissel, L. Bakke, D. Negrut,
Reducing the constrained multibody dynamics problem to the solution of a
system of ordinary differential equations via velocity partitioning and Lie
group integration,
Journal of Computational and Nonlinear Dynamics 19 (7) (2024).
doi:10.1115/1.4065254.

\bibitem{jay1996symplectic}
L.O. Jay,
Symplectic partitioned Runge--Kutta methods for constrained Hamiltonian
systems,
SIAM Journal on Numerical Analysis 33 (1) (1996) 368--387.
doi:10.1137/0733019.

\bibitem{marsden2001discrete}
J.E. Marsden, M. West,
Discrete mechanics and variational integrators,
Acta Numerica 10 (2001) 357--514.
doi:10.1017/S096249290100006X.

\bibitem{leyendecker2008variational}
S. Leyendecker, J.E. Marsden, M. Ortiz,
Variational integrators for constrained dynamical systems,
ZAMM -- Journal of Applied Mathematics and Mechanics 88 (9) (2008) 677--708.
doi:10.1002/zamm.200700173.

\bibitem{leok2012prolongation}
M. Leok, T. Shingel,
Prolongation--collocation variational integrators,
IMA Journal of Numerical Analysis 32 (3) (2012) 1194--1216.
doi:10.1093/imanum/drr042.

\bibitem{stewart1996implicit}
D.E. Stewart, J.C. Trinkle,
An implicit time-stepping scheme for rigid body dynamics with inelastic
collisions and Coulomb friction,
International Journal for Numerical Methods in Engineering 39 (15) (1996)
2673--2691.
doi:10.1002/(SICI)1097-0207(19960815)39:15<2673::AID-NME972>3.0.CO;2-I.

\bibitem{anitescu1997formulating}
M. Anitescu, F.A. Potra,
Formulating dynamic multi-rigid-body contact problems with friction as
solvable linear complementarity problems,
Nonlinear Dynamics 14 (3) (1997) 231--247.
doi:10.1023/A:1008292328909.

\bibitem{tasora2011matrixfree}
A. Tasora, M. Anitescu,
A matrix-free cone complementarity approach for solving large-scale,
nonsmooth, rigid body dynamics,
Computer Methods in Applied Mechanics and Engineering 200 (2011) 439--453.
doi:10.1016/j.cma.2010.06.030.

\bibitem{borri1988rigid}
M. Borri, S.N. Atluri,
Time-finite element method for the constrained dynamics of a rigid body,
in: Computational Mechanics, Springer, Berlin, 1988.

\bibitem{hughes1988spacetime}
T.J.R. Hughes, G.M. Hulbert,
Space-time finite element formulations for elastodynamics: formulations and
error estimates,
Computer Methods in Applied Mechanics and Engineering 66 (1988) 339--363.

\bibitem{hulbert1992time}
G.M. Hulbert,
Time finite element methods for structural dynamics,
International Journal for Numerical Methods in Engineering 33 (1992)
307--331.

\bibitem{jog2016robust}
C.S. Jog, M. Agrawal, A. Nandy,
The time finite element as a robust general scheme for solving nonlinear
dynamic equations including chaotic systems,
Applied Mathematics and Computation 279 (2016) 43--61.

\bibitem{bauchau2024time}
O.A. Bauchau,
The finite element method in time for multibody dynamics,
Journal of Computational and Nonlinear Dynamics 19 (7) (2024) 071001.
doi:10.1115/1.4063953.

\bibitem{browmcphee2016friction}
P. Brown, J. McPhee,
A continuous velocity-based friction model for dynamics and control with
physically meaningful parameters,
Journal of Computational and Nonlinear Dynamics 11 (2016) 054502.

\end{thebibliography}

\end{document}
